\chapter{Control Flow}
EDADSL provides several control flow constructs for conditional execution, iteration, and program flow control.

\section{Conditional Execution}
Conditional execution uses the \texttt{if} keyword followed by a boolean expression and a code block. Additional conditions can be tested with \texttt{otherwise}. A final \texttt{else} clause provides a fallback code block executed when no conditions match:
\begin{verbatim}
EBNF:
conditional = "if", boolean_expression, 
              "::", code_block,
              { "otherwise", boolean_expression, 
                "::", code_block },
              ["else", "::", code_block]

Example:
$flag [bool]= True
if $flag ::{
    f.util.print("Flag is set")
} otherwise $other ::{
    f.util.print("Other condition")
} else ::{
    f.util.print("Flag is not set")
}
\end{verbatim}

Basic \texttt{if} without \texttt{else}:
\begin{verbatim}
Example:
$flag [bool]= True
if $flag ::{
    f.util.print("Flag is set")
}
\end{verbatim}

\section{For Loops}
For loops iterate over a list or set:
\begin{verbatim}
EBNF:
iterator = "for", variable, "E", 
           (list | set), "::", code_block

Example:
$colors [list]= ["red", "green", "blue"]
for $c E $colors ::{
    f.util.print($c)
}
\end{verbatim}
The loop variable (\texttt{\$c}) takes on each value in the list/set in sequence.

\paragraph{Variable Scoping}
The loop variable is block-scoped: it is only visible inside the loop body and is destroyed when the loop exits. If the loop iterates over an empty list, the loop body never executes and the variable is never initialized; accessing it after the loop raises a \texttt{Runtime Error}.

\section{While Loops}
While loops repeat a code block as long as a boolean condition evaluates to \texttt{True}.

\subsection{Basic While Loop}
\begin{verbatim}
EBNF:
while_loop = "while", boolean_expression, 
             "::", code_block

Example:
$i [int]= 0
while $i < 5 ::{
    f.util.print($i)
    $i = $i + 1
}
\end{verbatim}

\paragraph{Variable Scoping}
While loops do not introduce any implicit variables. Any variables declared inside the loop body are block-scoped and destroyed when the loop exits. Variables used in the loop condition must be declared before the loop.

\subsection{While with Else Clause}
An optional \texttt{else} clause executes when the loop condition becomes \texttt{False}. If the loop exits via \texttt{break}, the \texttt{else} clause is \textbf{not} executed:
\begin{verbatim}
EBNF:
while_loop = "while", boolean_expression, 
             "::", code_block,
             ["else", "::", code_block]

Example:
$i [int]= 0
while $i < 5 ::{
    $i = $i + 1
} else ::{
    f.util.print("Loop completed normally")
}

Example with break:
$i [int]= 0
while $i < 10 ::{
    if $i == 3 ::{
        break
    }
    $i = $i + 1
} else ::{
    f.util.print("This will NOT be printed")
}
\end{verbatim}

\subsection{Do-While Loops}
Do-while loops execute the code block \textbf{at least once} before evaluating the condition:
\begin{verbatim}
EBNF:
do_while = "do", "::", code_block,
           "while", boolean_expression

Example:
$count [int]= 0
do ::{
    $count = $count + 1
    f.util.print($count)
} while $count < 5
\end{verbatim}
The code block executes first, then the condition is checked. If the condition is \texttt{True}, the loop repeats. The same block-scoping rules as while loops apply: variables declared inside the code block are destroyed when the block exits. If the code block raises an error, the loop terminates immediately and the condition is not evaluated. \texttt{break} and \texttt{continue} work as expected inside do-while loops.

\subsection{Break and Continue}
\texttt{break} exits the innermost loop immediately. \texttt{continue} skips the remainder of the current iteration and proceeds to the next evaluation of the loop condition:
\begin{verbatim}
EBNF:
break_stmt    = "break"
continue_stmt = "continue"

Example:
for $x E $values ::{
    if $x < 0 ::{
        continue
    }
    if $x > 100 ::{
        break
    }
    f.util.print($x)
}

Example with while:
$i [int]= 0
while $i < 10 ::{
    $i = $i + 1
    if $i % 2 == 0 ::{
        continue
    }
    f.util.print($i)
}
\end{verbatim}
\texttt{break} and \texttt{continue} apply only to the innermost enclosing loop.

\subsection{Nested While Loops}
While loops may be nested to arbitrary depth:
\begin{verbatim}
Example:
$rows [int]= 3
$cols [int]= 3
$i [int]= 0
while $i < $rows ::{
    $j [int]= 0
    while $j < $cols ::{
        f.util.print("Cell (" + f.str.str($i) + "," + f.str.str($j) + ")")
        $j = $j + 1
    }
    $i = $i + 1
}
\end{verbatim}

\subsection{Safety: Maximum Iteration Limit}
To prevent infinite loops, the interpreter enforces a maximum iteration count per loop. The limit is controlled by \texttt{syst.flag.max\_loop\_iter} (default: 100000). If a loop exceeds this limit, a \texttt{Value Error} is raised:
\begin{verbatim}
Example:
syst.flag.max_loop_iter 500
$i [int]= 0
while $i < 1000 ::{
    $i = $i + 1
}
f.util.print($i)  # never reached — Value Error at 501 iterations
\end{verbatim}
The limit applies independently to each loop. Nested loops each have their own counter.

\subsection{Infinite Loops}
An infinite loop can be written with \texttt{True} as the condition. Use \texttt{break} to exit:
\begin{verbatim}
Example:
$count [int]= 0
while True ::{
    $count = $count + 1
    if $count == 10 ::{
        break
    }
}
f.util.print($count)   #c# Prints: 10
\end{verbatim}

\section{Goto Statements}
Goto statements provide unconditional jumps to labeled positions:
\begin{verbatim}
EBNF:
gotomarker = "gotomarker", id
gotojump   = "gotojump", id

Example:
gotojump skip_section
f.util.print("This will not be printed")
gotomarker skip_section
f.util.print("Jumped here")
\end{verbatim}

\paragraph{Warning -- Deprecated Usage}
Goto statements are \textbf{strongly discouraged}. They make code difficult to read, debug, and maintain. Prefer structured control flow (\texttt{if}, \texttt{for}, \texttt{while}, \texttt{break}, \texttt{continue}) for all new code. Gotos are retained for backward compatibility only. Gotos may cross loop boundaries (e.g., jump from outside a loop into the loop body) but may not cross function boundaries.

\section{End of Program Statements}
\begin{verbatim}
EBNF:
halt = "halt"
noop = "skip"
\end{verbatim}
\begin{itemize}
\item \texttt{halt} -- Terminates program execution and exports all created files
\item \texttt{skip} -- No operation; program continues to next statement
\end{itemize}

\section{Output and Logging}
\begin{verbatim}
EBNF:
echo_statement = "echo", variable
writetolog     = "writetolog", string_expression

Example:
f.util.print("Hello, EDADSL!")
writetolog "Board design completed successfully"
\end{verbatim}
\begin{itemize}
\item \texttt{echo} -- Keyword that prints a variable reference to the console. Only accepts a variable, not an arbitrary expression. Use \texttt{f.util.print()} for expression output.
\item \texttt{writetolog} -- Keyword that appends a message to the log file
\end{itemize}

\chapter{Functions}
\section{Built-in Functions}

\textbf{Important:} All built-in function calls use a two-level namespace: \texttt{f.\textit{namespace}.\textit{function}}. This is defined in the formal grammar (Chapter~3):
\begin{verbatim}
function_call_built_in = "f", ".", namespace, ".", id, "(",
                        [{ data, "," }, data], ")"
namespace = "math" | "trig" | "hyp" | "cplx" | "log"
          | "special" | "str" | "list" | "arr"
          | "stat" | "linalg" | "ee" | "int" | "util"
\end{verbatim}
Correct usage: \texttt{f.math.abs(-3.2)}, \texttt{f.trig.sin(\$x)}, \texttt{f.arr.zeros([3,3])}\\
Incorrect: \texttt{f.math.abs(-3.2)}, \texttt{f.trig.sin(\$x)}, \texttt{abs(-3.2)}

\texttt{f.util.print()} takes a \textbf{single} argument. Use \texttt{f.str.str()} for string concatenation:
\begin{verbatim}
f.util.print("value = " + $x)              # OK
f.util.print(f.str.str("value = ", $x))    # OK
f.util.print("value = ", $x)               # WRONG - two arguments
\end{verbatim}

\subsection{Mathematical Functions}
\subsubsection{Absolute Value Function}
The Absolute Value Function ( denoted as \texttt{abs} ) returns the Absolute Value of the Argument or the Magnitude of a Vector. \\
Morphology of the Absolute Value Function:
\begin{verbatim}
EBNF:
"f.math.abs", "(", argument, [ ",", "vec", "=", 
       boolean_data_type ], ")"
\end{verbatim}
If the Type of the Argument does not match the expected Type (a Numeric Value or a List of Numeric Values) a Value Error is raised.\\
Example:
\begin{verbatim}
f.math.abs(-3.2)
\end{verbatim}
This Example would return a value of \texttt{3.2}.

The Argument may also be a List of Numeric Values. If This is the case the Function is iterated over each Element.\\
Example:
\begin{verbatim}
f.math.abs([-3, 0.2, 0, -7])
\end{verbatim}
This Example would return \texttt{[3, 0.2, 0, 7]}. \vspace{5 mm}\\
If \texttt{vec = True} is passed a List shall be passed as the Argument and the Function will return Magnitude of the List as an n-dimensional Vector, where n is the Length of the List.\\ %$\texttt{abs}(\vec{v}\texttt{, vec = True}) = \sqrt{\sum_{i=1}^nv_i^2}$. 
Example:
\begin{verbatim}
f.math.abs( [ 3, -4 ], vec = True )
\end{verbatim}
This Example would return \texttt{5.0}.\\
There are no Domain-Restrictions for the function inside $\mathbb{R}^n$.

\paragraph{Edge Cases -- Absolute Value}
\begin{itemize}
\item A \texttt{Type Error} is raised if the argument is not \texttt{int}, \texttt{real}, \texttt{complex}, or a list of numeric values.
\item A \texttt{Type Error} is raised if \texttt{vec = True} and the argument is not a list.
\item A \texttt{Type Error} is raised if \texttt{vec = True} and the list contains complex numbers (vector magnitude is only defined for real vectors).
\end{itemize}

\paragraph{Complex Overload}
When the Argument is a Complex Number $z = a + bi$, \texttt{abs} returns the Modulus (Magnitude) of the Complex Number, defined as:
\[
\lvert a + bi \rvert = \sqrt{a^2 + b^2}
\]
The Result is always a Real Number, even if the Argument is Complex.\\
Domain: All Complex Numbers ($\mathbb{C}$).\\
Example:
\begin{verbatim}
f.math.abs(3 + 4i)   #c# Returns 5.0
f.math.abs(1 + 1i)   #c# Returns approximately 1.41421
f.math.abs(0 - 5i)   #c# Returns 5.0
\end{verbatim}

\subsubsection{Parallel Equivalent Function}
The Parallel Equivalent Function ( denoted as \texttt{parallel} ) computes the parallel combination of a list of values. This is the standard formula for parallel resistances, conductances, or any set of admittances. \\
Morphology of the Parallel Equivalent Function:
\begin{verbatim}
EBNF:
"f.math.parallel", "(", argument, ")"
\end{verbatim}
If the Type of the Argument does not match the expected Type (a List of Numeric Values) a Value Error is raised.\\
The Argument must be a List of Numeric Values (\texttt{int}, \texttt{real}, or \texttt{complex}). The function computes:
\[
\text{parallel}([v_1, v_2, \ldots, v_n]) = \frac{1}{\frac{1}{v_1} + \frac{1}{v_2} + \cdots + \frac{1}{v_n}}
\]
If any Element in the List is Zero, the Function immediately Returns Zero.\\
If any Element is Infinite, it is Ignored (does not contribute to the Result).\\
The Return Type is determined by the Numerical Value of the Result:
\begin{itemize}
\item If the Result is Numerically Equal to an Integer, Return \texttt{int}
\item Otherwise, if the Result is Numerically Equal to a Real, Return \texttt{real}
\item Otherwise, Return \texttt{complex}
\end{itemize}
Examples:
\begin{verbatim}
f.math.parallel([100, 100])         #c# Returns 50 (int)
f.math.parallel([100, 200, 300])    #c# Returns 54 (int, softint)
f.math.parallel([1e3, 2.2e3])       #c# Returns 687.5 (real)
f.math.parallel([100, 0])           #c# Returns 0 (any zero)
f.math.parallel([100, 200, 0, 300]) #c# Returns 0 (any zero)
f.math.parallel([100, inf])           #c# Returns 100 (inf ignored)
f.math.parallel([100, 200, inf])      #c# Returns 66.666... (inf ignored)
\end{verbatim}

\paragraph{Complex Overload}
When the List contains Complex Numbers, the function computes the parallel combination using Complex Arithmetic. If the Result has a Zero Imaginary Part, it is Returned as \texttt{real}. \\
Example:
\begin{verbatim}
f.math.parallel([100 + 0i, 200 + 0i])
    #c# Returns 66.666... (real, imag=0)
f.math.parallel([3 + 4i, 1 - 2i])           #c# Returns complex result
\end{verbatim}
Domain: All Lists of Numeric Values. Empty Lists raise a Value Error. Lists with a single Element return that Element unchanged (after \texttt{softint}(\texttt{softreal}())).\\

\subsubsection{Signum Function}
The Signum Function ( denoted as \texttt{sign} ) returns the Sign of the Argument. \\
Morphology of the Signum Function:
\begin{verbatim}
EBNF:
"f.math.sign", "(", argument, ")"
\end{verbatim}
If the Type of the Argument does not match the expected Type (a Numeric Value or a List of Numeric Values) a Value Error is raised.\\
Example:
\begin{verbatim}
f.math.sign(-3.2)
\end{verbatim}
This Example would return a value of \texttt{-1}.

The Argument may also be a List of Numeric Values. If This is the case the Function is iterated over each Element.\\
Example:
\begin{verbatim}
f.math.sign([-3, 0.2, 0, -7])
\end{verbatim}
This Example would return \texttt{[-1, 1, 0, -1]}.\\
There are no Domain-Restrictions for the function inside $\mathbb{R}^n$.

\paragraph{Complex Overload}
When the Argument is a Complex Number $z = a + bi$, \texttt{sign} returns the unit complex $z / |z|$. For $z = 0$, \texttt{sign} returns \texttt{0} (int).
\[
\text{sign}(z) = \begin{cases} \frac{z}{|z|} & z \neq 0 \\ 0 & z = 0 \end{cases}
\]
Domain: All Complex Numbers ($\mathbb{C}$).
Examples:
\begin{verbatim}
f.math.sign(3 + 4i)     #c# Returns approximately 0.6 + 0.8i
f.math.sign(0 + 0i)     #c# Returns 0 (int)
f.math.sign(1 + 0i)     #c# Returns 1 + 0i
\end{verbatim}


\subsubsection{Sine Function}
Computes the sine of an angle. Angles are interpreted as radians by default. \\
Morphology:
\begin{verbatim}
EBNF:
"f.trig.sin", "(", ( int | real | complex ), [ ",", "unit", "=", string_type ], ")"
\end{verbatim}
\begin{itemize}
\item Operands: \texttt{int}, \texttt{real}, \texttt{complex}, or list of numeric values.
\item Optional \texttt{unit}: \texttt{"rad"} (default) or \texttt{"deg"}.
\item Result: \texttt{real} for \texttt{int}/\texttt{real} input, \texttt{complex} for \texttt{complex} input. For list input, the function is applied element-wise.
\item Domain: All $\mathbb{R}$ and all $\mathbb{C}$.
\item A \texttt{Type Error} is raised if the argument is not a numeric value or list of numeric values.
\end{itemize}
Examples:
\begin{verbatim}
f.trig.sin(1.570796327)              #c# Returns approximately 1.0
f.trig.sin([0, 30, 90], unit="deg")  #c# Returns [0, 0.5, 1]
f.trig.sin(1 + 2i)                  #c# Returns approximately 3.1658 - 1.9596i
\end{verbatim}

\subsubsection{Cosine Function}
Computes the cosine of an angle. Angles are interpreted as radians by default. \\
Morphology:
\begin{verbatim}
EBNF:
"f.trig.cos", "(", ( int | real | complex ), [ ",", "unit", "=", string_type ], ")"
\end{verbatim}
\begin{itemize}
\item Operands: \texttt{int}, \texttt{real}, \texttt{complex}, or list of numeric values.
\item Optional \texttt{unit}: \texttt{"rad"} (default) or \texttt{"deg"}.
\item Result: \texttt{real} for \texttt{int}/\texttt{real} input, \texttt{complex} for \texttt{complex} input. For list input, the function is applied element-wise.
\item Domain: All $\mathbb{R}$ and all $\mathbb{C}$.
\item A \texttt{Type Error} is raised if the argument is not a numeric value or list of numeric values.
\end{itemize}
Examples:
\begin{verbatim}
f.trig.cos(0)                       #c# Returns 1.0
f.trig.cos([0, 60, 90], unit="deg")  #c# Returns [1, 0.5, 0]
f.trig.cos(1 + 2i)                  #c# Returns approximately 2.0327 + 3.0519i
\end{verbatim}

\subsubsection{Tangent Function}
Computes the tangent of an angle. Angles are interpreted as radians by default. \\
Morphology:
\begin{verbatim}
EBNF:
"f.trig.tan", "(", ( int | real | complex ), [ ",", "unit", "=", string_type ], ")"
\end{verbatim}
\begin{itemize}
\item Operands: \texttt{int}, \texttt{real}, \texttt{complex}, or list of numeric values.
\item Optional \texttt{unit}: \texttt{"rad"} (default) or \texttt{"deg"}.
\item Result: \texttt{real} for \texttt{int}/\texttt{real} input, \texttt{complex} for \texttt{complex} input. For list input, the function is applied element-wise.
\item Domain for real: $\{x \in \mathbb{R} : \cos x \neq 0\}$. A \texttt{Division by Zero Error} is raised when $\cos x = 0$.
\item Domain for complex: all $\mathbb{C}$ where $\cos z \neq 0$.
\item A \texttt{Type Error} is raised if the argument is not a numeric value or list of numeric values.
\end{itemize}
Examples:
\begin{verbatim}
f.trig.tan(0)                       #c# Returns 0.0
f.trig.tan([0, 45, 90], unit="deg")  #c# Returns [0, 1, inf]
f.trig.tan(0 + 1i)                  #c# Returns approximately 0.0 + 0.7616i
\end{verbatim}

\subsubsection{Cotangent Function}
Computes the cotangent $\cot(x) = \frac{1}{\tan(x)} = \frac{\cos(x)}{\sin(x)}$. \\
Morphology:
\begin{verbatim}
EBNF:
"f.trig.cot", "(", ( int | real | complex ), [ ",", "unit", "=", string_type ], ")"
\end{verbatim}
\begin{itemize}
\item Operands: \texttt{int}, \texttt{real}, \texttt{complex}, or list of numeric values.
\item Optional \texttt{unit}: \texttt{"rad"} (default) or \texttt{"deg"}.
\item Result: \texttt{real} for \texttt{int}/\texttt{real} input, \texttt{complex} for \texttt{complex} input.
\item A \texttt{Division by Zero Error} is raised when $\sin x = 0$.
\end{itemize}
Examples:
\begin{verbatim}
f.trig.cot(0.7854)                  #c# Returns approximately 1.0  (pi/4)
f.trig.cot([45, 90], unit="deg")     #c# Returns [1.0, 0.0]
\end{verbatim}

\subsubsection{Secant Function}
Computes the secant $\sec(x) = \frac{1}{\cos(x)}$. \\
Morphology:
\begin{verbatim}
EBNF:
"f.trig.sec", "(", ( int | real | complex ), [ ",", "unit", "=", string_type ], ")"
\end{verbatim}
\begin{itemize}
\item Operands: \texttt{int}, \texttt{real}, \texttt{complex}, or list of numeric values.
\item Optional \texttt{unit}: \texttt{"rad"} (default) or \texttt{"deg"}.
\item Result: \texttt{real} for \texttt{int}/\texttt{real} input, \texttt{complex} for \texttt{complex} input.
\item A \texttt{Division by Zero Error} is raised when $\cos x = 0$.
\end{itemize}
Examples:
\begin{verbatim}
f.trig.sec(0)                       #c# Returns 1.0
f.trig.sec([0, 60], unit="deg")      #c# Returns [1.0, 2.0]
\end{verbatim}

\subsubsection{Cosecant Function}
Computes the cosecant $\csc(x) = \frac{1}{\sin(x)}$. \\
Morphology:
\begin{verbatim}
EBNF:
"f.trig.csc", "(", ( int | real | complex ), [ ",", "unit", "=", string_type ], ")"
\end{verbatim}
\begin{itemize}
\item Operands: \texttt{int}, \texttt{real}, \texttt{complex}, or list of numeric values.
\item Optional \texttt{unit}: \texttt{"rad"} (default) or \texttt{"deg"}.
\item Result: \texttt{real} for \texttt{int}/\texttt{real} input, \texttt{complex} for \texttt{complex} input.
\item A \texttt{Division by Zero Error} is raised when $\sin x = 0$.
\end{itemize}
Examples:
\begin{verbatim}
f.trig.csc(1.5708)                  #c# Returns approximately 1.0  (pi/2)
f.trig.csc([30, 90], unit="deg")     #c# Returns [2.0, 1.0]
\end{verbatim}

\subsubsection{Complex Trigonometric Formulas}
For complex $z = a + bi$, the trigonometric functions are computed via analytic continuation:
\[
\sin(a + bi) = \sin(a)\cosh(b) + i\cos(a)\sinh(b)
\]
\[
\cos(a + bi) = \cos(a)\cosh(b) - i\sin(a)\sinh(b)
\]
\[
\tan(a + bi) = \frac{\sin(a + bi)}{\cos(a + bi)}
\]

\subsubsection{Arcsine Function}
Computes the inverse sine. Angles are interpreted as radians by default. \\
Morphology:
\begin{verbatim}
EBNF:
("f.trig.asin"), "(", ( int | real | complex ), 
  [ ",", "unit", "=", string_type ], ")"
\end{verbatim}
\begin{itemize}
\item Operands: \texttt{int}, \texttt{real}, or list of numeric values.
\item Optional \texttt{unit}: \texttt{"rad"} (default) or \texttt{"deg"}.
\item Result: \texttt{real}.
\item Domain: $[-1, 1]$. A \texttt{Value Error} is raised for $|x| > 1$.
\item A \texttt{Type Error} is raised if the argument is not a numeric value or list of numeric values.
\end{itemize}
Examples:
\begin{verbatim}
f.trig.asin(0)                        #c# Returns 0.0
f.trig.asin(1)                        #c# Returns approximately 1.5708 (pi/2)
f.trig.asin([0, 0.5, 1], unit="deg")  #c# Returns [0, 30, 90]
\end{verbatim}

\subsubsection{Arccosine Function}
Computes the inverse cosine. Angles are interpreted as radians by default. \\
Morphology:
\begin{verbatim}
EBNF:
("f.trig.acos"), "(", ( int | real | complex ), 
  [ ",", "unit", "=", string_type ], ")"
\end{verbatim}
\begin{itemize}
\item Operands: \texttt{int}, \texttt{real}, or list of numeric values.
\item Optional \texttt{unit}: \texttt{"rad"} (default) or \texttt{"deg"}.
\item Result: \texttt{real}.
\item Domain: $[-1, 1]$. A \texttt{Value Error} is raised for $|x| > 1$.
\item A \texttt{Type Error} is raised if the argument is not a numeric value or list of numeric values.
\end{itemize}
Examples:
\begin{verbatim}
f.trig.acos(1)                        #c# Returns 0.0
f.trig.acos(0)                        #c# Returns approximately 1.5708 (pi/2)
f.trig.acos([0, 0.5, 1], unit="deg")  #c# Returns [90, 60, 0]
\end{verbatim}

\subsubsection{Arctangent Function}
Computes the inverse tangent. Angles are interpreted as radians by default. \\
Morphology:
\begin{verbatim}
EBNF:
("f.trig.atan"), "(", ( int | real | complex ), 
  [ ",", "unit", "=", string_type ], ")"
\end{verbatim}
\begin{itemize}
\item Operands: \texttt{int}, \texttt{real}, or list of numeric values.
\item Optional \texttt{unit}: \texttt{"rad"} (default) or \texttt{"deg"}.
\item Result: \texttt{real}.
\item Domain: all $\mathbb{R}$. No domain restrictions.
\item A \texttt{Type Error} is raised if the argument is not a numeric value or list of numeric values.
\end{itemize}
Examples:
\begin{verbatim}
f.trig.atan(0)                        #c# Returns 0.0
f.trig.atan(1)                        #c# Returns approximately 0.7854 (pi/4)
f.trig.atan([0, 1, inf], unit="deg")  #c# Returns [0, 45, 90]
\end{verbatim}

\subsubsection{Arccotangent Function}
Computes the inverse cotangent $\text{acot}(x) = \frac{\pi}{2} - \text{atan}(x)$. \\
Morphology:
\begin{verbatim}
EBNF:
"f.trig.acot", "(", ( int | real ), [ ",", "unit", "=", string_type ], ")"
\end{verbatim}
\begin{itemize}
\item Operands: \texttt{int} or \texttt{real}.
\item Optional \texttt{unit}: \texttt{"rad"} (default) or \texttt{"deg"}.
\item Result: \texttt{real}. Domain: all $\mathbb{R}$.
\end{itemize}
Examples:
\begin{verbatim}
f.trig.acot(1)                        #c# Returns approximately 0.7854 (pi/4)
f.trig.acot(inf)                      #c# Returns 0.0
\end{verbatim}

\subsubsection{Arcsecant Function}
Computes the inverse secant $\text{asec}(x) = \text{acos}(1/x)$. \\
Morphology:
\begin{verbatim}
EBNF:
"f.trig.asec", "(", ( int | real ), [ ",", "unit", "=", string_type ], ")"
\end{verbatim}
\begin{itemize}
\item Operands: \texttt{int} or \texttt{real}.
\item Optional \texttt{unit}: \texttt{"rad"} (default) or \texttt{"deg"}.
\item Result: \texttt{real}.
\item A \texttt{Value Error} is raised for $|x| < 1$ (out of domain).
\end{itemize}
Examples:
\begin{verbatim}
f.trig.asec(1)                        #c# Returns 0.0
f.trig.asec(2)                        #c# Returns approximately 1.0472 (pi/3)
\end{verbatim}

\subsubsection{Arccosecant Function}
Computes the inverse cosecant $\text{acsc}(x) = \text{asin}(1/x)$. \\
Morphology:
\begin{verbatim}
EBNF:
"f.trig.acsc", "(", ( int | real ), [ ",", "unit", "=", string_type ], ")"
\end{verbatim}
\begin{itemize}
\item Operands: \texttt{int} or \texttt{real}.
\item Optional \texttt{unit}: \texttt{"rad"} (default) or \texttt{"deg"}.
\item Result: \texttt{real}.
\item A \texttt{Value Error} is raised for $|x| < 1$ (out of domain).
\end{itemize}
Examples:
\begin{verbatim}
f.trig.acsc(1)                        #c# Returns approximately 1.5708 (pi/2)
f.trig.acsc(2)                        #c# Returns approximately 0.5236 (pi/6)
\end{verbatim}

\subsubsection{Hyperbolic Sine Function}
Computes the hyperbolic sine of the argument. \\
Morphology:
\begin{verbatim}
EBNF:
"f.hyp.sinh", "(", ( int | real | complex ), ")"
\end{verbatim}
\begin{itemize}
\item Operands: \texttt{int}, \texttt{real}, \texttt{complex}, or list of numeric values.
\item Result: \texttt{real} for \texttt{int}/\texttt{real} input, \texttt{complex} for \texttt{complex} input.
\item Domain: all $\mathbb{R}$ and all $\mathbb{C}$. No restrictions.
\item A \texttt{Type Error} is raised if the argument is not a numeric value or list of numeric values.
\end{itemize}
Examples:
\begin{verbatim}
f.hyp.sinh(0)          #c# Returns 0.0
f.hyp.sinh(1.443635475)  #c# Returns approximately 2.0
f.hyp.sinh(1 + 2i)     #c# Returns approximately -0.4891 + 0.9889i
\end{verbatim}

\subsubsection{Hyperbolic Cosine Function}
Computes the hyperbolic cosine of the argument. \\
Morphology:
\begin{verbatim}
EBNF:
"f.hyp.cosh", "(", ( int | real | complex ), ")"
\end{verbatim}
\begin{itemize}
\item Operands: \texttt{int}, \texttt{real}, \texttt{complex}, or list of numeric values.
\item Result: \texttt{real} for \texttt{int}/\texttt{real} input, \texttt{complex} for \texttt{complex} input.
\item Domain: all $\mathbb{R}$ and all $\mathbb{C}$. No restrictions.
\item A \texttt{Type Error} is raised if the argument is not a numeric value or list of numeric values.
\end{itemize}
Examples:
\begin{verbatim}
f.hyp.cosh(0)          #c# Returns 1.0
f.hyp.cosh(1)          #c# Returns approximately 1.5431
f.hyp.cosh(1 + 2i)     #c# Returns approximately -0.4891 + 0.9889i
\end{verbatim}

\subsubsection{Hyperbolic Tangent Function}
Computes the hyperbolic tangent of the argument. \\
Morphology:
\begin{verbatim}
EBNF:
"f.hyp.tanh", "(", ( int | real | complex ), ")"
\end{verbatim}
\begin{itemize}
\item Operands: \texttt{int}, \texttt{real}, \texttt{complex}, or list of numeric values.
\item Result: \texttt{real} for \texttt{int}/\texttt{real} input, \texttt{complex} for \texttt{complex} input.
\item Domain for real: all $\mathbb{R}$. For complex: all $\mathbb{C}$ where $\cosh z \neq 0$.
\item A \texttt{Type Error} is raised if the argument is not a numeric value or list of numeric values.
\end{itemize}
Examples:
\begin{verbatim}
f.hyp.tanh(0)          #c# Returns 0.0
f.hyp.tanh(1)          #c# Returns approximately 0.7616
f.hyp.tanh(0 + 1i)     #c# Returns approximately 0.0 + 0.5524i
\end{verbatim}

\subsubsection{Hyperbolic Cotangent Function}
Computes the hyperbolic cotangent $\coth(x) = \frac{\cosh(x)}{\sinh(x)}$. \\
Morphology:
\begin{verbatim}
EBNF:
"f.hyp.coth", "(", ( int | real | complex ), ")"
\end{verbatim}
\begin{itemize}
\item Operands: \texttt{int}, \texttt{real}, or \texttt{complex}.
\item Result: same type as input.
\item A \texttt{Division by Zero Error} is raised at $x = 0$.
\end{itemize}
Examples:
\begin{verbatim}
f.hyp.coth(1)          #c# Returns approximately 1.3130
\end{verbatim}

\subsubsection{Hyperbolic Secant Function}
Computes the hyperbolic secant $\text{sech}(x) = \frac{1}{\cosh(x)}$. \\
Morphology:
\begin{verbatim}
EBNF:
"f.hyp.sech", "(", ( int | real | complex ), ")"
\end{verbatim}
\begin{itemize}
\item Operands: \texttt{int}, \texttt{real}, or \texttt{complex}.
\item Result: same type as input.
\item No domain restrictions.
\end{itemize}
Examples:
\begin{verbatim}
f.hyp.sech(0)          #c# Returns 1.0
f.hyp.sech(1)          #c# Returns approximately 0.6481
\end{verbatim}

\subsubsection{Hyperbolic Cosecant Function}
Computes the hyperbolic cosecant $\text{csch}(x) = \frac{1}{\sinh(x)}$. \\
Morphology:
\begin{verbatim}
EBNF:
"f.hyp.csch", "(", ( int | real | complex ), ")"
\end{verbatim}
\begin{itemize}
\item Operands: \texttt{int}, \texttt{real}, or \texttt{complex}.
\item Result: same type as input.
\item A \texttt{Division by Zero Error} is raised at $x = 0$.
\end{itemize}
Examples:
\begin{verbatim}
f.hyp.csch(1)          #c# Returns approximately 0.8509
\end{verbatim}

\subsubsection{Complex Hyperbolic Formulas}
For complex $z = a + bi$, the hyperbolic functions are computed as:
\[
\sinh(a + bi) = \sinh(a)\cos(b) + i\cosh(a)\sin(b)
\]
\[
\cosh(a + bi) = \cosh(a)\cos(b) + i\sinh(a)\sin(b)
\]
\[
\tanh(a + bi) = \frac{\sinh(a + bi)}{\cosh(a + bi)}
\]
Equivalently: $\sinh(z) = -i \sin(iz)$, $\cosh(z) = \cos(iz)$, $\tanh(z) = \sinh(z) / \cosh(z)$.

\subsubsection{Hyperbolic Arcsine Function}
Computes the inverse hyperbolic sine. \\
Morphology:
\begin{verbatim}
EBNF:
("f.hyp.asinh"), "(", ( int | real | complex ), ")"
\end{verbatim}
\begin{itemize}
\item Operands: \texttt{int}, \texttt{real}, \texttt{complex}, or list of numeric values.
\item Result: \texttt{real} for \texttt{int}/\texttt{real} input, \texttt{complex} for \texttt{complex} input.
\item Domain: all $\mathbb{R}$ and all $\mathbb{C}$. No restrictions.
\item A \texttt{Type Error} is raised if the argument is not a numeric value or list of numeric values.
\end{itemize}
Examples:
\begin{verbatim}
f.hyp.asinh(0)          #c# Returns 0.0
f.hyp.asinh(1)          #c# Returns approximately 0.8814
f.hyp.asinh(1 + 2i)     #c# Returns approximately 1.4694 + 1.0987i
\end{verbatim}

\subsubsection{Hyperbolic Arccosine Function}
Computes the inverse hyperbolic cosine. \\
Morphology:
\begin{verbatim}
EBNF:
("f.hyp.acosh"), "(", ( int | real | complex ), ")"
\end{verbatim}
\begin{itemize}
\item Operands: \texttt{int}, \texttt{real}, \texttt{complex}, or list of numeric values.
\item Result: \texttt{real} for \texttt{int}/\texttt{real} input, \texttt{complex} for \texttt{complex} input.
\item Domain for real: $[1, \infty)$. A \texttt{Value Error} is raised for $x < 1$.
\item Domain for complex: all $\mathbb{C}$ (using principal branch).
\item A \texttt{Type Error} is raised if the argument is not a numeric value or list of numeric values.
\end{itemize}
Examples:
\begin{verbatim}
f.hyp.acosh(1)          #c# Returns 0.0
f.hyp.acosh(2)          #c# Returns approximately 1.3170
f.hyp.acosh(1 + 2i)     #c# Returns approximately 1.5286 + 1.1437i
\end{verbatim}

\subsubsection{Hyperbolic Arctangent Function}
Computes the inverse hyperbolic tangent. \\
Morphology:
\begin{verbatim}
EBNF:
("f.hyp.atanh"), "(", ( int | real | complex ), ")"
\end{verbatim}
\begin{itemize}
\item Operands: \texttt{int}, \texttt{real}, \texttt{complex}, or list of numeric values.
\item Result: \texttt{real} for \texttt{int}/\texttt{real} input, \texttt{complex} for \texttt{complex} input.
\item Domain for real: $(-1, 1)$. A \texttt{Value Error} is raised for $|x| \geq 1$.
\item Domain for complex: all $\mathbb{C}$ where $z \neq \pm 1$.
\item A \texttt{Type Error} is raised if the argument is not a numeric value or list of numeric values.
\end{itemize}
Examples:
\begin{verbatim}
f.hyp.atanh(0)          #c# Returns 0.0
f.hyp.atanh(0.5)        #c# Returns approximately 0.5493
f.hyp.atanh(0.5 + 0i)   #c# Returns approximately 0.5493 + 0.0i
\end{verbatim}

\subsubsection{Inverse Hyperbolic Cotangent Function}
Computes $\text{acoth}(x) = \frac{1}{2} \ln\!\left(\frac{x+1}{x-1}\right)$. \\
Morphology:
\begin{verbatim}
EBNF:
"f.hyp.acoth", "(", ( int | real | complex ), ")"
\end{verbatim}
\begin{itemize}
\item Operands: \texttt{int}, \texttt{real}, or \texttt{complex}.
\item Result: same type as input.
\item A \texttt{Value Error} is raised for real $|x| \leq 1$ (out of domain).
\end{itemize}
Examples:
\begin{verbatim}
f.hyp.acoth(2)         #c# Returns approximately 0.5493
\end{verbatim}

\subsubsection{Inverse Hyperbolic Secant Function}
Computes $\text{asech}(x) = \text{acosh}(1/x)$. \\
Morphology:
\begin{verbatim}
EBNF:
"f.hyp.asech", "(", ( int | real | complex ), ")"
\end{verbatim}
\begin{itemize}
\item Operands: \texttt{int}, \texttt{real}, or \texttt{complex}.
\item Result: same type as input.
\item A \texttt{Value Error} is raised for real $x \leq 0$ or $x > 1$ (out of domain).
\end{itemize}
Examples:
\begin{verbatim}
f.hyp.asech(1)         #c# Returns 0.0
\end{verbatim}

\subsubsection{Inverse Hyperbolic Cosecant Function}
Computes $\text{acsch}(x) = \text{asch}(1/x)$. \\
Morphology:
\begin{verbatim}
EBNF:
"f.hyp.acsch", "(", ( int | real | complex ), ")"
\end{verbatim}
\begin{itemize}
\item Operands: \texttt{int}, \texttt{real}, or \texttt{complex}.
\item Result: same type as input.
\item A \texttt{Value Error} is raised at $x = 0$ (singularity).
\end{itemize}
Examples:
\begin{verbatim}
f.hyp.acsch(1)         #c# Returns approximately 0.8814
\end{verbatim}

\subsubsection{Complex Hyperbolic Inverse Formulas}
For complex arguments, the inverse hyperbolic functions are computed via the complex logarithm:
\[
\operatorname{asinh}(z) = \ln\!\left(z + \sqrt{z^2 + 1}\right)
\]
\[
\operatorname{acosh}(z) = \ln\!\left(z + \sqrt{z^2 - 1}\right)
\]
\[
\operatorname{atanh}(z) = \frac{1}{2} \ln\!\left(\frac{1 + z}{1 - z}\right)
\]

\subsubsection{Exponential Function}
Returns $e$ raised to the power of the argument. \\
Morphology:
\begin{verbatim}
EBNF:
"f.log.exp", "(", ( int | real | complex ), ")"
\end{verbatim}
\begin{itemize}
\item Operands: \texttt{int}, \texttt{real}, \texttt{complex}, or list of numeric values.
\item Result: \texttt{real} for \texttt{int}/\texttt{real} input, \texttt{complex} for \texttt{complex} input.
\item Domain: all $\mathbb{R}$ and all $\mathbb{C}$.
\item A \texttt{Type Error} is raised if the argument is not a numeric value or list of numeric values.
\end{itemize}
For complex $z = a + bi$, computed via Euler's Formula: $e^{a+bi} = e^a(\cos b + i \sin b)$.
Examples:
\begin{verbatim}
f.log.exp(0)              #c# Returns 1.0
f.log.exp(1)              #c# Returns approximately 2.71828
f.log.exp(0 + 3.14159i)   #c# Returns approximately -1 + 0.0i (Euler's Identity)
f.log.exp(1 + 1i)         #c# Returns approximately 1.4687 + 2.2874i
\end{verbatim}

\subsubsection{Square Root Function}
Returns the principal square root of the argument. \\
Morphology:
\begin{verbatim}
EBNF:
"f.log.sqrt", "(", ( int | real | complex ), ")"
\end{verbatim}
\begin{itemize}
\item Operands: \texttt{int}, \texttt{real}, \texttt{complex}, or list of numeric values.
\item Result: \texttt{real} for non-negative \texttt{int}/\texttt{real} input, \texttt{complex} for negative real or \texttt{complex} input.
\item Domain for real: $[0, \infty)$ returns \texttt{real}. Negative real input returns \texttt{complex} (principal branch).
\item Domain for complex: all $\mathbb{C}$, principal branch $-\pi < \arg(\sqrt{z}) \leq \pi$.
\item A \texttt{Type Error} is raised if the argument is not a numeric value or list of numeric values.
\end{itemize}
Examples:
\begin{verbatim}
f.log.sqrt(9)          #c# Returns 3.0
f.log.sqrt(2)          #c# Returns approximately 1.41421
f.log.sqrt(-1)         #c# Returns 0 + 1i
f.log.sqrt(-4)         #c# Returns 0 + 2i
f.log.sqrt(3 + 4i)     #c# Returns approximately 2 + 1i
f.log.sqrt(0)          #c# Returns 0.0
\end{verbatim}

\subsubsection{Power Function}
Computes exponentiation and optional modular exponentiation. \\
Morphology:
\begin{verbatim}
EBNF:
"f.log.pow", "(", ( int | real | complex ), ",", ( int | real | complex ), [ ",", int ], ")"
\end{verbatim}
\begin{itemize}
\item \texttt{f.log.pow(a, b)}: returns $a^b$. Both arguments \texttt{int}, \texttt{real}, or \texttt{complex}.
\item \texttt{f.log.pow(a, b, c)}: returns $a^b \bmod c$ (modular exponentiation). \texttt{a}, \texttt{c} must be \texttt{int}, \texttt{c} must be $> 0$, \texttt{b} can be negative (computes modular inverse).
\item Return type: \texttt{int} when both arguments are \texttt{int} and the result is exact; \texttt{real} otherwise.
\item Without \texttt{c}: negative exponents on \texttt{int} produce \texttt{real}.
\item With \texttt{c}: negative \texttt{b} computes $a^b \bmod c$ via modular multiplicative inverse (exists when $\gcd(a, c) = 1$).
\item For complex $z_1, z_2$: $z_1^{z_2} = e^{z_2 \ln z_1}$ (principal branch). $0^z$ only defined for $\operatorname{Re}(z) > 0$.
\item A \texttt{Type Error} is raised if arguments are not numeric.
\end{itemize}
Examples:
\begin{verbatim}
f.log.pow(2, 3)         #c# Returns 8 (int)
f.log.pow(9, 0.5)       #c# Returns 3.0 (real)
f.log.pow(3, 4, 7)      #c# Returns 4  (81 mod 7)
f.log.pow(2, 32, 100)   #c# Returns 96  (4294967296 mod 100)
f.log.pow(3, -1, 7)     #c# Returns 5  (modular inverse: 3*5=15==1 mod 7)
f.log.pow(2, 1 + 1i)    #c# Returns approximately 1.538 + 0.999i
f.log.pow(1 + 1i, 2)    #c# Returns 0 + 2i
\end{verbatim}

\subsubsection{Natural Logarithm}
Returns the natural logarithm (base $e$) of the argument. \\
Morphology:
\begin{verbatim}
EBNF:
"f.log.ln", "(", ( int | real | complex ), ")"
\end{verbatim}
\begin{itemize}
\item Operands: \texttt{int} (must be $>$ 0), \texttt{real} (must be $>$ 0), \texttt{complex} ($\neq 0$).
\item Result: \texttt{real} for positive \texttt{int}/\texttt{real} input, \texttt{complex} for \texttt{complex} input.
\item A \texttt{Value Error} is raised for non-positive real arguments.
\item A \texttt{Type Error} is raised if the argument is not a numeric value or list of numeric values.
\end{itemize}
Examples:
\begin{verbatim}
f.log.ln(1)          #c# Returns 0.0
f.log.ln(c.math.e)        #c# Returns approximately 1.0
f.log.ln(100)        #c# Returns approximately 4.6052
f.log.ln(1 + 1i)     #c# Returns approximately 0.3466 + 0.7854i
\end{verbatim}

\subsubsection{Base-10 Logarithm}
Returns the base-10 logarithm of the argument. \\
Morphology:
\begin{verbatim}
EBNF:
"f.log.log10", "(", ( int | real | complex ), ")"
\end{verbatim}
\begin{itemize}
\item Operands: \texttt{int} (must be $>$ 0), \texttt{real} (must be $>$ 0), \texttt{complex} ($\neq 0$).
\item Result: \texttt{real} for positive \texttt{int}/\texttt{real} input, \texttt{complex} for \texttt{complex} input.
\item A \texttt{Value Error} is raised for non-positive real arguments.
\item A \texttt{Type Error} is raised if the argument is not a numeric value or list of numeric values.
\end{itemize}
Examples:
\begin{verbatim}
f.log10(1)          #c# Returns 0.0
f.log10(10)         #c# Returns 1.0
f.log10(1000)       #c# Returns 3.0
f.log10(1 + 1i)     #c# Returns approximately 0.1505 + 0.3411i
\end{verbatim}

\subsubsection{Expm1 Function}
Computes $e^x - 1$ with high precision for small $x$, avoiding catastrophic cancellation. \\
Morphology:
\begin{verbatim}
EBNF:
"f.math.expm1", "(", ( int | real ), ")"
\end{verbatim}
\begin{itemize}
\item Operand: \texttt{int} or \texttt{real}.
\item Result: \texttt{real}.
\end{itemize}
Examples:
\begin{verbatim}
f.math.expm1(0)       #c# Returns 0.0
f.math.expm1(1)       #c# Returns approximately 1.7183 (e - 1)
f.math.expm1(0.001)   #c# Returns approximately 0.0010005  (more precise than exp(0.001) - 1)
\end{verbatim}

\subsubsection{Log1p Function}
Computes $\ln(1 + x)$ with high precision for small $x$, avoiding catastrophic cancellation. \\
Morphology:
\begin{verbatim}
EBNF:
"f.math.log1p", "(", ( int | real ), ")"
\end{verbatim}
\begin{itemize}
\item Operand: \texttt{int} or \texttt{real}.
\item Result: \texttt{real}.
\item A \texttt{Value Error} is raised for $x \leq -1$.
\end{itemize}
Examples:
\begin{verbatim}
f.math.log1p(0)       #c# Returns 0.0
f.math.log1p(1)       #c# Returns approximately 0.6931 (ln 2)
f.math.log1p(0.001)   #c# Returns approximately 0.0009995  (more precise than ln(1.001))
\end{verbatim}

\label{sec:complex_functions}

\subsubsection{Real Part Function}
Returns the real part of a complex number. \\
Morphology:
\begin{verbatim}
EBNF:
"f.cplx.real", "(", ( int | real | complex ), ")"
\end{verbatim}
\begin{itemize}
\item Operands: \texttt{int} (promoted to \texttt{real}), \texttt{real}, \texttt{complex}.
\item Result: \texttt{real}.
\item For \texttt{int}/\texttt{real} input, the value is returned as \texttt{real}.
\item A \texttt{Type Error} is raised for non-numeric arguments.
\end{itemize}
Examples:
\begin{verbatim}
f.cplx.real(3 + 4i)   #c# Returns 3.0
f.cplx.real(-2i)      #c# Returns 0.0
f.cplx.real(5)        #c# Returns 5.0 (int promoted)
\end{verbatim}

\subsubsection{Imaginary Part Function}
Returns the imaginary part of a complex number. \\
Morphology:
\begin{verbatim}
EBNF:
"f.cplx.imag", "(", ( int | real | complex ), ")"
\end{verbatim}
\begin{itemize}
\item Operands: \texttt{int}, \texttt{real} (imaginary part is 0.0), \texttt{complex}.
\item Result: \texttt{real}.
\item A \texttt{Type Error} is raised for non-numeric arguments.
\end{itemize}
Examples:
\begin{verbatim}
f.cplx.imag(3 + 4i)   #c# Returns 4.0
f.cplx.imag(-2i)      #c# Returns -2.0
f.cplx.imag(5)        #c# Returns 0.0 (real has no imaginary part)
\end{verbatim}

\subsubsection{Angle Function}
Returns the angle (argument) of a complex number in radians, in the range $(-\pi, \pi]$. \\
Morphology:
\begin{verbatim}
EBNF:
"f.cplx.angle", "(", ( int | real | complex ), ")"
\end{verbatim}
\begin{itemize}
\item Operands: \texttt{int}, \texttt{real}, \texttt{complex}.
\item Result: \texttt{real}.
\item \texttt{angle(0)} returns \texttt{0.0}.
\item A \texttt{Type Error} is raised for non-numeric arguments.
\end{itemize}
Examples:
\begin{verbatim}
f.cplx.angle(1 + 0i)      #c# Returns 0.0
f.cplx.angle(0 + 1i)      #c# Returns 1.5708 (pi/2)
f.cplx.angle(-1 + 0i)     #c# Returns 3.14159 (pi)
f.cplx.angle(1 + 1i)      #c# Returns 0.7854 (pi/4)
\end{verbatim}

\subsubsection{Complex Conjugate Function}
Returns the complex conjugate: $\overline{a + bi} = a - bi$. \\
Morphology:
\begin{verbatim}
EBNF:
"f.cplx.conj", "(", ( int | real | complex ), ")"
\end{verbatim}
\begin{itemize}
\item Operands: \texttt{int}, \texttt{real} (self-conjugate), \texttt{complex}.
\item Result: same type as input. Real numbers are self-conjugate.
\item A \texttt{Type Error} is raised for non-numeric arguments.
\end{itemize}
Examples:
\begin{verbatim}
f.cplx.conj(3 + 4i)   #c# Returns 3 - 4i
f.cplx.conj(-2i)      #c# Returns 2i
f.cplx.conj(5)        #c# Returns 5 (real numbers are self-conjugate)
\end{verbatim}

\subsubsection{Polar Conversion Function}
Converts a complex number to polar form, returning \texttt{[magnitude, angle]}. \\
Morphology:
\begin{verbatim}
EBNF:
"f.cplx.polar", "(", ( int | real | complex ), ")"
\end{verbatim}
\begin{itemize}
\item Operands: \texttt{int}, \texttt{real}, \texttt{complex}.
\item Result: \texttt{list} of two \texttt{real} values: \texttt{[|z|, arg(z)]}.
\item The angle is in radians, in the range $(-\pi, \pi]$.
\item A \texttt{Type Error} is raised for non-numeric arguments.
\end{itemize}
Examples:
\begin{verbatim}
f.cplx.polar(3 + 4i)     #c# Returns [5.0, 0.9273]
f.cplx.polar(1 + 0i)     #c# Returns [1.0, 0.0]
f.cplx.polar(0 + 1i)     #c# Returns [1.0, 1.5708]
\end{verbatim}

\subsubsection{Rectangular Conversion Function}
Creates a complex number from polar form (magnitude and angle in radians). \\
Morphology:
\begin{verbatim}
EBNF:
"f.cplx.rect", "(", real, ",", real, ")"
\end{verbatim}
\begin{itemize}
\item Operands: magnitude (\texttt{int} or \texttt{real}), angle in radians (\texttt{int} or \texttt{real}).
\item Result: \texttt{complex}.
\item Computes $z = r(\cos\theta + i\sin\theta)$.
\item A \texttt{Type Error} is raised if arguments are not numeric.
\end{itemize}
Examples:
\begin{verbatim}
f.cplx.rect(5, 0.9273)      #c# Returns approximately 3 + 4i
f.cplx.rect(1, 0)           #c# Returns 1 + 0i
f.cplx.rect(1, 1.5708)      #c# Returns approximately 0 + 1i
\end{verbatim}

\subsubsection{Cis Function}
Computes the cis function: $\text{cis}(\theta) = \cos(\theta) + i \cdot \sin(\theta) = e^{i\theta}$. Returns a complex number on the unit circle at angle $\theta$ (radians). \\
Morphology:
\begin{verbatim}
EBNF:
"f.cplx.cis", "(", ( int | real ), ")"
\end{verbatim}
\begin{itemize}
\item Operands: angle in radians (\texttt{int} or \texttt{real}).
\item Result: \texttt{complex}.
\item Computes $\text{cis}(\theta) = \cos(\theta) + i \cdot \sin(\theta)$.
\item A \texttt{Type Error} is raised if the argument is not \texttt{int} or \texttt{real}.
\end{itemize}
Examples:
\begin{verbatim}
f.cplx.cis(0)           #c# Returns 1 + 0i
f.cplx.cis(1.5708)      #c# Returns approximately 0 + 1i  (pi/2)
f.cplx.cis(3.14159)     #c# Returns approximately -1 + 0i (pi)
f.cplx.cis(-1.5708)     #c# Returns approximately 0 - 1i  (-pi/2)
\end{verbatim}

\subsubsection{Radians to Degrees Function}
Converts an angle from radians to degrees. \\
Morphology:
\begin{verbatim}
EBNF:
"f.trig.rad2deg", "(", ( int | real ), ")"
\end{verbatim}
\begin{itemize}
\item Operands: \texttt{int} or \texttt{real}.
\item Result: \texttt{real}.
\item Computes $\text{degrees} = \text{radians} \times \frac{180}{\pi}$.
\item A \texttt{Type Error} is raised if the argument is not \texttt{int} or \texttt{real}.
\end{itemize}
Examples:
\begin{verbatim}
f.trig.rad2deg(3.14159)   #c# Returns 180.0
f.trig.rad2deg(1.5708)    #c# Returns 90.0
f.trig.rad2deg(0)         #c# Returns 0.0
\end{verbatim}

\subsubsection{Degrees to Radians Function}
Converts an angle from degrees to radians. \\
Morphology:
\begin{verbatim}
EBNF:
"f.trig.deg2rad", "(", ( int | real ), ")"
\end{verbatim}
\begin{itemize}
\item Operands: \texttt{int} or \texttt{real}.
\item Result: \texttt{real}.
\item Computes $\text{radians} = \text{degrees} \times \frac{\pi}{180}$.
\item A \texttt{Type Error} is raised if the argument is not \texttt{int} or \texttt{real}.
\end{itemize}
Examples:
\begin{verbatim}
f.trig.deg2rad(180)   #c# Returns 3.14159
f.trig.deg2rad(90)    #c# Returns 1.5708
f.trig.deg2rad(0)     #c# Returns 0.0
\end{verbatim}

\subsubsection{Floor Function}
Rounds the argument down to the nearest integer. \\
Morphology:
\begin{verbatim}
EBNF:
"f.math.floor", "(", ( int | real ), ")"
\end{verbatim}
\begin{itemize}
\item Operands: \texttt{int} (returned as-is), \texttt{real}.
\item Result: \texttt{int}.
\item Returns the largest integer $\leq$ the argument.
\item A \texttt{Type Error} is raised if the argument is not \texttt{int} or \texttt{real}.
\end{itemize}
Examples:
\begin{verbatim}
f.math.floor(2.7)   #c# Returns 2
f.math.floor(-1.3)  #c# Returns -2
f.math.floor(5)     #c# Returns 5
\end{verbatim}

\subsubsection{Ceiling Function}
Rounds the argument up to the nearest integer. \\
Morphology:
\begin{verbatim}
EBNF:
"f.math.ceil", "(", ( int | real ), ")"
\end{verbatim}
\begin{itemize}
\item Operands: \texttt{int} (returned as-is), \texttt{real}.
\item Result: \texttt{int}.
\item Returns the smallest integer $\geq$ the argument.
\item A \texttt{Type Error} is raised if the argument is not \texttt{int} or \texttt{real}.
\end{itemize}
Examples:
\begin{verbatim}
f.math.ceil(2.3)    #c# Returns 3
f.math.ceil(-2.7)   #c# Returns -2
f.math.ceil(5)      #c# Returns 5
\end{verbatim}

\subsubsection{Round Function}
Rounds the argument to the nearest value. Without \texttt{digits}, rounds to the nearest integer. With \texttt{digits}, rounds to the specified number of decimal places. Negative \texttt{digits} rounds to the nearest $10^n$: \texttt{digits=-1} rounds to nearest 10, \texttt{digits=-2} to nearest 100, etc. \\
Morphology:
\begin{verbatim}
EBNF:
"f.math.round", "(", ( int | real ), [",", "digits", "=", int], ")"
\end{verbatim}
\begin{itemize}
\item Operands: value (\texttt{int} or \texttt{real}), optional \texttt{digits} (\texttt{int}, default 0).
\item Result: \texttt{int} when no \texttt{digits} specified and input is \texttt{int}; \texttt{int} when \texttt{digits} $\geq 0$; \texttt{real} when \texttt{digits} $< 0$.
\item Uses banker's rounding (round half to even).
\item A \texttt{Type Error} is raised if the first argument is not \texttt{int} or \texttt{real}.
\item A \texttt{Type Error} is raised if \texttt{digits} is not \texttt{int}.
\end{itemize}
Examples:
\begin{verbatim}
f.math.round(2.5)                #c# Returns 2 (banker's rounding)
f.math.round(2.4)                #c# Returns 2
f.math.round(3.14159, digits=2)  #c# Returns 3.14
f.math.round(3.14159, digits=0)  #c# Returns 3
f.math.round(3.14159, digits=-1) #c# Returns 0
f.math.round(155, digits=-2)     #c# Returns 200
\end{verbatim}

\subsubsection{E-Series Rounding Function}
Rounds a value to the nearest standard E-series value. The E-series are logarithmically spaced value series used for resistors, capacitors, and inductors. \\
Morphology:
\begin{verbatim}
EBNF:
"f.math.eseries", "(", ( int | real ), ",", string, [",", "rm", "=", string], ")"
\end{verbatim}
\begin{itemize}
\item Operands: value (\texttt{int} or \texttt{real}), series name (\texttt{string}: \texttt{"E3"}, \texttt{"E6"}, \texttt{"E12"}, \texttt{"E24"}, \texttt{"E48"}, \texttt{"E96"}, or \texttt{"E192"}, case-insensitive), optional \texttt{rm} (\texttt{string}: \texttt{"nearest"} (default), \texttt{"floor"}, or \texttt{"ceil"}).
\item Result: \texttt{real} in the selected series, scaled to the correct decade.
\item Zero returns \texttt{0} directly ($0\Omega$ resistors are valid components).
\item A \texttt{Type Error} is raised if the first argument is not \texttt{int} or \texttt{real}.
\item A \texttt{Type Error} is raised if the second argument is not a \texttt{string}.
\item A \texttt{Value Error} is raised if the value is negative.
\item A \texttt{Value Error} is raised if the series name is not one of the valid E-series.
\item A \texttt{Value Error} is raised if the rounding mode is not one of \texttt{"nearest"}, \texttt{"floor"}, or \texttt{"ceil"}.
\end{itemize}
Examples:
\begin{verbatim}
f.math.eseries(330, "E24")              #c# Returns 330
f.math.eseries(350, "E24")              #c# Returns 330
f.math.eseries(350, "E12")              #c# Returns 330
f.math.eseries(350, "E6")               #c# Returns 330
f.math.eseries(390, "E24")              #c# Returns 390
f.math.eseries(4.7e3, "E12")            #c# Returns 4700
f.math.eseries(1.5e3, "E3")             #c# Returns 2200
f.math.eseries(330, "e24")              #c# Returns 330 (case-insensitive)
f.math.eseries(0, "E24")               #c# Returns 0
f.math.eseries(330, "E5")              #c# Value Error
f.math.eseries(-100, "E24")            #c# Value Error
f.math.eseries("abc", "E24")           #c# Type Error
f.math.eseries(330, "E24", rm="floor") #c# Returns 330
f.math.eseries(350, "E24", rm="ceil")  #c# Returns 390
\end{verbatim}

\paragraph{E-Series Reference}
All values are base values within one decade (1.0 to 10). Multiply by $10^n$ for other decades. Every value in E($n$) is also a member of E($2n$).\\
\textbf{E3} (20\%): 1.0, 2.2, 4.7\\
\textbf{E6} (20\%): 1.0, 1.5, 2.2, 3.3, 4.7, 6.8\\
\textbf{E12} (10\%): 1.0, 1.2, 1.5, 1.8, 2.2, 2.7, 3.3, 3.9, 4.7, 5.6, 6.8, 8.2\\
\textbf{E24} (5\%): 1.0, 1.1, 1.2, 1.3, 1.5, 1.6, 1.8, 2.0, 2.2, 2.4, 2.7, 3.0, 3.3, 3.6, 3.9, 4.3, 4.7, 5.1, 5.6, 6.2, 6.8, 7.5, 8.2, 9.1\\
\textbf{E48} (2\%): 1.00, 1.05, 1.10, 1.15, 1.21, 1.27, 1.33, 1.40, 1.47, 1.54, 1.62, 1.69, 1.78, 1.87, 1.96, 2.05, 2.15, 2.26, 2.37, 2.49, 2.61, 2.74, 2.87, 3.01, 3.16, 3.32, 3.48, 3.65, 3.83, 4.02, 4.22, 4.42, 4.64, 4.87, 5.11, 5.36, 5.62, 5.90, 6.19, 6.49, 6.81, 7.15, 7.50, 7.87, 8.25, 8.66, 9.09, 9.53\\
\textbf{E96} (1\%): 1.00, 1.02, 1.05, 1.07, 1.10, 1.13, 1.15, 1.18, 1.21, 1.24, 1.27, 1.30, 1.33, 1.37, 1.40, 1.43, 1.47, 1.50, 1.54, 1.58, 1.62, 1.65, 1.69, 1.74, 1.78, 1.82, 1.87, 1.91, 1.96, 2.00, 2.05, 2.10, 2.15, 2.21, 2.26, 2.32, 2.37, 2.43, 2.49, 2.55, 2.61, 2.67, 2.74, 2.80, 2.87, 2.94, 3.01, 3.09, 3.16, 3.24, 3.32, 3.40, 3.48, 3.57, 3.65, 3.74, 3.83, 3.92, 4.02, 4.12, 4.22, 4.32, 4.42, 4.53, 4.64, 4.75, 4.87, 4.99, 5.11, 5.23, 5.36, 5.49, 5.62, 5.76, 5.90, 6.04, 6.19, 6.34, 6.49, 6.65, 6.81, 6.98, 7.15, 7.32, 7.50, 7.68, 7.87, 8.06, 8.25, 8.45, 8.66, 8.87, 9.09, 9.31, 9.53, 9.76\\
\textbf{E192} (0.1\%): 1.00, 1.01, 1.02, 1.04, 1.05, 1.06, 1.07, 1.09, 1.10, 1.11, 1.13, 1.14, 1.15, 1.17, 1.18, 1.20, 1.21, 1.23, 1.24, 1.26, 1.27, 1.29, 1.30, 1.32, 1.33, 1.35, 1.37, 1.38, 1.40, 1.42, 1.43, 1.45, 1.47, 1.49, 1.50, 1.52, 1.54, 1.56, 1.58, 1.60, 1.62, 1.64, 1.65, 1.67, 1.69, 1.72, 1.74, 1.76, 1.78, 1.80, 1.82, 1.84, 1.87, 1.89, 1.91, 1.93, 1.96, 1.98, 2.00, 2.03, 2.05, 2.08, 2.10, 2.13, 2.15, 2.18, 2.21, 2.23, 2.26, 2.29, 2.32, 2.34, 2.37, 2.40, 2.43, 2.46, 2.49, 2.52, 2.55, 2.58, 2.61, 2.64, 2.67, 2.71, 2.74, 2.77, 2.80, 2.84, 2.87, 2.91, 2.94, 2.98, 3.01, 3.05, 3.09, 3.12, 3.16, 3.20, 3.24, 3.28, 3.32, 3.36, 3.40, 3.44, 3.48, 3.52, 3.57, 3.61, 3.65, 3.70, 3.74, 3.79, 3.83, 3.88, 3.92, 3.97, 4.02, 4.07, 4.12, 4.17, 4.22, 4.27, 4.32, 4.37, 4.42, 4.48, 4.53, 4.59, 4.64, 4.70, 4.75, 4.81, 4.87, 4.93, 4.99, 5.05, 5.11, 5.17, 5.23, 5.30, 5.36, 5.42, 5.49, 5.56, 5.62, 5.69, 5.76, 5.83, 5.90, 5.97, 6.04, 6.12, 6.19, 6.26, 6.34, 6.42, 6.49, 6.57, 6.65, 6.73, 6.81, 6.90, 6.98, 7.06, 7.15, 7.23, 7.32, 7.41, 7.50, 7.59, 7.68, 7.77, 7.87, 7.96, 8.06, 8.16, 8.25, 8.35, 8.45, 8.56, 8.66, 8.76, 8.87, 8.98, 9.09, 9.20, 9.31, 9.42, 9.53, 9.65, 9.76, 9.88


\subsubsection{Minimum Function}
Returns the minimum value in a list. For ndarrays, an optional \texttt{axis} parameter reduces along that axis. \\
Morphology:
\begin{verbatim}
EBNF:
"f.math.min", "(", list, [ ",", "axis", "=", int ], ")"
\end{verbatim}
\begin{itemize}
\item Operands: \texttt{list} of comparable elements, or \texttt{ndarray}. Optional \texttt{axis}: \texttt{int} (ndarray only).
\item Result: the smallest element in the list/array.
\item For ndarrays with \texttt{axis}, reduces along the specified axis, returning an ndarray.
\item A \texttt{Value Error} is raised if the list is empty.
\item A \texttt{Type Error} is raised if the argument is not a list or ndarray.
\item A \texttt{Type Error} is raised if elements are not comparable (e.g., mixing \texttt{int} and \texttt{string}).
\end{itemize}
Examples:
\begin{verbatim}
f.math.min([3, 1, 4, 1, 5])              #c# Returns 1
f.math.min([[3,1],[4,2]], axis=0)        #c# Returns [3, 1]
f.math.min([[3,1],[4,2]], axis=1)        #c# Returns [1, 2]
\end{verbatim}

\subsubsection{Maximum Function}
Returns the maximum value in a list. For ndarrays, an optional \texttt{axis} parameter reduces along that axis. \\
Morphology:
\begin{verbatim}
EBNF:
"f.math.max", "(", list, [ ",", "axis", "=", int ], ")"
\end{verbatim}
\begin{itemize}
\item Operands: \texttt{list} of comparable elements, or \texttt{ndarray}. Optional \texttt{axis}: \texttt{int} (ndarray only).
\item Result: the largest element in the list/array.
\item For ndarrays with \texttt{axis}, reduces along the specified axis, returning an ndarray.
\item A \texttt{Value Error} is raised if the list is empty.
\item A \texttt{Type Error} is raised if the argument is not a list or ndarray.
\item A \texttt{Type Error} is raised if elements are not comparable.
\end{itemize}
Examples:
\begin{verbatim}
f.math.max([3, 1, 4, 1, 5])              #c# Returns 5
f.math.max([[3,1],[4,2]], axis=0)        #c# Returns [4, 2]
f.math.max([[3,1],[4,2]], axis=1)        #c# Returns [3, 4]
\end{verbatim}

\subsubsection{Inverse Tangent Function based on a Vector Argument}
The \texttt{atan2} function computes the arctangent of $y/x$ using the signs of both arguments to determine the correct quadrant. \\
Morphology:
\begin{verbatim}
EBNF:
"f.trig.atan2", "(", argument, ",", argument, ")"
\end{verbatim}
\paragraph{Edge Cases}
\begin{itemize}
\item Operands: both arguments must be \texttt{int} or \texttt{real}.
\item \texttt{atan2(0, 0)} returns \texttt{0.0}.
\item A \texttt{Type Error} is raised if either argument is not \texttt{int} or \texttt{real}.
\end{itemize}
Example:
\begin{verbatim}
f.atan2(1, 0)    #c# Returns approximately 1.5708 (pi/2)
f.atan2(0, 1)    #c# Returns 0.0
\end{verbatim}

\subsubsection{Norm Function}
Computes the p-norm of a vector or matrix. \\
Morphology:
\begin{verbatim}
EBNF:
"f.arr.norm", "(", ( list | list[list] ), [",", "p", "=", ( real | "fro" ) ], ")"
\end{verbatim}
\begin{itemize}
\item Operands: \texttt{list} of \texttt{int}/\texttt{real} (vector), \texttt{list[list]} of \texttt{int}/\texttt{real} (matrix), optional \texttt{p} (\texttt{int}, \texttt{real}, or \texttt{"fro"}, default 2).
\item Result: \texttt{real}.
\end{itemize}

\textbf{Vector norms:}
\begin{itemize}
\item \texttt{p = 0} -- count of non-zero elements.
\item \texttt{p = 1} -- L1/Manhattan norm.
\item \texttt{p = 2} -- L2/Euclidean norm (default).
\item Any real $p \geq 2$ -- Lp norm.
\item \texttt{p = -1} -- L-infinity/Chebyshev norm (max absolute value).
\item A \texttt{Value Error} is raised for $p < -1$ or $-1 < p < 0$.
\end{itemize}

\textbf{Matrix norms:}
\begin{itemize}
\item \texttt{p = 1} -- maximum absolute column sum.
\item \texttt{p = -1} -- minimum absolute column sum.
\item \texttt{p = 2} (default) -- spectral norm (largest singular value).
\item \texttt{p = -2} -- smallest singular value.
\item \texttt{p = "fro"} -- Frobenius norm ($\sqrt{\sum_{i,j} a_{ij}^2}$).
\item A \texttt{Value Error} is raised for unsupported matrix \texttt{p} values.
\end{itemize}

\begin{itemize}
\item A \texttt{Type Error} is raised if the argument is not a list of numeric values or a list of lists of numeric values.
\end{itemize}
Examples:
\begin{verbatim}
# Vector norms
f.arr.norm([3, 4])            #c# Returns 5.0 (L2 norm)
f.arr.norm([3, 4], p=1)       #c# Returns 7 (L1 norm)
f.arr.norm([3, 4], p=0)       #c# Returns 2 (L0 norm)
f.arr.norm([-3, 4, -2], p=-1) #c# Returns 4 (L-inf norm)

# Matrix norms
f.arr.norm([[1,2],[3,4]])              #c# Returns approximately 5.4649 (spectral norm)
f.arr.norm([[1,2],[3,4]], p="fro")     #c# Returns approximately 5.4772 (Frobenius)
f.arr.norm([[1,2],[3,4]], p=1)         #c# Returns 7 (max column sum: |3|+|1|=4, |4|+|2|=6)
f.arr.norm([[1,2],[3,4]], p=-1)        #c# Returns 3 (min column sum)
f.arr.norm([[1,2],[3,4]], p=-2)        #c# Returns approximately 0.1010 (smallest singular value)
\end{verbatim}

\subsubsection{Normalize Function}
Returns a unit vector in the direction of the input vector: $\hat{\mathbf{x}} = \mathbf{x} / \|\mathbf{x}\|_p$. \\
Morphology:
\begin{verbatim}
EBNF:
"f.arr.normalize", "(", list, [",", "p", "=", real], ")"
\end{verbatim}
\begin{itemize}
\item Operands: \texttt{list} of \texttt{int} or \texttt{real} values, optional \texttt{p} (\texttt{int} or \texttt{real}, default 2).
\item Result: \texttt{list} of \texttt{real} values.
\item Uses the same norm as \texttt{f.arr.norm} for the denominator.
\item A \texttt{Division by Zero Error} is raised if the norm of the vector is zero.
\item A \texttt{Type Error} is raised if the argument is not a list of numeric values.
\end{itemize}
Examples:
\begin{verbatim}
f.arr.normalize([3, 4])        #c# Returns [0.6, 0.8] (L2)
f.arr.normalize([3, 4], p=1)   #c# Returns [0.43, 0.57] (L1)
f.arr.normalize([6, 8], p=2)   #c# Returns [0.6, 0.8] (L2)
\end{verbatim}

\paragraph{Relationship between abs(vec=True) and norm}
The function \texttt{abs(v, vec=True)} is equivalent to \texttt{norm(v)} (L2 norm by default).

\subsubsection{Clamp}
The \texttt{clamp} function limits a value to a specified range.\\
Morphology:
\begin{verbatim}
EBNF:
"f.math.clamp", "(", ( int | real ), ",", ( int | real ), ",", ( int | real ), ")"
\end{verbatim}
\texttt{f.math.clamp(x, min, max)} returns \texttt{min} if \texttt{x $<$ min}, \texttt{max} if \texttt{x $>$ max}, otherwise \texttt{x}.
Operates element-wise on lists.\\
\paragraph{Edge Cases}
\begin{itemize}
\item A \texttt{Value Error} is raised if \texttt{min $>$ max}.
\item A \texttt{Type Error} is raised if all three arguments are not the same numeric type (\texttt{int}, \texttt{real}).
\end{itemize}
Examples:
\begin{verbatim}
f.math.clamp(5, 0, 10)        #c# Returns 5
f.math.clamp(-3, 0, 10)       #c# Returns 0
f.math.clamp(15, 0, 10)       #c# Returns 10
f.math.clamp([-2, 5, 12], 0, 10)  #c# Returns [0, 5, 10]
\end{verbatim}

\subsubsection{Linear Interpolation}
The \texttt{lerp} function performs linear interpolation between two values. \\
Morphology:
\begin{verbatim}
EBNF:
"f.math.lerp", "(", ( int | real ), ",", ( int | real ), ",", real, [",", "clamp", "=", bool], ")"
\end{verbatim}
\begin{itemize}
\item Operands: \texttt{a} (start value), \texttt{b} (end value), \texttt{t} (interpolation factor), optional \texttt{clamp} (\texttt{bool}, default \texttt{True}).
\item Result: \texttt{real} (always \texttt{real}, even when \texttt{a} and \texttt{b} are \texttt{int}).
\item Computes $a + (b - a) \cdot t$.
\item When \texttt{clamp = True} (default), \texttt{t} is clamped to $[0, 1]$ and the result is always between \texttt{a} and \texttt{b}.
\item When \texttt{clamp = False}, extrapolation is allowed for any \texttt{t}.
\item A \texttt{Type Error} is raised if arguments are not numeric or \texttt{clamp} is not \texttt{bool}.
\end{itemize}
Examples:
\begin{verbatim}
f.math.lerp(0, 10, 0.5)              #c# Returns 5.0
f.math.lerp(0, 10, 0.0)              #c# Returns 0.0
f.math.lerp(0, 10, 1.0)              #c# Returns 10.0
f.math.lerp(0, 10, 1.5)              #c# Returns 10.0 (clamped by default)
f.math.lerp(0, 10, 1.5, clamp=False) #c# Returns 15.0 (extrapolated)
f.math.lerp(0, 10, -0.5, clamp=False) #c# Returns -5.0 (extrapolated)
\end{verbatim}

\subsubsection{Is Finite Function}
The \texttt{f.math.isFinite} function checks whether a value is finite (not NaN and not $\pm\infty$). \\
Morphology:
\begin{verbatim}
EBNF:
"f.math.isFinite", "(", ( int | real | complex ), ")"
\end{verbatim}
\begin{itemize}
\item Result: \texttt{bool} --- \texttt{True} if the value is finite, \texttt{False} otherwise.
\item \texttt{int} values are always finite.
\end{itemize}
Examples:
\begin{verbatim}
f.math.isFinite(42)       #c# True
f.math.isFinite(3.14)     #c# True
f.math.isFinite(0.0 / 0.0) #c# False (NaN)
f.math.isFinite(1.0 / 0.0) #c# False (inf)
\end{verbatim}

\subsubsection{Is NaN Function}
The \texttt{f.math.isNaN} function checks whether a value is NaN (Not a Number). \\
Morphology:
\begin{verbatim}
EBNF:
"f.math.isNaN", "(", ( int | real | complex ), ")"
\end{verbatim}
\begin{itemize}
\item Result: \texttt{bool} --- \texttt{True} if the value is NaN, \texttt{False} otherwise.
\item \texttt{int} values are never NaN.
\end{itemize}
Examples:
\begin{verbatim}
f.math.isNaN(42)          #c# False
f.math.isNaN(0.0 / 0.0)   #c# True
f.math.isNaN(3.14)        #c# False
\end{verbatim}

\subsubsection{Unit in the Last Place Function}
The \texttt{f.math.ulp} function returns the unit in the last place of a floating-point number --- the gap between \texttt{x} and its next representable neighbor. Useful for numerical precision analysis. \\
Morphology:
\begin{verbatim}
EBNF:
"f.math.ulp", "(", real, ")"
\end{verbatim}
\begin{itemize}
\item Operands: \texttt{real}.
\item Result: \texttt{real}.
\item A \texttt{Value Error} is raised if the argument is not finite.
\end{itemize}
Examples:
\begin{verbatim}
f.math.ulp(1.0)      #c# 2.220446049250313e-16 (for f64)
f.math.ulp(1000.0)   #c# 1.1368683772161603e-13
f.math.ulp(0.0)      #c# 5e-324 (smallest subnormal)
\end{verbatim}

\subsubsection{Machine Epsilon Function}
The \texttt{f.math.epsilon} function returns the machine epsilon --- the smallest positive value $\epsilon$ such that $1.0 + \epsilon \neq 1.0$. \\
Morphology:
\begin{verbatim}
EBNF:
"f.math.epsilon", "(", ")"
\end{verbatim}
\begin{itemize}
\item Result: \texttt{real} --- approximately $2.22 \times 10^{-16}$ for IEEE 754 f64.
\end{itemize}
Examples:
\begin{verbatim}
$f eps = f.math.epsilon()
#c# 2.220446049250313e-16
1.0 + $eps / 2 == 1.0   #c# True (below precision threshold)
1.0 + $eps == 1.0       #c# False (above precision threshold)
\end{verbatim}

\subsubsection{Truncate}
The \texttt{trunc} function truncates the decimal part of a number, returning the integer part toward zero.\\
Morphology:
\begin{verbatim}
EBNF:
"f.math.trunc", "(", ( int | real ), ")"
\end{verbatim}
\begin{itemize}
\item Operands: \texttt{int} (returned unchanged), \texttt{real}, \texttt{complex}
\item Result: \texttt{int}
\item For complex numbers, truncates real and imaginary parts separately and returns complex.
\end{itemize}
Examples:
\begin{verbatim}
f.math.trunc(3.7)        #c# Returns 3
f.math.trunc(-3.7)       #c# Returns -3
f.math.trunc(3.2 + 4.8i) #c# Returns 3 + 4i
\end{verbatim}

\subsubsection{Cube Root}
The \texttt{cbrt} function returns the real cube root of a number.\\
Morphology:
\begin{verbatim}
EBNF:
"f.math.cbrt", "(", ( int | real ), ")"
\end{verbatim}
\begin{itemize}
\item Operands: \texttt{int}, \texttt{real}
\item Result: \texttt{real}
\item Returns the real cube root (negative inputs yield negative results, unlike \texttt{\^{}}).
\end{itemize}
Examples:
\begin{verbatim}
f.math.cbrt(27)         #c# Returns 3.0
f.math.cbrt(-8)         #c# Returns -2.0
f.math.cbrt(0)          #c# Returns 0.0
\end{verbatim}

\subsubsection{Nth Root}
The \texttt{nthrt} function returns the $n$-th root of a value, computed as $x^{1/n}$. Supports negative and complex $n$ (reciprocal and complex roots), though \textbf{integer $n \geq 1$ is encouraged}.\\
Morphology:
\begin{verbatim}
EBNF:
"f.math.nthrt", "(", ( int | real | complex ), ",", ( int | real | complex ), ")"
\end{verbatim}
\begin{itemize}
\item Operands: \texttt{int}, \texttt{real}, or \texttt{complex} for both $n$ and $x$.
\item Result: \texttt{complex} if either operand is \texttt{complex}, otherwise \texttt{real}.
\item Positive $n$: returns $x^{1/n}$. For odd integer $n$ with negative real $x$, returns the negative real root (like \texttt{cbrt}). For even $n$ with negative $x$, returns \texttt{complex}.
\item Negative $n$: returns $x^{1/n} = 1 / x^{1/|n|}$ (reciprocal root).
\item $n = 0$: raises a \texttt{Division by Zero Error}.
\item $x = 0$, $n > 0$: returns $0.0$.
\item $x = 0$, $n \leq 0$: raises a \texttt{Division by Zero Error}.
\item A \texttt{Type Error} is raised if either argument is not numeric.
\end{itemize}
Examples:
\begin{verbatim}
f.math.nthrt(2, 9)        #c# Returns 3.0 (square root)
f.math.nthrt(3, 27)       #c# Returns 3.0 (cube root)
f.math.nthrt(4, 16)       #c# Returns 2.0 (fourth root)
f.math.nthrt(-2, 9)       #c# Returns 0.333... (reciprocal: 1/sqrt(9))
f.math.nthrt(3, -8)       #c# Returns -2.0 (real odd root)
f.math.nthrt(2, -4)       #c# Returns 0 + 2i (complex even root)
f.math.nthrt(1, 42)       #c# Returns 42.0 (first root)
f.math.nthrt(0.5, 4)      #c# Returns 16.0 (real n)
f.math.nthrt(2+0i, 9)     #c# Returns 3 + 0i (complex n)
\end{verbatim}

\subsubsection{Base-2 Logarithm}
The \texttt{log2} function returns the base-2 logarithm.\\
Morphology:
\begin{verbatim}
EBNF:
"f.log.log2", "(", ( int | real ), ")"
\end{verbatim}
\begin{itemize}
\item Operands: \texttt{int} (must be $>$ 0), \texttt{real} (must be $>$ 0)
\item Result: \texttt{real}
\item A Value Error is raised for arguments $\leq$ 0.
\end{itemize}
Examples:
\begin{verbatim}
f.log.log2(8)          #c# Returns 3.0
f.log.log2(1)          #c# Returns 0.0
f.log.log2(1024)       #c# Returns 10.0
\end{verbatim}

\subsubsection{General Logarithm}
Returns the logarithm of the argument to an arbitrary base, computed as $\log_b(x) = \frac{\ln(x)}{\ln(b)}$. \\
Morphology:
\begin{verbatim}
EBNF:
"f.log.log", "(", ( int | real ), ",", ( int | real | complex ), ")"
\end{verbatim}
\begin{itemize}
\item \texttt{base}: \texttt{int} or \texttt{real} (must be $> 0$ and $\neq 1$).
\item \texttt{x}: \texttt{int} (must be $> 0$), \texttt{real} (must be $> 0$), \texttt{complex} ($\neq 0$).
\item Result: \texttt{real} for positive \texttt{int}/\texttt{real} input, \texttt{complex} for \texttt{complex} input.
\item A \texttt{Value Error} is raised if \texttt{base} $\leq 0$ or \texttt{base} $= 1$.
\item A \texttt{Value Error} is raised for non-positive real arguments or zero complex argument.
\end{itemize}
Examples:
\begin{verbatim}
f.log.log(10, 100)    #c# Returns 2.0
f.log.log(2, 8)       #c# Returns 3.0
f.log.log(c.math.e, 1)       #c# Returns 0.0
f.log.log(10, 1 + 1i) #c# Returns approximately 0.6505 + 0.3411i
\end{verbatim}

\subsubsection{Hypotenuse}
The \texttt{hypot} function computes $\sqrt{a^2 + b^2}$, the Euclidean distance / hypotenuse.\\
Morphology:
\begin{verbatim}
EBNF:
"f.math.hypot", "(", ( int | real ), ",", ( int | real ), ")"
\end{verbatim}
\begin{itemize}
\item Operands: two numeric values (\texttt{int} or \texttt{real})
\item Result: \texttt{real}
\item Avoids intermediate overflow/underflow for large/small values.
\end{itemize}
Examples:
\begin{verbatim}
f.math.hypot(3, 4)      #c# Returns 5.0
f.math.hypot(5, 12)     #c# Returns 13.0
f.math.hypot(0, 1)      #c# Returns 1.0
\end{verbatim}

\subsubsection{Sinc Function}
The \texttt{sinc} function computes the normalized sinc function $\text{sinc}(x) = \frac{\sin(\pi x)}{\pi x}$ for $x \neq 0$, and $1$ for $x = 0$. \\
Morphology:
\begin{verbatim}
EBNF:
"f.math.sinc", "(", ( int | real ), ")"
\end{verbatim}
\begin{itemize}
\item Operands: \texttt{int} or \texttt{real}
\item Result: \texttt{real}
\item $\text{sinc}(0) = 1.0$. Zeros at all non-zero integers.
\end{itemize}
Examples:
\begin{verbatim}
f.math.sinc(0)          #c# Returns 1.0
f.math.sinc(0.5)        #c# Returns approximately 0.6366
f.math.sinc(1)          #c# Returns 0.0
f.math.sinc(2)          #c# Returns 0.0
\end{verbatim}

\subsubsection{Tanc Function}
The \texttt{tanc} function computes $\text{tanc}(x) = \frac{\tan(x)}{x}$ for $x \neq 0$, and $1$ for $x = 0$. \\
Morphology:
\begin{verbatim}
EBNF:
"f.math.tanc", "(", ( int | real ), ")"
\end{verbatim}
\begin{itemize}
\item Operands: \texttt{int} or \texttt{real}
\item Result: \texttt{real}
\item $\text{tanc}(0) = 1.0$. Raises a Division by Zero Error at $x = \pm\pi/2$ (where $\tan$ is undefined).
\end{itemize}
Examples:
\begin{verbatim}
f.math.tanc(0)          #c# Returns 1.0
f.math.tanc(0.5)        #c# Returns approximately 1.0450
f.math.tanc(1)          #c# Returns approximately 1.5574
\end{verbatim}

\subsubsection{Factorial}
The \texttt{factorial} function returns $n! = n \times (n-1) \times \cdots \times 1$.\\
Morphology:
\begin{verbatim}
EBNF:
"f.math.factorial", "(", int, ")"
\end{verbatim}
\begin{itemize}
\item Operands: \texttt{int} (must be $\geq$ 0)
\item Result: \texttt{int}
\item $0! = 1$. A Value Error is raised for negative arguments.
\end{itemize}
Examples:
\begin{verbatim}
f.math.factorial(0)     #c# Returns 1
f.math.factorial(5)     #c# Returns 120
f.math.factorial(10)    #c# Returns 3628800
\end{verbatim}

\subsubsection{Greatest Common Divisor}
Computes the greatest common divisor of two integers. \\
Morphology:
\begin{verbatim}
EBNF:
"f.math.gcd", "(", int, ",", int, ")"
\end{verbatim}
\begin{itemize}
\item Operands: two \texttt{int} values.
\item Result: \texttt{int} ($\geq 0$).
\item $\gcd(0, 0) = 0$. $\gcd(a, 0) = |a|$. $\gcd(a, b) = \gcd(b, a \bmod b)$.
\end{itemize}
Examples:
\begin{verbatim}
f.math.gcd(12, 8)      #c# Returns 4
f.math.gcd(0, 5)       #c# Returns 5
f.math.gcd(7, 13)      #c# Returns 1
\end{verbatim}

\subsubsection{Least Common Multiple}
Computes the least common multiple of two integers. \\
Morphology:
\begin{verbatim}
EBNF:
"f.math.lcm", "(", int, ",", int, ")"
\end{verbatim}
\begin{itemize}
\item Operands: two \texttt{int} values.
\item Result: \texttt{int} ($\geq 0$).
\item $\text{lcm}(0, 0) = 0$. $\text{lcm}(a, 0) = 0$. $\text{lcm}(a, b) = |a \cdot b| / \gcd(a, b)$.
\end{itemize}
Examples:
\begin{verbatim}
f.math.lcm(4, 6)       #c# Returns 12
f.math.lcm(0, 5)       #c# Returns 0
f.math.lcm(7, 13)      #c# Returns 91
\end{verbatim}

\subsubsection{Double Factorial}
Computes the double factorial $n!! = n \times (n-2) \times (n-4) \times \cdots$. \\
Morphology:
\begin{verbatim}
EBNF:
"f.math.doubleFactorial", "(", int, ")"
\end{verbatim}
\begin{itemize}
\item Operand: \texttt{int} ($\geq 0$).
\item Result: \texttt{int}.
\item $0!! = 1$, $1!! = 1$, $5!! = 15$, $6!! = 48$.
\item A \texttt{Value Error} is raised for negative arguments.
\end{itemize}
Examples:
\begin{verbatim}
f.math.doubleFactorial(0)    #c# Returns 1
f.math.doubleFactorial(5)    #c# Returns 15  (5 * 3 * 1)
f.math.doubleFactorial(6)    #c# Returns 48  (6 * 4 * 2)
\end{verbatim}

\subsubsection{Falling Factorial}
Computes the falling factorial $n^{\underline{k}} = n \times (n-1) \times \cdots \times (n-k+1)$. \\
Morphology:
\begin{verbatim}
EBNF:
"f.math.fallingFactorial", "(", int, ",", int, ")"
\end{verbatim}
\begin{itemize}
\item Operands: $n$ (\texttt{int}), $k$ (\texttt{int} $\geq 0$).
\item Result: \texttt{int}.
\item $n^{\underline{0}} = 1$. $n^{\underline{1}} = n$. $n^{\underline{n}} = n!$.
\item A \texttt{Value Error} is raised if $k < 0$.
\end{itemize}
Examples:
\begin{verbatim}
f.math.fallingFactorial(5, 3)   #c# Returns 60  (5 * 4 * 3)
f.math.fallingFactorial(5, 0)   #c# Returns 1
\end{verbatim}

\subsubsection{Rising Factorial}
Computes the rising factorial (Pochhammer symbol) $n^{\overline{k}} = n \times (n+1) \times \cdots \times (n+k-1)$. \\
Morphology:
\begin{verbatim}
EBNF:
"f.math.risingFactorial", "(", int, ",", int, ")"
\end{verbatim}
\begin{itemize}
\item Operands: $n$ (\texttt{int}), $k$ (\texttt{int} $\geq 0$).
\item Result: \texttt{int}.
\item $n^{\overline{0}} = 1$. $n^{\overline{1}} = n$. $n^{\overline{k}} = (n+k-1)! / (n-1)!$.
\item A \texttt{Value Error} is raised if $k < 0$.
\end{itemize}
Examples:
\begin{verbatim}
f.math.risingFactorial(3, 4)   #c# Returns 360  (3 * 4 * 5 * 6)
f.math.risingFactorial(1, 5)   #c# Returns 120  (1 * 2 * 3 * 4 * 5)
\end{verbatim}

\subsubsection{Binomial Coefficient}
Computes $\binom{n}{k} = \frac{n!}{k!(n-k)!}$. Same as \texttt{f.math.comb}. \\
Morphology:
\begin{verbatim}
EBNF:
"f.math.binomial", "(", int, ",", int, ")"
\end{verbatim}
\begin{itemize}
\item Operands: $n$ (\texttt{int} $\geq 0$), $k$ (\texttt{int} $\geq 0$).
\item Result: \texttt{int}.
\item $\binom{n}{0} = 1$, $\binom{n}{1} = n$, $\binom{n}{n} = 1$.
\item A \texttt{Value Error} is raised if $k > n$ or $n < 0$ or $k < 0$.
\end{itemize}
Examples:
\begin{verbatim}
f.math.binomial(5, 3)    #c# Returns 10
f.math.binomial(10, 2)   #c# Returns 45
\end{verbatim}

\subsubsection{Multinomial Coefficient}
Computes $\binom{n}{k_1, k_2, \ldots, k_m} = \frac{n!}{k_1!\, k_2!\, \cdots\, k_m!}$ where $n = \sum k_i$. \\
Morphology:
\begin{verbatim}
EBNF:
"f.math.multinomial", "(", list, ")"
\end{verbatim}
\begin{itemize}
\item Operand: \texttt{list[int]} ($\geq 0$).
\item Result: \texttt{int}.
\item The sum of all elements is $n$.
\item A \texttt{Value Error} is raised if any element is negative.
\end{itemize}
Examples:
\begin{verbatim}
f.math.multinomial([1, 2, 3])    #c# Returns 60  (6! / (1!*2!*3!))
f.math.multinomial([2, 2])       #c# Returns 6   (4! / (2!*2!))
f.math.multinomial([5])          #c# Returns 1
\end{verbatim}

\subsubsection{Modular Inverse}
Computes the modular multiplicative inverse $x$ such that $a \cdot x \equiv 1 \pmod{m}$. \\
Morphology:
\begin{verbatim}
EBNF:
"f.math.modInverse", "(", int, ",", int, ")"
\end{verbatim}
\begin{itemize}
\item Operands: $a$ (\texttt{int}), $m$ (\texttt{int}, $> 0$).
\item Result: \texttt{int} in range $[0, m)$.
\item A \texttt{Value Error} is raised if $\gcd(a, m) \neq 1$ (no inverse exists) or $m \leq 0$.
\end{itemize}
Examples:
\begin{verbatim}
f.math.modInverse(3, 7)     #c# Returns 5  (3*5 = 15 ≡ 1 mod 7)
f.math.modInverse(10, 17)   #c# Returns 12  (10*12 = 120 ≡ 1 mod 17)
\end{verbatim}

\subsubsection{Divides}
Tests whether $a$ divides $b$ evenly ($a \mid b$). \\
Morphology:
\begin{verbatim}
EBNF:
"f.math.divides", "(", int, ",", int, ")"
\end{verbatim}
\begin{itemize}
\item Operands: $a$ (\texttt{int}), $b$ (\texttt{int}).
\item Result: \texttt{bool}. Returns \texttt{True} if $b \bmod a = 0$.
\item A \texttt{Value Error} is raised if $a = 0$.
\end{itemize}
Examples:
\begin{verbatim}
f.math.divides(3, 9)     #c# Returns True
f.math.divides(3, 10)    #c# Returns False
f.math.divides(1, 7)     #c# Returns True  (1 divides everything)
\end{verbatim}

\subsubsection{Permutations}
Computes the number of ordered selections: $P(n, k) = \frac{n!}{(n-k)!}$. \\
Morphology:
\begin{verbatim}
EBNF:
"f.math.perm", "(", int, ",", int, ")"
\end{verbatim}
\begin{itemize}
\item Operands: $n$ (\texttt{int} $\geq 0$), $k$ (\texttt{int} $\geq 0$).
\item Result: \texttt{int}.
\item $P(n, 0) = 1$, $P(n, 1) = n$.
\item A \texttt{Value Error} is raised if $k > n$ or $n < 0$ or $k < 0$.
\end{itemize}
Examples:
\begin{verbatim}
f.math.perm(5, 3)     #c# Returns 60
f.math.perm(10, 2)    #c# Returns 90
f.math.perm(5, 0)     #c# Returns 1
\end{verbatim}

\subsubsection{Combinations}
Computes the number of unordered selections: $C(n, k) = \frac{n!}{k!(n-k)!}$. \\
Morphology:
\begin{verbatim}
EBNF:
"f.math.comb", "(", int, ",", int, ")"
\end{verbatim}
\begin{itemize}
\item Operands: $n$ (\texttt{int} $\geq 0$), $k$ (\texttt{int} $\geq 0$).
\item Result: \texttt{int}.
\item $C(n, 0) = 1$, $C(n, 1) = n$, $C(n, n) = 1$.
\item A \texttt{Value Error} is raised if $k > n$ or $n < 0$ or $k < 0$.
\end{itemize}
Examples:
\begin{verbatim}
f.math.comb(5, 3)     #c# Returns 10
f.math.comb(10, 2)    #c# Returns 45
f.math.comb(5, 5)     #c# Returns 1
\end{verbatim}

\subsubsection{Partitions}
Computes the number of integer partitions of $n$, i.e.\ the number of ways to write $n$ as a sum of positive integers (order does not matter). \\
Morphology:
\begin{verbatim}
EBNF:
"f.math.partitions", "(", int, ")"
\end{verbatim}
\begin{itemize}
\item Operand: $n$ (\texttt{int} $\geq 0$).
\item Result: \texttt{int}.
\item $p(0) = 1$, $p(1) = 1$, $p(5) = 7$.
\item A \texttt{Value Error} is raised if $n < 0$.
\end{itemize}
Examples:
\begin{verbatim}
f.math.partitions(0)   #c# Returns 1
f.math.partitions(5)   #c# Returns 7  (5, 4+1, 3+2, 3+1+1, 2+2+1, 2+1+1+1, 1+1+1+1+1)
f.math.partitions(10)  #c# Returns 42
\end{verbatim}

\subsubsection{Compositions}
Computes the number of compositions of $n$, i.e.\ the number of ways to write $n$ as an ordered sum of positive integers. Equal to $2^{n-1}$ for $n \geq 1$. \\
Morphology:
\begin{verbatim}
EBNF:
"f.math.compositions", "(", int, ")"
\end{verbatim}
\begin{itemize}
\item Operand: $n$ (\texttt{int} $\geq 0$).
\item Result: \texttt{int}.
\item $c(0) = 1$, $c(1) = 1$, $c(n) = 2^{n-1}$ for $n \geq 1$.
\item A \texttt{Value Error} is raised if $n < 0$.
\end{itemize}
Examples:
\begin{verbatim}
f.math.compositions(0)   #c# Returns 1
f.math.compositions(1)   #c# Returns 1
f.math.compositions(4)   #c# Returns 8  (2^3)
\end{verbatim}

\subsubsection{Gray Code}
Computes the $n$-th value in the binary reflected Gray code sequence. \\
Morphology:
\begin{verbatim}
EBNF:
"f.math.grayCode", "(", int, ")"
\end{verbatim}
\begin{itemize}
\item Operand: $n$ (\texttt{int} $\geq 0$).
\item Result: \texttt{int}. Formula: $G(n) = n \oplus \lfloor n / 2 \rfloor$.
\item Adjacent Gray code values differ by exactly one bit.
\item A \texttt{Value Error} is raised if $n < 0$.
\end{itemize}
Examples:
\begin{verbatim}
f.math.grayCode(0)    #c# Returns 0
f.math.grayCode(1)    #c# Returns 1
f.math.grayCode(2)    #c# Returns 3
f.math.grayCode(3)    #c# Returns 2
f.math.grayCode(4)    #c# Returns 6
\end{verbatim}

\subsubsection{Derangements}
Computes the number of derangements $!n$ (permutations with no fixed points). \\
Morphology:
\begin{verbatim}
EBNF:
"f.math.derangements", "(", int, ")"
\end{verbatim}
\begin{itemize}
\item Operand: $n$ (\texttt{int} $\geq 0$).
\item Result: \texttt{int}.
\item $!0 = 1$, $!1 = 0$, $!2 = 1$, $!3 = 2$, $!4 = 9$.
\item A \texttt{Value Error} is raised if $n < 0$.
\end{itemize}
Examples:
\begin{verbatim}
f.math.derangements(0)   #c# Returns 1
f.math.derangements(1)   #c# Returns 0
f.math.derangements(4)   #c# Returns 9
f.math.derangements(10)  #c# Returns 1334961
\end{verbatim}

\subsubsection{Gamma Function}
The \texttt{gamma} function returns $\Gamma(x) = (x-1)!$ for positive integers, extended to real and complex numbers.\\
Morphology:
\begin{verbatim}
EBNF:
"f.special.gamma", "(", ( int | real ), ")"
\end{verbatim}
\begin{itemize}
\item Operands: \texttt{int}, \texttt{real}, \texttt{complex}
\item Result: \texttt{real} or \texttt{complex}
\item $\Gamma(1) = 1$, $\Gamma(0.5) = \sqrt{\pi}$. Poles at non-positive integers raise a Value Error.
\end{itemize}
Examples:
\begin{verbatim}
f.special.gamma(1)         #c# Returns 1.0
f.special.gamma(0.5)       #c# Returns approximately 1.77245
f.special.gamma(5)         #c# Returns 24.0  (= 4!)
\end{verbatim}

\subsubsection{Log-Gamma Function}
The \texttt{lgamma} function returns $\ln|\Gamma(x)|$, the natural logarithm of the absolute value of the gamma function. Useful for large factorials without overflow.\\
Morphology:
\begin{verbatim}
EBNF:
"f.special.lgamma", "(", ( int | real ), ")"
\end{verbatim}
\begin{itemize}
\item Operands: \texttt{int}, \texttt{real}
\item Result: \texttt{real}
\item Always returns a non-negative value.
\end{itemize}
Examples:
\begin{verbatim}
f.special.lgamma(1)        #c# Returns 0.0
f.special.lgamma(5)        #c# Returns approximately 3.17805  (= ln(24))
f.special.lgamma(10)       #c# Returns approximately 12.8018  (= ln(362880))
\end{verbatim}

\subsection{String Functions}

\subsubsection{String Conversion}
The \texttt{str} function converts its arguments to a string. \\
Morphology:
\begin{verbatim}
EBNF:
"f.str.str", "(", { argument }, ")"
\end{verbatim}
Example:
\begin{verbatim}
f.str.str("Hello")           #c# Returns "Hello"
f.str.str(42)                #c# Returns "42"
f.str.str("Value: ", 3.14)   #c# Returns "Value: 3.14"
f.str.str()                  #c# Returns "" (empty string)
\end{verbatim}

With zero arguments, \texttt{f.str.str()} returns an empty string.

\subsubsection{String Formatting}
The \texttt{format} function formats a value using printf-style format specifiers. \\
Morphology:
\begin{verbatim}
EBNF:
"f.str.format", "(", argument, ",", string_type, ")"
\end{verbatim}

Format specifiers:
\begin{itemize}
\item \texttt{\%d} -- Integer decimal
\item \texttt{\%f} -- Fixed-point real (default 6 decimal places)
\item \texttt{\%e} -- Scientific notation
\item \texttt{\%g} -- Shortest of \texttt{\%f} or \texttt{\%e}
\item \texttt{\%s} -- String
\item \texttt{\%x} -- Hexadecimal (integers)
\item \texttt{\%o} -- Octal (integers)
\item \texttt{\%b} -- Binary (integers)
\end{itemize}

Width and precision: \texttt{\%<width>f}, \texttt{\%<width>.<precision>f}

\paragraph{Edge Cases}
\begin{itemize}
\item A \texttt{Type Error} is raised if the value type does not match the format specifier (e.g. \texttt{f.str.format("abc", "\%d")}).
\end{itemize}

Examples:
\begin{verbatim}
f.str.format(3.14159, "%.2f")     #c# Returns "3.14"
f.str.format(42, "%d")            #c# Returns "42"
f.str.format(42, "%5d")           #c# Returns "   42"
f.str.format(0xff, "%x")          #c# Returns "ff"
f.str.format(3.14159, "%4.1f")    #c# Returns " 3.1"
f.str.format(1234567, "%e")       #c# Returns "1.234567e+06"
\end{verbatim}

\subsubsection{Uppercase Function}
Converts a string to uppercase. \\
Morphology:
\begin{verbatim}
EBNF:
"f.str.upperc", "(", string, ")"
\end{verbatim}
\begin{itemize}
\item Operands: \texttt{string}.
\item Result: \texttt{string}.
\item Empty string returns empty string.
\item A \texttt{Type Error} is raised if the argument is not a \texttt{string}.
\end{itemize}
Examples:
\begin{verbatim}
f.str.upperc("hello")   #c# Returns "HELLO"
f.str.upperc("")        #c# Returns ""
f.str.upperc("abc123")  #c# Returns "ABC123"
\end{verbatim}

\subsubsection{Lowercase Function}
Converts a string to lowercase. \\
Morphology:
\begin{verbatim}
EBNF:
"f.str.lowerc", "(", string, ")"
\end{verbatim}
\begin{itemize}
\item Operands: \texttt{string}.
\item Result: \texttt{string}.
\item Empty string returns empty string.
\item A \texttt{Type Error} is raised if the argument is not a \texttt{string}.
\end{itemize}
Examples:
\begin{verbatim}
f.str.lowerc("HELLO")   #c# Returns "hello"
f.str.lowerc("")        #c# Returns ""
f.str.lowerc("ABC123")  #c# Returns "abc123"
\end{verbatim}

\subsubsection{String Character Access}
Returns a single character at the given index. Strings support index and slice operations, similar to lists. \\
Syntax:
\begin{verbatim}
str[i]              Returns the character at index i
\end{verbatim}
\begin{itemize}
\item Index is zero-based.
\item Negative indices count from the end: \texttt{str[-1]} is the last character.
\item \texttt{str[-0]} equals \texttt{str[0]} (negative zero equals zero).
\item A \texttt{Value Error} is raised if the index is out of bounds.
\end{itemize}
Examples:
\begin{verbatim}
"hello"[0]           #c# Returns "h"
"hello"[-1]          #c# Returns "o"
\end{verbatim}

\subsubsection{String Slice}
Returns a substring using start:end:step slicing. \\
Syntax:
\begin{verbatim}
str[start]          Single character access
str[start:end]      Returns str[start] through str[end-1]
str[start:end:step] Returns every step-th character from start to end
\end{verbatim}
\begin{itemize}
\item \texttt{start} defaults to 0, \texttt{end} defaults to string length, \texttt{step} defaults to 1.
\item Negative \texttt{step} reverses the direction.
\item Out-of-range indices are clamped (not errors).
\end{itemize}
Examples:
\begin{verbatim}
"hello"[1:4]         #c# Returns "ell"
"hello"[0:5:2]       #c# Returns "hlo"
"hello"[::-1]        #c# Returns "olleh"
\end{verbatim}

\subsubsection{String Length Function}
Returns the number of characters in a string. \\
Morphology:
\begin{verbatim}
EBNF:
"f.str.length", "(", string, ")"
\end{verbatim}
\begin{itemize}
\item Operands: \texttt{string}.
\item Result: \texttt{int}.
\item \texttt{f.str.length("")} returns 0.
\item A \texttt{Type Error} is raised if the argument is not a \texttt{string}.
\end{itemize}
Examples:
\begin{verbatim}
f.str.length("")        #c# Returns 0
f.str.length("hello")   #c# Returns 5
\end{verbatim}

\subsubsection{Unicode Code Point Function}
Returns the Unicode code point of the character at the given index. \\
Morphology:
\begin{verbatim}
EBNF:
"f.str.charCodeAt", "(", string, ",", int, ")"
\end{verbatim}
\begin{itemize}
\item Operands: \texttt{string}, \texttt{int} (index).
\item Result: \texttt{int} (Unicode code point).
\item A \texttt{Value Error} is raised if the index is out of bounds.
\item A \texttt{Type Error} is raised if the first argument is not a \texttt{string} or the second is not an \texttt{int}.
\end{itemize}
Examples:
\begin{verbatim}
f.str.charCodeAt("A", 0)    #c# Returns 65
f.str.charCodeAt("hello", 4) #c# Returns 111
\end{verbatim}

\subsubsection{Contains Function}
Returns \texttt{True} if the string contains the substring. \\
Morphology:
\begin{verbatim}
EBNF:
"f.str.contains", "(", string, ",", string, ")"
\end{verbatim}
\begin{itemize}
\item Operands: \texttt{string}, \texttt{string} (substring).
\item Result: \texttt{bool}.
\item \texttt{f.str.contains("hello", "")} returns \texttt{True} (empty string is contained in any string).
\item A \texttt{Type Error} is raised if either argument is not a \texttt{string}.
\end{itemize}
Examples:
\begin{verbatim}
f.str.contains("hello world", "world")  #c# Returns True
f.str.contains("hello", "xyz")          #c# Returns False
f.str.contains("hello", "")             #c# Returns True
\end{verbatim}

\subsubsection{IndexOf Function}
Returns the index of the first occurrence of the substring, or $-1$ if not found. \\
Morphology:
\begin{verbatim}
EBNF:
"f.str.indexOf", "(", string, ",", string, ")"
\end{verbatim}
\begin{itemize}
\item Operands: \texttt{string}, \texttt{string} (substring).
\item Result: \texttt{int}.
\item Returns $-1$ if the substring is not found.
\item A \texttt{Type Error} is raised if either argument is not a \texttt{string}.
\end{itemize}
Examples:
\begin{verbatim}
f.str.indexOf("hello", "l")         #c# Returns 2
f.str.indexOf("hello", "xyz")       #c# Returns -1
\end{verbatim}

\subsubsection{LastIndexOf Function}
Returns the index of the last occurrence of the substring, or $-1$ if not found. \\
Morphology:
\begin{verbatim}
EBNF:
"f.str.lastIndexOf", "(", string, ",", string, ")"
\end{verbatim}
\begin{itemize}
\item Operands: \texttt{string}, \texttt{string} (substring).
\item Result: \texttt{int}.
\item Returns $-1$ if the substring is not found.
\item A \texttt{Type Error} is raised if either argument is not a \texttt{string}.
\end{itemize}
Examples:
\begin{verbatim}
f.str.lastIndexOf("hello", "l")     #c# Returns 3
f.str.lastIndexOf("hello", "xyz")   #c# Returns -1
\end{verbatim}

\subsubsection{Count Function}
Returns the number of non-overlapping occurrences of the substring. \\
Morphology:
\begin{verbatim}
EBNF:
"f.str.count", "(", string, ",", string, ")"
\end{verbatim}
\begin{itemize}
\item Operands: \texttt{string}, \texttt{string} (substring).
\item Result: \texttt{int}.
\item Occurrences do not overlap: \texttt{f.util.count("aaa", "aa")} returns 1.
\item A \texttt{Type Error} is raised if either argument is not a \texttt{string}.
\end{itemize}
Examples:
\begin{verbatim}
f.util.count("hello", "l")           #c# Returns 2
f.util.count("aaa", "aa")            #c# Returns 1 (non-overlapping)
\end{verbatim}

\subsubsection{StartsWith Function}
Returns \texttt{True} if the string starts with the given prefix. \\
Morphology:
\begin{verbatim}
EBNF:
"f.str.startsWith", "(", string, ",", string, ")"
\end{verbatim}
\begin{itemize}
\item Operands: \texttt{string}, \texttt{string} (prefix).
\item Result: \texttt{bool}.
\item A \texttt{Type Error} is raised if either argument is not a \texttt{string}.
\end{itemize}
Examples:
\begin{verbatim}
f.str.startsWith("hello", "hel")    #c# Returns True
f.str.startsWith("hello", "xyz")    #c# Returns False
\end{verbatim}

\subsubsection{EndsWith Function}
Returns \texttt{True} if the string ends with the given suffix. \\
Morphology:
\begin{verbatim}
EBNF:
"f.str.endsWith", "(", string, ",", string, ")"
\end{verbatim}
\begin{itemize}
\item Operands: \texttt{string}, \texttt{string} (suffix).
\item Result: \texttt{bool}.
\item A \texttt{Type Error} is raised if either argument is not a \texttt{string}.
\end{itemize}
Examples:
\begin{verbatim}
f.str.endsWith("hello", "llo")      #c# Returns True
f.str.endsWith("hello", "xyz")      #c# Returns False
\end{verbatim}

\subsubsection{Replace Function}
Replaces the first occurrence of the search string with the replacement string. \\
Morphology:
\begin{verbatim}
EBNF:
"f.str.replace", "(", string, ",", string, ",", string, ")"
\end{verbatim}
\begin{itemize}
\item Operands: \texttt{string} (haystack), \texttt{string} (search), \texttt{string} (replacement).
\item Result: \texttt{string}.
\item Only the first occurrence is replaced. Use \texttt{f.replaceAll} for all occurrences.
\item A \texttt{Type Error} is raised if any argument is not a \texttt{string}.
\end{itemize}
Examples:
\begin{verbatim}
f.str.replace("hello world", "world", "there")  #c# Returns "hello there"
f.str.replace("aabbcc", "a", "x")              #c# Returns "xabbcc"
\end{verbatim}

\subsubsection{ReplaceAll Function}
Replaces all occurrences of the search string with the replacement string. \\
Morphology:
\begin{verbatim}
EBNF:
"f.str.replaceAll", "(", string, ",", string, ",", string, ")"
\end{verbatim}
\begin{itemize}
\item Operands: \texttt{string} (haystack), \texttt{string} (search), \texttt{string} (replacement).
\item Result: \texttt{string}.
\item A \texttt{Type Error} is raised if any argument is not a \texttt{string}.
\end{itemize}
Examples:
\begin{verbatim}
f.str.replaceAll("aabbcc", "a", "x")  #c# Returns "xxbbcc"
\end{verbatim}

\subsubsection{Insert Function}
Inserts a string before the given index. \\
Morphology:
\begin{verbatim}
EBNF:
"f.str.insert", "(", string, ",", int, ",", string, ")"
\end{verbatim}
\begin{itemize}
\item Operands: \texttt{string}, \texttt{int} (index), \texttt{string} (insertion).
\item Result: \texttt{string}.
\item If the index is less than 0, it is clamped to 0 (insert at beginning).
\item If the index exceeds the string length, it is clamped to the string length (insert at end).
\item A \texttt{Type Error} is raised if the first or third argument is not a \texttt{string} or the second is not an \texttt{int}.
\end{itemize}
Examples:
\begin{verbatim}
f.str.insert("ace", 1, "b")   #c# Returns "abce"
f.str.insert("abc", 0, "x")   #c# Returns "xabc"
f.str.insert("abc", 3, "x")   #c# Returns "abcx"
f.str.insert("abc", 99, "x")  #c# Returns "abcx" (clamped to end)
\end{verbatim}

\subsubsection{Remove Function}
Removes characters from a string by index range or by value. \\
Morphology:
\begin{verbatim}
EBNF:
"f.str.remove", "(", string, ",", data, ",", data, [",", "mode", "=", string], ")"
\end{verbatim}
\begin{itemize}
\item Operands: \texttt{string}, second/third operands depend on \texttt{mode}, optional \texttt{mode} (\texttt{string}: \texttt{"index"} (default) or \texttt{"value"}).
\item \texttt{mode="index"}: \texttt{f.str.remove(str, start, end)} removes characters from index \texttt{start} up to (not including) \texttt{end}.
\item \texttt{mode="value"}: \texttt{f.str.remove(str, search, pos, mode="value")} removes the first occurrence of \texttt{search} starting at position \texttt{pos}.
\item Out-of-bounds indices are clamped to the string length.
\item A \texttt{Type Error} is raised if the first argument is not a \texttt{string}.
\end{itemize}
Examples:
\begin{verbatim}
f.str.remove("abcdef", 2, 4)                      #c# Returns "abef"
f.str.remove("hello world", 5, 6)                 #c# Returns "helloworld"
f.str.remove("aabbcc", "bb", 0, mode="value")     #c# Returns "aacc"
\end{verbatim}

\subsubsection{ToUpperCase Function}
Converts a string to uppercase. Identical in behavior to \texttt{f.str.upperc}. \\
Morphology:
\begin{verbatim}
EBNF:
"f.str.toUpperCase", "(", string, ")"
\end{verbatim}
\begin{itemize}
\item Operands: \texttt{string}.
\item Result: \texttt{string}.
\item A \texttt{Type Error} is raised if the argument is not a \texttt{string}.
\end{itemize}
Examples:
\begin{verbatim}
f.str.toUpperCase("hello")  #c# Returns "HELLO"
\end{verbatim}

\subsubsection{ToLowerCase Function}
Converts a string to lowercase. Identical in behavior to \texttt{f.str.lowerc}. \\
Morphology:
\begin{verbatim}
EBNF:
"f.str.toLowerCase", "(", string, ")"
\end{verbatim}
\begin{itemize}
\item Operands: \texttt{string}.
\item Result: \texttt{string}.
\item A \texttt{Type Error} is raised if the argument is not a \texttt{string}.
\end{itemize}
Examples:
\begin{verbatim}
f.str.toLowerCase("HELLO")  #c# Returns "hello"
\end{verbatim}

\subsubsection{Capitalize Function}
Returns the string with the first character uppercased and the rest lowercased. \\
Morphology:
\begin{verbatim}
EBNF:
"f.str.capitalize", "(", string, ")"
\end{verbatim}
\begin{itemize}
\item Operands: \texttt{string}.
\item Result: \texttt{string}.
\item Empty string returns empty string.
\item A \texttt{Type Error} is raised if the argument is not a \texttt{string}.
\end{itemize}
Examples:
\begin{verbatim}
f.str.capitalize("hello")       #c# Returns "Hello"
f.str.capitalize("hELLO")       #c# Returns "Hello"
f.str.capitalize("")            #c# Returns ""
\end{verbatim}

\subsubsection{Title Case Function}
Returns the string with the first character of each word uppercased and the rest lowercased. Words are delimited by whitespace. \\
Morphology:
\begin{verbatim}
EBNF:
"f.str.titleCase", "(", string, ")"
\end{verbatim}
\begin{itemize}
\item Operands: \texttt{string}.
\item Result: \texttt{string}.
\item A \texttt{Type Error} is raised if the argument is not a \texttt{string}.
\end{itemize}
Returns the string with the first character of each word uppercased.
\begin{verbatim}
f.str.titleCase("hello world")       #c# Returns "Hello World"
f.str.titleCase("foo bar baz")       #c# Returns "Foo Bar Baz"
\end{verbatim}

\subsubsection{Swap Case Function}
Swaps the case of each character in the string (uppercase $\leftrightarrow$ lowercase). \\
Morphology:
\begin{verbatim}
EBNF:
"f.str.swapCase", "(", string, ")"
\end{verbatim}
\begin{itemize}
\item Operands: \texttt{string}.
\item Result: \texttt{string}.
\item Non-alphabetic characters are unchanged.
\item A \texttt{Type Error} is raised if the argument is not a \texttt{string}.
\end{itemize}
Examples:
\begin{verbatim}
f.str.swapCase("Hello World")   #c# Returns "hELLO wORLD"
f.str.swapCase("ABC123")        #c# Returns "abc123"
\end{verbatim}

\subsubsection{Trim Function}
Removes leading and/or trailing whitespace from a string. \\
Morphology:
\begin{verbatim}
EBNF:
"f.str.trim", "(", string, ")"
\end{verbatim}
\begin{itemize}
\item Operands: \texttt{string}.
\item Result: \texttt{string}.
\item Removes spaces, tabs, and newlines from both ends.
\item A \texttt{Type Error} is raised if the argument is not a \texttt{string}.
\end{itemize}
Examples:
\begin{verbatim}
f.str.trim("  hello  ")      #c# Returns "hello"
\end{verbatim}

\subsubsection{TrimStart Function}
Removes leading whitespace from a string. \\
Morphology:
\begin{verbatim}
EBNF:
"f.str.trimStart", "(", string, ")"
\end{verbatim}
\begin{itemize}
\item Operands: \texttt{string}.
\item Result: \texttt{string}.
\item Removes spaces, tabs, and newlines from the beginning.
\item A \texttt{Type Error} is raised if the argument is not a \texttt{string}.
\end{itemize}
Examples:
\begin{verbatim}
f.trimStart("  hello  ") #c# Returns "hello  "
\end{verbatim}

\subsubsection{TrimEnd Function}
Removes trailing whitespace from a string. \\
Morphology:
\begin{verbatim}
EBNF:
"f.str.trimEnd", "(", string, ")"
\end{verbatim}
\begin{itemize}
\item Operands: \texttt{string}.
\item Result: \texttt{string}.
\item Removes spaces, tabs, and newlines from the end.
\item A \texttt{Type Error} is raised if the argument is not a \texttt{string}.
\end{itemize}
Examples:
\begin{verbatim}
f.trimEnd("  hello  ")   #c# Returns "  hello"
\end{verbatim}

\subsubsection{Split Function}
Splits a string into a list of substrings. If no delimiter is given, splits on whitespace (one or more whitespace characters are treated as a single delimiter). \\
Morphology:
\begin{verbatim}
EBNF:
"f.str.split", "(", string, [",", string], ")"
\end{verbatim}
\begin{itemize}
\item Operands: \texttt{string}, optional delimiter (\texttt{string}).
\item Result: \texttt{list} of \texttt{string}.
\item \texttt{f.str.split("")} returns \texttt{[""]} (a list containing one empty string).
\item If the delimiter is not found, returns a list containing the entire string.
\item A \texttt{Type Error} is raised if the argument is not a \texttt{string}.
\end{itemize}
Examples:
\begin{verbatim}
f.str.split("hello world")          #c# Returns ["hello", "world"]
f.str.split("a,b,c", ",")           #c# Returns ["a", "b", "c"]
f.str.split("one  two  three")      #c# Returns ["one", "two", "three"]
\end{verbatim}

\subsubsection{SplitLines Function}
Splits a string on newline boundaries. \\
Morphology:
\begin{verbatim}
EBNF:
"f.str.splitLines", "(", string, ")"
\end{verbatim}
\begin{itemize}
\item Operands: \texttt{string}.
\item Result: \texttt{list} of \texttt{string}.
\item Handles both \texttt{\\n} and \texttt{\\r\\n} line endings.
\item A \texttt{Type Error} is raised if the argument is not a \texttt{string}.
\end{itemize}
Examples:
\begin{verbatim}
f.str.splitLines("line1\nline2\nline3")  #c# Returns ["line1", "line2", "line3"]
f.str.splitLines("single")               #c# Returns ["single"]
\end{verbatim}

\subsubsection{Join Function}
Joins a list of strings with the given separator. \\
Morphology:
\begin{verbatim}
EBNF:
"f.str.join", "(", list, ",", string, ")"
\end{verbatim}
\begin{itemize}
\item Operands: \texttt{list} of \texttt{string}, separator (\texttt{string}).
\item Result: \texttt{string}.
\item A \texttt{Type Error} is raised if the list contains non-string elements.
\item A \texttt{Type Error} is raised if the separator is not a \texttt{string}.
\end{itemize}
Examples:
\begin{verbatim}
f.str.join(["a", "b", "c"], ",")    #c# Returns "a,b,c"
f.str.join(["hello", "world"], " ")  #c# Returns "hello world"
f.str.join(["a", "b"], "")           #c# Returns "ab"
\end{verbatim}

\subsubsection{Partition Function}
Splits at the first occurrence of the separator, returning a 3-element list: \texttt{[before, separator, after]}. \\
Morphology:
\begin{verbatim}
EBNF:
"f.str.partition", "(", string, ",", string, ")"
\end{verbatim}
\begin{itemize}
\item Operands: \texttt{string}, \texttt{string} (separator).
\item Result: \texttt{list} of 3 \texttt{string} elements.
\item If the separator is not found, returns \texttt{[original, "", ""]}.
\item A \texttt{Type Error} is raised if either argument is not a \texttt{string}.
\end{itemize}
Examples:
\begin{verbatim}
f.str.partition("hello world", " ")   #c# Returns ["hello", " ", "world"]
f.str.partition("abc", "x")           #c# Returns ["abc", "", ""]
\end{verbatim}

\subsubsection{Format Hex Function}
Formats an integer as a zero-padded hexadecimal string of the given width. \\
Morphology:
\begin{verbatim}
EBNF:
"f.str.formatHex", "(", int, ",", int, ")"
\end{verbatim}
\begin{itemize}
\item Operands: value (\texttt{int}), width (\texttt{int}).
\item Result: \texttt{string}.
\item Output is lowercase hex, zero-padded to the specified width.
\item A \texttt{Value Error} is raised for negative values.
\item A \texttt{Type Error} is raised if arguments are not \texttt{int}.
\end{itemize}
Examples:
\begin{verbatim}
f.str.formatHex(255, 4)    #c# Returns "00ff"
f.str.formatHex(10, 2)     #c# Returns "0a"
\end{verbatim}

\subsubsection{Format Binary Function}
Formats an integer as a zero-padded binary string of the given width. \\
Morphology:
\begin{verbatim}
EBNF:
"f.str.formatBinary", "(", int, ",", int, ")"
\end{verbatim}
\begin{itemize}
\item Operands: value (\texttt{int}), width (\texttt{int}).
\item Result: \texttt{string}.
\item Output is binary digits, zero-padded to the specified width.
\item A \texttt{Value Error} is raised for negative values.
\item A \texttt{Type Error} is raised if arguments are not \texttt{int}.
\end{itemize}
Examples:
\begin{verbatim}
f.str.formatBinary(10, 8)    #c# Returns "00001010"
f.str.formatBinary(255, 8)   #c# Returns "11111111"
\end{verbatim}

\subsubsection{Is Empty Function}
Returns \texttt{True} if the string has length 0. \\
Morphology:
\begin{verbatim}
EBNF:
"f.str.isEmpty", "(", string, ")"
\end{verbatim}
\begin{itemize}
\item Operands: \texttt{string}.
\item Result: \texttt{bool}.
\item A \texttt{Type Error} is raised if the argument is not a \texttt{string}.
\end{itemize}
Examples:
\begin{verbatim}
f.str.isEmpty("")        #c# Returns True
f.str.isEmpty("hello")   #c# Returns False
\end{verbatim}

\subsubsection{Is Blank Function}
Returns \texttt{True} if the string is empty or contains only whitespace. \\
Morphology:
\begin{verbatim}
EBNF:
"f.str.isBlank", "(", string, ")"
\end{verbatim}
\begin{itemize}
\item Operands: \texttt{string}.
\item Result: \texttt{bool}.
\item A \texttt{Type Error} is raised if the argument is not a \texttt{string}.
\end{itemize}
Examples:
\begin{verbatim}
f.str.isBlank("")        #c# Returns True
f.str.isBlank("   ")     #c# Returns True
f.str.isBlank("hello")   #c# Returns False
\end{verbatim}

\subsubsection{Is Numeric Function}
Returns \texttt{True} if the string represents any number (int, real, or complex). \\
Morphology:
\begin{verbatim}
EBNF:
"f.str.isNumeric", "(", string, ")"
\end{verbatim}
\begin{itemize}
\item Operands: \texttt{string}.
\item Result: \texttt{bool}.
\item Returns \texttt{False} for empty strings and whitespace-only strings.
\item \texttt{"1+2i"} returns \texttt{True} (complex), \texttt{"3.14"} returns \texttt{True} (real), \texttt{"42"} returns \texttt{True} (int).
\item A \texttt{Type Error} is raised if the argument is not a \texttt{string}.
\end{itemize}
Examples:
\begin{verbatim}
f.str.isNumeric("42")       #c# Returns True
f.str.isNumeric("3.14")     #c# Returns True
f.str.isNumeric("1+2i")     #c# Returns True
f.str.isNumeric("abc")      #c# Returns False
f.str.isNumeric("")         #c# Returns False
\end{verbatim}

\subsubsection{Is Integer Function}
Returns \texttt{True} if the string represents an integer. \\
Morphology:
\begin{verbatim}
EBNF:
"f.str.isInteger", "(", string, ")"
\end{verbatim}
\begin{itemize}
\item Operands: \texttt{string}.
\item Result: \texttt{bool}.
\item Returns \texttt{False} for empty strings, reals, and complex numbers.
\item A \texttt{Type Error} is raised if the argument is not a \texttt{string}.
\end{itemize}
Examples:
\begin{verbatim}
f.str.isInteger("42")       #c# Returns True
f.str.isInteger("3.14")     #c# Returns False
f.str.isInteger("1+2i")     #c# Returns False
\end{verbatim}

\subsubsection{Is Real Function}
Returns \texttt{True} if the string represents a real number (not complex). \\
Morphology:
\begin{verbatim}
EBNF:
"f.str.isReal", "(", string, ")"
\end{verbatim}
\begin{itemize}
\item Operands: \texttt{string}.
\item Result: \texttt{bool}.
\item \texttt{"3.0"} returns \texttt{True}, \texttt{"1+2i"} returns \texttt{False}.
\item Returns \texttt{False} for empty strings.
\item A \texttt{Type Error} is raised if the argument is not a \texttt{string}.
\end{itemize}
Examples:
\begin{verbatim}
f.str.isReal("3.14")        #c# Returns True
f.str.isReal("1+2i")        #c# Returns False
\end{verbatim}

\subsubsection{Is Identifier Function}
Returns \texttt{True} if the string is a valid EDADSL identifier (starts with letter or underscore, followed by letters, digits, or underscores). Keywords (e.g., \texttt{"if"}, \texttt{"return"}, \texttt{"function"}) are syntactically valid identifiers and return \texttt{True}. \\
Morphology:
\begin{verbatim}
EBNF:
"f.str.isIdentifier", "(", string, ")"
\end{verbatim}
\begin{itemize}
\item Operands: \texttt{string}.
\item Result: \texttt{bool}.
\item A \texttt{Type Error} is raised if the argument is not a \texttt{string}.
\end{itemize}
Examples:
\begin{verbatim}
f.str.isIdentifier("hello")    #c# Returns True
f.str.isIdentifier("_foo")     #c# Returns True
f.str.isIdentifier("if")       #c# Returns True  (keyword is syntactically valid)
f.str.isIdentifier("123abc")   #c# Returns False
f.str.isIdentifier("my-var")   #c# Returns False
\end{verbatim}

\subsubsection{Is Valid UTF-8 Function}
Returns \texttt{True} if the string is valid UTF-8. \\
Morphology:
\begin{verbatim}
EBNF:
"f.str.isValidUTF8", "(", string, ")"
\end{verbatim}
\begin{itemize}
\item Operands: \texttt{string}.
\item Result: \texttt{bool}.
\item All string literals in EDADSL are UTF-8; this function validates whether a byte sequence (if accessible) forms valid UTF-8.
\item A \texttt{Type Error} is raised if the argument is not a \texttt{string}.
\end{itemize}
Examples:
\begin{verbatim}
f.str.isValidUTF8("hello")     #c# Returns True
\end{verbatim}

\subsubsection{Regex Functions}
\begin{itemize}
\item A \texttt{Value Error} is raised by all four regex functions if the pattern is not a valid regular expression.
\end{itemize}

\paragraph{Regex Match}
\begin{verbatim}
EBNF:
"f.str.regexMatch", "(", string, ",", string, ")"
\end{verbatim}
Returns \texttt{True} if the string matches the regular expression pattern (full match).
\begin{verbatim}
f.str.regexMatch("hello123", "[a-z]+[0-9]+")  #c# Returns True
f.str.regexMatch("hello", "[0-9]+")            #c# Returns False
\end{verbatim}

\paragraph{Regex Find}
\begin{verbatim}
EBNF:
"f.str.regexFind", "(", string, ",", string, ")"
\end{verbatim}
Returns the first match as a string, or \texttt{null} if no match.
\begin{verbatim}
f.str.regexFind("abc123def", "[0-9]+")   #c# Returns "123"
f.str.regexFind("hello", "[0-9]+")        #c# Returns null
\end{verbatim}

\paragraph{Regex FindAll}
\begin{verbatim}
EBNF:
"f.str.regexFindAll", "(", string, ",", string, ")"
\end{verbatim}
Returns a list of all non-overlapping matches.
\begin{verbatim}
f.str.regexFindAll("a1b2c3", "[0-9]+")  #c# Returns ["1", "2", "3"]
f.str.regexFindAll("hello", "[0-9]+")    #c# Returns []
\end{verbatim}

\paragraph{Regex Replace}
\begin{verbatim}
EBNF:
"f.str.regexReplace", "(", string, ",", string, ",", string, ")"
\end{verbatim}
Replaces all matches of the pattern with the replacement string.
\begin{verbatim}
f.str.regexReplace("abc123def456", "[0-9]+", "X")  #c# Returns "abcXdefX"
f.str.regexReplace("hello", "[0-9]+", "X")          #c# Returns "hello"
\end{verbatim}

\subsubsection{Encoding Functions}
\paragraph{Bytes Conversion}
\begin{verbatim}
EBNF:
"f.str.toBytes", "(", string, ")"
\end{verbatim}
\begin{itemize}
\item Converts a string to a list of integer byte values using UTF-8 encoding. Multi-byte characters produce multiple byte values (e.g. \texttt{f.str.toBytes("€")} returns \texttt{[0xE2, 0x82, 0xAC]}).
\end{itemize}
Examples:
\begin{verbatim}
f.str.toBytes("AB")        #c# Returns [65, 66]
f.str.toBytes("€")         #c# Returns [0xE2, 0x82, 0xAC]
\end{verbatim}

\paragraph{FromBytes}
\begin{verbatim}
EBNF:
"f.str.fromBytes", "(", list, ")"
\end{verbatim}
\begin{itemize}
\item Converts a list of integer byte values (0--255) back to a string via UTF-8 decoding. A \texttt{Value Error} is raised if any byte value is outside 0--255 or if the resulting bytes do not form valid UTF-8.
\end{itemize}
\begin{verbatim}
f.str.fromBytes([65, 66])  #c# Returns "AB"
f.str.fromBytes([0xE2, 0x82, 0xAC])  #c# Returns "€"
\end{verbatim}

\paragraph{Base64}
\begin{verbatim}
EBNF:
("f.str.base64Encode" | "f.str.base64Decode"), "(", string, ")"
\end{verbatim}
\begin{itemize}
\item A \texttt{Value Error} is raised by \texttt{f.base64Decode} if the input is not valid base64.
\end{itemize}
\begin{verbatim}
f.str.base64Encode("hello")       #c# Returns "aGVsbG8="
f.str.base64Decode("aGVsbG8=")   #c# Returns "hello"
\end{verbatim}

\paragraph{Hex Encode/Decode}
\begin{verbatim}
EBNF:
("f.str.hexEncode" | "f.str.hexDecode"), "(", string, ")"
\end{verbatim}
\begin{itemize}
\item \texttt{f.str.hexEncode(string)} -- converts each byte to its two-digit hex representation
\item \texttt{f.str.hexDecode(string)} -- converts hex string back to bytes/string. A \texttt{Value Error} is raised for odd-length strings or invalid hex characters.
\end{itemize}
\begin{verbatim}
f.str.hexEncode("AB")       #c# Returns "4142"
f.str.hexDecode("4142")     #c# Returns "AB"
\end{verbatim}

\subsubsection{Identifier Functions}
\paragraph{Case Checks}
\begin{verbatim}
EBNF:
("f.str.isUpper" | "f.str.isLower" | "f.str.isAlpha" | "f.str.isDigit" | "f.str.isAlphanumeric"),
"(", string, ")"
\end{verbatim}
\begin{itemize}
\item \texttt{f.isUpper} -- \texttt{True} if all cased characters are uppercase
\item \texttt{f.isLower} -- \texttt{True} if all cased characters are lowercase
\item \texttt{f.isAlpha} -- \texttt{True} if all characters are letters
\item \texttt{f.isDigit} -- \texttt{True} if all characters are digits
\item \texttt{f.isAlphanumeric} -- \texttt{True} if all characters are letters or digits
\item All five functions return \texttt{True} for empty strings (vacuously true).
\end{itemize}
\begin{verbatim}
f.str.isUpper("HELLO")       #c# Returns True
f.str.isUpper("Hello")       #c# Returns False
f.str.isUpper("")            #c# Returns True  (empty string)
f.str.isLower("hello")       #c# Returns True
f.str.isAlpha("abc")         #c# Returns True
f.str.isAlpha("abc123")      #c# Returns False
f.str.isDigit("123")         #c# Returns True
f.str.isDigit("12.3")        #c# Returns False
f.str.isAlphanumeric("abc123") #c# Returns True
\end{verbatim}

\paragraph{Sanitize Identifier}
\begin{verbatim}
EBNF:
"f.str.sanitizeIdentifier", "(", string, ")"
\end{verbatim}
Converts a string to a valid identifier by replacing invalid characters with underscores and collapsing consecutive underscores.
\begin{itemize}
\item If the result starts with a digit, an underscore is prepended.
\end{itemize}
\begin{verbatim}
f.str.sanitizeIdentifier("my-var")      #c# Returns "my_var"
f.str.sanitizeIdentifier("  spaces  ")  #c# Returns "spaces"
f.str.sanitizeIdentifier("a!!b")        #c# Returns "a_b"
f.str.sanitizeIdentifier("123abc")      #c# Returns "_123abc"
\end{verbatim}


\subsection{List and Set Functions}

\subsubsection{Concatenation}
The \texttt{concat} function combines all arguments into a single list. \\
Morphology:
\begin{verbatim}
EBNF:
"f.list.concat", "(", { argument }, ")"
\end{verbatim}
Example:
\begin{verbatim}
f.list.concat([1, 2], [3, 4])        #c# Returns [1, 2, 3, 4]
f.list.concat("a", "b", "c")         #c# Returns ["a", "b", "c"]
\end{verbatim}

\subsubsection{Push}
The \texttt{push} function appends an element to the end of a list.\\
Morphology:
\begin{verbatim}
EBNF:
"f.list.push", "(", list, ",", data, ")"
\end{verbatim}
\begin{itemize}
\item Operands: list, any data type matching the list element type
\item Returns the modified list (the original list is also mutated)
\end{itemize}
Examples:
\begin{verbatim}
f.list.push([1, 2, 3], 4)        #c# Returns [1, 2, 3, 4]
$f = [10, 20]
f.list.push($f, 30)              #c# $f is now [10, 20, 30]
\end{verbatim}

\subsubsection{Pop}
The \texttt{pop} function removes and returns the last element of a list.\\
Morphology:
\begin{verbatim}
EBNF:
"f.list.pop", "(", list, ")"
\end{verbatim}
\begin{itemize}
\item Operands: list (must not be empty)
\item Returns the removed element
\item A Value Error is raised if the list is empty
\end{itemize}
Examples:
\begin{verbatim}
f.list.pop([1, 2, 3])             #c# Returns 3, list is now [1, 2]
$f = [10, 20, 30]
f.list.pop($f)                    #c# Returns 30, $f is now [10, 20]
\end{verbatim}

\subsubsection{Head}
The \texttt{head} function returns the first \texttt{n} elements of a list.\\
Morphology:
\begin{verbatim}
EBNF:
"f.list.head", "(", list, ",", int, ")"
\end{verbatim}
\begin{itemize}
\item Operands: list, int (n $\geq$ 0)
\item Returns a new list containing the first \texttt{n} elements
\item If \texttt{n} exceeds the list length, returns the entire list (no error)
\item A \texttt{Value Error} is raised if \texttt{n $<$ 0}.
\end{itemize}
Examples:
\begin{verbatim}
f.list.head([1, 2, 3, 4, 5], 3)  #c# Returns [1, 2, 3]
f.list.head([1, 2, 3], 10)       #c# Returns [1, 2, 3]
f.list.head([1, 2, 3], 0)        #c# Returns []
\end{verbatim}

\subsubsection{Tail}
The \texttt{tail} function returns the last \texttt{n} elements of a list.\\
Morphology:
\begin{verbatim}
EBNF:
"f.list.tail", "(", list, ",", int, ")"
\end{verbatim}
\begin{itemize}
\item Operands: list, int (n $\geq$ 0)
\item Returns a new list containing the last \texttt{n} elements
\item If \texttt{n} exceeds the list length, returns the entire list (no error)
\item A \texttt{Value Error} is raised if \texttt{n $<$ 0}.
\end{itemize}
Examples:
\begin{verbatim}
f.list.tail([1, 2, 3, 4, 5], 3)  #c# Returns [3, 4, 5]
f.list.tail([1, 2, 3], 10)       #c# Returns [1, 2, 3]
f.list.tail([1, 2, 3], 0)        #c# Returns []
\end{verbatim}

\subsubsection{Type Conversion}
\paragraph{List to Set}
The \texttt{set} function converts a list to a set. All unique elements are preserved; duplicate elements are removed and order is lost. \\
Morphism:
\begin{verbatim}
EBNF:
"f.list.set", "(", list, ")"
\end{verbatim}
Example:
\begin{verbatim}
f.list.set([Q1, Q2, R5])   #c# Returns {Q1, Q2, R5}
\end{verbatim}

\paragraph{Set to List}
The \texttt{list} function converts a set to a list. The order is determined by the lexicographic order of the canonical string representations of the elements. \\
Morphology:
\begin{verbatim}
EBNF:
"f.list.fromSet", "(", set, ")"
\end{verbatim}
Example:
\begin{verbatim}
f.list.fromSet({Q1, Q2})   #c# Returns [Q1, Q2] (order from *e)
\end{verbatim}

\subsubsection{Copy and Clone}
\paragraph{Shallow Copy}
The \texttt{copy} function creates a shallow copy of a value. For lists, a new list is created containing the same elements (but nested structures are shared by reference). For scalars (int, real, complex, bool, string), the value is returned unchanged (scalars are immutable).\\
Morphology:
\begin{verbatim}
EBNF:
"f.list.copy", "(", data, ")"
\end{verbatim}
\begin{itemize}
\item Operands: any
\item Result: shallow copy of the input
\end{itemize}
Examples:
\begin{verbatim}
$a = [1, 2, 3]
$b = f.list.copy($a)       #c# $b is a new list [1, 2, 3]
$b[0] = 99            #c# $a is still [1, 2, 3]
\end{verbatim}

\paragraph{Deep Clone}
The \texttt{clone} function creates a deep copy of a value. All nested structures are recursively copied, so modifications to the clone never affect the original. For scalars (int, real, complex, bool, string), the value is returned unchanged (scalars are immutable).\\
Morphology:
\begin{verbatim}
EBNF:
"f.list.clone", "(", data, ")"
\end{verbatim}
\begin{itemize}
\item Operands: any
\item Result: deep copy of the input
\end{itemize}
Examples:
\begin{verbatim}
$a = [[1, 2], [3, 4]]
$b = f.list.clone($a)      #c# $b is a fully independent copy
$b[0][0] = 99         #c# $a is still [[1, 2], [3, 4]]
\end{verbatim}

\subsection{Ndarray Functions}
Functions for operating on ndarrays (and lists). These functions extend the built-in math and list functions with array-aware behavior.

\subsubsection{Array Creation}

\paragraph{Zeros}
The \texttt{zeros} function creates an array filled with zeros.\\
Morphology:
\begin{verbatim}
EBNF:
"f.arr.zeros", "(", shape, ")"
\end{verbatim}
\begin{itemize}
\item Operands: int (length for 1D), or list of ints (dimensions for nD)
\item Result: ndarray filled with 0
\item A \texttt{Value Error} is raised for negative dimensions.
\end{itemize}
Examples:
\begin{verbatim}
f.arr.zeros(5)           #c# Returns [0, 0, 0, 0, 0]
f.arr.zeros([2, 3])      #c# Returns [[0,0,0],[0,0,0]]
\end{verbatim}

\paragraph{Ones}
The \texttt{ones} function creates an array filled with ones.\\
Morphology:
\begin{verbatim}
EBNF:
"f.arr.ones", "(", shape, ")"
\end{verbatim}
\begin{itemize}
\item Operands: int (length for 1D), or list of ints (dimensions for nD)
\item Result: ndarray filled with 1
\item A \texttt{Value Error} is raised for negative dimensions.
\end{itemize}
Examples:
\begin{verbatim}
f.arr.ones(4)            #c# Returns [1, 1, 1, 1]
f.arr.ones([2, 2])       #c# Returns [[1,1],[1,1]]
\end{verbatim}

\paragraph{Identity Matrix}
The \texttt{identity} function creates an $n \times n$ identity matrix.\\
Morphology:
\begin{verbatim}
EBNF:
"f.arr.identity", "(", int, ")"
\end{verbatim}
\begin{itemize}
\item Operands: int (n $\geq$ 0)
\item Result: ndarray (n$\times$n) with 1s on the diagonal and 0s elsewhere
\item \texttt{f.arr.identity(0)} returns a 0$\times$0 matrix \texttt{[[]]}.
\end{itemize}
Examples:
\begin{verbatim}
f.arr.identity(1)        #c# Returns [[1]]
f.arr.identity(3)        #c# Returns [[1,0,0],[0,1,0],[0,0,1]]
\end{verbatim}

\paragraph{Edge Cases -- Reduction Functions}
\begin{itemize}
\item \texttt{f.stat.sum([])} returns \texttt{0} (additive identity).
\item \texttt{f.stat.prod([])} returns \texttt{1} (multiplicative identity).
\item \texttt{f.stat.mean([])}, \texttt{f.stat.std([])}, \texttt{f.stat.var([])}, \texttt{f.stat.median([])} raise a \texttt{Value Error} on empty lists.
\item \texttt{f.util.sort([])}, \texttt{f.util.reverse([])}, \texttt{f.util.unique([])} return \texttt{[]} on empty lists.
\item A \texttt{Type Error} is raised if the argument is not a list.
\end{itemize}

\subsubsection{Reduction Functions}
Reduction functions collapse an ndarray along an axis. When \texttt{axis} is omitted, the operation is applied to the flattened array (all elements). When \texttt{axis = 0}, the operation is applied row-wise (collapses rows). When \texttt{axis = 1}, the operation is applied column-wise (collapses columns).

\paragraph{Sum}
\begin{verbatim}
EBNF:
"f.stat.sum", "(", argument, [ ",", "axis", "=", int ], ")"
\end{verbatim}
\begin{verbatim}
f.stat.sum([1, 2, 3])              #c# Returns 6
f.stat.sum([[1,2],[3,4]])          #c# Returns 10 (flattened)
f.stat.sum([[1,2],[3,4]], axis=0)  #c# Returns [4, 6]  (row-wise)
f.stat.sum([[1,2],[3,4]], axis=1)  #c# Returns [3, 7]  (column-wise)
\end{verbatim}

\paragraph{Product}
\begin{verbatim}
EBNF:
"f.stat.prod", "(", argument, [ ",", "axis", "=", int ], ")"
\end{verbatim}
\begin{verbatim}
f.stat.prod([2, 3, 4])              #c# Returns 24
f.stat.prod([[2,3],[4,5]], axis=0)  #c# Returns [8, 15]
f.stat.prod([[2,3],[4,5]], axis=1)  #c# Returns [6, 20]
\end{verbatim}

\paragraph{Mean}
\begin{verbatim}
EBNF:
"f.stat.mean", "(", argument, [ ",", "axis", "=", int ], ")"
\end{verbatim}
\begin{verbatim}
f.stat.mean([1, 2, 3, 4])            #c# Returns 2.5
f.stat.mean([[1,2],[3,4]], axis=0)   #c# Returns [2.0, 3.0]
f.stat.mean([[1,2],[3,4]], axis=1)   #c# Returns [1.5, 3.5]
\end{verbatim}

\paragraph{Standard Deviation}
\begin{verbatim}
EBNF:
"f.stat.std", "(", argument, [ ",", "axis", "=", int ], ")"
\end{verbatim}
\begin{verbatim}
f.stat.std([1, 2, 3, 4, 5])         #c# Returns ~1.4142
f.stat.std([1, 1, 1, 1])            #c# Returns 0.0
\end{verbatim}

\paragraph{Variance}
Computes the population variance (divides by $N$, not $N-1$).
\begin{verbatim}
EBNF:
"f.stat.var", "(", argument, [ ",", "axis", "=", int ], ")"
\end{verbatim}
\begin{verbatim}
f.stat.var([1, 2, 3, 4, 5])         #c# Returns 2.0
f.stat.var([1, 1, 1, 1])            #c# Returns 0.0
\end{verbatim}

\paragraph{Median}
\begin{verbatim}
EBNF:
"f.stat.median", "(", argument, ")"
\end{verbatim}
\begin{verbatim}
f.stat.median([3, 1, 4, 1, 5])   #c# Returns 3
f.stat.median([1, 2, 3, 4])      #c# Returns 2.5 (average of middle two)
\end{verbatim}

\paragraph{Percentile}
\begin{verbatim}
EBNF:
"f.stat.percentile", "(", argument, ",", real, ")"
\end{verbatim}
The second argument is a percentage (0--100).

\paragraph{Edge Cases -- Percentile}
\begin{itemize}
\item A \texttt{Value Error} is raised if the list is empty.
\item A \texttt{Value Error} is raised if the percentage is outside 0--100.
\end{itemize}
\begin{verbatim}
f.stat.percentile([10, 20, 30, 40, 50], 50)  #c# Returns 30
f.stat.percentile([10, 20, 30, 40, 50], 0)   #c# Returns 10
f.stat.percentile([10, 20, 30, 40, 50], 100) #c# Returns 50
\end{verbatim}

\paragraph{Mode}
Returns the most frequently occurring value(s) in a list. If multiple values tie for the highest frequency, returns a sorted list of all modes.
\begin{verbatim}
EBNF:
"f.stat.mode", "(", list, ")"
\end{verbatim}

\paragraph{Edge Cases -- Mode}
\begin{itemize}
\item A \texttt{Value Error} is raised if the list is empty.
\item A \texttt{Type Error} is raised if the argument is not a list.
\item If all values occur equally often, returns a sorted list of all unique values.
\end{itemize}
\begin{verbatim}
f.stat.mode([1, 2, 2, 3, 3, 3])     #c# Returns 3 (single mode)
f.stat.mode([1, 1, 2, 2, 3])        #c# Returns [1, 2] (bimodal)
f.stat.mode([1, 2, 3])              #c# Returns [1, 2, 3] (all equally frequent)
\end{verbatim}

\paragraph{Quantile}
Returns the value at quantile $q$ (in $[0, 1]$). Equivalent to \texttt{f.stat.percentile(x, q * 100)}.
\begin{verbatim}
EBNF:
"f.stat.quantile", "(", list, ",", real, ")"
\end{verbatim}

\paragraph{Edge Cases -- Quantile}
\begin{itemize}
\item A \texttt{Value Error} is raised if the list is empty.
\item A \texttt{Value Error} is raised if the quantile is outside 0--1.
\item A \texttt{Type Error} is raised if the first argument is not a list.
\end{itemize}
\begin{verbatim}
f.stat.quantile([10, 20, 30, 40, 50], 0.5)  #c# Returns 30
f.stat.quantile([10, 20, 30, 40, 50], 0.0)  #c# Returns 10
f.stat.quantile([10, 20, 30, 40, 50], 1.0)  #c# Returns 50
f.stat.quantile([10, 20, 30, 40, 50], 0.25) #c# Returns 20
\end{verbatim}

\paragraph{Root Mean Square}
\begin{verbatim}
EBNF:
"f.stat.rms", "(", argument, ")"
\end{verbatim}
Computes $\sqrt{\frac{1}{n}\sum_{i=1}^{n} x_i^2}$. Useful for AC signal analysis and power calculations.
\begin{itemize}
\item A \texttt{Value Error} is raised if the list is empty.
\end{itemize}
\begin{verbatim}
f.stat.rms([1, -1, 1, -1])        #c# Returns 1.0
f.stat.rms([3, 4])                 #c# Returns 3.5355...
f.stat.rms([0, 0, 0])             #c# Returns 0.0
\end{verbatim}

\paragraph{Peak}
\begin{verbatim}
EBNF:
"f.stat.peak", "(", argument, [ ",", "mode", "=", string ], ")"
\end{verbatim}
Modes:
\begin{itemize}
\item \texttt{"p2p"} (default) -- Peak-to-peak: $\max(x) - \min(x)$
\item \texttt{"abs"} -- Maximum absolute value: $\max(|x|)$
\item \texttt{"pos"} -- Maximum positive value: $\max(x)$
\item \texttt{"neg"} -- Minimum negative value: $\min(x)$
\end{itemize}
\begin{verbatim}
f.stat.peak([1, -3, 2, -5, 4])                  #c# Returns 9 (p2p: 4 - (-5))
f.stat.peak([1, -3, 2, -5, 4], mode="abs")      #c# Returns 5
f.stat.peak([1, -3, 2, -5, 4], mode="pos")      #c# Returns 4
f.stat.peak([1, -3, 2, -5, 4], mode="neg")      #c# Returns -5
\end{verbatim}

\paragraph{Edge Cases -- Peak}
\begin{itemize}
\item A \texttt{Value Error} is raised if the list is empty.
\item A \texttt{Value Error} is raised if an invalid mode is specified.
\end{itemize}

\subsubsection{Search Functions}

\paragraph{Edge Cases -- Search Functions}
\begin{itemize}
\item A \texttt{Value Error} is raised if the list is empty.
\end{itemize}

\paragraph{Index of Minimum}
\begin{verbatim}
EBNF:
"f.stat.argmin", "(", argument, ")"
\end{verbatim}
\begin{verbatim}
f.stat.argmin([5, 3, 1, 4])   #c# Returns 2 (index of 1)
\end{verbatim}

\paragraph{Index of Maximum}
\begin{verbatim}
EBNF:
"f.stat.argmax", "(", argument, ")"
\end{verbatim}
\begin{verbatim}
f.stat.argmax([5, 3, 1, 4])   #c# Returns 0 (index of 5)
\end{verbatim}

\subsubsection{Shape Functions}
\paragraph{Flatten}
\begin{verbatim}
EBNF:
"f.arr.flatten", "(", argument, ")"
\end{verbatim}
Flattens one level of nesting. Already-flat lists are returned unchanged. Empty lists return \texttt{[]}.
\begin{verbatim}
f.arr.flatten([[1,2],[3,4],[5,6]])   #c# Returns [1, 2, 3, 4, 5, 6]
f.arr.flatten([[1,2,3]])             #c# Returns [1, 2, 3]
f.arr.flatten([1, 2, 3])            #c# Returns [1, 2, 3]  (already flat)
f.arr.flatten([])                   #c# Returns []
\end{verbatim}

\paragraph{Sort}
\begin{verbatim}
EBNF:
"f.util.sort", "(", argument, [ ",", lambda ], ")"
\end{verbatim}
Returns a new sorted list (does not modify the original). An optional lambda comparator can be provided for custom ordering (see Higher-Order Functions section).
\begin{verbatim}
f.util.sort([5, 3, 1, 4])           #c# Returns [1, 3, 4, 5]
f.util.sort([3.1, 2.7, 1.4])       #c# Returns [1.4, 2.7, 3.1]
f.util.sort([5, 1, 3], L::($a [int], $b [int])->{ $b - $a })
#c# Returns [5, 3, 1]  (descending)
\end{verbatim}

\paragraph{Reverse}
\begin{verbatim}
EBNF:
"f.util.reverse", "(", argument, ")"
\end{verbatim}
Returns a new reversed list (does not modify the original).
\begin{verbatim}
f.util.reverse([1, 2, 3])           #c# Returns [3, 2, 1]
f.util.reverse([[1,2],[3,4]])       #c# Returns [[3,4],[1,2]]
\end{verbatim}

\paragraph{Unique}
\begin{verbatim}
EBNF:
"f.util.unique", "(", argument, ")"
\end{verbatim}
\begin{verbatim}
f.util.unique([1, 2, 2, 3, 3, 3])   #c# Returns [1, 2, 3]
\end{verbatim}

\subsubsection{Linear Algebra Functions}

\paragraph{Edge Cases -- Linear Algebra Functions}
\begin{itemize}
\item A \texttt{Type Error} is raised if the argument is not an \texttt{ndarray} (matrix).
\item A \texttt{Value Error} is raised for \texttt{f.linalg.det}, \texttt{f.linalg.inv}, \texttt{f.linalg.eig}, \texttt{f.linalg.eigvec}, \texttt{f.linalg.lu}, \texttt{f.linalg.qr}, \texttt{f.linalg.trace}, \texttt{f.linalg.cond} if the matrix is not square.
\item A \texttt{Value Error} is raised for \texttt{f.linalg.cond} if the matrix is singular.
\item A \texttt{Value Error} is raised for \texttt{f.arr.normalize} if the vector has zero magnitude.
\end{itemize}

\paragraph{Determinant}
\begin{verbatim}
EBNF:
"f.linalg.det", "(", argument, ")"
\end{verbatim}
\begin{verbatim}
f.linalg.det([[1,2],[3,4]])   #c# Returns -2.0
f.linalg.det([[2,0],[0,3]])   #c# Returns 6.0
\end{verbatim}

\paragraph{Matrix Inverse}
\begin{verbatim}
EBNF:
"f.linalg.inv", "(", argument, ")"
\end{verbatim}
Returns the inverse matrix as a list of lists. A Division by Zero Error is raised if the matrix is singular.
\begin{verbatim}
f.linalg.inv([[1,2],[3,4]])   #c# Returns [[-2.0, 1.0],[1.5, -0.5]]
\end{verbatim}

\paragraph{Eigenvalues}
\begin{verbatim}
EBNF:
"f.linalg.eig", "(", argument, ")"
\end{verbatim}
Returns a list of eigenvalues.
\begin{verbatim}
f.linalg.eig([[2,0],[0,3]])   #c# Returns [2, 3]
f.linalg.eig([[1,2],[2,1]])   #c# Returns [3, -1]
\end{verbatim}

\paragraph{Eigenvectors}
\begin{verbatim}
EBNF:
"f.linalg.eigvec", "(", argument, ")"
\end{verbatim}
Returns a list of eigenvectors (one per eigenvalue, as rows).
\begin{verbatim}
f.linalg.eigvec([[2,0],[0,3]])   #c# Returns [[1,0],[0,1]]
\end{verbatim}

\paragraph{Singular Value Decomposition}
\begin{verbatim}
EBNF:
"f.linalg.svd", "(", argument, ")"
\end{verbatim}
Returns a list of 3 lists: \texttt{[U, S, V]} where \texttt{U} and \texttt{V} are orthogonal matrices and \texttt{S} is the list of singular values. Supports non-square matrices: for an $m \times n$ input, \texttt{U} is $m \times m$, \texttt{S} has $\min(m, n)$ values, and \texttt{V} is $n \times n$.
\begin{verbatim}
f.linalg.svd([[1,0],[0,2]])
    #c# Returns [[[1,0],[0,1]], [1, 2], [[1,0],[0,1]]]
\end{verbatim}

\paragraph{LU Decomposition}
\begin{verbatim}
EBNF:
"f.linalg.lu", "(", argument, ")"
\end{verbatim}
Returns a list of 2 lists: \texttt{[L, U]} where \texttt{L} is lower triangular and \texttt{U} is upper triangular. Supports non-square matrices.
\begin{verbatim}
f.linalg.lu([[2,1],[4,3]])   #c# Returns L and U such that L .* U = original
\end{verbatim}

\paragraph{QR Decomposition}
\begin{verbatim}
EBNF:
"f.linalg.qr", "(", argument, ")"
\end{verbatim}
Returns a list of 2 lists: \texttt{[Q, R]} where \texttt{Q} is orthogonal and \texttt{R} is upper triangular. Supports non-square matrices.
\begin{verbatim}
f.linalg.qr([[1,1],[1,-1]])   #c# Returns Q and R such that Q .* R = original
\end{verbatim}

\paragraph{Cholesky Decomposition}
\begin{verbatim}
EBNF:
"f.linalg.chol", "(", argument, ")"
\end{verbatim}
Returns the lower triangular matrix \texttt{L} such that \texttt{L .* L.T = original}. \\

\paragraph{Edge Cases}
\begin{itemize}
\item A \texttt{Type Error} is raised if the argument is not a list of lists (ndarray).
\item A \texttt{Value Error} is raised if the matrix is not square.
\item A \texttt{Value Error} is raised if the matrix is not symmetric.
\item A \texttt{Value Error} is raised if the matrix is not positive definite.
\end{itemize}

\begin{verbatim}
f.linalg.chol([[4,2],[2,3]])   #c# Returns [[2,0],[1,1.7321]]
f.linalg.chol([[1,2],[3,4]])   #c# Value Error (not symmetric)
f.linalg.chol([[1,2,3]])       #c# Value Error (not square)
\end{verbatim}

\paragraph{Trace}
\begin{verbatim}
EBNF:
"f.linalg.trace", "(", argument, ")"
\end{verbatim}
\begin{verbatim}
f.linalg.trace([[1,2],[3,4]])   #c# Returns 5
f.linalg.trace([[2,0],[0,3]])   #c# Returns 5
\end{verbatim}

\paragraph{Matrix Rank}
\begin{verbatim}
EBNF:
"f.linalg.rank", "(", argument, ")"
\end{verbatim}
\begin{verbatim}
f.linalg.rank([[1,2],[3,4]])   #c# Returns 2
f.linalg.rank([[1,2],[2,4]])   #c# Returns 1 (rows are linearly dependent)
\end{verbatim}

\paragraph{Condition Number}
\begin{verbatim}
EBNF:
"f.linalg.cond", "(", argument, ")"
\end{verbatim}
Returns the ratio of the largest to smallest singular value. A high condition number indicates numerical instability.
\begin{verbatim}
f.linalg.cond([[1,0],[0,1]])   #c# Returns 1.0 (well-conditioned)
f.linalg.cond([[1,0],[0,1e-10]]) #c# Returns 1e10 (ill-conditioned)
\end{verbatim}

\paragraph{Pseudo-Inverse}
\begin{verbatim}
EBNF:
"f.linalg.pinv", "(", argument, ")"
\end{verbatim}
Returns the Moore--Penrose pseudo-inverse. Works for non-square or singular matrices.
\begin{verbatim}
f.linalg.pinv([[1,2],[3,4]])   #c# Returns [[-2, 1],[1.5, -0.5]]
\end{verbatim}

\subsubsection{Statistical Functions}
\paragraph{Histogram}
\begin{verbatim}
EBNF:
"f.stat.histogram", "(", argument, ",", int, ")"
\end{verbatim}
Returns a list of 2 lists: \texttt{[edges, counts]} where \texttt{edges} has \texttt{n + 1} values and \texttt{counts} has \texttt{n} values.

\paragraph{Edge Cases -- Histogram}
\begin{itemize}
\item A \texttt{Type Error} is raised if the first argument is not a list.
\item A \texttt{Value Error} is raised if \texttt{n <= 0}.
\item A \texttt{Value Error} is raised if the list is empty.
\end{itemize}

\begin{verbatim}
f.stat.histogram([1, 2, 3, 4, 5], 3)
#c# Returns [[1.0, 2.333, 3.667, 5.0], [2, 1, 2]]
\end{verbatim}

\subsubsection{Linear System Solver}
The \texttt{solve} function solves the linear system $A\mathbf{x} = \mathbf{b}$ and returns $\mathbf{x}$.\\
Morphology:
\begin{verbatim}
EBNF:
"f.linalg.solve", "(", ndarray, ",", ndarray, ")"
\end{verbatim}
\begin{itemize}
\item Operands: A (ndarray, square coefficient matrix), b (ndarray, right-hand side vector or matrix)
\item Result: ndarray $\mathbf{x}$ such that $A\mathbf{x} = \mathbf{b}$
\item A must be square and non-singular (determinant $\neq$ 0)
\item A Division by Zero Error is raised if A is singular
\item A Value Error is raised if A and b have incompatible dimensions
\end{itemize}
Examples:
\begin{verbatim}
f.linalg.solve([[2,1],[5,7]], [11,13])     #c# Returns [7.333, -3.667]
f.linalg.solve([[1,0],[0,1]], [3,4])       #c# Returns [3, 4]
\end{verbatim}


\subsection{Higher-Order Functions}
Higher-order functions take function values (lambdas or function references) as arguments. They operate on lists and ndarrays. See Section~\ref{sec:lambda} for lambda syntax and Section~\ref{ch:types} for the function data type.

\paragraph{Edge Cases -- Higher-Order Functions}
\begin{itemize}
\item A \texttt{Type Error} is raised if the second argument is not a function value (lambda or function reference).
\item A \texttt{Type Error} is raised if the first argument is not a list.
\item A \texttt{Runtime Error} is raised if the function value is \texttt{null}.
\item \texttt{f.util.reduce} raises a \texttt{Value Error} if the list is empty and no initial value is provided (third argument).
\end{itemize}

\subsubsection{Map}
The \texttt{map} function applies a lambda to each element of a list and returns a new list with the results. The input list is not modified.

\paragraph{Morphology}
\begin{verbatim}
EBNF:
"f.util.map", "(", list, ",", lambda, ")"
\end{verbatim}

\paragraph{Type Overloading}
\begin{itemize}
\item \texttt{map(list[T], L::(\$x [T])->\{...\})} $\rightarrow$ \texttt{list[U]}
  (T and U can be different types)
\item Works on any list: flat, nested, mixed types
\item The lambda parameter type must match the element type
\end{itemize}

\paragraph{Examples}
\begin{verbatim}
f.util.map([1, 2, 3], L::($x [int])->{ $x * 2 })
#c# Returns [2, 4, 6]

f.util.map([1, 2, 3], L::($x [int])->{ $x > 2 })
#c# Returns [False, False, True]  (type change: int -> bool)

f.util.map(["a", "bb", "ccc"], L::($x [string])->{ f.util.len($x) })
#c# Returns [1, 2, 3]  (type change: string -> int)

f.util.map([1, 2, 3], L::($x [int])->{
    $y [int]= $x * 10
    $y + 1
})
#c# Returns [11, 21, 31]  (multi-line lambda)

f.util.map([], L::($x [int])->{ $x * 2 })
#c# Returns []  (empty list)
\end{verbatim}

\paragraph{Error Conditions}
\begin{itemize}
\item Type mismatch: lambda parameter type does not match element type raises a Runtime Error
\end{itemize}

\subsubsection{Filter}
The \texttt{filter} function returns a new list containing only elements for which the lambda returns \texttt{True}. The input list is not modified.

\paragraph{Morphology}
\begin{verbatim}
EBNF:
"f.util.filter", "(", list, ",", lambda, ")"
\end{verbatim}

\paragraph{Type Overloading}
\begin{itemize}
\item \texttt{filter(list[T], L::(\$x [T])->\{bool\})} $\rightarrow$ \texttt{list[T]}
  (return type is the same as input)
\item The lambda must return \texttt{bool}
\end{itemize}

\paragraph{Examples}
\begin{verbatim}
f.util.filter([1, 2, 3, 4, 5], L::($x [int])->{ $x % 2 == 0 })
#c# Returns [2, 4]

f.util.filter([1, 2, 3, 4, 5], L::($x [int])->{ $x > 10 })
#c# Returns []  (no matches)

f.util.filter(["a", "bb", "ccc", "dd"], L::($x [string])->{ f.util.len($x) > 2 })
#c# Returns ["ccc"]

f.util.filter([-1, -2, 3, -4, 5], L::($x [int])->{ $x > 0 })
#c# Returns [3, 5]

f.util.filter([1, 2, 3], L::($x [int])->{ $x > 5 })
#c# Returns []  (empty result)

f.util.filter([], L::($x [int])->{ $x > 0 })
#c# Returns []  (empty input)
\end{verbatim}

\paragraph{Error Conditions}
\begin{itemize}
\item Type mismatch: lambda parameter type does not match element type raises a Runtime Error
\item Lambda does not return \texttt{bool} raises a Type Error
\end{itemize}

\subsubsection{Reduce}
The \texttt{reduce} function folds a list to a single value by applying a binary lambda cumulatively. The lambda receives an accumulator and the current element. The \texttt{init} value is always required.

\paragraph{Morphology}
\begin{verbatim}
EBNF:
"f.util.reduce", "(", list, ",", lambda, ",", expression, ")"
\end{verbatim}

\paragraph{Type Overloading}
\begin{itemize}
\item \texttt{reduce(list[T], L::(\$acc [U], \$x [T])->\{U\}, init[U])} $\rightarrow$ \texttt{U}
  (accumulator type U can differ from element type T)
\item The \texttt{init} value type must match the accumulator type
\end{itemize}

\paragraph{Examples}
\begin{verbatim}
#c# Sum integers
f.util.reduce([1, 2, 3, 4], L::($acc [int], $x [int])->{ $acc + $x }, 0)
#c# Returns 10

#c# Product integers
f.util.reduce([2, 3, 4], L::($acc [int], $x [int])->{ $acc * $x }, 1)
#c# Returns 24

#c# Find maximum
f.util.reduce([3, 1, 4, 1, 5], L::($acc [int], $x [int])->{ 
    ($acc > $x)?->($acc)!($x) 
}, 0)
#c# Returns 5

#c# Reduce with string concatenation
f.util.reduce(["a", "b", "c"],
    L::($acc [string], $x [string])->{ $acc + $x }, "")
#c# Returns "abc"

#c# Reduce on empty list returns init
f.util.reduce([], L::($acc [int], $x [int])->{ $acc + $x }, 42)
#c# Returns 42

#c# Flatten with reduce (using concat)
f.util.reduce([[1,2],[3,4]],
    L::($acc [list], $x [list])->{ f.list.concat($acc, $x) }, [])
#c# Returns [1, 2, 3, 4]
\end{verbatim}

\paragraph{Error Conditions}
\begin{itemize}
\item Type mismatch: lambda parameter types do not match accumulator/element types raises a Runtime Error
\item Lambda does not return a type coercible to the accumulator type raises a Type Error
\end{itemize}

\subsubsection{Where}
The \texttt{where} function returns a list of \textbf{indices} (ints) for which the lambda returns \texttt{True}. Indices are 0-based.

\paragraph{Morphology}
\begin{verbatim}
EBNF:
"f.util.where", "(", list, ",", lambda, ")"
\end{verbatim}

\paragraph{Type Overloading}
\begin{itemize}
\item \texttt{where(list[T], L::(\$x [T])->\{bool\})} $\rightarrow$ \texttt{list[int]}
\item The lambda must return \texttt{bool}
\end{itemize}

\paragraph{Examples}
\begin{verbatim}
f.util.where([-1, 3, -2, 5], L::($x [int])->{ $x > 0 })
#c# Returns [1, 3]

f.util.where([10, 20, 30], L::($x [int])->{ $x >= 20 })
#c# Returns [1, 2]

f.util.where([1, 2, 3], L::($x [int])->{ $x > 10 })
#c# Returns []  (no matches)

f.util.where(["a", "bb", "ccc"], L::($x [string])->{ f.util.len($x) >= 2 })
#c# Returns [1, 2]

f.util.where([], L::($x [int])->{ $x > 0 })
#c# Returns []  (empty input)
\end{verbatim}

\paragraph{Error Conditions}
\begin{itemize}
\item Type mismatch: lambda parameter type does not match element type raises a Runtime Error
\item Lambda does not return \texttt{bool} raises a Type Error
\end{itemize}

\subsubsection{Count}
The \texttt{count} function returns the number of elements for which the lambda returns \texttt{True}.

\paragraph{Morphology}
\begin{verbatim}
EBNF:
"f.util.count", "(", list, ",", lambda, ")"
\end{verbatim}

\paragraph{Type Overloading}
\begin{itemize}
\item \texttt{count(list[T], L::(\$x [T])->\{bool\})} $\rightarrow$ \texttt{int}
\item The lambda must return \texttt{bool}
\end{itemize}

\paragraph{Examples}
\begin{verbatim}
f.util.count([1, 2, 3, 4, 5], L::($x [int])->{ $x > 3 })
#c# Returns 2

f.util.count([True, False, True, True], L::($x [bool])->{ $x })
#c# Returns 3

f.util.count([1, 2, 3], L::($x [int])->{ $x > 10 })
#c# Returns 0

f.util.count([], L::($x [int])->{ $x > 0 })
#c# Returns 0
\end{verbatim}

\paragraph{Error Conditions}
\begin{itemize}
\item Type mismatch: lambda parameter type does not match element type raises a Runtime Error
\item Lambda does not return \texttt{bool} raises a Type Error
\end{itemize}

\subsubsection{Find}
The \texttt{find} function returns the first element for which the lambda returns \texttt{True}. If no element matches, a Value Error is raised.

\paragraph{Morphology}
\begin{verbatim}
EBNF:
"f.util.find", "(", list, ",", lambda, ")"
\end{verbatim}

\paragraph{Type Overloading}
\begin{itemize}
\item \texttt{find(list[T], L::(\$x [T])->\{bool\})} $\rightarrow$ \texttt{T}
  (return type is the same as element type)
\item The lambda must return \texttt{bool}
\end{itemize}

\paragraph{Examples}
\begin{verbatim}
f.util.find([1, 2, 3, 4, 5], L::($x [int])->{ $x > 2 })
#c# Returns 3

f.util.find([1, 2, 3, 4, 5], L::($x [int])->{ $x == 4 })
#c# Returns 4

f.util.find(["a", "bb", "ccc"], L::($x [string])->{ f.util.len($x) > 1 })
#c# Returns "bb"

f.util.find([1, 2, 3, 4, 5], L::($x [int])->{ $x > 10 })
#c# Value Error: no matching element

f.util.find([], L::($x [int])->{ $x > 0 })
#c# Value Error: no matching element
\end{verbatim}

\paragraph{Error Conditions}
\begin{itemize}
\item Value Error: no element matches the lambda condition
\item Type mismatch: lambda parameter type does not match element type raises a Runtime Error
\item Lambda does not return \texttt{bool} raises a Type Error
\end{itemize}

\subsubsection{Sort with Comparator}
The \texttt{sort} function returns a new sorted list. Without a comparator, default sort order is ascending. An optional lambda comparator can be provided for custom ordering (see Sort paragraph above for EBNF).

\paragraph{Comparator Lambda}
The comparator lambda receives two elements (\texttt{\$a} and \texttt{\$b}) and must return a numeric value:
\begin{itemize}
\item \textbf{Negative} value (e.g., \texttt{\$a - \$b}): \texttt{\$a} comes before \texttt{\$b}
\item \textbf{Zero}: order unchanged (stable sort)
\item \textbf{Positive} value (e.g., \texttt{\$b - \$a}): \texttt{\$b} comes before \texttt{\$a}
\end{itemize}

This is equivalent to the \texttt{cmp} convention used in many languages. The simplest ascending comparator is \texttt{\$a - \$b}, and descending is \texttt{\$b - \$a}.

\paragraph{Type Overloading}
\begin{itemize}
\item \texttt{sort(list[T])} $\rightarrow$ \texttt{list[T]} (default ascending)
\item \texttt{sort(list[T], L::(\$a [T], \$b [T])->\{int\})} $\rightarrow$ \texttt{list[T]}
  (custom comparator, must return \texttt{int} or \texttt{real})
\end{itemize}

\paragraph{Examples}
\begin{verbatim}
#c# Default sort (ascending)
f.util.sort([5, 1, 3])
#c# Returns [1, 3, 5]

#c# Ascending with comparator
f.util.sort([5, 1, 3], L::($a [int], $b [int])->{ $a - $b })
#c# Returns [1, 3, 5]

#c# Descending with comparator
f.util.sort([5, 1, 3], L::($a [int], $b [int])->{ $b - $a })
#c# Returns [5, 3, 1]

#c# Sort strings by length
f.util.sort(["bb", "a", "ccc"],
    L::($a [string], $b [string])->{ f.util.len($a) - f.util.len($b) })
#c# Returns ["a", "bb", "ccc"]

#c# Sort reals ascending
f.util.sort([3.14, 1.41, 2.72], L::($a [real], $b [real])->{ $a - $b })
#c# Returns [1.41, 2.72, 3.14]

#c# Sort with nested sort key (sort by absolute value)
f.util.sort([-5, 3, -1, 4], L::($a [int], $b [int])->{ f.math.abs($a) - f.math.abs($b) })
#c# Returns [-1, 3, 4, -5]

#c# Sort preserves original (returns new list)
$list [list]= [5, 3, 1]
$sorted [list]= f.util.sort($list)
#c# $list is still [5, 3, 1], $sorted is [1, 3, 5]
\end{verbatim}

\paragraph{Error Conditions}
\begin{itemize}
\item Type mismatch: comparator parameter types do not match element types raises a Runtime Error
\item Comparator does not return a numeric type raises a Type Error
\end{itemize}

\subsubsection{Arity}
The \texttt{arity} function returns the number of parameters a function value accepts. This is useful for introspection and validation before calling a function value stored in a variable.
\begin{verbatim}
EBNF:
"f.util.arity", "(", func_value, ")"
\end{verbatim}

\paragraph{Type Overloading}
\begin{itemize}
\item \texttt{arity(func)} $\rightarrow$ \texttt{int}
\item Works on any function value: built-in references, user-defined references, and lambdas
\end{itemize}

\paragraph{Examples}
\begin{verbatim}
f.util.arity(f.math.sin)                    #c# Returns 1
f.util.arity(f.math.pow)                    #c# Returns 2
f.util.arity(f.math.convert)                #c# Returns 4
f.util.arity(L::($x [int])->{ $x * 2 })    #c# Returns 1
f.util.arity(L::($a [int], $b [int])->{ $a + $b })  #c# Returns 2
\end{verbatim}

\paragraph{Error Conditions}
\begin{itemize}
\item \texttt{Type Error}: argument is not a function value
\item \texttt{Runtime Error}: argument is \texttt{null}
\end{itemize}

\subsubsection{Assert}
\label{sec:assert}
The \texttt{assert} statement checks a boolean condition and raises an error if the condition is \texttt{False}. It is available both as a keyword statement and as a built-in function \texttt{f.util.assert}.

\paragraph{Keyword Form}
\begin{verbatim}
EBNF:
assertion = "assert", boolean_expression, "::", string_type
          , [";", error_type]

error_type = "Assertion Error" | "aerr"
           | "Type Error"       | "terr"
           | "Value Error"      | "verr"
           | "Division by Zero Error" | "0div"
           | "Runtime Error"    | "rerr"
\end{verbatim}

\paragraph{Function Form}
\begin{verbatim}
EBNF:
"f.util.assert", "(", boolean_expression, ","
                , string_type, [",", error_type], ")"
\end{verbatim}

\paragraph{Semantics}
\begin{itemize}
\item The boolean expression is evaluated
\item If \texttt{True}: execution continues (no-op)
\item If \texttt{False}: raises the specified error type with the message string
\item Default error type (when omitted): \texttt{Assertion Error}
\item The \texttt{error\_type} accepts both full names and short aliases
\end{itemize}

\paragraph{Error Type Aliases}
\begin{center}
\begin{tabular}{ll}
Full Name & Alias \\ \hline
\texttt{"Assertion Error"} & \texttt{"aerr"} \\
\texttt{"Type Error"} & \texttt{"terr"} \\
\texttt{"Value Error"} & \texttt{"verr"} \\
\texttt{"Division by Zero Error"} & \texttt{"0div"} \\
\texttt{"Runtime Error"} & \texttt{"rerr"} \\
\end{tabular}
\end{center}

\paragraph{Examples}
\begin{verbatim}
# Keyword form
assert $x > 0 :: "x must be positive"
assert $x > 0 :: "x must be positive" ; "terr"
assert $freq > 0 :: "freq must be positive" ; "0div"

# Function form
f.util.assert($x > 0, "x must be positive")
f.util.assert($x > 0, "x must be positive", "terr")
f.util.assert($freq > 0, "freq must be positive", "0div")

# Short aliases
assert $val != 0 :: "val must be nonzero" ; "aerr"
assert $len > 0 :: "list must not be empty" ; "rerr"
\end{verbatim}

\paragraph{Error Conditions}
\begin{itemize}
\item \texttt{Assertion Error} (default): condition is \texttt{False}
\item The error message is the string provided as the second argument
\item The error type string must be one of the 5 valid types or their aliases; an invalid error type string raises a \texttt{Value Error}
\end{itemize}

\subsubsection{Try/Except}
\label{sec:try-except}
The \texttt{try/except} construct allows structured error handling. Code in the \texttt{try} block is executed; if an error occurs, matching \texttt{otherwise} blocks catch it. An optional \texttt{except} block runs always, regardless of success or failure.

\paragraph{Syntax}
\begin{verbatim}
EBNF:
try_except = "try", "::", code_block
           , { "otherwise", [error_type], "::", code_block }
           , ["except", "::", code_block]

error_type = "Assertion Error" | "aerr"
           | "Type Error"       | "terr"
           | "Value Error"      | "verr"
           | "Division by Zero Error" | "0div"
           | "Runtime Error"    | "rerr"
\end{verbatim}

Reuses the \texttt{error\_type} rule defined for \texttt{assert}.

\paragraph{Semantics}
\begin{itemize}
\item Execute the \texttt{try} block sequentially
\item If no error: skip all \texttt{otherwise} blocks, execute \texttt{except} if present, continue
\item If error: find the first matching \texttt{otherwise} (by error type or catch-all), execute its block, then execute \texttt{except} if present, continue
\item If error occurs and no matching \texttt{otherwise}: the error propagates (halts as normal), but \texttt{except} still runs first
\item Bare \texttt{otherwise} (no error type) catches all errors
\item Multiple \texttt{otherwise} blocks are checked in order; first match wins
\end{itemize}

\paragraph{Variable Scope}
Variables declared in the \texttt{try} block are visible in all \texttt{otherwise} and \texttt{except} blocks. Variables declared in an \texttt{otherwise} block are visible in the \texttt{except} block.

\paragraph{Examples}
\begin{verbatim}
# Basic try/except
try ::{
    $result [real]= f.math.sqrt($x)
} otherwise "Value Error" ::{
    $result [real]= 0.0
}

# Multiple error types
try ::{
    $inv [real]= 1.0 / $x
} otherwise "Division by Zero Error" ::{
    f.util.print("division by zero")
    $inv [real]= 0.0
} otherwise "Value Error" ::{
    f.util.print("invalid value")
    $inv [real]= 0.0
}

# Catch-all otherwise
try ::{
    $data [list]= f.list.fromFile("input.csv")
} otherwise ::{
    f.util.print("failed to load, using defaults")
    $data [list]= []
}

# With except (always runs)
try ::{
    $result [real]= f.math.sqrt($x)
} otherwise "Value Error" ::{
    $result [real]= 0.0
} except ::{
    f.util.print("attempted sqrt calculation")
}

# Nested try
try ::{
    try ::{
        $val [real]= 1.0 / $x
    } otherwise "Division by Zero Error" ::{
        $val [real]= 0.0
    }
    $result [real]= f.math.sqrt($val)
} otherwise "Value Error" ::{
    $result [real]= 0.0
}
\end{verbatim}

\paragraph{Error Conditions}
\begin{itemize}
\item The \texttt{error\_type} string must be one of the 5 valid types or their aliases; an invalid error type string raises a \texttt{Value Error}
\item If an error occurs and no matching \texttt{otherwise} exists, the error propagates after the \texttt{except} block (if present)
\end{itemize}

\subsection{Signal Processing Functions}
Functions for frequency-domain analysis of signals.

\subsubsection{Fast Fourier Transform}
The \texttt{fft} function computes the Discrete Fourier Transform of a real-valued signal.
\begin{verbatim}
EBNF:
"f.ee.fft", "(", list, ")"
\end{verbatim}

\paragraph{Type Overloading}
\begin{itemize}
\item \texttt{fft(list[int])} $\rightarrow$ \texttt{list[complex]}
\item \texttt{fft(list[real])} $\rightarrow$ \texttt{list[complex]}
\item Output length equals input length
\item Works on any signal length (not restricted to powers of 2)
\end{itemize}

\paragraph{Examples}
\begin{verbatim}
#c# FFT of a simple signal
f.ee.fft([1, 0, 1, 0])   #c# Returns [2+0i, 0+0i, 2+0i, 0+0i]

#c# FFT of a single impulse
f.ee.fft([1, 0, 0, 0])   #c# Returns [1+0i, 1+0i, 1+0i, 1+0i]

#c# FFT of a DC signal (constant value)
f.ee.fft([5, 5, 5, 5])   #c# Returns [20+0i, 0+0i, 0+0i, 0+0i]

#c# Extract magnitude spectrum
$freq [list]= f.ee.fft([1, 0, 1, 0])
$mags [list]= f.util.map($freq, L::($z [complex])->{ f.math.abs($z) })
#c# $mags = [2.0, 0.0, 2.0, 0.0]
\end{verbatim}

\paragraph{Error Conditions}
\begin{itemize}
\item Empty list raises a Value Error
\end{itemize}

\subsubsection{Inverse Fast Fourier Transform}
The \texttt{ifft} function computes the Inverse Discrete Fourier Transform, reconstructing a real signal from complex frequency components.
\begin{verbatim}
EBNF:
"f.ee.ifft", "(", list, ")"
\end{verbatim}

\paragraph{Type Overloading}
\begin{itemize}
\item \texttt{ifft(list[complex])} $\rightarrow$ \texttt{list[real]}
\item Output length equals input length
\item Imaginary parts are discarded (reconstructed signal is real)
\end{itemize}

\paragraph{Examples}
\begin{verbatim}
#c# Round-trip: fft then ifft
$signal [list]= [1, 0, 1, 0]
$freq [list]= f.ee.fft($signal)
$reconstructed [list]= f.ee.ifft($freq)
#c# $reconstructed is approximately [1, 0, 1, 0]

#c# ifft of DC component
f.ee.ifft([4+0i, 0+0i, 0+0i, 0+0i])   #c# Returns [1, 1, 1, 1]

#c# ifft of impulse in frequency domain
f.ee.ifft([1+0i, 1+0i, 1+0i, 1+0i])   #c# Returns [1, 0, 0, 0]
\end{verbatim}

\paragraph{Error Conditions}
\begin{itemize}
\item Empty list raises a Value Error
\end{itemize}

\subsection{Debug Functions}
\label{sec:debug-functions}
Debug functions (\texttt{debg.*}) provide runtime inspection, tracing, and profiling. All debug functions are no-ops when \texttt{syst.flag.debug\_mode} is \texttt{False} (default). Enable with:
\begin{verbatim}
syst.flag.debug_mode True
\end{verbatim}

\subsubsection{Debug Print Function}
The \texttt{debg.print} function outputs a variable's value and type. The destination is controlled by the \texttt{mode} parameter or the global \texttt{debug\_output} flag.
\begin{verbatim}
EBNF:
"debg.print", "(", variable, [",", string_type], ")"
\end{verbatim}

\paragraph{Modes}
\begin{itemize}
\item \texttt{"console"} (default) -- Output to console
\item \texttt{"log"} -- Output to debug log file
\item \texttt{"both"} -- Output to console and debug log file
\end{itemize}

\paragraph{Examples}
\begin{verbatim}
debg.print($voltage)
debg.print($voltage, "log")
debg.print($voltage, "both")
\end{verbatim}

\subsubsection{Breakpoint Function}
The \texttt{debg.breakpoint} function pauses execution and prints the current state (variables, stack). An optional condition can be provided.
\begin{verbatim}
EBNF:
"debg.breakpoint", "(", [boolean_expression], ")"
\end{verbatim}

\paragraph{Examples}
\begin{verbatim}
debg.breakpoint()
debg.breakpoint($voltage > 5.0)
debg.breakpoint($count == 10)
\end{verbatim}

If a condition is provided, execution only pauses when the condition is \texttt{True}. When \texttt{debug\_mode} is \texttt{False}, breakpoints are no-ops.

\subsubsection{Graph Dump Function}
The \texttt{debg.graph} function dumps the full electrical graph state (parts, nets, connections, pins).
\begin{verbatim}
EBNF:
"debg.graph", "(", ")"
\end{verbatim}

\paragraph{Examples}
\begin{verbatim}
debg.graph()
\end{verbatim}

Output includes all parts with models and parameters, all nets with tags, all connections, and all pin assignments.

\subsubsection{Resource Monitor Function}
The \texttt{debg.resources} function prints current resource usage: memory, CPU time elapsed, and object counts (parts, nets, connections, pins).
\begin{verbatim}
EBNF:
"debg.resources", "(", ")"
\end{verbatim}

\paragraph{Examples}
\begin{verbatim}
debg.resources()
\end{verbatim}

\subsubsection{Execution Trace Function}
The \texttt{debg.trace} function logs an execution trace entry with a timestamp and optional message. Trace entries are written to the debug log file.
\begin{verbatim}
EBNF:
"debg.trace", "(", [string_type], ")"
\end{verbatim}

\paragraph{Examples}
\begin{verbatim}
debg.trace("Entering main loop")
debg.trace("Before constraint solve")
debg.trace()
\end{verbatim}

\subsubsection{Profiling Function}
The \texttt{debg.profile} function measures the execution time of a function call. It returns the elapsed time in seconds and optionally prints the result.
\begin{verbatim}
EBNF:
"debg.profile", "(", function_call, [",", string_type], ")"
\end{verbatim}

\paragraph{Examples}
\begin{verbatim}
debg.profile(f.util.sort($data))
debg.profile(f.linalg.inv($matrix), "log")
debg.profile(f.ee.fft($signal), "both")
\end{verbatim}

Output includes the function name, execution time in seconds, and the return value.

\subsubsection{Memory Inspection Function}
The \texttt{debg.memory} function reports detailed memory usage: total heap, live objects, and per-type breakdown.
\begin{verbatim}
EBNF:
"debg.memory", "(", ")"
\end{verbatim}

\paragraph{Examples}
\begin{verbatim}
debg.memory()
\end{verbatim}

Output includes total allocated bytes, live object count by type, and peak memory usage.

\subsubsection{CPU Usage Function}
The \texttt{debg.cpu} function returns the current CPU usage as a percentage. Unlike \texttt{debg.resources} which reports cumulative CPU time, this reports instantaneous usage.
\begin{verbatim}
EBNF:
"debg.cpu", "(", ")"
\end{verbatim}

\paragraph{Examples}
\begin{verbatim}
debg.cpu()            #c# Returns 42.5 (percent)
\end{verbatim}

\subsubsection{Variable Watch Function}
The \texttt{debg.watch} function registers a variable for watching. Changes to the watched variable are logged to the debug log file at each breakpoint.
\begin{verbatim}
EBNF:
"debg.watch", "(", variable, ")"
\end{verbatim}

\paragraph{Examples}
\begin{verbatim}
debg.watch($voltage)
debg.watch($current)
debg.breakpoint($count == 10)   #c# Logs all watched variable changes
\end{verbatim}

Multiple variables can be watched. Watch entries persist until the program terminates.

\subsubsection{Call Stack Function}
The \texttt{debg.stack} function prints the current call stack, showing function names and line numbers.
\begin{verbatim}
EBNF:
"debg.stack", "(", ")"
\end{verbatim}

\paragraph{Examples}
\begin{verbatim}
debg.stack()
\end{verbatim}

Output shows the chain of nested function calls from the current execution point back to the entry point.

\subsubsection{Variable Inspector Function}
The \texttt{debg.inspect} function outputs detailed information about a variable: its type, value, memory footprint, and metadata. Richer than \texttt{debg.print}.
\begin{verbatim}
EBNF:
"debg.inspect", "(", variable, ")"
\end{verbatim}

\paragraph{Examples}
\begin{verbatim}
debg.inspect($matrix)
debg.inspect($signal)
\end{verbatim}

Output includes:
\begin{itemize}
\item Type and dimensions (for lists/ndarrays)
\item Element types and count
\item Memory footprint in bytes
\item Value summary (truncated for large structures)
\end{itemize}

\subsection{System Functions}

System functions (\texttt{f.sys.*}) provide runtime introspection and configuration. Unlike \texttt{syst.flag} assignments (which are declarative and appear before the program body), system functions are callable at runtime.

\subsubsection{Get Configuration Function}
The \texttt{f.sys.getConfig} function retrieves a named configuration value. Configuration keys are strings referencing runtime settings.
\\ Morphology:
\begin{verbatim}
"f.sys.getConfig", "(", string, ")"
\end{verbatim}
\\ Returns: The value associated with the given key (type depends on the key). \\
\\ Raises: \texttt{Runtime Error} if the key is not recognized.

\begin{verbatim}
$f max_iter = f.sys.getConfig("max_loop_iter")
$f dbg = f.sys.getConfig("debug_mode")
\end{verbatim}

\subsubsection{Set Configuration Function}
The \texttt{f.sys.setConfig} function assigns a value to a named configuration key at runtime.
\\ Morphology:
\begin{verbatim}
"f.sys.setConfig", "(", string, ",", data, ")"
\end{verbatim}
\\ Returns: \texttt{None}. \\
\\ Raises: \texttt{Runtime Error} if the key is read-only or not recognized.

\begin{verbatim}
f.sys.setConfig("max_loop_iter", 5000)
f.sys.setConfig("debug_mode", True)
\end{verbatim}

\subsubsection{Environment Variable Function}
The \texttt{f.sys.environment} function retrieves the value of an environment variable by name.
\\ Morphology:
\begin{verbatim}
"f.sys.environment", "(", string, ")"
\end{verbatim}
\\ Returns: \texttt{string} --- the value of the environment variable. \\
\\ Raises: \texttt{Runtime Error} if the variable is not set.

\begin{verbatim}
$f home = f.sys.environment("HOME")
$f path = f.sys.environment("PATH")
\end{verbatim}

\subsubsection{Set Environment Function}
The \texttt{f.sys.setEnvironment} function sets an environment variable for the current runtime session.
\\ Morphology:
\begin{verbatim}
"f.sys.setEnvironment", "(", string, ",", string, ")"
\end{verbatim}
\\ Returns: \texttt{None}.

\begin{verbatim}
f.sys.setEnvironment("EDA_LIB", "/opt/eda/lib")
\end{verbatim}

\subsubsection{Version Function}
The \texttt{f.sys.version} function returns the language or runtime version string.
\\ Morphology:
\begin{verbatim}
"f.sys.version", "(", ")"
\end{verbatim}
\\ Returns: \texttt{string} --- version identifier (e.g.\ \texttt{"0.1.0"}).

\begin{verbatim}
$f ver = f.sys.version()
\end{verbatim}

\subsubsection{Features Function}
The \texttt{f.sys.features} function returns a list of supported runtime feature identifiers.
\\ Morphology:
\begin{verbatim}
"f.sys.features", "(", ")"
\end{verbatim}
\\ Returns: \texttt{list[string]} --- feature identifiers enabled in the current runtime.

\begin{verbatim}
$f feats = f.sys.features()
\end{verbatim}

\subsection{Time Functions}

Time functions (\texttt{f.time.*}) provide timestamps, sleep, and named timers. Timers are useful for profiling solver performance or measuring execution time of code sections.

\subsubsection{Epoch Nanoseconds Function}
The \texttt{f.time.now} function returns the current time as nanoseconds since the Unix epoch (1970-01-01T00:00:00Z). \\
Morphology:
\begin{verbatim}
EBNF:
"f.time.now", "(", ")"
\end{verbatim}
\begin{itemize}
\item Result: \texttt{int} --- nanoseconds since epoch.
\end{itemize}
Examples:
\begin{verbatim}
$f t = f.time.now()
#c# 1753036200000000000 (varies)
\end{verbatim}

\subsubsection{ISO 8601 Timestamp Function}
The \texttt{f.time.nowISO} function returns the current time as an ISO 8601 string. Convenience wrapper for \texttt{f.time.nsToISO8601(f.time.now())}. \\
Morphology:
\begin{verbatim}
EBNF:
"f.time.nowISO", "(", ")"
\end{verbatim}
\begin{itemize}
\item Result: \texttt{string} --- ISO 8601 formatted timestamp (e.g.\ \texttt{"2025-07-20T14:30:00Z"}).
\end{itemize}
Examples:
\begin{verbatim}
$f ts = f.time.nowISO()
#c# "2025-07-20T14:30:00Z" (varies)
\end{verbatim}

\subsubsection{Nanoseconds to ISO 8601 Function}
The \texttt{f.time.nsToISO8601} function converts a nanosecond timestamp to an ISO 8601 string. \\
Morphology:
\begin{verbatim}
EBNF:
"f.time.nsToISO8601", "(", int, ")"
\end{verbatim}
\begin{itemize}
\item Operands: \texttt{int} --- nanoseconds since epoch.
\item Result: \texttt{string} --- ISO 8601 formatted timestamp.
\item A \texttt{Value Error} is raised if the argument is negative.
\end{itemize}
Examples:
\begin{verbatim}
f.time.nsToISO8601(1753036200000000000) #c# "2025-07-20T14:30:00Z"
f.time.nsToISO8601(f.time.now())        #c# "2025-07-20T14:30:00Z" (current time)
\end{verbatim}

\subsubsection{Sleep Function}
The \texttt{f.time.sleep} function pauses execution for the given number of milliseconds. \\
Morphology:
\begin{verbatim}
EBNF:
"f.time.sleep", "(", ( int | real ), ")"
\end{verbatim}
\begin{itemize}
\item Operands: \texttt{int} or \texttt{real} --- milliseconds to sleep.
\item Result: \texttt{None}.
\item A \texttt{Value Error} is raised if the argument is negative.
\end{itemize}
Examples:
\begin{verbatim}
f.time.sleep(100)   #c# Pauses for 100 ms
f.time.sleep(0.5)   #c# Pauses for 0.5 ms
\end{verbatim}

\subsubsection{Timer Start Function}
The \texttt{f.time.timerStart} function starts or resets a named timer. If no name is given, the default timer is used. \\
Morphology:
\begin{verbatim}
EBNF:
"f.time.timerStart", "(", [string], ")"
\end{verbatim}
\begin{itemize}
\item Operands: optional \texttt{string} --- timer name (default: \texttt{"default"}).
\item Result: \texttt{None}.
\end{itemize}
Examples:
\begin{verbatim}
f.time.timerStart()              #c# Starts default timer
f.time.timerStart("placement")   #c# Starts timer named "placement"
f.time.timerStart("routing")     #c# Starts timer named "routing"
\end{verbatim}

\subsubsection{Timer Stop Function}
The \texttt{f.time.timerStop} function stops a named timer. \\
Morphology:
\begin{verbatim}
EBNF:
"f.time.timerStop", "(", [string], ")"
\end{verbatim}
\begin{itemize}
\item Operands: optional \texttt{string} --- timer name (default: \texttt{"default"}).
\item Result: \texttt{None}.
\item A \texttt{Runtime Error} is raised if the timer has not been started.
\end{itemize}
Examples:
\begin{verbatim}
f.time.timerStop()              #c# Stops default timer
f.time.timerStop("placement")   #c# Stops timer named "placement"
\end{verbatim}

\subsubsection{Timer Read Function}
The \texttt{f.time.timerRead} function returns the elapsed time in milliseconds since the named timer was started. \\
Morphology:
\begin{verbatim}
EBNF:
"f.time.timerRead", "(", [string], ")"
\end{verbatim}
\begin{itemize}
\item Operands: optional \texttt{string} --- timer name (default: \texttt{"default"}).
\item Result: \texttt{real} --- elapsed milliseconds.
\item A \texttt{Runtime Error} is raised if the timer has not been started.
\end{itemize}
Examples:
\begin{verbatim}
f.time.timerStart("placement")
# ... work ...
f.time.timerStop("placement")
$f elapsed = f.time.timerRead("placement")
#c# 1234.56 (varies)

# Multiple timers are independent
f.time.timerStart("a")
f.time.timerStart("b")
# ... work ...
f.time.timerStop("a")
f.time.timerRead("a")   #c# 500.0 (varies)
f.time.timerRead("b")   #c# Runtime Error (not stopped yet)
\end{verbatim}

\subsection{Random Number Functions}

Random number functions (\texttt{f.rng.*}) provide pseudo-random number generation for Monte Carlo analysis, randomized test vectors, and solver seeding. Use \texttt{f.rng.seed} for reproducible results.

\subsubsection{Seed Function}
The \texttt{f.rng.seed} function sets the random number generator seed. Calling with the same seed produces identical sequences. \\
Morphology:
\begin{verbatim}
EBNF:
"f.rng.seed", "(", int, ")"
\end{verbatim}
\begin{itemize}
\item Operands: \texttt{int} --- seed value.
\item Result: \texttt{None}.
\end{itemize}
Examples:
\begin{verbatim}
f.rng.seed(42)
\end{verbatim}

\subsubsection{Random Function}
The \texttt{f.rng.random} function returns a random real in $[0, 1)$. \\
Morphology:
\begin{verbatim}
EBNF:
"f.rng.random", "(", ")"
\end{verbatim}
\begin{itemize}
\item Result: \texttt{real} --- value in $[0, 1)$.
\end{itemize}
Examples:
\begin{verbatim}
f.rng.seed(42)
f.rng.random()          #c# 0.37454 (reproducible with seed 42)
\end{verbatim}

\subsubsection{Random Integer Function}
The \texttt{f.rng.randInt} function returns a random integer in $[\text{min}, \text{max}]$ (inclusive both ends). \\
Morphology:
\begin{verbatim}
EBNF:
"f.rng.randInt", "(", int, ",", int, ")"
\end{verbatim}
\begin{itemize}
\item Operands: \texttt{min} (\texttt{int}), \texttt{max} (\texttt{int}).
\item Result: \texttt{int} --- value in $[\text{min}, \text{max}]$.
\item A \texttt{Value Error} is raised if \texttt{min > max}.
\end{itemize}
Examples:
\begin{verbatim}
f.rng.seed(42)
$f v = f.rng.randInt(1, 100)   #c# 73 (reproducible with seed 42)
\end{verbatim}

\subsubsection{Random Real Function}
The \texttt{f.rng.randReal} function returns a random real in $[\text{min}, \text{max})$ (half-open interval). \\
Morphology:
\begin{verbatim}
EBNF:
"f.rng.randReal", "(", real, ",", real, ")"
\end{verbatim}
\begin{itemize}
\item Operands: \texttt{min} (\texttt{real}), \texttt{max} (\texttt{real}).
\item Result: \texttt{real} --- value in $[\text{min}, \text{max})$.
\item A \texttt{Value Error} is raised if \texttt{min > max}.
\end{itemize}
Examples:
\begin{verbatim}
f.rng.seed(42)
$f r = f.rng.randReal(0.0, 10.0)   #c# 3.7454 (reproducible with seed 42)
\end{verbatim}

\subsubsection{Uniform Function}
The \texttt{f.rng.uniform} function returns a random real in $[\text{min}, \text{max}]$ (closed interval, inclusive both ends). \\
Morphology:
\begin{verbatim}
EBNF:
"f.rng.uniform", "(", real, ",", real, ")"
\end{verbatim}
\begin{itemize}
\item Operands: \texttt{min} (\texttt{real}), \texttt{max} (\texttt{real}).
\item Result: \texttt{real} --- value in $[\text{min}, \text{max}]$.
\item A \texttt{Value Error} is raised if \texttt{min > max}.
\end{itemize}
Examples:
\begin{verbatim}
f.rng.seed(42)
$f u = f.rng.uniform(0.0, 1.0)    #c# 0.37454 (reproducible with seed 42)
\end{verbatim}

\subsubsection{Normal Distribution Function}
The \texttt{f.rng.normal} function returns a Gaussian-distributed random sample. \\
Morphology:
\begin{verbatim}
EBNF:
"f.rng.normal", "(", [real], [real], ")"
\end{verbatim}
\begin{itemize}
\item Operands: optional \texttt{mean} (\texttt{real}, default 0.0), optional \texttt{stddev} (\texttt{real}, default 1.0).
\item Result: \texttt{real}.
\item A \texttt{Value Error} is raised if \texttt{stddev < 0}.
\end{itemize}
Examples:
\begin{verbatim}
f.rng.seed(42)
$f n = f.rng.normal()              #c# -0.5823 (mean=0, stddev=1)
$f n = f.rng.normal(100.0, 15.0)   #c# 91.234 (mean=100, stddev=15)
\end{verbatim}

\subsubsection{Choice Function}
The \texttt{f.rng.choice} function returns a random element from a list. \\
Morphology:
\begin{verbatim}
EBNF:
"f.rng.choice", "(", list, ")"
\end{verbatim}
\begin{itemize}
\item Operands: \texttt{list} (non-empty).
\item Result: element from the list.
\item A \texttt{Value Error} is raised if the list is empty.
\end{itemize}
Examples:
\begin{verbatim}
f.rng.seed(42)
$f c = f.rng.choice(["R", "C", "L"])   #c# "C" (reproducible with seed 42)
\end{verbatim}

\subsubsection{Shuffle Function}
The \texttt{f.rng.shuffle} function returns a shuffled copy of a list. The original list is unchanged. \\
Morphology:
\begin{verbatim}
EBNF:
"f.rng.shuffle", "(", list, ")"
\end{verbatim}
\begin{itemize}
\item Operands: \texttt{list}.
\item Result: \texttt{list} --- shuffled copy.
\end{itemize}
Examples:
\begin{verbatim}
f.rng.seed(42)
$f s = f.rng.shuffle([1, 2, 3, 4, 5])   #c# [3, 1, 5, 2, 4]
\end{verbatim}

\subsubsection{Random Boolean Function}
The \texttt{f.rng.bool} function returns a random boolean with 50/50 probability. \\
Morphology:
\begin{verbatim}
EBNF:
"f.rng.bool", "(", ")"
\end{verbatim}
\begin{itemize}
\item Result: \texttt{bool} --- \texttt{True} or \texttt{False}.
\end{itemize}
Examples:
\begin{verbatim}
f.rng.seed(42)
$f b = f.rng.bool()   #c# True (reproducible with seed 42)
\end{verbatim}

\subsubsection{Cryptographically Secure Random Function}
The \texttt{f.rng.csprng} function returns a cryptographically secure pseudo-random integer of the specified bit width. Unlike other \texttt{f.rng.*} functions, this is not affected by \texttt{f.rng.seed} and produces non-deterministic output. \\
Morphology:
\begin{verbatim}
EBNF:
"f.rng.csprng", "(", int, ")"
\end{verbatim}
\begin{itemize}
\item Operands: \texttt{int} --- bit width (must be $> 0$).
\item Result: \texttt{int} --- random integer with the specified number of bits.
\item A \texttt{Value Error} is raised if the argument is $\leq 0$.
\end{itemize}
Examples:
\begin{verbatim}
$f uid = f.rng.csprng(128)   #c# 2894839274893728472983748923748923 (varies, 128-bit)
$f key = f.rng.csprng(256)   #c# (varies, 256-bit)
\end{verbatim}

\subsubsection{Random Prime Function}
The \texttt{f.rng.randomPrime} function returns a random prime number in $[\text{min}, \text{max}]$. Useful for cryptographic key generation. \\
Morphology:
\begin{verbatim}
EBNF:
"f.rng.randomPrime", "(", int, ",", int, ")"
\end{verbatim}
\begin{itemize}
\item Operands: \texttt{min} (\texttt{int}, must be $\geq 2$), \texttt{max} (\texttt{int}).
\item Result: \texttt{int} --- random prime in $[\text{min}, \text{max}]$.
\item A \texttt{Value Error} is raised if \texttt{min $<$ 2} or if no prime exists in the range.
\end{itemize}
Examples:
\begin{verbatim}
f.rng.seed(42)
$f p = f.rng.randomPrime(100, 200)   #c# 173 (reproducible with seed 42)
$f p = f.rng.randomPrime(2, 10)      #c# 5 (one of: 2, 3, 5, 7)
\end{verbatim}

\subsection{Hash Functions}

Hash functions (\texttt{f.hash.*}) provide integrity checking, tamper detection, deduplication, and data fingerprinting. Useful for PCB design verification, supply chain integrity, and secure hardware validation.

\subsubsection{General Hash Function}
The \texttt{f.hash.hash} function computes a hash value using the algorithm specified by \texttt{syst.flag.default\_hash} (default: \texttt{"sha512"}). \\
Morphology:
\begin{verbatim}
EBNF:
"f.hash.hash", "(", ( string | list[int] ), [",", ( "hex" | "bytes" | "int" ) ], ")"
\end{verbatim}
\begin{itemize}
\item Operands: \texttt{string} or \texttt{list[int]} (raw bytes), optional output format (\texttt{"hex"} (default), \texttt{"bytes"}, \texttt{"int"}).
\item Result: \texttt{string} (hex), \texttt{list[int]} (bytes), or \texttt{int} depending on format.
\end{itemize}
Examples:
\begin{verbatim}
f.hash.hash("Hello World")               #c# "a591a6d40bf420404a011733cfb7b190d62c65bf0bcda32b57b277d9ad9f146e"
f.hash.hash("Hello World", "hex")        #c# same as above
f.hash.hash("Hello World", "int")        #c# 10102030405060708091011... (integer)
f.hash.hash("Hello World", "bytes")      #c# [165, 145, 166, 212, ...]
\end{verbatim}

\subsubsection{Hash Bytes Function}
The \texttt{f.hash.hashBytes} function computes a hash of raw byte data. \\
Morphology:
\begin{verbatim}
EBNF:
"f.hash.hashBytes", "(", list[int], ")"
\end{verbatim}
\begin{itemize}
\item Operands: \texttt{list[int]} --- byte values (0--255).
\item Result: \texttt{string} --- hex-encoded hash.
\end{itemize}
Examples:
\begin{verbatim}
f.hash.hashBytes([0x48, 0x65, 0x6c, 0x6c, 0x6f])   #c# same hash as "Hello"
\end{verbatim}

\subsubsection{FNV-1a Hash Function}
The \texttt{f.hash.fnv1a} function computes a non-cryptographic FNV-1a hash. Fast, good distribution. \\
Morphology:
\begin{verbatim}
EBNF:
"f.hash.fnv1a", "(", ( string | list[int] ), [",", ( "hex" | "bytes" | "int" ) ], ")"
\end{verbatim}
\begin{itemize}
\item Operands: \texttt{string} or \texttt{list[int]}, optional output format (\texttt{"int"} (default), \texttt{"hex"}, \texttt{"bytes"}).
\item Result: \texttt{int} (default), \texttt{string} (hex), or \texttt{list[int]} (bytes).
\end{itemize}
Examples:
\begin{verbatim}
f.hash.fnv1a("Hello World")                  #c# 2550413245 (int, default)
f.hash.fnv1a("Hello World", "hex")           #c# "982414e5" (hex string)
f.hash.fnv1a("Hello World", "bytes")         #c# [152, 36, 20, 229]
\end{verbatim}

\subsubsection{MurmurHash Function}
The \texttt{f.hash.murmurHash} function computes a non-cryptographic MurmurHash. Fast, good distribution for hash tables. \\
Morphology:
\begin{verbatim}
EBNF:
"f.hash.murmurHash", "(", ( string | list[int] ), [",", ( "hex" | "bytes" | "int" ) ], ")"
\end{verbatim}
\begin{itemize}
\item Operands: \texttt{string} or \texttt{list[int]}, optional output format (\texttt{"int"} (default), \texttt{"hex"}, \texttt{"bytes"}).
\item Result: \texttt{int} (default), \texttt{string} (hex), or \texttt{list[int]} (bytes).
\end{itemize}
Examples:
\begin{verbatim}
f.hash.murmurHash("Hello World")                  #c# (int, varies)
f.hash.murmurHash("Hello World", "hex")           #c# (hex string, varies)
f.hash.murmurHash("Hello World", "bytes")         #c# (byte array, varies)
\end{verbatim}

\subsubsection{xxHash Function}
The \texttt{f.hash.xxHash} function computes a non-cryptographic xxHash. Extremely fast. \\
Morphology:
\begin{verbatim}
EBNF:
"f.hash.xxHash", "(", ( string | list[int] ), [",", ( "hex" | "bytes" | "int" ) ], ")"
\end{verbatim}
\begin{itemize}
\item Operands: \texttt{string} or \texttt{list[int]}, optional output format (\texttt{"int"} (default), \texttt{"hex"}, \texttt{"bytes"}).
\item Result: \texttt{int} (default), \texttt{string} (hex), or \texttt{list[int]} (bytes).
\end{itemize}
Examples:
\begin{verbatim}
f.hash.xxHash("Hello World")                  #c# (int, varies)
f.hash.xxHash("Hello World", "hex")           #c# (hex string, varies)
f.hash.xxHash("Hello World", "bytes")         #c# (byte array, varies)
\end{verbatim}

\subsubsection{CRC-32 Checksum Function}
The \texttt{f.hash.crc32} function computes a CRC-32 checksum. Common in hardware data integrity verification. \\
Morphology:
\begin{verbatim}
EBNF:
"f.hash.crc32", "(", ( string | list[int] ), [",", ( "hex" | "bytes" | "int" ) ], ")"
\end{verbatim}
\begin{itemize}
\item Operands: \texttt{string} or \texttt{list[int]}, optional output format (\texttt{"int"} (default), \texttt{"hex"}, \texttt{"bytes"}).
\item Result: \texttt{int} (32-bit unsigned, default), \texttt{string} (hex), or \texttt{list[int]} (bytes).
\end{itemize}
Examples:
\begin{verbatim}
f.hash.crc32("Hello World")                  #c# 3603590121 (int, default)
f.hash.crc32("Hello World", "hex")           #c# "d5663949" (hex string)
f.hash.crc32("Hello World", "bytes")         #c# [213, 102, 57, 73]
\end{verbatim}

\subsubsection{Adler-32 Checksum Function}
The \texttt{f.hash.adler32} function computes an Adler-32 checksum. Faster than CRC-32, weaker error detection. \\
Morphology:
\begin{verbatim}
EBNF:
"f.hash.adler32", "(", ( string | list[int] ), [",", ( "hex" | "bytes" | "int" ) ], ")"
\end{verbatim}
\begin{itemize}
\item Operands: \texttt{string} or \texttt{list[int]}, optional output format (\texttt{"int"} (default), \texttt{"hex"}, \texttt{"bytes"}).
\item Result: \texttt{int} (32-bit unsigned, default), \texttt{string} (hex), or \texttt{list[int]} (bytes).
\end{itemize}
Examples:
\begin{verbatim}
f.hash.adler32("Hello World")                  #c# 466698496 (int, default)
f.hash.adler32("Hello World", "hex")           #c# "1bd42d50" (hex string)
f.hash.adler32("Hello World", "bytes")         #c# [27, 212, 45, 80]
\end{verbatim}

\subsubsection{MD5 Hash Function}
The \texttt{f.hash.md5} function computes an MD5 digest. Legacy --- use for compatibility only, not for security. \\
Morphology:
\begin{verbatim}
EBNF:
"f.hash.md5", "(", ( string | list[int] ), [",", ( "hex" | "bytes" | "int" ) ], ")"
\end{verbatim}
\begin{itemize}
\item Operands: \texttt{string} or \texttt{list[int]}, optional output format (\texttt{"hex"} (default), \texttt{"bytes"}, \texttt{"int"}).
\item Result: \texttt{string} (hex, default), \texttt{list[int]} (bytes), or \texttt{int}.
\end{itemize}
Examples:
\begin{verbatim}
f.hash.md5("Hello World")                  #c# "b10a8db164e0754105b7a99be72e3fe5" (hex, default)
f.hash.md5("Hello World", "int")           #c# (integer value)
f.hash.md5("Hello World", "bytes")         #c# [177, 10, 141, 177, ...]
\end{verbatim}

\subsubsection{SHA-1 Hash Function}
The \texttt{f.hash.sha1} function computes a SHA-1 digest. Legacy --- use for compatibility only, not for security. \\
Morphology:
\begin{verbatim}
EBNF:
"f.hash.sha1", "(", ( string | list[int] ), [",", ( "hex" | "bytes" | "int" ) ], ")"
\end{verbatim}
\begin{itemize}
\item Operands: \texttt{string} or \texttt{list[int]}, optional output format (\texttt{"hex"} (default), \texttt{"bytes"}, \texttt{"int"}).
\item Result: \texttt{string} (hex, default), \texttt{list[int]} (bytes), or \texttt{int}.
\end{itemize}
Examples:
\begin{verbatim}
f.hash.sha1("Hello World")                  #c# "0a4d55a8d778e5022fab701977c5d840bbc486d0" (hex, default)
f.hash.sha1("Hello World", "int")           #c# (integer value)
f.hash.sha1("Hello World", "bytes")         #c# [10, 77, 85, 168, ...]
\end{verbatim}

\subsubsection{SHA-224 Hash Function}
The \texttt{f.hash.sha224} function computes a SHA-224 digest. \\
Morphology:
\begin{verbatim}
EBNF:
"f.hash.sha224", "(", ( string | list[int] ), [",", ( "hex" | "bytes" | "int" ) ], ")"
\end{verbatim}
\begin{itemize}
\item Operands: \texttt{string} or \texttt{list[int]}, optional output format (\texttt{"hex"} (default), \texttt{"bytes"}, \texttt{"int"}).
\item Result: \texttt{string} (hex, default), \texttt{list[int]} (bytes), or \texttt{int}.
\end{itemize}
Examples:
\begin{verbatim}
f.hash.sha224("Hello World")                  #c# "c8b073b548cd38f5e0eea5727bea4657" (hex, default)
f.hash.sha224("Hello World", "int")           #c# (integer value)
f.hash.sha224("Hello World", "bytes")         #c# [200, 176, 115, 181, ...]
\end{verbatim}

\subsubsection{SHA-256 Hash Function}
The \texttt{f.hash.sha256} function computes a SHA-256 digest. \\
Morphology:
\begin{verbatim}
EBNF:
"f.hash.sha256", "(", ( string | list[int] ), [",", ( "hex" | "bytes" | "int" ) ], ")"
\end{verbatim}
\begin{itemize}
\item Operands: \texttt{string} or \texttt{list[int]}, optional output format (\texttt{"hex"} (default), \texttt{"bytes"}, \texttt{"int"}).
\item Result: \texttt{string} (hex, default), \texttt{list[int]} (bytes), or \texttt{int}.
\end{itemize}
Examples:
\begin{verbatim}
f.hash.sha256("Hello World")                  #c# "a591a6d40bf420404a011733cfb7b190d62c65bf0bcda32b57b277d9ad9f146e" (hex, default)
f.hash.sha256("Hello World", "int")           #c# (integer value)
f.hash.sha256("Hello World", "bytes")         #c# [165, 145, 166, 212, ...]
\end{verbatim}

\subsubsection{SHA-384 Hash Function}
The \texttt{f.hash.sha384} function computes a SHA-384 digest. \\
Morphology:
\begin{verbatim}
EBNF:
"f.hash.sha384", "(", ( string | list[int] ), [",", ( "hex" | "bytes" | "int" ) ], ")"
\end{verbatim}
\begin{itemize}
\item Operands: \texttt{string} or \texttt{list[int]}, optional output format (\texttt{"hex"} (default), \texttt{"bytes"}, \texttt{"int"}).
\item Result: \texttt{string} (hex, default), \texttt{list[int]} (bytes), or \texttt{int}.
\end{itemize}
Examples:
\begin{verbatim}
f.hash.sha384("Hello World")                  #c# "14fdb4453b92a0c18562b5a9705826b45a3ef8de5..." (hex, default)
f.hash.sha384("Hello World", "int")           #c# (integer value)
f.hash.sha384("Hello World", "bytes")         #c# [20, 253, 180, 69, ...]
\end{verbatim}

\subsubsection{SHA-512 Hash Function}
The \texttt{f.hash.sha512} function computes a SHA-512 digest. \\
Morphology:
\begin{verbatim}
EBNF:
"f.hash.sha512", "(", ( string | list[int] ), [",", ( "hex" | "bytes" | "int" ) ], ")"
\end{verbatim}
\begin{itemize}
\item Operands: \texttt{string} or \texttt{list[int]}, optional output format (\texttt{"hex"} (default), \texttt{"bytes"}, \texttt{"int"}).
\item Result: \texttt{string} (hex, default), \texttt{list[int]} (bytes), or \texttt{int}.
\end{itemize}
Examples:
\begin{verbatim}
f.hash.sha512("Hello World")                  #c# "861844d6704e8573fec34d967e20bcfef3d424cf4..." (hex, default)
f.hash.sha512("Hello World", "int")           #c# (integer value)
f.hash.sha512("Hello World", "bytes")         #c# [134, 24, 68, 214, ...]
\end{verbatim}

\subsubsection{SHA-3 256 Hash Function}
The \texttt{f.hash.sha3\_256} function computes a SHA-3 256-bit digest. \\
Morphology:
\begin{verbatim}
EBNF:
"f.hash.sha3_256", "(", ( string | list[int] ), [",", ( "hex" | "bytes" | "int" ) ], ")"
\end{verbatim}
\begin{itemize}
\item Operands: \texttt{string} or \texttt{list[int]}, optional output format (\texttt{"hex"} (default), \texttt{"bytes"}, \texttt{"int"}).
\item Result: \texttt{string} (hex, default), \texttt{list[int]} (bytes), or \texttt{int}.
\end{itemize}
Examples:
\begin{verbatim}
f.hash.sha3_256("Hello World")                  #c# "644bcc7e56197e90906997b5657bd..." (hex, default)
f.hash.sha3_256("Hello World", "int")           #c# (integer value)
f.hash.sha3_256("Hello World", "bytes")         #c# [100, 75, 204, 126, ...]
\end{verbatim}

\subsubsection{SHA-3 512 Hash Function}
The \texttt{f.hash.sha3\_512} function computes a SHA-3 512-bit digest. \\
Morphology:
\begin{verbatim}
EBNF:
"f.hash.sha3_512", "(", ( string | list[int] ), [",", ( "hex" | "bytes" | "int" ) ], ")"
\end{verbatim}
\begin{itemize}
\item Operands: \texttt{string} or \texttt{list[int]}, optional output format (\texttt{"hex"} (default), \texttt{"bytes"}, \texttt{"int"}).
\item Result: \texttt{string} (hex, default), \texttt{list[int]} (bytes), or \texttt{int}.
\end{itemize}
Examples:
\begin{verbatim}
f.hash.sha3_512("Hello World")                  #c# (128-char hex, default)
f.hash.sha3_512("Hello World", "int")           #c# (integer value)
f.hash.sha3_512("Hello World", "bytes")         #c# (byte array)
\end{verbatim}

\subsubsection{BLAKE2b Hash Function}
The \texttt{f.hash.blake2b} function computes a BLAKE2b digest. High performance, cryptographic. \\
Morphology:
\begin{verbatim}
EBNF:
"f.hash.blake2b", "(", ( string | list[int] ), [",", ( "hex" | "bytes" | "int" ) ], ")"
\end{verbatim}
\begin{itemize}
\item Operands: \texttt{string} or \texttt{list[int]}, optional output format (\texttt{"hex"} (default), \texttt{"bytes"}, \texttt{"int"}).
\item Result: \texttt{string} (hex, default), \texttt{list[int]} (bytes), or \texttt{int}.
\end{itemize}
Examples:
\begin{verbatim}
f.hash.blake2b("Hello World")                  #c# (128-char hex, default)
f.hash.blake2b("Hello World", "int")           #c# (integer value)
f.hash.blake2b("Hello World", "bytes")         #c# (byte array)
\end{verbatim}

\subsubsection{BLAKE2s Hash Function}
The \texttt{f.hash.blake2s} function computes a BLAKE2s digest. High performance, cryptographic. \\
Morphology:
\begin{verbatim}
EBNF:
"f.hash.blake2s", "(", ( string | list[int] ), [",", ( "hex" | "bytes" | "int" ) ], ")"
\end{verbatim}
\begin{itemize}
\item Operands: \texttt{string} or \texttt{list[int]}, optional output format (\texttt{"hex"} (default), \texttt{"bytes"}, \texttt{"int"}).
\item Result: \texttt{string} (hex, default), \texttt{list[int]} (bytes), or \texttt{int}.
\end{itemize}
Examples:
\begin{verbatim}
f.hash.blake2s("Hello World")                  #c# (64-char hex, default)
f.hash.blake2s("Hello World", "int")           #c# (integer value)
f.hash.blake2s("Hello World", "bytes")         #c# (byte array)
\end{verbatim}

\subsubsection{BLAKE3 Hash Function}
The \texttt{f.hash.blake3} function computes a BLAKE3 digest. Extremely fast, cryptographic. \\
Morphology:
\begin{verbatim}
EBNF:
"f.hash.blake3", "(", ( string | list[int] ), [",", ( "hex" | "bytes" | "int" ) ], ")"
\end{verbatim}
\begin{itemize}
\item Operands: \texttt{string} or \texttt{list[int]}, optional output format (\texttt{"hex"} (default), \texttt{"bytes"}, \texttt{"int"}).
\item Result: \texttt{string} (hex, default), \texttt{list[int]} (bytes), or \texttt{int}.
\end{itemize}
Examples:
\begin{verbatim}
f.hash.blake3("Hello World")                  #c# (64-char hex, default)
f.hash.blake3("Hello World", "int")           #c# (integer value)
f.hash.blake3("Hello World", "bytes")         #c# (byte array)
\end{verbatim}

\subsection{Probability Distributions}

Probability distribution functions (\texttt{f.stat.*}) provide PDF/PMF, CDF, inverse CDF, and random sampling for common statistical distributions. Useful for Monte Carlo tolerance analysis, reliability modeling, and statistical design validation.

\subsubsection{Normal PDF}
The \texttt{f.stat.normalPDF} function computes the probability density of the normal (Gaussian) distribution at a given point.
\begin{verbatim}
EBNF:
"f.stat.normalPDF", "(", real, ",", real, ",", real, ")"
\end{verbatim}
\begin{itemize}
\item Operands: \texttt{x} (evaluation point), \texttt{mean}, \texttt{stddev} (must be $> 0$).
\item Result: \texttt{real}.
\item A \texttt{Value Error} is raised if \texttt{stddev $\leq$ 0}.
\end{itemize}
\begin{verbatim}
f.stat.normalPDF(0, 0, 1)    #c# 0.39894
f.stat.normalPDF(1, 0, 1)    #c# 0.24197
\end{verbatim}

\subsubsection{Normal CDF}
The \texttt{f.stat.normalCDF} function computes the cumulative probability $P(X \leq x)$ for the normal distribution.
\begin{verbatim}
EBNF:
"f.stat.normalCDF", "(", real, ",", real, ",", real, ")"
\end{verbatim}
\begin{itemize}
\item Operands: \texttt{x}, \texttt{mean}, \texttt{stddev} (must be $> 0$).
\item Result: \texttt{real} (in $[0, 1]$).
\item A \texttt{Value Error} is raised if \texttt{stddev $\leq$ 0}.
\end{itemize}
\begin{verbatim}
f.stat.normalCDF(0, 0, 1)    #c# 0.5
f.stat.normalCDF(1.96, 0, 1) #c# 0.975
\end{verbatim}

\subsubsection{Normal Inverse CDF}
The \texttt{f.stat.normalInvCDF} function computes the quantile (inverse CDF) for the normal distribution.
\begin{verbatim}
EBNF:
"f.stat.normalInvCDF", "(", real, ",", real, ",", real, ")"
\end{verbatim}
\begin{itemize}
\item Operands: \texttt{p} (probability, in $(0, 1)$), \texttt{mean}, \texttt{stddev} (must be $> 0$).
\item Result: \texttt{real}.
\item A \texttt{Value Error} is raised if \texttt{p $\leq$ 0}, \texttt{p $\geq$ 1}, or \texttt{stddev $\leq$ 0}.
\end{itemize}
\begin{verbatim}
f.stat.normalInvCDF(0.5, 0, 1)    #c# 0.0
f.stat.normalInvCDF(0.975, 0, 1)  #c# 1.96
\end{verbatim}

\subsubsection{Normal Random}
The \texttt{f.stat.normalRandom} function returns a random sample from the normal distribution.
\begin{verbatim}
EBNF:
"f.stat.normalRandom", "(", real, ",", real, ")"
\end{verbatim}
\begin{itemize}
\item Operands: \texttt{mean}, \texttt{stddev} (must be $> 0$).
\item Result: \texttt{real}.
\item A \texttt{Value Error} is raised if \texttt{stddev $\leq$ 0}.
\end{itemize}
\begin{verbatim}
f.rng.seed(42)
$f v = f.stat.normalRandom(0, 1)   #c# -0.5823 (varies)
\end{verbatim}

\subsubsection{Exponential PDF}
The \texttt{f.stat.exponentialPDF} function computes the probability density of the exponential distribution.
\begin{verbatim}
EBNF:
"f.stat.exponentialPDF", "(", real, ",", real, ")"
\end{verbatim}
\begin{itemize}
\item Operands: \texttt{x} (must be $\geq 0$), \texttt{lambda} (rate, must be $> 0$).
\item Result: \texttt{real}.
\item Returns 0 for $x < 0$.
\item A \texttt{Value Error} is raised if \texttt{lambda $\leq$ 0}.
\end{itemize}
\begin{verbatim}
f.stat.exponentialPDF(0, 1)   #c# 1.0
f.stat.exponentialPDF(1, 1)   #c# 0.36788
\end{verbatim}

\subsubsection{Exponential CDF}
The \texttt{f.stat.exponentialCDF} function computes the cumulative probability for the exponential distribution.
\begin{verbatim}
EBNF:
"f.stat.exponentialCDF", "(", real, ",", real, ")"
\end{verbatim}
\begin{itemize}
\item Operands: \texttt{x}, \texttt{lambda} (must be $> 0$).
\item Result: \texttt{real} (in $[0, 1]$).
\item A \texttt{Value Error} is raised if \texttt{lambda $\leq$ 0}.
\end{itemize}
\begin{verbatim}
f.stat.exponentialCDF(0, 1)   #c# 0.0
f.stat.exponentialCDF(1, 1)   #c# 0.63212
\end{verbatim}

\subsubsection{Exponential Inverse CDF}
The \texttt{f.stat.exponentialInvCDF} function computes the quantile for the exponential distribution.
\begin{verbatim}
EBNF:
"f.stat.exponentialInvCDF", "(", real, ",", real, ")"
\end{verbatim}
\begin{itemize}
\item Operands: \texttt{p} (probability, in $(0, 1)$), \texttt{lambda} (must be $> 0$).
\item Result: \texttt{real}.
\item A \texttt{Value Error} is raised if \texttt{p $\leq$ 0}, \texttt{p $\geq$ 1}, or \texttt{lambda $\leq$ 0}.
\end{itemize}
\begin{verbatim}
f.stat.exponentialInvCDF(0.5, 1)   #c# 0.69315
f.stat.exponentialInvCDF(0.9, 1)   #c# 2.30259
\end{verbatim}

\subsubsection{Exponential Random}
The \texttt{f.stat.exponentialRandom} function returns a random sample from the exponential distribution.
\begin{verbatim}
EBNF:
"f.stat.exponentialRandom", "(", real, ")"
\end{verbatim}
\begin{itemize}
\item Operands: \texttt{lambda} (must be $> 0$).
\item Result: \texttt{real}.
\item A \texttt{Value Error} is raised if \texttt{lambda $\leq$ 0}.
\end{itemize}
\begin{verbatim}
f.rng.seed(42)
$f v = f.stat.exponentialRandom(1)   #c# 0.4869 (varies)
\end{verbatim}

\subsubsection{Gamma PDF}
The \texttt{f.stat.gammaPDF} function computes the probability density of the gamma distribution.
\begin{verbatim}
EBNF:
"f.stat.gammaPDF", "(", real, ",", real, ",", real, ")"
\end{verbatim}
\begin{itemize}
\item Operands: \texttt{x} (must be $> 0$), \texttt{shape} (must be $> 0$), \texttt{rate} (must be $> 0$).
\item Result: \texttt{real}.
\item Returns 0 for $x \leq 0$.
\item A \texttt{Value Error} is raised if \texttt{shape $\leq$ 0} or \texttt{rate $\leq$ 0}.
\end{itemize}
\begin{verbatim}
f.stat.gammaPDF(1, 1, 1)   #c# 1.0
f.stat.gammaPDF(2, 2, 1)   #c# 0.27067
\end{verbatim}

\subsubsection{Gamma CDF}
The \texttt{f.stat.gammaCDF} function computes the cumulative probability for the gamma distribution.
\begin{verbatim}
EBNF:
"f.stat.gammaCDF", "(", real, ",", real, ",", real, ")"
\end{verbatim}
\begin{itemize}
\item Operands: \texttt{x}, \texttt{shape} (must be $> 0$), \texttt{rate} (must be $> 0$).
\item Result: \texttt{real} (in $[0, 1]$).
\item A \texttt{Value Error} is raised if \texttt{shape $\leq$ 0} or \texttt{rate $\leq$ 0}.
\end{itemize}
\begin{verbatim}
f.stat.gammaCDF(1, 1, 1)   #c# 0.63212
f.stat.gammaCDF(2, 2, 1)   #c# 0.59399
\end{verbatim}

\subsubsection{Gamma Inverse CDF}
The \texttt{f.stat.gammaInvCDF} function computes the quantile for the gamma distribution.
\begin{verbatim}
EBNF:
"f.stat.gammaInvCDF", "(", real, ",", real, ",", real, ")"
\end{verbatim}
\begin{itemize}
\item Operands: \texttt{p} (probability, in $(0, 1)$), \texttt{shape} (must be $> 0$), \texttt{rate} (must be $> 0$).
\item Result: \texttt{real}.
\item A \texttt{Value Error} is raised if \texttt{p $\leq$ 0}, \texttt{p $\geq$ 1}, \texttt{shape $\leq$ 0}, or \texttt{rate $\leq$ 0}.
\end{itemize}
\begin{verbatim}
f.stat.gammaInvCDF(0.5, 1, 1)   #c# 0.69315
\end{verbatim}

\subsubsection{Gamma Random}
The \texttt{f.stat.gammaRandom} function returns a random sample from the gamma distribution.
\begin{verbatim}
EBNF:
"f.stat.gammaRandom", "(", real, ",", real, ")"
\end{verbatim}
\begin{itemize}
\item Operands: \texttt{shape} (must be $> 0$), \texttt{rate} (must be $> 0$).
\item Result: \texttt{real}.
\item A \texttt{Value Error} is raised if \texttt{shape $\leq$ 0} or \texttt{rate $\leq$ 0}.
\end{itemize}
\begin{verbatim}
f.rng.seed(42)
$f v = f.stat.gammaRandom(2, 1)   #c# 1.2345 (varies)
\end{verbatim}

\subsubsection{Beta PDF}
The \texttt{f.stat.betaPDF} function computes the probability density of the beta distribution.
\begin{verbatim}
EBNF:
"f.stat.betaPDF", "(", real, ",", real, ",", real, ")"
\end{verbatim}
\begin{itemize}
\item Operands: \texttt{x} (in $[0, 1]$), \texttt{alpha} (must be $> 0$), \texttt{beta} (must be $> 0$).
\item Result: \texttt{real}.
\item Returns 0 for $x < 0$ or $x > 1$.
\item A \texttt{Value Error} is raised if \texttt{alpha $\leq$ 0} or \texttt{beta $\leq$ 0}.
\end{itemize}
\begin{verbatim}
f.stat.betaPDF(0.5, 1, 1)   #c# 1.0
f.stat.betaPDF(0.5, 2, 2)   #c# 1.5
\end{verbatim}

\subsubsection{Beta CDF}
The \texttt{f.stat.betaCDF} function computes the cumulative probability for the beta distribution.
\begin{verbatim}
EBNF:
"f.stat.betaCDF", "(", real, ",", real, ",", real, ")"
\end{verbatim}
\begin{itemize}
\item Operands: \texttt{x}, \texttt{alpha} (must be $> 0$), \texttt{beta} (must be $> 0$).
\item Result: \texttt{real} (in $[0, 1]$).
\item A \texttt{Value Error} is raised if \texttt{alpha $\leq$ 0} or \texttt{beta $\leq$ 0}.
\end{itemize}
\begin{verbatim}
f.stat.betaCDF(0.5, 1, 1)   #c# 0.5
f.stat.betaCDF(0.5, 2, 2)   #c# 0.5
\end{verbatim}

\subsubsection{Beta Inverse CDF}
The \texttt{f.stat.betaInvCDF} function computes the quantile for the beta distribution.
\begin{verbatim}
EBNF:
"f.stat.betaInvCDF", "(", real, ",", real, ",", real, ")"
\end{verbatim}
\begin{itemize}
\item Operands: \texttt{p} (probability, in $(0, 1)$), \texttt{alpha} (must be $> 0$), \texttt{beta} (must be $> 0$).
\item Result: \texttt{real}.
\item A \texttt{Value Error} is raised if \texttt{p $\leq$ 0}, \texttt{p $\geq$ 1}, \texttt{alpha $\leq$ 0}, or \texttt{beta $\leq$ 0}.
\end{itemize}
\begin{verbatim}
f.stat.betaInvCDF(0.5, 1, 1)   #c# 0.5
\end{verbatim}

\subsubsection{Beta Random}
The \texttt{f.stat.betaRandom} function returns a random sample from the beta distribution.
\begin{verbatim}
EBNF:
"f.stat.betaRandom", "(", real, ",", real, ")"
\end{verbatim}
\begin{itemize}
\item Operands: \texttt{alpha} (must be $> 0$), \texttt{beta} (must be $> 0$).
\item Result: \texttt{real} (in $[0, 1]$).
\item A \texttt{Value Error} is raised if \texttt{alpha $\leq$ 0} or \texttt{beta $\leq$ 0}.
\end{itemize}
\begin{verbatim}
f.rng.seed(42)
$f v = f.stat.betaRandom(2, 5)   #c# 0.3456 (varies)
\end{verbatim}

\subsubsection{Weibull PDF}
The \texttt{f.stat.weibullPDF} function computes the probability density of the Weibull distribution.
\begin{verbatim}
EBNF:
"f.stat.weibullPDF", "(", real, ",", real, ",", real, ")"
\end{verbatim}
\begin{itemize}
\item Operands: \texttt{x} (must be $\geq 0$), \texttt{shape} (must be $> 0$), \texttt{scale} (must be $> 0$).
\item Result: \texttt{real}.
\item Returns 0 for $x < 0$.
\item A \texttt{Value Error} is raised if \texttt{shape $\leq$ 0} or \texttt{scale $\leq$ 0}.
\end{itemize}
\begin{verbatim}
f.stat.weibullPDF(1, 1, 1)   #c# 0.36788
f.stat.weibullPDF(1, 2, 1)   #c# 0.73576
\end{verbatim}

\subsubsection{Weibull CDF}
The \texttt{f.stat.weibullCDF} function computes the cumulative probability for the Weibull distribution.
\begin{verbatim}
EBNF:
"f.stat.weibullCDF", "(", real, ",", real, ",", real, ")"
\end{verbatim}
\begin{itemize}
\item Operands: \texttt{x}, \texttt{shape} (must be $> 0$), \texttt{scale} (must be $> 0$).
\item Result: \texttt{real} (in $[0, 1]$).
\item A \texttt{Value Error} is raised if \texttt{shape $\leq$ 0} or \texttt{scale $\leq$ 0}.
\end{itemize}
\begin{verbatim}
f.stat.weibullCDF(1, 1, 1)   #c# 0.63212
f.stat.weibullCDF(1, 2, 1)   #c# 0.39347
\end{verbatim}

\subsubsection{Weibull Inverse CDF}
The \texttt{f.stat.weibullInvCDF} function computes the quantile for the Weibull distribution.
\begin{verbatim}
EBNF:
"f.stat.weibullInvCDF", "(", real, ",", real, ",", real, ")"
\end{verbatim}
\begin{itemize}
\item Operands: \texttt{p} (probability, in $(0, 1)$), \texttt{shape} (must be $> 0$), \texttt{scale} (must be $> 0$).
\item Result: \texttt{real}.
\item A \texttt{Value Error} is raised if \texttt{p $\leq$ 0}, \texttt{p $\geq$ 1}, \texttt{shape $\leq$ 0}, or \texttt{scale $\leq$ 0}.
\end{itemize}
\begin{verbatim}
f.stat.weibullInvCDF(0.5, 1, 1)   #c# 0.69315
\end{verbatim}

\subsubsection{Weibull Random}
The \texttt{f.stat.weibullRandom} function returns a random sample from the Weibull distribution.
\begin{verbatim}
EBNF:
"f.stat.weibullRandom", "(", real, ",", real, ")"
\end{verbatim}
\begin{itemize}
\item Operands: \texttt{shape} (must be $> 0$), \texttt{scale} (must be $> 0$).
\item Result: \texttt{real}.
\item A \texttt{Value Error} is raised if \texttt{shape $\leq$ 0} or \texttt{scale $\leq$ 0}.
\end{itemize}
\begin{verbatim}
f.rng.seed(42)
$f v = f.stat.weibullRandom(2, 1)   #c# 0.8765 (varies)
\end{verbatim}

\subsubsection{Logistic PDF}
The \texttt{f.stat.logisticPDF} function computes the probability density of the logistic distribution.
\begin{verbatim}
EBNF:
"f.stat.logisticPDF", "(", real, ",", real, ",", real, ")"
\end{verbatim}
\begin{itemize}
\item Operands: \texttt{x}, \texttt{mu} (location), \texttt{s} (scale, must be $> 0$).
\item Result: \texttt{real}.
\item A \texttt{Value Error} is raised if \texttt{s $\leq$ 0}.
\end{itemize}
\begin{verbatim}
f.stat.logisticPDF(0, 0, 1)   #c# 0.25
f.stat.logisticPDF(1, 0, 1)   #c# 0.19661
\end{verbatim}

\subsubsection{Logistic CDF}
The \texttt{f.stat.logisticCDF} function computes the cumulative probability for the logistic distribution.
\begin{verbatim}
EBNF:
"f.stat.logisticCDF", "(", real, ",", real, ",", real, ")"
\end{verbatim}
\begin{itemize}
\item Operands: \texttt{x}, \texttt{mu}, \texttt{s} (must be $> 0$).
\item Result: \texttt{real} (in $(0, 1)$).
\item A \texttt{Value Error} is raised if \texttt{s $\leq$ 0}.
\end{itemize}
\begin{verbatim}
f.stat.logisticCDF(0, 0, 1)   #c# 0.5
f.stat.logisticCDF(2, 0, 1)   #c# 0.88080
\end{verbatim}

\subsubsection{Logistic Inverse CDF}
The \texttt{f.stat.logisticInvCDF} function computes the quantile for the logistic distribution.
\begin{verbatim}
EBNF:
"f.stat.logisticInvCDF", "(", real, ",", real, ",", real, ")"
\end{verbatim}
\begin{itemize}
\item Operands: \texttt{p} (probability, in $(0, 1)$), \texttt{mu}, \texttt{s} (must be $> 0$).
\item Result: \texttt{real}.
\item A \texttt{Value Error} is raised if \texttt{p $\leq$ 0}, \texttt{p $\geq$ 1}, or \texttt{s $\leq$ 0}.
\end{itemize}
\begin{verbatim}
f.stat.logisticInvCDF(0.5, 0, 1)   #c# 0.0
f.stat.logisticInvCDF(0.9, 0, 1)   #c# 2.19722
\end{verbatim}

\subsubsection{Logistic Random}
The \texttt{f.stat.logisticRandom} function returns a random sample from the logistic distribution.
\begin{verbatim}
EBNF:
"f.stat.logisticRandom", "(", real, ",", real, ")"
\end{verbatim}
\begin{itemize}
\item Operands: \texttt{mu}, \texttt{s} (must be $> 0$).
\item Result: \texttt{real}.
\item A \texttt{Value Error} is raised if \texttt{s $\leq$ 0}.
\end{itemize}
\begin{verbatim}
f.rng.seed(42)
$f v = f.stat.logisticRandom(0, 1)   #c# -0.3456 (varies)
\end{verbatim}

\subsubsection{Cauchy PDF}
The \texttt{f.stat.cauchyPDF} function computes the probability density of the Cauchy distribution.
\begin{verbatim}
EBNF:
"f.stat.cauchyPDF", "(", real, ",", real, ",", real, ")"
\end{verbatim}
\begin{itemize}
\item Operands: \texttt{x}, \texttt{x0} (location), \texttt{gamma} (scale, must be $> 0$).
\item Result: \texttt{real}.
\item A \texttt{Value Error} is raised if \texttt{gamma $\leq$ 0}.
\end{itemize}
\begin{verbatim}
f.stat.cauchyPDF(0, 0, 1)   #c# 0.31831
f.stat.cauchyPDF(1, 0, 1)   #c# 0.15915
\end{verbatim}

\subsubsection{Cauchy CDF}
The \texttt{f.stat.cauchyCDF} function computes the cumulative probability for the Cauchy distribution.
\begin{verbatim}
EBNF:
"f.stat.cauchyCDF", "(", real, ",", real, ",", real, ")"
\end{verbatim}
\begin{itemize}
\item Operands: \texttt{x}, \texttt{x0}, \texttt{gamma} (must be $> 0$).
\item Result: \texttt{real} (in $(0, 1)$).
\item A \texttt{Value Error} is raised if \texttt{gamma $\leq$ 0}.
\end{itemize}
\begin{verbatim}
f.stat.cauchyCDF(0, 0, 1)   #c# 0.5
f.stat.cauchyCDF(1, 0, 1)   #c# 0.75
\end{verbatim}

\subsubsection{Cauchy Inverse CDF}
The \texttt{f.stat.cauchyInvCDF} function computes the quantile for the Cauchy distribution.
\begin{verbatim}
EBNF:
"f.stat.cauchyInvCDF", "(", real, ",", real, ",", real, ")"
\end{verbatim}
\begin{itemize}
\item Operands: \texttt{p} (probability, in $(0, 1)$), \texttt{x0}, \texttt{gamma} (must be $> 0$).
\item Result: \texttt{real}.
\item A \texttt{Value Error} is raised if \texttt{p $\leq$ 0}, \texttt{p $\geq$ 1}, or \texttt{gamma $\leq$ 0}.
\end{itemize}
\begin{verbatim}
f.stat.cauchyInvCDF(0.5, 0, 1)   #c# 0.0
f.stat.cauchyInvCDF(0.75, 0, 1)  #c# 1.0
\end{verbatim}

\subsubsection{Cauchy Random}
The \texttt{f.stat.cauchyRandom} function returns a random sample from the Cauchy distribution.
\begin{verbatim}
EBNF:
"f.stat.cauchyRandom", "(", real, ",", real, ")"
\end{verbatim}
\begin{itemize}
\item Operands: \texttt{x0}, \texttt{gamma} (must be $> 0$).
\item Result: \texttt{real}.
\item A \texttt{Value Error} is raised if \texttt{gamma $\leq$ 0}.
\end{itemize}
\begin{verbatim}
f.rng.seed(42)
$f v = f.stat.cauchyRandom(0, 1)   #c# -1.2345 (varies)
\end{verbatim}

\subsubsection{Student's t PDF}
The \texttt{f.stat.studentTPDF} function computes the probability density of the Student's t distribution.
\begin{verbatim}
EBNF:
"f.stat.studentTPDF", "(", real, ",", real, ")"
\end{verbatim}
\begin{itemize}
\item Operands: \texttt{x}, \texttt{df} (degrees of freedom, must be $> 0$).
\item Result: \texttt{real}.
\item A \texttt{Value Error} is raised if \texttt{df $\leq$ 0}.
\end{itemize}
\begin{verbatim}
f.stat.studentTPDF(0, 1)   #c# 0.31831 (Cauchy)
f.stat.studentTPDF(0, 10)  #c# 0.37589
\end{verbatim}

\subsubsection{Student's t CDF}
The \texttt{f.stat.studentTCDF} function computes the cumulative probability for the Student's t distribution.
\begin{verbatim}
EBNF:
"f.stat.studentTCDF", "(", real, ",", real, ")"
\end{verbatim}
\begin{itemize}
\item Operands: \texttt{x}, \texttt{df} (must be $> 0$).
\item Result: \texttt{real} (in $(0, 1)$).
\item A \texttt{Value Error} is raised if \texttt{df $\leq$ 0}.
\end{itemize}
\begin{verbatim}
f.stat.studentTCDF(0, 10)    #c# 0.5
f.stat.studentTCDF(1.812, 10) #c# 0.95
\end{verbatim}

\subsubsection{Student's t Inverse CDF}
The \texttt{f.stat.studentTInvCDF} function computes the quantile for the Student's t distribution.
\begin{verbatim}
EBNF:
"f.stat.studentTInvCDF", "(", real, ",", real, ")"
\end{verbatim}
\begin{itemize}
\item Operands: \texttt{p} (probability, in $(0, 1)$), \texttt{df} (must be $> 0$).
\item Result: \texttt{real}.
\item A \texttt{Value Error} is raised if \texttt{p $\leq$ 0}, \texttt{p $\geq$ 1}, or \texttt{df $\leq$ 0}.
\end{itemize}
\begin{verbatim}
f.stat.studentTInvCDF(0.5, 10)    #c# 0.0
f.stat.studentTInvCDF(0.95, 10)   #c# 1.812
\end{verbatim}

\subsubsection{Student's t Random}
The \texttt{f.stat.studentTRandom} function returns a random sample from the Student's t distribution.
\begin{verbatim}
EBNF:
"f.stat.studentTRandom", "(", real, ")"
\end{verbatim}
\begin{itemize}
\item Operands: \texttt{df} (must be $> 0$).
\item Result: \texttt{real}.
\item A \texttt{Value Error} is raised if \texttt{df $\leq$ 0}.
\end{itemize}
\begin{verbatim}
f.rng.seed(42)
$f v = f.stat.studentTRandom(5)   #c# -0.7890 (varies)
\end{verbatim}

\subsubsection{Chi-Square PDF}
The \texttt{f.stat.chiSquarePDF} function computes the probability density of the chi-square distribution.
\begin{verbatim}
EBNF:
"f.stat.chiSquarePDF", "(", real, ",", real, ")"
\end{verbatim}
\begin{itemize}
\item Operands: \texttt{x} (must be $\geq 0$), \texttt{df} (degrees of freedom, must be $> 0$).
\item Result: \texttt{real}.
\item Returns 0 for $x < 0$.
\item A \texttt{Value Error} is raised if \texttt{df $\leq$ 0}.
\end{itemize}
\begin{verbatim}
f.stat.chiSquarePDF(1, 1)   #c# 0.24197
f.stat.chiSquarePDF(2, 2)   #c# 0.18394
\end{verbatim}

\subsubsection{Chi-Square CDF}
The \texttt{f.stat.chiSquareCDF} function computes the cumulative probability for the chi-square distribution.
\begin{verbatim}
EBNF:
"f.stat.chiSquareCDF", "(", real, ",", real, ")"
\end{verbatim}
\begin{itemize}
\item Operands: \texttt{x}, \texttt{df} (must be $> 0$).
\item Result: \texttt{real} (in $[0, 1]$).
\item A \texttt{Value Error} is raised if \texttt{df $\leq$ 0}.
\end{itemize}
\begin{verbatim}
f.stat.chiSquareCDF(1, 1)   #c# 0.68269
f.stat.chiSquareCDF(2, 2)   #c# 0.63212
\end{verbatim}

\subsubsection{Chi-Square Inverse CDF}
The \texttt{f.stat.chiSquareInvCDF} function computes the quantile for the chi-square distribution.
\begin{verbatim}
EBNF:
"f.stat.chiSquareInvCDF", "(", real, ",", real, ")"
\end{verbatim}
\begin{itemize}
\item Operands: \texttt{p} (probability, in $(0, 1)$), \texttt{df} (must be $> 0$).
\item Result: \texttt{real}.
\item A \texttt{Value Error} is raised if \texttt{p $\leq$ 0}, \texttt{p $\geq$ 1}, or \texttt{df $\leq$ 0}.
\end{itemize}
\begin{verbatim}
f.stat.chiSquareInvCDF(0.5, 1)   #c# 0.45494
\end{verbatim}

\subsubsection{Chi-Square Random}
The \texttt{f.stat.chiSquareRandom} function returns a random sample from the chi-square distribution.
\begin{verbatim}
EBNF:
"f.stat.chiSquareRandom", "(", real, ")"
\end{verbatim}
\begin{itemize}
\item Operands: \texttt{df} (must be $> 0$).
\item Result: \texttt{real}.
\item A \texttt{Value Error} is raised if \texttt{df $\leq$ 0}.
\end{itemize}
\begin{verbatim}
f.rng.seed(42)
$f v = f.stat.chiSquareRandom(3)   #c# 2.3456 (varies)
\end{verbatim}

\subsubsection{F-Distribution PDF}
The \texttt{f.stat.fDistributionPDF} function computes the probability density of the F-distribution.
\begin{verbatim}
EBNF:
"f.stat.fDistributionPDF", "(", real, ",", real, ",", real, ")"
\end{verbatim}
\begin{itemize}
\item Operands: \texttt{x} (must be $\geq 0$), \texttt{df1} (numerator df, must be $> 0$), \texttt{df2} (denominator df, must be $> 0$).
\item Result: \texttt{real}.
\item Returns 0 for $x < 0$.
\item A \texttt{Value Error} is raised if \texttt{df1 $\leq$ 0} or \texttt{df2 $\leq$ 0}.
\end{itemize}
\begin{verbatim}
f.stat.fDistributionPDF(1, 5, 10)   #c# 0.31831 (approx)
\end{verbatim}

\subsubsection{F-Distribution CDF}
The \texttt{f.stat.fDistributionCDF} function computes the cumulative probability for the F-distribution.
\begin{verbatim}
EBNF:
"f.stat.fDistributionCDF", "(", real, ",", real, ",", real, ")"
\end{verbatim}
\begin{itemize}
\item Operands: \texttt{x}, \texttt{df1}, \texttt{df2} (both must be $> 0$).
\item Result: \texttt{real} (in $[0, 1]$).
\item A \texttt{Value Error} is raised if \texttt{df1 $\leq$ 0} or \texttt{df2 $\leq$ 0}.
\end{itemize}
\begin{verbatim}
f.stat.fDistributionCDF(1, 5, 10)   #c# 0.5 (approx)
\end{verbatim}

\subsubsection{F-Distribution Inverse CDF}
The \texttt{f.stat.fDistributionInvCDF} function computes the quantile for the F-distribution.
\begin{verbatim}
EBNF:
"f.stat.fDistributionInvCDF", "(", real, ",", real, ",", real, ")"
\end{verbatim}
\begin{itemize}
\item Operands: \texttt{p} (probability, in $(0, 1)$), \texttt{df1}, \texttt{df2} (both must be $> 0$).
\item Result: \texttt{real}.
\item A \texttt{Value Error} is raised if \texttt{p $\leq$ 0}, \texttt{p $\geq$ 1}, \texttt{df1 $\leq$ 0}, or \texttt{df2 $\leq$ 0}.
\end{itemize}
\begin{verbatim}
f.stat.fDistributionInvCDF(0.5, 5, 10)   #c# 0.976 (approx)
\end{verbatim}

\subsubsection{F-Distribution Random}
The \texttt{f.stat.fDistributionRandom} function returns a random sample from the F-distribution.
\begin{verbatim}
EBNF:
"f.stat.fDistributionRandom", "(", real, ",", real, ")"
\end{verbatim}
\begin{itemize}
\item Operands: \texttt{df1} (must be $> 0$), \texttt{df2} (must be $> 0$).
\item Result: \texttt{real}.
\item A \texttt{Value Error} is raised if \texttt{df1 $\leq$ 0} or \texttt{df2 $\leq$ 0}.
\end{itemize}
\begin{verbatim}
f.rng.seed(42)
$f v = f.stat.fDistributionRandom(5, 10)   #c# 1.2345 (varies)
\end{verbatim}

\subsubsection{Uniform PDF}
The \texttt{f.stat.uniformPDF} function computes the probability density of the continuous uniform distribution.
\begin{verbatim}
EBNF:
"f.stat.uniformPDF", "(", real, ",", real, ",", real, ")"
\end{verbatim}
\begin{itemize}
\item Operands: \texttt{x}, \texttt{a} (lower bound), \texttt{b} (upper bound, must be $> a$).
\item Result: \texttt{real} --- $1/(b-a)$ if $a \leq x \leq b$, else 0.
\item A \texttt{Value Error} is raised if \texttt{a $\geq$ b}.
\end{itemize}
\begin{verbatim}
f.stat.uniformPDF(0.5, 0, 1)   #c# 1.0
f.stat.uniformPDF(2, 0, 1)     #c# 0.0
\end{verbatim}

\subsubsection{Uniform CDF}
The \texttt{f.stat.uniformCDF} function computes the cumulative probability for the continuous uniform distribution.
\begin{verbatim}
EBNF:
"f.stat.uniformCDF", "(", real, ",", real, ",", real, ")"
\end{verbatim}
\begin{itemize}
\item Operands: \texttt{x}, \texttt{a}, \texttt{b} (must be $> a$).
\item Result: \texttt{real} (in $[0, 1]$).
\item A \texttt{Value Error} is raised if \texttt{a $\geq$ b}.
\end{itemize}
\begin{verbatim}
f.stat.uniformCDF(0.5, 0, 1)   #c# 0.5
f.stat.uniformCDF(1, 0, 1)     #c# 1.0
\end{verbatim}

\subsubsection{Uniform Inverse CDF}
The \texttt{f.stat.uniformInvCDF} function computes the quantile for the continuous uniform distribution.
\begin{verbatim}
EBNF:
"f.stat.uniformInvCDF", "(", real, ",", real, ",", real, ")"
\end{verbatim}
\begin{itemize}
\item Operands: \texttt{p} (probability, in $(0, 1)$), \texttt{a}, \texttt{b} (must be $> a$).
\item Result: \texttt{real}.
\item A \texttt{Value Error} is raised if \texttt{p $\leq$ 0}, \texttt{p $\geq$ 1}, or \texttt{a $\geq$ b}.
\end{itemize}
\begin{verbatim}
f.stat.uniformInvCDF(0.5, 0, 1)   #c# 0.5
f.stat.uniformInvCDF(0.9, 0, 10)  #c# 9.0
\end{verbatim}

\subsubsection{Uniform Random}
The \texttt{f.stat.uniformRandom} function returns a random sample from the continuous uniform distribution.
\begin{verbatim}
EBNF:
"f.stat.uniformRandom", "(", real, ",", real, ")"
\end{verbatim}
\begin{itemize}
\item Operands: \texttt{a} (lower bound), \texttt{b} (upper bound, must be $> a$).
\item Result: \texttt{real}.
\item A \texttt{Value Error} is raised if \texttt{a $\geq$ b}.
\end{itemize}
\begin{verbatim}
f.rng.seed(42)
$f v = f.stat.uniformRandom(0, 1)   #c# 0.3745 (varies)
\end{verbatim}

\subsubsection{Binomial PMF}
The \texttt{f.stat.binomialPMF} function computes the probability mass of the binomial distribution.
\begin{verbatim}
EBNF:
"f.stat.binomialPMF", "(", int, ",", int, ",", real, ")"
\end{verbatim}
\begin{itemize}
\item Operands: \texttt{k} (number of successes, $0 \leq k \leq n$), \texttt{n} (trials, must be $\geq 0$), \texttt{p} (probability, in $[0, 1]$).
\item Result: \texttt{real}.
\item Returns 0 if $k < 0$ or $k > n$.
\item A \texttt{Value Error} is raised if \texttt{n $<$ 0} or \texttt{p $<$ 0} or \texttt{p $>$ 1}.
\end{itemize}
\begin{verbatim}
f.stat.binomialPMF(3, 10, 0.5)   #c# 0.11719
f.stat.binomialPMF(0, 10, 0.5)   #c# 0.00098
\end{verbatim}

\subsubsection{Binomial CDF}
The \texttt{f.stat.binomialCDF} function computes the cumulative probability for the binomial distribution.
\begin{verbatim}
EBNF:
"f.stat.binomialCDF", "(", int, ",", int, ",", real, ")"
\end{verbatim}
\begin{itemize}
\item Operands: \texttt{k}, \texttt{n} (must be $\geq 0$), \texttt{p} (in $[0, 1]$).
\item Result: \texttt{real} (in $[0, 1]$).
\item A \texttt{Value Error} is raised if \texttt{n $<$ 0} or \texttt{p $<$ 0} or \texttt{p $>$ 1}.
\end{itemize}
\begin{verbatim}
f.stat.binomialCDF(3, 10, 0.5)   #c# 0.17188
f.stat.binomialCDF(5, 10, 0.5)   #c# 0.62305
\end{verbatim}

\subsubsection{Binomial Inverse CDF}
The \texttt{f.stat.binomialInvCDF} function computes the quantile for the binomial distribution.
\begin{verbatim}
EBNF:
"f.stat.binomialInvCDF", "(", real, ",", int, ",", real, ")"
\end{verbatim}
\begin{itemize}
\item Operands: \texttt{p} (probability, in $(0, 1)$), \texttt{n} (must be $\geq 0$), \texttt{p\_succ} (in $[0, 1]$).
\item Result: \texttt{int}.
\item A \texttt{Value Error} is raised if \texttt{p $\leq$ 0}, \texttt{p $\geq$ 1}, \texttt{n $<$ 0}, or \texttt{p\_succ $<$ 0} or \texttt{p\_succ $>$ 1}.
\end{itemize}
\begin{verbatim}
f.stat.binomialInvCDF(0.5, 10, 0.5)   #c# 5
\end{verbatim}

\subsubsection{Binomial Random}
The \texttt{f.stat.binomialRandom} function returns a random sample from the binomial distribution.
\begin{verbatim}
EBNF:
"f.stat.binomialRandom", "(", int, ",", real, ")"
\end{verbatim}
\begin{itemize}
\item Operands: \texttt{n} (trials, must be $\geq 0$), \texttt{p} (probability, in $[0, 1]$).
\item Result: \texttt{int}.
\item A \texttt{Value Error} is raised if \texttt{n $<$ 0} or \texttt{p $<$ 0} or \texttt{p $>$ 1}.
\end{itemize}
\begin{verbatim}
f.rng.seed(42)
$f v = f.stat.binomialRandom(10, 0.5)   #c# 6 (varies)
\end{verbatim}

\subsubsection{Poisson PMF}
The \texttt{f.stat.poissonPMF} function computes the probability mass of the Poisson distribution.
\begin{verbatim}
EBNF:
"f.stat.poissonPMF", "(", int, ",", real, ")"
\end{verbatim}
\begin{itemize}
\item Operands: \texttt{k} (events, must be $\geq 0$), \texttt{lambda} (rate, must be $> 0$).
\item Result: \texttt{real}.
\item Returns 0 for $k < 0$.
\item A \texttt{Value Error} is raised if \texttt{k $<$ 0} or \texttt{lambda $\leq$ 0}.
\end{itemize}
\begin{verbatim}
f.stat.poissonPMF(3, 5)   #c# 0.14037
f.stat.poissonPMF(0, 5)   #c# 0.00674
\end{verbatim}

\subsubsection{Poisson CDF}
The \texttt{f.stat.poissonCDF} function computes the cumulative probability for the Poisson distribution.
\begin{verbatim}
EBNF:
"f.stat.poissonCDF", "(", int, ",", real, ")"
\end{verbatim}
\begin{itemize}
\item Operands: \texttt{k}, \texttt{lambda} (must be $> 0$).
\item Result: \texttt{real} (in $[0, 1]$).
\item A \texttt{Value Error} is raised if \texttt{k $<$ 0} or \texttt{lambda $\leq$ 0}.
\end{itemize}
\begin{verbatim}
f.stat.poissonCDF(3, 5)   #c# 0.26503
f.stat.poissonCDF(5, 5)   #c# 0.61596
\end{verbatim}

\subsubsection{Poisson Inverse CDF}
The \texttt{f.stat.poissonInvCDF} function computes the quantile for the Poisson distribution.
\begin{verbatim}
EBNF:
"f.stat.poissonInvCDF", "(", real, ",", real, ")"
\end{verbatim}
\begin{itemize}
\item Operands: \texttt{p} (probability, in $(0, 1)$), \texttt{lambda} (must be $> 0$).
\item Result: \texttt{int}.
\item A \texttt{Value Error} is raised if \texttt{p $\leq$ 0}, \texttt{p $\geq$ 1}, or \texttt{lambda $\leq$ 0}.
\end{itemize}
\begin{verbatim}
f.stat.poissonInvCDF(0.5, 5)   #c# 5
\end{verbatim}

\subsubsection{Poisson Random}
The \texttt{f.stat.poissonRandom} function returns a random sample from the Poisson distribution.
\begin{verbatim}
EBNF:
"f.stat.poissonRandom", "(", real, ")"
\end{verbatim}
\begin{itemize}
\item Operands: \texttt{lambda} (must be $> 0$).
\item Result: \texttt{int}.
\item A \texttt{Value Error} is raised if \texttt{lambda $\leq$ 0}.
\end{itemize}
\begin{verbatim}
f.rng.seed(42)
$f v = f.stat.poissonRandom(5)   #c# 4 (varies)
\end{verbatim}

\subsubsection{Geometric PMF}
The \texttt{f.stat.geometricPMF} function computes the probability mass of the geometric distribution (number of failures before first success).
\begin{verbatim}
EBNF:
"f.stat.geometricPMF", "(", int, ",", real, ")"
\end{verbatim}
\begin{itemize}
\item Operands: \texttt{k} (failures, must be $\geq 0$), \texttt{p} (success probability, in $(0, 1]$).
\item Result: \texttt{real}.
\item A \texttt{Value Error} is raised if \texttt{k $<$ 0}, \texttt{p $\leq$ 0}, or \texttt{p $>$ 1}.
\end{itemize}
\begin{verbatim}
f.stat.geometricPMF(0, 0.5)   #c# 0.5
f.stat.geometricPMF(2, 0.5)   #c# 0.125
\end{verbatim}

\subsubsection{Geometric CDF}
The \texttt{f.stat.geometricCDF} function computes the cumulative probability for the geometric distribution.
\begin{verbatim}
EBNF:
"f.stat.geometricCDF", "(", int, ",", real, ")"
\end{verbatim}
\begin{itemize}
\item Operands: \texttt{k}, \texttt{p} (in $(0, 1]$).
\item Result: \texttt{real} (in $[0, 1]$).
\item A \texttt{Value Error} is raised if \texttt{k $<$ 0}, \texttt{p $\leq$ 0}, or \texttt{p $>$ 1}.
\end{itemize}
\begin{verbatim}
f.stat.geometricCDF(0, 0.5)   #c# 0.5
f.stat.geometricCDF(2, 0.5)   #c# 0.875
\end{verbatim}

\subsubsection{Geometric Inverse CDF}
The \texttt{f.stat.geometricInvCDF} function computes the quantile for the geometric distribution.
\begin{verbatim}
EBNF:
"f.stat.geometricInvCDF", "(", real, ",", real, ")"
\end{verbatim}
\begin{itemize}
\item Operands: \texttt{p\_prob} (probability, in $(0, 1)$), \texttt{p\_succ} (success probability, in $(0, 1]$).
\item Result: \texttt{int}.
\item A \texttt{Value Error} is raised if \texttt{p\_prob $\leq$ 0}, \texttt{p\_prob $\geq$ 1}, \texttt{p\_succ $\leq$ 0}, or \texttt{p\_succ $>$ 1}.
\end{itemize}
\begin{verbatim}
f.stat.geometricInvCDF(0.5, 0.5)   #c# 1
\end{verbatim}

\subsubsection{Geometric Random}
The \texttt{f.stat.geometricRandom} function returns a random sample from the geometric distribution.
\begin{verbatim}
EBNF:
"f.stat.geometricRandom", "(", real, ")"
\end{verbatim}
\begin{itemize}
\item Operands: \texttt{p} (success probability, in $(0, 1]$).
\item Result: \texttt{int}.
\item A \texttt{Value Error} is raised if \texttt{p $\leq$ 0} or \texttt{p $>$ 1}.
\end{itemize}
\begin{verbatim}
f.rng.seed(42)
$f v = f.stat.geometricRandom(0.5)   #c# 1 (varies)
\end{verbatim}

\subsubsection{Negative Binomial PMF}
The \texttt{f.stat.negativeBinomialPMF} function computes the probability mass of the negative binomial distribution (number of failures before $r$ successes).
\begin{verbatim}
EBNF:
"f.stat.negativeBinomialPMF", "(", int, ",", int, ",", real, ")"
\end{verbatim}
\begin{itemize}
\item Operands: \texttt{k} (failures, must be $\geq 0$), \texttt{r} (target successes, must be $> 0$), \texttt{p} (success probability, in $(0, 1]$).
\item Result: \texttt{real}.
\item A \texttt{Value Error} is raised if \texttt{k $<$ 0}, \texttt{r $\leq$ 0}, \texttt{p $\leq$ 0}, or \texttt{p $>$ 1}.
\end{itemize}
\begin{verbatim}
f.stat.negativeBinomialPMF(3, 5, 0.5)   #c# 0.15625
\end{verbatim}

\subsubsection{Negative Binomial CDF}
The \texttt{f.stat.negativeBinomialCDF} function computes the cumulative probability for the negative binomial distribution.
\begin{verbatim}
EBNF:
"f.stat.negativeBinomialCDF", "(", int, ",", int, ",", real, ")"
\end{verbatim}
\begin{itemize}
\item Operands: \texttt{k}, \texttt{r} (must be $> 0$), \texttt{p} (in $(0, 1]$).
\item Result: \texttt{real} (in $[0, 1]$).
\item A \texttt{Value Error} is raised if \texttt{k $<$ 0}, \texttt{r $\leq$ 0}, \texttt{p $\leq$ 0}, or \texttt{p $>$ 1}.
\end{itemize}
\begin{verbatim}
f.stat.negativeBinomialCDF(3, 5, 0.5)   #c# 0.5
\end{verbatim}

\subsubsection{Negative Binomial Inverse CDF}
The \texttt{f.stat.negativeBinomialInvCDF} function computes the quantile for the negative binomial distribution.
\begin{verbatim}
EBNF:
"f.stat.negativeBinomialInvCDF", "(", real, ",", int, ",", real, ")"
\end{verbatim}
\begin{itemize}
\item Operands: \texttt{p\_prob} (probability, in $(0, 1)$), \texttt{r} (must be $> 0$), \texttt{p\_succ} (in $(0, 1]$).
\item Result: \texttt{int}.
\item A \texttt{Value Error} is raised if \texttt{p\_prob $\leq$ 0}, \texttt{p\_prob $\geq$ 1}, \texttt{r $\leq$ 0}, \texttt{p\_succ $\leq$ 0}, or \texttt{p\_succ $>$ 1}.
\end{itemize}
\begin{verbatim}
f.stat.negativeBinomialInvCDF(0.5, 5, 0.5)   #c# 4
\end{verbatim}

\subsubsection{Negative Binomial Random}
The \texttt{f.stat.negativeBinomialRandom} function returns a random sample from the negative binomial distribution.
\begin{verbatim}
EBNF:
"f.stat.negativeBinomialRandom", "(", int, ",", real, ")"
\end{verbatim}
\begin{itemize}
\item Operands: \texttt{r} (target successes, must be $> 0$), \texttt{p} (success probability, in $(0, 1]$).
\item Result: \texttt{int}.
\item A \texttt{Value Error} is raised if \texttt{r $\leq$ 0}, \texttt{p $\leq$ 0}, or \texttt{p $>$ 1}.
\end{itemize}
\begin{verbatim}
f.rng.seed(42)
$f v = f.stat.negativeBinomialRandom(5, 0.5)   #c# 3 (varies)
\end{verbatim}

\subsection{Numerical Analysis}

Numerical analysis functions (\texttt{f.math.*}) provide root finding, interpolation, calculus, and ODE solving for engineering computations. All functions in this section accept function references (lambdas or named functions) as arguments.

\paragraph{When to Use}
\begin{itemize}
\item \textbf{Root finding} --- solving $f(x) = 0$ for circuit operating points, bias calculations, threshold detection.
\item \textbf{Interpolation} --- lookup tables, component model approximation, sensor calibration curves.
\item \textbf{Calculus} --- filter response integration, sensitivity analysis, power/energy calculations.
\item \textbf{ODE solving} --- transient circuit analysis, SPICE-like simulation, thermal modeling, PLL loop dynamics.
\end{itemize}

\subsubsection{Newton's Method}
The \texttt{f.math.newton} function finds a root of $f(x) = 0$ using Newton's method (also known as the Newton-Raphson method). Requires the derivative $f'(x)$ to be computed automatically via automatic differentiation.

\textbf{Algorithm:} Starting from initial guess $x_0$, iterate $x_{n+1} = x_n - f(x_n) / f'(x_n)$ until $|f(x_n)| < \text{tol}$ or \texttt{maxIter} is reached.

\textbf{Convergence:} Quadratic convergence when near a root and $f'(x) \neq 0$. May diverge if the initial guess is far from a root, or if $f'(x)$ is near zero.

\textbf{Behavior:} Finds ONE root. The root found depends on the initial guess $x_0$. Different initial guesses may converge to different roots. If $f$ has multiple roots, the user must try different starting points.

\begin{verbatim}
EBNF:
"f.math.newton", "(", func, ",", real, [",", real], [",", int], ")"
\end{verbatim}
\begin{itemize}
\item Operands: \texttt{\$f} (function $f(x)$), \texttt{x0} (initial guess), optional \texttt{tol} (\texttt{real}, default $10^{-10}$), optional \texttt{maxIter} (\texttt{int}, default 100).
\item Result: \texttt{real} --- root approximation where $|f(x)| < \text{tol}$.
\item A \texttt{Runtime Error} is raised if $f'(x) = 0$ during any iteration (division by zero).
\item A \texttt{Value Error} is raised if no convergence within \texttt{maxIter} iterations, if \texttt{tol} $\leq 0$, or if \texttt{maxIter} $\leq 0$.
\item A \texttt{Type Error} is raised if \texttt{\$f} is not a function reference.
\end{itemize}

\paragraph{When to Prefer Newton's Method}
Fast convergence when you have a good initial guess and the function is smooth. Not suitable for functions with discontinuities or zero derivatives near the root.

\begin{verbatim}
$f sq = func($x) $x * $x - 2
f.math.newton($sq, 1.0)         #c# 1.41421 (sqrt(2), found from guess 1.0)
f.math.newton($sq, -1.0)        #c# -1.41421 (sqrt(2), found from guess -1.0)
f.math.newton($sq, 1.0, 1e-8)   #c# 1.41421 (custom tolerance)
\end{verbatim}

\subsubsection{Bisection Method}
The \texttt{f.math.bisection} function finds a root of $f(x) = 0$ using the bisection method. Requires $f(a)$ and $f(b)$ to have opposite signs (Intermediate Value Theorem guarantees a root exists in $[a, b]$).

\textbf{Algorithm:} Repeatedly halve the interval $[a, b]$ by checking the sign of $f$ at the midpoint. Converges to a root within $\text{tol}$.

\textbf{Convergence:} Linear convergence. Always converges if $f(a)$ and $f(b)$ have opposite signs and $f$ is continuous. Maximum $\lceil \log_2((b-a)/\text{tol}) \rceil$ iterations.

\textbf{Behavior:} Finds ONE root in the interval $[a, b]$. If multiple roots exist in the interval, returns one of them (not specified which). The result is always within $\text{tol}$ of the true root.

\begin{verbatim}
EBNF:
"f.math.bisection", "(", func, ",", real, ",", real, [",", real], ")"
\end{verbatim}
\begin{itemize}
\item Operands: \texttt{\$f} (function $f(x)$), \texttt{a} (lower bound), \texttt{b} (upper bound), optional \texttt{tol} (\texttt{real}, default $10^{-10}$).
\item Result: \texttt{real} --- root approximation within $\text{tol}$ of the true root.
\item A \texttt{Value Error} is raised if $f(a)$ and $f(b)$ have the same sign (no guaranteed root), if $a = b$, if no convergence within internal limit ($\sim$1000 iterations), or if \texttt{tol} $\leq 0$.
\item A \texttt{Type Error} is raised if \texttt{\$f} is not a function reference.
\end{itemize}

\paragraph{When to Prefer Bisection}
Most robust method. Guaranteed to converge if a sign change exists. Use when you know the root is bracketed but the function is not smooth, or when Newton's method fails.

\begin{verbatim}
$f sq = func($x) $x * $x - 2
f.math.bisection($sq, 0, 2)       #c# 1.41421 (sqrt(2))
f.math.bisection($sq, 0, 2, 1e-8) #c# 1.41421 (custom tolerance)
\end{verbatim}

\subsubsection{Secant Method}
The \texttt{f.math.secant} function finds a root of $f(x) = 0$ using the secant method. A finite-difference approximation of Newton's method that does not require computing $f'(x)$.

\textbf{Algorithm:} Starting from two points $x_0, x_1$, iterate $x_{n+1} = x_n - f(x_n) \cdot (x_n - x_{n-1}) / (f(x_n) - f(x_{n-1}))$.

\textbf{Convergence:} Superlinear convergence (order $\approx 1.618$). Faster than bisection, slower than Newton. May fail if $f(x_n) = f(x_{n-1})$ (zero denominator).

\textbf{Behavior:} Finds ONE root. The root found depends on both initial points $x_0$ and $x_1$. Not guaranteed to converge (unlike bisection).

\begin{verbatim}
EBNF:
"f.math.secant", "(", func, ",", real, ",", real, [",", real], [",", int], ")"
\end{verbatim}
\begin{itemize}
\item Operands: \texttt{\$f} (function $f(x)$), \texttt{x0}, \texttt{x1} (two initial points), optional \texttt{tol} (\texttt{real}, default $10^{-10}$), optional \texttt{maxIter} (\texttt{int}, default 100).
\item Result: \texttt{real} --- root approximation.
\item A \texttt{Runtime Error} is raised if $f(x_0) = f(x_1)$ (zero denominator during first step).
\item A \texttt{Value Error} is raised if no convergence within \texttt{maxIter} iterations, if \texttt{tol} $\leq 0$, or if \texttt{maxIter} $\leq 0$.
\item A \texttt{Type Error} is raised if \texttt{\$f} is not a function reference.
\end{itemize}

\paragraph{When to Prefer Secant Method}
When you cannot compute $f'(x)$ analytically but need faster convergence than bisection. Good general-purpose method when Newton's automatic differentiation is not available or too expensive.

\begin{verbatim}
$f sq = func($x) $x * $x - 2
f.math.secant($sq, 0, 2)         #c# 1.41421 (sqrt(2))
f.math.secant($sq, 0, 2, 1e-8)   #c# 1.41421 (custom tolerance)
\end{verbatim}

\subsubsection{Brent's Method}
The \texttt{f.math.brent} function finds a root of $f(x) = 0$ using Brent's method. Combines the robustness of bisection with the speed of secant and inverse quadratic interpolation.

\textbf{Algorithm:} Maintains a bracketing interval $[a, b]$. Attempts inverse quadratic interpolation when safe, falls back to secant, and uses bisection as a last resort. Guarantees progress at every step.

\textbf{Convergence:} Superlinear convergence in practice. Guaranteed to converge (like bisection) if the root is bracketed. Typically faster than bisection, sometimes faster than secant.

\textbf{Behavior:} Finds ONE root in the interval $[a, b]$. Combines the best properties of bisection (guaranteed convergence) and secant (fast convergence). Recommended as the default root-finding method when a bracket is known.

\begin{verbatim}
EBNF:
"f.math.brent", "(", func, ",", real, ",", real, [",", real], ")"
\end{verbatim}
\begin{itemize}
\item Operands: \texttt{\$f} (function $f(x)$), \texttt{a} (lower bound), \texttt{b} (upper bound), optional \texttt{tol} (\texttt{real}, default $10^{-10}$).
\item Result: \texttt{real} --- root approximation.
\item A \texttt{Value Error} is raised if $f(a)$ and $f(b)$ have the same sign (no guaranteed root), if $a = b$, if no convergence within internal limit ($\sim$1000 iterations), or if \texttt{tol} $\leq 0$.
\item A \texttt{Type Error} is raised if \texttt{\$f} is not a function reference.
\end{itemize}

\paragraph{When to Prefer Brent's Method}
Recommended default when you have a bracketed interval. More robust than secant, faster than bisection. Use this unless you have a specific reason to prefer another method.

\begin{verbatim}
$f sq = func($x) $x * $x - 2
f.math.brent($sq, 0, 2)       #c# 1.41421 (sqrt(2))
f.math.brent($sq, 0, 2, 1e-8) #c# 1.41421 (custom tolerance)
\end{verbatim}

\paragraph{Root-Finding Method Comparison}

\begin{center}
\begin{tabular}{|l|c|c|c|c|}
\hline
\textbf{Method} & \textbf{Convergence} & \textbf{Needs \$f'} & \textbf{Needs Bracket} & \textbf{Guaranteed} \\
\hline
Newton & Quadratic & Yes & No & No \\
Bisection & Linear & No & Yes & Yes \\
Secant & Superlinear ($\phi$) & No & No & No \\
Brent & Superlinear & No & Yes & Yes \\
\hline
\end{tabular}
\end{center}

\subsubsection{Interpolation}
The \texttt{f.math.interpolate} function estimates $y$ values between (or beyond) known data points.

\textbf{Linear interpolation:} For each interval $[x_i, x_{i+1}]$, computes the straight-line segment between $(x_i, y_i)$ and $(x_{i+1}, y_{i+1})$. Continuous but not smooth (kinks at data points). Always uses the two nearest data points.

\textbf{Polynomial interpolation:} Fits a single polynomial of degree $n-1$ through all $n$ data points using Lagrange interpolation. Exact at all data points. Smooth everywhere. May exhibit Runge's oscillation for equally-spaced points with large $n$ (use spline instead for $n > 10$).

\textbf{Extrapolation:} Both methods extrapolate beyond the data range. Linear extrapolation extends the last segment. Polynomial extrapolation follows the fitted polynomial. No error is raised for out-of-range $x$, but extrapolation accuracy may be poor.

\begin{verbatim}
EBNF:
"f.math.interpolate", "(", list[real], ",", list[real], ",", real, [",", string], ")"
\end{verbatim}
\begin{itemize}
\item Operands: \texttt{xs} (x-coordinates, must be sorted ascending, unique), \texttt{ys} (y-coordinates), \texttt{x} (evaluation point), optional \texttt{method} (\texttt{"linear"} (default) or \texttt{"polynomial"}).
\item Result: \texttt{real} --- interpolated or extrapolated y-value.
\item A \texttt{Value Error} is raised if \texttt{len(xs)} $\neq$ \texttt{len(ys)}, if fewer than 2 data points, if fewer than 4 points with \texttt{"polynomial"} method, if x values are not unique, or if x values are not sorted.
\item A \texttt{Type Error} is raised if data lists contain non-numeric values.
\end{itemize}

\paragraph{Runge's Phenomenon}
For equally-spaced data points, high-degree polynomial interpolation can oscillate wildly near the edges. Use \texttt{"linear"} or \texttt{f.math.spline} instead for more than 10 data points.

\begin{verbatim}
# Linear interpolation
f.math.interpolate([0, 1, 2], [0, 1, 4], 1.5)                     #c# 2.25
f.math.interpolate([0, 1, 2], [0, 1, 4], 0.5)                     #c# 0.5
f.math.interpolate([0, 1, 2], [0, 1, 4], 3.0)                     #c# 7.0 (extrapolated)

# Polynomial interpolation (exact for quadratic data)
f.math.interpolate([0, 1, 2], [0, 1, 4], 1.5, "polynomial")       #c# 2.25
f.math.interpolate([0, 1, 2, 3], [0, 1, 4, 9], 2.5, "polynomial") #c# 6.25
\end{verbatim}

\subsubsection{Spline Interpolation}
The \texttt{f.math.spline} function performs cubic spline interpolation. Fits piecewise cubic polynomials between data points such that the result is continuous and has continuous first and second derivatives at all interior data points.

\textbf{Advantages over polynomial interpolation:} No Runge's oscillation. Smooth (continuous $f$, $f'$, $f''$). Stable for large numbers of data points.

\textbf{Boundary conditions:}
\begin{itemize}
\item \texttt{"natural"} (default) --- Zero second derivative at the endpoints: $f''(x_0) = 0$, $f''(x_n) = 0$. Produces the ``smoothest'' curve.
\item \texttt{"clamp"} --- Zero first derivative at the endpoints: $f'(x_0) = 0$, $f'(x_n) = 0$. Curve is flat at the boundaries.
\item \texttt{"periodic"} --- The function wraps around: $f(x_0) = f(x_n)$, $f'(x_0) = f'(x_n)$, $f''(x_0) = f''(x_n)$. For periodic data.
\end{itemize}

\textbf{Extrapolation:} Beyond the data range, the spline extends using the boundary slope of the nearest segment. No error is raised, but accuracy may be poor.

\begin{verbatim}
EBNF:
"f.math.spline", "(", list[real], ",", list[real], ",", real, [",", string], ")"
\end{verbatim}
\begin{itemize}
\item Operands: \texttt{xs} (x-coordinates, must be sorted ascending, unique), \texttt{ys} (y-coordinates), \texttt{x} (evaluation point), optional \texttt{boundary} (\texttt{"natural"} (default), \texttt{"clamp"}, or \texttt{"periodic"}).
\item Result: \texttt{real} --- interpolated or extrapolated y-value.
\item A \texttt{Value Error} is raised if \texttt{len(xs)} $\neq$ \texttt{len(ys)}, if fewer than 3 data points, if x values are not unique, if x values are not sorted, or if \texttt{boundary} is unknown.
\item A \texttt{Type Error} is raised if data lists contain non-numeric values.
\end{itemize}

\begin{verbatim}
# Natural spline
f.math.spline([0, 1, 2], [0, 1, 4], 1.5)                      #c# 2.25

# Clamped spline (flat at boundaries)
f.math.spline([0, 1, 2], [0, 1, 4], 1.5, "clamp")             #c# 2.25

# Periodic spline
f.math.spline([0, 1, 2, 3], [0, 1, 0, -1], 4.0, "periodic")   #c# 0.0 (wraps to start)
\end{verbatim}

\subsubsection{Adaptive Integration}
The \texttt{f.math.integrateAdaptive} function computes the definite integral $\int_a^b f(x)\,dx$ using adaptive Simpson quadrature. The algorithm subdivides intervals where the integrand changes rapidly, concentrating computation where it matters most.

\textbf{Algorithm:} Recursively apply Simpson's rule to each subinterval. If the estimated error exceeds $\text{tol}$, split the subinterval and recurse. Stops when the estimated total error is below $\text{tol}$.

\textbf{Accuracy:} Typically achieves the requested tolerance. The actual error is bounded by the estimated error. For smooth functions, the error is much smaller than the tolerance.

\textbf{Behavior:}
\begin{itemize}
\item If $a = b$, returns 0.0 (zero-width interval).
\item If $a > b$, computes $-\int_b^a f(x)\,dx$ (reversed bounds).
\item The function may be called many times (adaptive subdivision). Ensure $f$ is efficient.
\end{itemize}

\begin{verbatim}
EBNF:
"f.math.integrateAdaptive", "(", func, ",", real, ",", real, [",", real], ")"
\end{verbatim}
\begin{itemize}
\item Operands: \texttt{\$f} (function $f(x)$), \texttt{a} (lower bound), \texttt{b} (upper bound), optional \texttt{tol} (\texttt{real}, default $10^{-8}$).
\item Result: \texttt{real} --- approximate integral $\int_a^b f(x)\,dx$.
\item A \texttt{Runtime Error} is raised if the integral fails to converge (e.g., oscillatory or singular integrand).
\item A \texttt{Value Error} is raised if \texttt{tol} $\leq 0$.
\item A \texttt{Type Error} is raised if \texttt{\$f} is not a function reference.
\end{itemize}

\begin{verbatim}
$f sq = func($x) $x * $x
f.math.integrateAdaptive($sq, 0, 1)           #c# 0.33333 (1/3, exact)
f.math.integrateAdaptive($sq, 0, 1, 1e-6)    #c# 0.33333 (custom tolerance)
f.math.integrateAdaptive($sq, 1, 0)           #c# -0.33333 (reversed bounds)

# Integral of sin(x) from 0 to pi = 2.0
$f sinx = func($x) f.math.sin($x)
f.math.integrateAdaptive($sinx, 0, c.math.pi) #c# 2.0
\end{verbatim}

\subsubsection{Numerical Differentiation}
The \texttt{f.math.differentiateNumerically} function computes the first derivative $f'(x)$ using central differences: $f'(x) \approx (f(x+h) - f(x-h)) / (2h)$.

\textbf{Accuracy:} Central differences have error $O(h^2)$, which is more accurate than forward differences ($O(h)$). The optimal step size balances truncation error ($O(h^2)$) against floating-point roundoff error ($O(\epsilon/h)$), typically $h \approx \sqrt{\epsilon} \cdot |x| \approx 10^{-5}$ for unit-magnitude $x$.

\textbf{Behavior:}
\begin{itemize}
\item If $h < 0$, uses $|h|$ (no error, documented).
\item If $x = 0$, uses $h$ as the absolute step (not relative).
\item The function is called twice per evaluation ($f(x+h)$ and $f(x-h)$).
\end{itemize}

\begin{verbatim}
EBNF:
"f.math.differentiateNumerically", "(", func, ",", real, [",", real], ")"
\end{verbatim}
\begin{itemize}
\item Operands: \texttt{\$f} (function $f(x)$), \texttt{x} (evaluation point), optional \texttt{h} (step size, default $10^{-5}$).
\item Result: \texttt{real} --- approximate first derivative $f'(x)$.
\item A \texttt{Value Error} is raised if $h = 0$.
\item A \texttt{Type Error} is raised if \texttt{\$f} is not a function reference.
\end{itemize}

\begin{verbatim}
$f sq = func($x) $x * $x
f.math.differentiateNumerically($sq, 3)       #c# 6.0 (f'(3) = 2*3 = 6)
f.math.differentiateNumerically($sq, 3, 1e-8) #c# 6.0 (smaller step, same result)
f.math.differentiateNumerically($sq, 0)       #c# 0.0 (f'(0) = 0)
\end{verbatim}

\subsubsection{ODE Solver}
The \texttt{f.math.odeSolve} function solves a system of ordinary differential equations $\frac{dy}{dt} = f(t, y)$ over a time span $[t_0, t_{End}]$ with initial conditions $y(t_0) = y_0$.

\textbf{Output:} Returns a list of $[t, y]$ pairs sampled at regular intervals. For non-adaptive methods, the step count is 1000 (plus the start point), giving 1001 points. For the adaptive \texttt{"rkf45"} method, the step count and spacing vary based on the local error estimate.

\textbf{Stiff systems:} Stiff ODEs have widely separated time scales and require implicit methods for stability. Explicit methods (Euler, Heun, RK2/3/4) will either require extremely small step sizes or produce unstable (divergent) results when applied to stiff systems. Use implicit methods (\texttt{"eulerBack"}, \texttt{"midpointImp"}, \texttt{"glrk*"}) for stiff problems.

\paragraph{\$dydt Signature}
\begin{itemize}
\item For 1D (\texttt{y0} is \texttt{real}): \texttt{\$dydt(t, y)} $\rightarrow$ \texttt{real}.
\item For systems (\texttt{y0} is \texttt{list[real]}): \texttt{\$dydt(t, y)} $\rightarrow$ \texttt{list[real]}.
\item Return type must match \texttt{y0} type. A \texttt{Type Error} is raised on mismatch.
\end{itemize}

\begin{verbatim}
EBNF:
"f.math.odeSolve", "(", func, ",", real, ",", real, ",", ( real | list[real] ), [",", string], ")"
\end{verbatim}
\begin{itemize}
\item Operands: \texttt{\$dydt} (derivative function), \texttt{t0} (start time), \texttt{tEnd} (end time), \texttt{y0} (initial conditions: \texttt{real} for 1D, \texttt{list[real]} for systems), optional \texttt{method} (string, default \texttt{"rk4"}).
\item Result: \texttt{list[list[real]]} --- trajectory as \texttt{[[t0, y0], [t1, y1], ...]}. Always includes the start point.
\item If $t_0 = t_{End}$, returns \texttt{[[t0, y0]]} (single point).
\item If $t_0 > t_{End}$, integrates backwards in time (no error).
\item A \texttt{Runtime Error} is raised if the solver fails (e.g., stiff system with explicit method, or step failure).
\item A \texttt{Value Error} is raised if \texttt{y0} is empty or if \texttt{method} is unknown.
\item A \texttt{Type Error} is raised if \texttt{\$dydt} is not a function reference.
\end{itemize}

\paragraph{Explicit Methods}
Explicit methods compute $y_{n+1}$ directly from known values. Fast per step but may require very small step sizes for stiff systems.

\begin{itemize}
\item \texttt{"euler"} --- Forward Euler, order 1. Simplest method. $y_{n+1} = y_n + h \cdot f(t_n, y_n)$. First-order accurate: error $O(h)$. Not suitable for stiff systems. Use only for quick estimates or non-stiff problems.
\item \texttt{"heun"} --- Heun's method (improved Euler), order 2. Predictor-corrector: predicts with Euler, corrects with average slope. Second-order accurate: error $O(h^2)$. Better than Euler for same step size.
\item \texttt{"rk2"} --- Midpoint method (RK2), order 2. Evaluates the slope at the midpoint of the interval. Second-order accurate: error $O(h^2)$. More stable than Heun for oscillatory problems.
\item \texttt{"rk3"} --- Kutta's third-order method (RK3), order 3. Three-stage method. Third-order accurate: error $O(h^3)$. Good balance of accuracy and cost.
\item \texttt{"rk4"} --- Classic fourth-order Runge-Kutta (RK4), order 4. (default) Four-stage method. Fourth-order accurate: error $O(h^4)$. The standard choice for non-stiff ODEs. Good accuracy for moderate computational cost.
\item \texttt{"rkf45"} --- Runge-Kutta-Fehlberg, order 4(5) with adaptive step size. Computes both 4th and 5th order solutions and uses the difference to estimate local error. Adjusts step size to keep the error below a target tolerance (default $10^{-6}$). Step size bounds: minimum $10^{-12}$, maximum $(t_{End} - t_0) / 10$. May take fewer or more than 1000 steps depending on the problem.
\end{itemize}

\paragraph{Implicit Methods}
Implicit methods solve an equation at each step (typically using Newton iteration). More expensive per step but unconditionally stable for stiff systems.

\begin{itemize}
\item \texttt{"eulerBack"} --- Backward Euler, order 1. $y_{n+1} = y_n + h \cdot f(t_{n+1}, y_{n+1})$. First-order accurate. A-stable (stable for all step sizes). The simplest stiff-safe method. Recommended for very stiff systems where accuracy is not critical.
\item \texttt{"midpointImp"} --- Implicit Midpoint, order 2. Symplectic method (preserves energy for Hamiltonian systems). Second-order accurate. A-stable. Good for long-time integration of conservative systems.
\item \texttt{"glrk2"} --- Gauss-Legendre RK of order 2 (implicit midpoint rule). Second-order accurate. A-stable. Symplectic. One-stage implicit method.
\item \texttt{"glrk4"} --- Gauss-Legendre RK of order 4. Fourth-order accurate. A-stable. Symplectic. Two-stage implicit method. Good balance of accuracy and stability for moderate-stiffness systems.
\item \texttt{"glrk6"} --- Gauss-Legendre RK of order 6. Sixth-order accurate. A-stable. Symplectic. Three-stage implicit method. High accuracy for smooth solutions.
\item \texttt{"glrk8"} --- Gauss-Legendre RK of order 8. Eighth-order accurate. A-stable. Symplectic. Four-stage implicit method. Highest available accuracy. Use when extreme precision is required.
\end{itemize}

\paragraph{Method Selection Guide}

\begin{center}
\begin{tabular}{|l|c|c|c|}
\hline
\textbf{Method} & \textbf{Order} & \textbf{Adaptive} & \textbf{Stiff-safe} \\
\hline
\texttt{euler} & 1 & No & No \\
\texttt{heun} & 2 & No & No \\
\texttt{rk2} & 2 & No & No \\
\texttt{rk3} & 3 & No & No \\
\texttt{rk4} & 4 & No & No \\
\texttt{rkf45} & 4(5) & Yes & No \\
\hline
\texttt{eulerBack} & 1 & No & Yes \\
\texttt{midpointImp} & 2 & No & Yes \\
\texttt{glrk2} & 2 & No & Yes \\
\texttt{glrk4} & 4 & No & Yes \\
\texttt{glrk6} & 6 & No & Yes \\
\texttt{glrk8} & 8 & No & Yes \\
\hline
\end{tabular}
\end{center}

\textbf{Guidelines:}
\begin{itemize}
\item Start with \texttt{"rk4"} for non-stiff problems. It is reliable and well-understood.
\item Use \texttt{"rkf45"} when you need automatic step size control or when the solution has varying time scales.
\item Switch to \texttt{"eulerBack"} or \texttt{"glrk4"} if \texttt{"rk4"} produces unstable (divergent) results.
\item Use \texttt{"glrk8"} only when high accuracy is required and the function evaluations are cheap.
\item For Hamiltonian/conservative systems, prefer \texttt{"midpointImp"} or \texttt{"glrk*"} (symplectic methods preserve energy).
\end{itemize}

\paragraph{Examples}
\begin{verbatim}
# 1D: Exponential decay dy/dt = -y, y(0) = 1
$f decay = func($t, $y) -$y
$f sol = f.math.odeSolve($decay, 0, 5, 1.0)
# sol = [[0, 1.0], [0.005, 0.995], ..., [5.0, 0.00674]]

# Same with backward Euler (stiff-safe)
$f sol = f.math.odeSolve($decay, 0, 5, 1.0, "eulerBack")

# Same with adaptive RK45
$f sol = f.math.odeSolve($decay, 0, 5, 1.0, "rkf45")
# sol may have more or fewer than 1001 points

# System: Lotka-Volterra (predator-prey)
$f lotka = func($t, $y) [
    1.5 * $y[0] - 1.0 * $y[0] * $y[1],
    -3.0 * $y[1] + 1.0 * $y[0] * $y[1]
]
$f sol = f.math.odeSolve($lotka, 0, 10, [1.0, 1.0])
# sol = [[0, [1.0, 1.0]], [0.01, [1.015, 0.970]], ..., [10, [...]]]

# Backwards in time
$f sol = f.math.odeSolve($decay, 5, 0, 0.01)
# Integrates from t=5 to t=0
\end{verbatim}

\subsection{Shared Functions}

\subsubsection{Assert Function}
The \texttt{f.util.assert()} function checks a boolean condition and raises an error if \texttt{False}. See Section~\ref{sec:assert} for the full documentation (keyword and function forms, error type aliases, and examples).

\subsubsection{Length/Size Function}
The \texttt{len} function returns the length of a list or string, or the size of a set. \\
Morphology:
\begin{verbatim}
EBNF:
"f.util.len", "(", ( list | string | set ), ")"
\end{verbatim}

\paragraph{Edge Cases -- Collection Utilities}
\begin{itemize}
\item A \texttt{Type Error} is raised if the argument to \texttt{f.util.len} is not a list, string, or set.
\item A \texttt{Type Error} is raised if the arguments to \texttt{f.list.concat} are not all lists or all strings.
\item A \texttt{Type Error} is raised if the argument to \texttt{f.list.fromSet} is not a set.
\item A \texttt{Type Error} is raised if the argument to \texttt{f.list.set} is not a list.
\end{itemize}

Examples:
\begin{verbatim}
f.util.len([1, 2, 3])     #c# Returns 3
f.util.len("hello")       #c# Returns 5
f.util.len({Q1, Q2, R5})  #c# Returns 3
\end{verbatim}


\subsection{Special Functions}

\subsubsection{Print Function}
The \texttt{f.util.print()} function outputs its argument to the console. Unlike \texttt{echo}, it can evaluate expressions and function calls directly. \\
Morphology:
\begin{verbatim}
EBNF:
"f.util.print", "(", argument, ")"
\end{verbatim}
Example:
\begin{verbatim}
f.util.print("Hello, World!")
f.util.print(3 + 2 * 5)   #c# Outputs: 13
\end{verbatim}

\subsubsection{Type Function}
The \texttt{type} function returns the type of an object as a string. \\
Morphology:
\begin{verbatim}
EBNF:
"f.util.type", "(", argument, ")"
\end{verbatim}
Example:
\begin{verbatim}
f.util.type(42)       #c# Returns "int"
f.util.type(3.14)     #c# Returns "real"
f.util.type("hello")  #c# Returns "string"
f.util.type(Q1)       #c# Returns "part"
\end{verbatim}

\subsubsection{Reference Function}
The \texttt{ref} function converts a string containing a part reference into the referenced part. \\
Morphology:
\begin{verbatim}
EBNF:
"f.util.ref", "(", string, ")"
\end{verbatim}
Example:
\begin{verbatim}
f.util.ref("R1")   #c# Returns the part referenced as R1
\end{verbatim}

\subsubsection{No-Operation Function}
The \texttt{noop} function is the identity function. It returns its argument unchanged. \\
Morphology:
\begin{verbatim}
EBNF:
"f.util.noop", "(", argument, ")"
\end{verbatim}
Example:
\begin{verbatim}
f.util.noop(42)       #c# Returns 42
f.util.noop("hello")  #c# Returns "hello"
f.util.noop(Q1)       #c# Returns Q1
\end{verbatim}
Use \texttt{noop} as a placeholder in queries where no transformation is needed.

\subsubsection{Soft Integer Conversion}
The \texttt{softint} function converts its argument to an integer if it is either already an integer or is a real number numerically equal to an integer. Otherwise it returns the argument unchanged. \\
Morphology:
\begin{verbatim}
EBNF:
"f.util.softint", "(", argument, ")"
\end{verbatim}
Examples:
\begin{verbatim}
f.util.softint(42)       #c# Returns 42 (already int)
f.util.softint(5.0)      #c# Returns 5 (real equal to int)
f.util.softint(3.7)      #c# Returns 3.7 (not equal to int, unchanged)
f.util.softint("hello")  #c# Returns "hello" (not numeric, unchanged)
\end{verbatim}

\subsubsection{Soft Real Conversion}
The \texttt{softreal} function converts its argument to a real if it is either already a real, an integer, or a complex number with zero imaginary part. Otherwise it returns the argument unchanged. \\
Morphology:
\begin{verbatim}
EBNF:
"f.util.softreal", "(", argument, ")"
\end{verbatim}
Examples:
\begin{verbatim}
f.util.softreal(42)        #c# Returns 42.0 (int promoted to real)
f.util.softreal(3.14)      #c# Returns 3.14 (already real)
f.util.softreal(5 + 0i)    #c# Returns 5.0 (complex with zero imag part)
f.util.softreal(3 + 4i)
    #c# Returns 3 + 4i (complex with non-zero imag, unchanged)
f.util.softreal("hello")   #c# Returns "hello" (not numeric, unchanged)
\end{verbatim}

\subsection{Engineering Math Functions}

\subsubsection{Magnitude to dB Function}
Computes $10 \cdot \log_{10}(x)$, converting a power ratio to decibels. \\
Morphology:
\begin{verbatim}
EBNF:
"f.ee.db", "(", ( int | real ), ")"
\end{verbatim}
\begin{itemize}
\item Operands: \texttt{int} or \texttt{real} (must be $> 0$).
\item Result: \texttt{real}.
\item A \texttt{Value Error} is raised for $x \leq 0$.
\item A \texttt{Type Error} is raised if the argument is not numeric.
\end{itemize}
Examples:
\begin{verbatim}
f.ee.db(1)           #c# Returns 0.0
f.ee.db(10)          #c# Returns 10.0
f.ee.db(100)         #c# Returns 20.0
f.ee.db(0.5)         #c# Returns approximately -3.0103
\end{verbatim}

\subsubsection{dB10 Function}
Alias for \texttt{f.ee.db}: computes $10 \cdot \log_{10}(x)$. Same semantics, provided for clarity in power contexts. \\
Morphology:
\begin{verbatim}
EBNF:
"f.ee.db10", "(", ( int | real ), ")"
\end{verbatim}
\begin{itemize}
\item Operands: \texttt{int} or \texttt{real} (must be $> 0$).
\item Result: \texttt{real}.
\item A \texttt{Value Error} is raised for $x \leq 0$.
\end{itemize}
Examples:
\begin{verbatim}
f.db10(10)        #c# Returns 10.0
f.db10(100)       #c# Returns 20.0
\end{verbatim}

\subsubsection{dB20 Function}
Computes $20 \cdot \log_{10}(x)$, converting a voltage/current ratio to decibels. \\
Morphology:
\begin{verbatim}
EBNF:
"f.ee.db20", "(", ( int | real ), ")"
\end{verbatim}
\begin{itemize}
\item Operands: \texttt{int} or \texttt{real} (must be $> 0$).
\item Result: \texttt{real}.
\item A \texttt{Value Error} is raised for $x \leq 0$.
\end{itemize}
Examples:
\begin{verbatim}
f.db20(1)         #c# Returns 0.0
f.db20(10)        #c# Returns 20.0
f.db20(0.001)     #c# Returns -60.0
\end{verbatim}

\subsubsection{dB to Magnitude Function}
Converts a decibel value back to a magnitude (inverse of \texttt{f.ee.db}). Computes $10^{x/10}$. \\
Morphology:
\begin{verbatim}
EBNF:
"f.ee.fromDb", "(", ( int | real ), ")"
\end{verbatim}
\begin{itemize}
\item Operands: \texttt{int} or \texttt{real}.
\item Result: \texttt{real}.
\item A \texttt{Type Error} is raised if the argument is not numeric.
\end{itemize}
Examples:
\begin{verbatim}
f.fromDb(0)       #c# Returns 1.0
f.fromDb(20)      #c# Returns 100.0
f.fromDb(-3)      #c# Returns approximately 0.5012
\end{verbatim}

\subsubsection{Voltage Divider Function}
Computes $V_{\text{out}} = V_{\text{in}} \cdot \frac{r_2}{r_1 + r_2}$. \\
Morphology:
\begin{verbatim}
EBNF:
"f.ee.voltageDivider", "(", ( int | real ), ",", ( int | real ), ",", ( int | real ), ")"
\end{verbatim}
\begin{itemize}
\item Operands: $r_1$ (\texttt{int} or \texttt{real}), $r_2$ (\texttt{int} or \texttt{real}), $V_{\text{in}}$ (\texttt{int} or \texttt{real}).
\item Result: \texttt{real}.
\item A \texttt{Division by Zero Error} is raised if $r_1 + r_2 = 0$.
\item A \texttt{Type Error} is raised if arguments are not numeric.
\end{itemize}
Examples:
\begin{verbatim}
f.ee.voltageDivider(1000, 1000, 5)    #c# Returns 2.5
f.ee.voltageDivider(1000, 2000, 3.3)  #c# Returns 2.2
\end{verbatim}

\subsubsection{Current Divider Function}
Computes $I_1 = I \cdot \frac{r_2}{r_1 + r_2}$ (current through $r_1$ in a parallel circuit). \\
Morphology:
\begin{verbatim}
EBNF:
"f.ee.currentDivider", "(", ( int | real ), ",", ( int | real ), ",", ( int | real ), ")"
\end{verbatim}
\begin{itemize}
\item Operands: $r_1$ (\texttt{int} or \texttt{real}), $r_2$ (\texttt{int} or \texttt{real}), $I$ (\texttt{int} or \texttt{real}).
\item Result: \texttt{real}.
\item A \texttt{Division by Zero Error} is raised if $r_1 + r_2 = 0$.
\item A \texttt{Type Error} is raised if arguments are not numeric.
\end{itemize}
Examples:
\begin{verbatim}
f.ee.currentDivider(100, 100, 1)     #c# Returns 0.5
f.ee.currentDivider(100, 300, 0.01)  #c# Returns 0.0075
\end{verbatim}

\subsubsection{LC Resonance Frequency Function}
Computes $f = \frac{1}{2\pi\sqrt{LC}}$. \\
Morphology:
\begin{verbatim}
EBNF:
"f.ee.lcFrequency", "(", ( int | real ), ",", ( int | real ), ")"
\end{verbatim}
\begin{itemize}
\item Operands: $L$ (inductance, \texttt{int} or \texttt{real}), $C$ (capacitance, \texttt{int} or \texttt{real}).
\item Result: \texttt{real} (frequency in Hz).
\item A \texttt{Value Error} is raised if either argument is $\leq 0$.
\item A \texttt{Type Error} is raised if arguments are not numeric.
\end{itemize}
Examples:
\begin{verbatim}
f.ee.lcFrequency(0.001, 0.000001)    #c# Returns approximately 5032.9
f.ee.lcFrequency(10e-6, 100e-12)     #c# Returns approximately 503292.1
\end{verbatim}

\subsubsection{Reactance Function}
Computes the reactance $X = \omega L - \frac{1}{\omega C}$ for a series LC circuit. \\
Morphology:
\begin{verbatim}
EBNF:
"f.ee.reactance", "(", ( int | real ), ",", ( int | real ), ",", ( int | real ), ")"
\end{verbatim}
\begin{itemize}
\item Operands: frequency (\texttt{int} or \texttt{real}), $L$ (inductance, \texttt{int} or \texttt{real}), $C$ (capacitance, \texttt{int} or \texttt{real}).
\item Result: \texttt{real}.
\item A \texttt{Division by Zero Error} is raised if frequency $= 0$ or $C = 0$.
\item A \texttt{Type Error} is raised if arguments are not numeric.
\end{itemize}
Examples:
\begin{verbatim}
f.ee.reactance(1000, 0.001, 0.000001)  #c# Returns approximately 5.36
\end{verbatim}

\subsubsection{Unit Conversion Function}
Converts a numeric value from one unit to another within a specified physical dimension. The dimension parameter scopes the conversion, making cross-dimension conversions (e.g.\ Farad to °F) impossible by design. Temperature conversions are offset-based. Logarithmic power units (dBm, dBW) are supported within the \texttt{"power"} dimension. \\
Morphology:
\begin{verbatim}
EBNF:
"f.math.convert", "(", ( int | real ), ",", string, ",", string, ",", string, ")"
\end{verbatim}
\begin{itemize}
\item Operands: value (\texttt{int} or \texttt{real}), dimension (\texttt{string}), from\_unit (\texttt{string}), to\_unit (\texttt{string}).
\item Result: \texttt{real}.
\item Supported dimensions and their unit strings are listed in Table~\ref{tab:unit-strings}.
\item A \texttt{Value Error} is raised if the dimension string is not recognized, or if either unit string does not belong to the specified dimension.
\item A \texttt{Type Error} is raised if the value is not numeric.
\end{itemize}

Supported dimensions and unit strings:

\begin{table}[h]
\centering
\footnotesize
\begin{tabular}{|l|p{10cm}|l|}
\hline
\textbf{Dimension} & \textbf{Unit Strings} & \textbf{Base} \\
\hline
\texttt{"length"} & \texttt{"pl"} (Planck), \texttt{"fm"} (fermi), \texttt{"pm"}, \texttt{"A"} (angstrom), \texttt{"nm"}, \texttt{"um"} (micron), \texttt{"mm"}, \texttt{"cm"}, \texttt{"m"}, \texttt{"km"}, \texttt{"mil"}, \texttt{"pt"} / \texttt{"point"}, \texttt{"in"}, \texttt{"ft"}, \texttt{"survey\_foot"}, \texttt{"yd"}, \texttt{"mi"}, \texttt{"survey\_mile"}, \texttt{"nmi"}, \texttt{"au"} / \texttt{"astronomical\_unit"}, \texttt{"ly"} / \texttt{"light\_year"}, \texttt{"pc"} / \texttt{"parsec"} & m \\
\hline
\texttt{"mass"} & \texttt{"ug"}, \texttt{"mg"}, \texttt{"g"} / \texttt{"gram"}, \texttt{"grain"}, \texttt{"kg"}, \texttt{"carat"}, \texttt{"t"} / \texttt{"metric\_ton"}, \texttt{"oz"} / \texttt{"ounce"}, \texttt{"troy\_ounce"}, \texttt{"lb"} / \texttt{"pound"}, \texttt{"troy\_pound"}, \texttt{"stone"}, \texttt{"blob"} / \texttt{"slinch"}, \texttt{"slug"}, \texttt{"short\_ton"}, \texttt{"long\_ton"}, \texttt{"u"} / \texttt{"m\_u"} / \texttt{"amu"} / \texttt{"atomic\_mass"} & kg \\
\hline
\texttt{"time"} & \texttt{"ps"}, \texttt{"ns"}, \texttt{"us"}, \texttt{"ms"}, \texttt{"s"}, \texttt{"min"}, \texttt{"h"}, \texttt{"d"}, \texttt{"wk"}, \texttt{"yr"}, \texttt{"Julian\_year"} & s \\
\hline
\texttt{"frequency"} & \texttt{"mHz"}, \texttt{"Hz"}, \texttt{"kHz"}, \texttt{"MHz"}, \texttt{"GHz"}, \texttt{"THz"} & Hz \\
\hline
\texttt{"angle"} & \texttt{"deg"} / \texttt{"degree"}, \texttt{"rad"}, \texttt{"grad"}, \texttt{"arcmin"} / \texttt{"arcminute"}, \texttt{"arcsec"} / \texttt{"arcsecond"}, \texttt{"rev"} & rad \\
\hline
\texttt{"temperature"} & \texttt{"C"}, \texttt{"F"}, \texttt{"K"}, \texttt{"R"} (Rankine), \texttt{"zero\_Celsius"}, \texttt{"degree\_Fahrenheit"} & K (offset) \\
\hline
\texttt{"voltage"} & \texttt{"uV"}, \texttt{"mV"}, \texttt{"V"}, \texttt{"kV"}, \texttt{"MV"}, \texttt{"dBV"} & V \\
\hline
\texttt{"current"} & \texttt{"pA"}, \texttt{"nA"}, \texttt{"uA"}, \texttt{"mA"}, \texttt{"A"}, \texttt{"kA"} & A \\
\hline
\texttt{"resistance"} & \texttt{"uohm"}, \texttt{"mohm"}, \texttt{"ohm"}, \texttt{"kohm"}, \texttt{"Mohm"}, \texttt{"Gohm"} & $\Omega$ \\
\hline
\texttt{"capacitance"} & \texttt{"pF"}, \texttt{"nF"}, \texttt{"uF"}, \texttt{"mF"}, \texttt{"F"} & F \\
\hline
\texttt{"inductance"} & \texttt{"nH"}, \texttt{"uH"}, \texttt{"mH"}, \texttt{"H"} & H \\
\hline
\texttt{"charge"} & \texttt{"pC"}, \texttt{"nC"}, \texttt{"uC"}, \texttt{"mC"}, \texttt{"C"}, \texttt{"Ah"}, \texttt{"mAh"} & C \\
\hline
\texttt{"power"} & \texttt{"uW"}, \texttt{"mW"}, \texttt{"W"}, \texttt{"kW"}, \texttt{"MW"}, \texttt{"hp"} / \texttt{"horsepower"}, \texttt{"dBm"}, \texttt{"dBW"} & W \\
\hline
\texttt{"energy"} & \texttt{"eV"} / \texttt{"electron\_volt"}, \texttt{"keV"}, \texttt{"MeV"}, \texttt{"erg"}, \texttt{"J"}, \texttt{"kJ"}, \texttt{"cal"} / \texttt{"calorie\_th"}, \texttt{"calorie\_IT"}, \texttt{"kcal"}, \texttt{"BTU"} / \texttt{"Btu\_th"}, \texttt{"Btu\_IT"}, \texttt{"ton\_TNT"} & J \\
\hline
\texttt{"force"} & \texttt{"dyn"} / \texttt{"dyne"}, \texttt{"N"}, \texttt{"kN"}, \texttt{"lbf"} / \texttt{"pound\_force"}, \texttt{"kgf"} / \texttt{"kilogram\_force"} & N \\
\hline
\texttt{"pressure"} & \texttt{"Pa"}, \texttt{"kPa"}, \texttt{"MPa"}, \texttt{"bar"}, \texttt{"mbar"}, \texttt{"atm"} / \texttt{"atmosphere"}, \texttt{"psi"}, \texttt{"torr"} / \texttt{"mmHg"} & Pa \\
\hline
\texttt{"area"} & \texttt{"mm2"}, \texttt{"cm2"}, \texttt{"m2"}, \texttt{"km2"}, \texttt{"in2"}, \texttt{"ft2"}, \texttt{"ac"} / \texttt{"acre"}, \texttt{"ha"} / \texttt{"hectare"}, \texttt{"mi2"} & m$^2$ \\
\hline
\texttt{"volume"} & \texttt{"mm3"}, \texttt{"cm3"}, \texttt{"m3"}, \texttt{"mL"}, \texttt{"L"} / \texttt{"liter"} / \texttt{"litre"}, \texttt{"in3"}, \texttt{"ft3"}, \texttt{"gal"} / \texttt{"gallon\_US"}, \texttt{"gallon\_imp"}, \texttt{"fl\_oz"} / \texttt{"fluid\_ounce\_US"}, \texttt{"fluid\_ounce\_imp"}, \texttt{"bbl"} / \texttt{"barrel"} & m$^3$ \\
\hline
\texttt{"velocity"} & \texttt{"m/s"}, \texttt{"km/h"} / \texttt{"kmh"}, \texttt{"mph"}, \texttt{"kn"} / \texttt{"knot"}, \texttt{"mach"} / \texttt{"speed\_of\_sound"}, \texttt{"c"} & m/s \\
\hline
\texttt{"acceleration"} & \texttt{"m/s2"}, \texttt{"g"} & m/s$^2$ \\
\hline
\texttt{"voltage*time"} & \texttt{"Vs"}, \texttt{"mVs"}, \texttt{"uVs"}, \texttt{"Wb"} & V$\cdot$s \\
\hline
\end{tabular}
\caption{Supported dimensions and unit strings for \texttt{f.math.convert}.}
\label{tab:unit-strings}
\end{table}

Examples:
\begin{verbatim}
# Length: astronomical
f.math.convert(1, "length", "parsec", "light_year")   #c# Returns approximately 3.2616
f.math.convert(1, "length", "light_year", "km")      #c# Returns approximately 9.461e12
f.math.convert(1, "length", "au", "m")               #c# Returns approximately 1.496e11

# Length: imperial / PCB
f.math.convert(1, "length", "mile", "km")            #c# Returns approximately 1.609
f.math.convert(1, "length", "angstrom", "nm")        #c# Returns 0.1
f.math.convert(25.4, "length", "mil", "mm")          #c# Returns 25.4
f.math.convert(1, "length", "inch", "mm")            #c# Returns 25.4
f.math.convert(1, "length", "survey_foot", "ft")     #c# Returns approximately 1.000002

# Mass
f.math.convert(1, "mass", "pound", "kg")             #c# Returns approximately 0.4536
f.math.convert(1, "mass", "stone", "lb")             #c# Returns 14.0
f.math.convert(1, "mass", "slug", "kg")              #c# Returns approximately 14.594
f.math.convert(1, "mass", "carat", "mg")             #c# Returns 200.0
f.math.convert(1, "mass", "troy_ounce", "g")         #c# Returns approximately 31.103
f.math.convert(1, "mass", "long_ton", "kg")          #c# Returns approximately 1016.0
f.math.convert(1, "mass", "short_ton", "kg")         #c# Returns approximately 907.2

# Time
f.math.convert(1, "time", "hour", "s")              #c# Returns 3600.0
f.math.convert(1, "time", "day", "hour")            #c# Returns 24.0
f.math.convert(1, "time", "Julian_year", "day")     #c# Returns 365.25

# Electrical
f.math.convert(100, "resistance", "ohm", "kohm")    #c# Returns 0.1
f.math.convert(1, "frequency", "MHz", "Hz")         #c# Returns 1000000.0
f.math.convert(30, "power", "dBm", "mW")            #c# Returns 1000.0
f.math.convert(1, "power", "horsepower", "W")       #c# Returns approximately 745.7

# Force
f.math.convert(1, "force", "lbf", "N")             #c# Returns approximately 4.448
f.math.convert(1, "force", "dyne", "N")            #c# Returns 1e-5

# Energy
f.math.convert(1, "energy", "electron_volt", "J")   #c# Returns approximately 1.602e-19
f.math.convert(1, "energy", "calorie_IT", "J")      #c# Returns approximately 4.1868
f.math.convert(1, "energy", "Btu_IT", "J")          #c# Returns approximately 1055.06
f.math.convert(1, "energy", "erg", "J")             #c# Returns 1e-7
f.math.convert(1, "energy", "ton_TNT", "kJ")        #c# Returns 4184.0

# Temperature
f.math.convert(100, "temperature", "C", "F")        #c# Returns 212.0
f.math.convert(0, "temperature", "C", "K")          #c# Returns 273.15

# Pressure
f.math.convert(1, "pressure", "atm", "psi")         #c# Returns approximately 14.696
f.math.convert(1, "pressure", "torr", "Pa")         #c# Returns approximately 133.322

# Velocity
f.math.convert(1, "velocity", "mach", "m/s")        #c# Returns approximately 340.3
f.math.convert(1, "velocity", "knot", "mph")        #c# Returns approximately 1.151

# Volume
f.math.convert(1, "volume", "gallon_imp", "L")      #c# Returns approximately 4.546
f.math.convert(1, "volume", "barrel", "L")          #c# Returns approximately 159.0

# Angle
f.math.convert(1, "angle", "rev", "deg")            #c# Returns 360.0
f.math.convert(90, "angle", "degree", "arcminute")  #c# Returns 5400.0
\end{verbatim}

\subsubsection{Wavelength to Frequency}
Converts a wavelength (in meters) to optical frequency (in Hz) using $\nu = c / \lambda$, where $c = 299\,792\,458$ m/s. \\
Morphology:
\begin{verbatim}
EBNF:
"f.ee.lambda2nu", "(", ( int | real ), ")"
\end{verbatim}
\begin{itemize}
\item Operands: wavelength in meters (\texttt{int} or \texttt{real}).
\item Result: frequency in Hz (\texttt{real}).
\item A \texttt{Division by Zero Error} is raised if the wavelength is $0$.
\item A \texttt{Type Error} is raised if the argument is not numeric.
\end{itemize}
Examples:
\begin{verbatim}
f.ee.lambda2nu(500e-9)      #c# Returns approximately 5.996e14 (green light)
f.ee.lambda2nu(1.55e-6)     #c# Returns approximately 1.934e14 (C-band telecom)
f.ee.lambda2nu(10.6e-6)     #c# Returns approximately 2.828e13 (CO2 laser)
\end{verbatim}

\subsubsection{Frequency to Wavelength}
Converts an optical frequency (in Hz) to wavelength (in meters) using $\lambda = c / \nu$, where $c = 299\,792\,458$ m/s. \\
Morphology:
\begin{verbatim}
EBNF:
"f.ee.nu2lambda", "(", ( int | real ), ")"
\end{verbatim}
\begin{itemize}
\item Operands: frequency in Hz (\texttt{int} or \texttt{real}).
\item Result: wavelength in meters (\texttt{real}).
\item A \texttt{Division by Zero Error} is raised if the frequency is $0$.
\item A \texttt{Type Error} is raised if the argument is not numeric.
\end{itemize}
Examples:
\begin{verbatim}
f.ee.nu2lambda(5.996e14)    #c# Returns approximately 5.000e-07 (green light)
f.ee.nu2lambda(1.934e14)    #c# Returns approximately 1.550e-06 (C-band telecom)
f.ee.nu2lambda(1e9)         #c# Returns approximately 0.2998 (1 GHz radio)
\end{verbatim}

\subsubsection{Error Function}
Computes the Gauss error function $\text{erf}(x) = \frac{2}{\sqrt{\pi}} \int_0^x e^{-t^2} dt$. \\
Morphology:
\begin{verbatim}
EBNF:
"f.special.erf", "(", ( int | real ), ")"
\end{verbatim}
\begin{itemize}
\item Operands: \texttt{int} or \texttt{real}.
\item Result: \texttt{real}.
\item Odd function: $\text{erf}(-x) = -\text{erf}(x)$.
\item A \texttt{Type Error} is raised if the argument is not numeric.
\end{itemize}
Examples:
\begin{verbatim}
f.special.erf(0)         #c# Returns 0.0
f.special.erf(1)         #c# Returns approximately 0.8427
f.special.erf(-1)        #c# Returns approximately -0.8427
\end{verbatim}

\subsubsection{Complementary Error Function}
Computes $\text{erfc}(x) = 1 - \text{erf}(x)$. \\
Morphology:
\begin{verbatim}
EBNF:
"f.special.erfc", "(", ( int | real ), ")"
\end{verbatim}
\begin{itemize}
\item Operands: \texttt{int} or \texttt{real}.
\item Result: \texttt{real}.
\item A \texttt{Type Error} is raised if the argument is not numeric.
\end{itemize}
Examples:
\begin{verbatim}
f.special.erfc(0)        #c# Returns 1.0
f.special.erfc(1)        #c# Returns approximately 0.1573
\end{verbatim}

\subsubsection{Imaginary Error Function}
Computes $\text{erfi}(x) = -i \cdot \text{erf}(ix)$. \\
Morphology:
\begin{verbatim}
EBNF:
"f.special.erfi", "(", ( int | real ), ")"
\end{verbatim}
\begin{itemize}
\item Operands: \texttt{int} or \texttt{real}.
\item Result: \texttt{real}.
\item Odd function: $\text{erfi}(-x) = -\text{erfi}(x)$.
\item A \texttt{Type Error} is raised if the argument is not numeric.
\end{itemize}
Examples:
\begin{verbatim}
f.special.erfi(0)        #c# Returns 0.0
f.special.erfi(1)        #c# Returns approximately 1.6504
\end{verbatim}

\subsubsection{Bessel J Function}
Bessel function of the first kind $J_n(x)$. \\
Morphology:
\begin{verbatim}
EBNF:
"f.special.besselJ", "(", int, ",", real, ")"
\end{verbatim}
\begin{itemize}
\item Operands: order $n$ (\texttt{int}), argument $x$ (\texttt{int} or \texttt{real}).
\item Result: \texttt{real}.
\item A \texttt{Type Error} is raised if $n$ is not \texttt{int} or $x$ is not numeric.
\end{itemize}
Examples:
\begin{verbatim}
f.special.besselJ(0, 0)   #c# Returns 1.0
f.special.besselJ(0, 1)   #c# Returns approximately 0.7652
f.special.besselJ(1, 1)   #c# Returns approximately 0.4401
\end{verbatim}

\subsubsection{Bessel Y Function}
Bessel function of the second kind $Y_n(x)$ (Weber/Neumann function). \\
Morphology:
\begin{verbatim}
EBNF:
"f.special.besselY", "(", int, ",", real, ")"
\end{verbatim}
\begin{itemize}
\item Operands: order $n$ (\texttt{int}), argument $x$ (\texttt{int} or \texttt{real}).
\item Result: \texttt{real}.
\item A \texttt{Value Error} is raised for $x = 0$ (singularity).
\item A \texttt{Type Error} is raised if $n$ is not \texttt{int} or $x$ is not numeric.
\end{itemize}
Examples:
\begin{verbatim}
f.special.besselY(0, 1)   #c# Returns approximately 0.0883
\end{verbatim}

\subsubsection{Bessel I Function}
Modified Bessel function of the first kind $I_n(x)$. \\
Morphology:
\begin{verbatim}
EBNF:
"f.special.besselI", "(", int, ",", real, ")"
\end{verbatim}
\begin{itemize}
\item Operands: order $n$ (\texttt{int}), argument $x$ (\texttt{int} or \texttt{real}).
\item Result: \texttt{real}.
\item A \texttt{Type Error} is raised if $n$ is not \texttt{int} or $x$ is not numeric.
\end{itemize}
Examples:
\begin{verbatim}
f.special.besselI(0, 1)   #c# Returns approximately 1.2661
\end{verbatim}

\subsubsection{Bessel K Function}
Modified Bessel function of the second kind $K_n(x)$. \\
Morphology:
\begin{verbatim}
EBNF:
"f.special.besselK", "(", int, ",", real, ")"
\end{verbatim}
\begin{itemize}
\item Operands: order $n$ (\texttt{int}), argument $x$ (\texttt{int} or \texttt{real}).
\item Result: \texttt{real}.
\item A \texttt{Value Error} is raised for $x = 0$ (singularity).
\item A \texttt{Type Error} is raised if $n$ is not \texttt{int} or $x$ is not numeric.
\end{itemize}
Examples:
\begin{verbatim}
f.special.besselK(0, 1)   #c# Returns approximately 0.4210
\end{verbatim}

\subsubsection{Lambert W Function}
Computes the Lambert W function, the inverse of $x e^x$. \\
Morphology:
\begin{verbatim}
EBNF:
"f.special.lambertW", "(", real, [",", int], ")"
\end{verbatim}
\begin{itemize}
\item Operands: $x$ (\texttt{int} or \texttt{real}), optional branch $k$ (\texttt{int}, default 0).
\item Result: \texttt{real}.
\item \texttt{f.special.lambertW(x)} returns the principal branch $W_0(x)$.
\item \texttt{f.special.lambertW(x, k)} returns the branch $W_k(x)$.
\item For the principal branch ($k = 0$), a \texttt{Value Error} is raised for $x < -1/e$ (domain restriction).
\item A \texttt{Type Error} is raised if arguments are not numeric.
\end{itemize}
Examples:
\begin{verbatim}
f.special.lambertW(0)       #c# Returns 0.0
f.special.lambertW(1)       #c# Returns approximately 0.5671
f.special.lambertW(10, 0)   #c# Returns approximately 1.7455
\end{verbatim}

\subsubsection{Dilogarithm Function}
Computes the dilogarithm $\text{Li}_2(x) = -\int_0^x \frac{\ln(1-t)}{t} dt$. \\
Morphology:
\begin{verbatim}
EBNF:
"f.special.Li2", "(", ( int | real | complex ), ")"
\end{verbatim}
\begin{itemize}
\item Operands: \texttt{int}, \texttt{real}, or \texttt{complex}.
\item Result: \texttt{real} for real input, \texttt{complex} for complex input.
\item A \texttt{Type Error} is raised for non-numeric arguments.
\end{itemize}
Examples:
\begin{verbatim}
f.special.Li2(0)          #c# Returns 0.0
f.special.Li2(1)          #c# Returns approximately 1.6449  (= pi^2/6)
\end{verbatim}

\subsubsection{Trilogarithm Function}
Computes the trilogarithm $\text{Li}_3(x) = \sum_{k=1}^{\infty} \frac{x^k}{k^3}$. \\
Morphology:
\begin{verbatim}
EBNF:
"f.special.Li3", "(", ( int | real | complex ), ")"
\end{verbatim}
\begin{itemize}
\item Operands: \texttt{int}, \texttt{real}, or \texttt{complex}.
\item Result: \texttt{real} for real input, \texttt{complex} for complex input.
\item A \texttt{Type Error} is raised for non-numeric arguments.
\end{itemize}
Examples:
\begin{verbatim}
f.special.Li3(1)          #c# Returns approximately 1.2021  (= zeta(3))
\end{verbatim}

\subsubsection{General Polylogarithm Function}
Computes the polylogarithm $\text{Li}_s(x) = \sum_{k=1}^{\infty} \frac{x^k}{k^s}$ for arbitrary order $s$. \\
Morphology:
\begin{verbatim}
EBNF:
"f.special.Li", "(", ( int | real | complex ), ",", ( int | real | complex ), ")"
\end{verbatim}
\begin{itemize}
\item Operands: order $s$ (\texttt{int} or \texttt{real}), argument $x$ (\texttt{int}, \texttt{real}, or \texttt{complex}).
\item Result: \texttt{real} for real input, \texttt{complex} for complex input.
\item When $s = 2$, equivalent to \texttt{f.Li2}. When $s = 3$, equivalent to \texttt{f.Li3}.
\item A \texttt{Type Error} is raised for non-numeric arguments.
\end{itemize}
Examples:
\begin{verbatim}
f.special.Li(2, 0.5)      #c# Returns approximately 0.5822
f.special.Li(2, 1)        #c# Returns approximately 1.6449  (= f.special.Li2(1))
f.special.Li(3, 1)        #c# Returns approximately 1.2021  (= f.special.Li3(1))
\end{verbatim}
\begin{itemize}
\item \texttt{f.Li(s, x)} -- polylogarithm $\text{Li}_s(x) = \sum_{k=1}^{\infty} \frac{x^k}{k^s}$
\item \texttt{s} must be \texttt{int} ($\geq$ 1), \texttt{x} can be \texttt{int}, \texttt{real}, or \texttt{complex}.
\item \texttt{f.Li(2, x)} is equivalent to \texttt{f.Li2(x)}, \texttt{f.Li(3, x)} to \texttt{f.Li3(x)}.
\end{itemize}
\begin{verbatim}
f.special.Li(1, 0.5)      #c# Returns approximately 0.6931  (= -ln(1 - 0.5))
f.special.Li(2, 1)        #c# Returns approximately 1.6449  (= pi^2/6)
f.special.Li(4, 1)        #c# Returns approximately 1.0823  (= pi^4/90)
\end{verbatim}

\subsubsection{Digamma Function}
Computes the digamma function $\psi(x) = \frac{d}{dx} \ln(\Gamma(x))$. \\
Morphology:
\begin{verbatim}
EBNF:
"f.special.digamma", "(", ( int | real ), ")"
\end{verbatim}
\begin{itemize}
\item Operands: \texttt{int} or \texttt{real}.
\item Result: \texttt{real}.
\item A \texttt{Value Error} is raised for non-positive integers (poles of the gamma function).
\item A \texttt{Type Error} is raised if the argument is not numeric.
\end{itemize}
Examples:
\begin{verbatim}
f.special.digamma(1)       #c# Returns approximately -0.5772  (= -gamma)
f.special.digamma(2)       #c# Returns approximately 0.4228
\end{verbatim}

\subsubsection{Trigamma Function}
Computes the trigamma function $\psi'(x) = \frac{d^2}{dx^2} \ln(\Gamma(x))$. \\
Morphology:
\begin{verbatim}
EBNF:
"f.special.trigamma", "(", ( int | real ), ")"
\end{verbatim}
\begin{itemize}
\item Operands: \texttt{int} or \texttt{real}.
\item Result: \texttt{real}.
\item A \texttt{Value Error} is raised for non-positive integers (poles).
\item A \texttt{Type Error} is raised if the argument is not numeric.
\end{itemize}
Examples:
\begin{verbatim}
f.special.trigamma(1)      #c# Returns approximately 1.6449  (= pi^2/6)
\end{verbatim}

\subsubsection{Beta Function}
Computes the beta function $B(a,b) = \frac{\Gamma(a)\Gamma(b)}{\Gamma(a+b)}$. \\
Morphology:
\begin{verbatim}
EBNF:
"f.special.beta", "(", ( int | real ), ",", ( int | real ), ")"
\end{verbatim}
\begin{itemize}
\item Operands: $a$ (\texttt{int} or \texttt{real}), $b$ (\texttt{int} or \texttt{real}).
\item Result: \texttt{real}.
\item Symmetric: $B(a, b) = B(b, a)$.
\item A \texttt{Value Error} is raised if either argument is a non-positive integer (poles).
\item A \texttt{Type Error} is raised if arguments are not numeric.
\end{itemize}
Examples:
\begin{verbatim}
f.special.beta(1, 1)       #c# Returns 1.0
f.special.beta(2, 3)       #c# Returns approximately 0.0833  (= 1/12)
\end{verbatim}

\subsubsection{Legendre P Function}
Computes the Legendre polynomial $P_n(x)$. \\
Morphology:
\begin{verbatim}
EBNF:
"f.special.legendreP", "(", int, ",", real, ")"
\end{verbatim}
\begin{itemize}
\item Operands: degree $n$ (\texttt{int} $\geq 0$), argument $x$ (\texttt{int} or \texttt{real}).
\item Result: \texttt{real}.
\item A \texttt{Type Error} is raised if arguments are not the correct types.
\end{itemize}
Examples:
\begin{verbatim}
f.special.legendreP(0, 0.5)   #c# Returns 1.0
f.special.legendreP(1, 0.5)   #c# Returns 0.5
f.special.legendreP(2, 0.5)   #c# Returns -0.125
\end{verbatim}

\subsubsection{Legendre Q Function}
Computes the Legendre function of the second kind $Q_n(x)$. \\
Morphology:
\begin{verbatim}
EBNF:
"f.special.legendreQ", "(", int, ",", real, ")"
\end{verbatim}
\begin{itemize}
\item Operands: degree $n$ (\texttt{int} $\geq 0$), argument $x$ (\texttt{int} or \texttt{real}).
\item Result: \texttt{real}.
\item A \texttt{Value Error} is raised for $|x| \geq 1$ (singular at $\pm 1$).
\item A \texttt{Type Error} is raised if arguments are not the correct types.
\end{itemize}
Examples:
\begin{verbatim}
f.special.legendreQ(0, 2.0)   #c# Returns approximately 0.5493
\end{verbatim}

\subsubsection{Associated Legendre Function}
Computes the associated Legendre function $P_l^m(x)$. \\
Morphology:
\begin{verbatim}
EBNF:
"f.special.associatedLegendre", "(", int, ",", int, ",", real, ")"
\end{verbatim}
\begin{itemize}
\item Operands: degree $l$ (\texttt{int} $\geq 0$), order $m$ (\texttt{int}), argument $x$ (\texttt{int} or \texttt{real}).
\item Result: \texttt{real}.
\item If $|m| > l$, the result is \texttt{0.0} (mathematical definition).
\item A \texttt{Type Error} is raised if arguments are not the correct types.
\end{itemize}
Examples:
\begin{verbatim}
f.special.associatedLegendre(1, 0, 0.5)   #c# Returns 0.5
f.special.associatedLegendre(2, 1, 0.5)   #c# Returns approximately -0.9682
f.special.associatedLegendre(1, 2, 0.5)   #c# Returns 0.0 (|m| > l)
\end{verbatim}

\subsubsection{Hermite Polynomial Function}
Computes the physicist's Hermite polynomial $H_n(x)$. \\
Morphology:
\begin{verbatim}
EBNF:
"f.special.hermite", "(", int, ",", real, ")"
\end{verbatim}
\begin{itemize}
\item Operands: degree $n$ (\texttt{int} $\geq 0$), argument $x$ (\texttt{int} or \texttt{real}).
\item Result: \texttt{real}.
\item A \texttt{Type Error} is raised if arguments are not the correct types.
\end{itemize}
Examples:
\begin{verbatim}
f.special.hermite(0, 1)     #c# Returns 1.0
f.special.hermite(1, 1)     #c# Returns 2.0
f.special.hermite(2, 1)     #c# Returns 2.0
\end{verbatim}

\subsubsection{Laguerre Polynomial Function}
Computes the Laguerre polynomial $L_n(x)$. \\
Morphology:
\begin{verbatim}
EBNF:
"f.special.laguerre", "(", int, ",", real, ")"
\end{verbatim}
\begin{itemize}
\item Operands: degree $n$ (\texttt{int} $\geq 0$), argument $x$ (\texttt{int} or \texttt{real}).
\item Result: \texttt{real}.
\item A \texttt{Type Error} is raised if arguments are not the correct types.
\end{itemize}
Examples:
\begin{verbatim}
f.special.laguerre(0, 1)    #c# Returns 1.0
f.special.laguerre(1, 1)    #c# Returns 0.0
f.special.laguerre(2, 1)    #c# Returns 0.5
\end{verbatim}

\subsubsection{Chebyshev Polynomial Function}
Computes the Chebyshev polynomial of the first kind $T_n(x)$. \\
Morphology:
\begin{verbatim}
EBNF:
"f.special.chebyshev", "(", int, ",", real, ")"
\end{verbatim}
\begin{itemize}
\item Operands: degree $n$ (\texttt{int} $\geq 0$), argument $x$ (\texttt{int} or \texttt{real}).
\item Result: \texttt{real}.
\item A \texttt{Type Error} is raised if arguments are not the correct types.
\end{itemize}
Examples:
\begin{verbatim}
f.special.chebyshev(0, 0.5)  #c# Returns 1.0
f.special.chebyshev(1, 0.5)  #c# Returns 0.5
f.special.chebyshev(2, 0.5)  #c# Returns -0.5
\end{verbatim}

\subsubsection{Chebyshev Polynomial of the Second Kind}
Computes the Chebyshev polynomial of the second kind $U_n(x)$. \\
Morphology:
\begin{verbatim}
EBNF:
"f.special.chebyshevU", "(", int, ",", real, ")"
\end{verbatim}
\begin{itemize}
\item Operands: degree $n$ (\texttt{int} $\geq 0$), argument $x$ (\texttt{int} or \texttt{real}).
\item Result: \texttt{real}.
\item $U_0(x) = 1$, $U_1(x) = 2x$.
\end{itemize}
Examples:
\begin{verbatim}
f.special.chebyshevU(0, 0.5)  #c# Returns 1.0
f.special.chebyshevU(1, 0.5)  #c# Returns 1.0
f.special.chebyshevU(2, 0.5)  #c# Returns 3.0
\end{verbatim}

\subsubsection{Gegenbauer Polynomial}
Computes the Gegenbauer (ultraspherical) polynomial $C_n^{(\alpha)}(x)$. For $\alpha = 0.5$ this is the Legendre polynomial; for $\alpha = 1$ this is the Chebyshev polynomial of the second kind. \\
Morphology:
\begin{verbatim}
EBNF:
"f.special.gegenbauer", "(", int, ",", ( int | real ), ",", ( int | real ), ")"
\end{verbatim}
\begin{itemize}
\item Operands: degree $n$ (\texttt{int} $\geq 0$), parameter $\alpha$ (\texttt{int} or \texttt{real}), argument $x$ (\texttt{int} or \texttt{real}).
\item Result: \texttt{real}.
\end{itemize}
Examples:
\begin{verbatim}
f.special.gegenbauer(0, 1, 0.5)  #c# Returns 1.0
f.special.gegenbauer(1, 1, 0.5)  #c# Returns 1.0
\end{verbatim}

\subsubsection{Jacobi Polynomial}
Computes the Jacobi polynomial $P_n^{(\alpha, \beta)}(x)$. Generalizes Legendre, Chebyshev, and Gegenbauer polynomials. \\
Morphology:
\begin{verbatim}
EBNF:
"f.special.jacobiP", "(", int, ",", ( int | real ), ",", ( int | real ), ",", ( int | real ), ")"
\end{verbatim}
\begin{itemize}
\item Operands: degree $n$ (\texttt{int} $\geq 0$), parameters $\alpha$, $\beta$ (\texttt{int} or \texttt{real}), argument $x$ (\texttt{int} or \texttt{real}).
\item Result: \texttt{real}.
\end{itemize}
Examples:
\begin{verbatim}
f.special.jacobiP(0, 0, 0, 0.5)  #c# Returns 1.0
f.special.jacobiP(1, 1, 1, 0.5)  #c# Returns 0.75
\end{verbatim}

\subsubsection{Zernike Polynomial}
Computes the Zernike polynomial $R_n^m(r)$, used in optics for wavefront analysis over circular apertures. \\
Morphology:
\begin{verbatim}
EBNF:
"f.special.zernike", "(", int, ",", int, ",", ( int | real ), ")"
\end{verbatim}
\begin{itemize}
\item Operands: radial order $n$ (\texttt{int} $\geq 0$), azimuthal frequency $m$ (\texttt{int}, $|m| \leq n$), radius $r$ (\texttt{int} or \texttt{real}).
\item Result: \texttt{real}.
\item A \texttt{Value Error} is raised if $|m| > n$ or $(n - m)$ is odd and negative.
\end{itemize}
Examples:
\begin{verbatim}
f.special.zernike(0, 0, 0.5)  #c# Returns 1.0
f.special.zernike(1, 1, 0.5)  #c# Returns 0.5
f.special.zernike(2, 0, 0.5)  #c# Returns -0.125
\end{verbatim}

\subsubsection{Sine Integral Function}
The \texttt{si} function computes the sine integral $\text{Si}(x) = \int_0^x \frac{\sin(t)}{t}\,dt$. \\
Morphology:
\begin{verbatim}
EBNF:
"f.special.si", "(", ( int | real ), ")"
\end{verbatim}
\begin{itemize}
\item Operands: \texttt{int} or \texttt{real}
\item Result: \texttt{real}
\item $\text{Si}(0) = 0$. $\text{Si}(\infty) = \pi/2$. Odd function: $\text{Si}(-x) = -\text{Si}(x)$.
\end{itemize}
Examples:
\begin{verbatim}
f.special.si(0)          #c# Returns 0.0
f.special.si(1)          #c# Returns approximately 0.9461
f.special.si(3.14159)    #c# Returns approximately 1.8519 (near pi/2)
f.special.si(-1)         #c# Returns approximately -0.9461
\end{verbatim}

\subsubsection{Airy Function}
Computes the Airy functions $\text{Ai}(x)$ and $\text{Bi}(x)$, solutions to the differential equation $y'' - xy = 0$. Airy functions arise in optics (Airy disk diffraction patterns), quantum mechanics (linear potential wells), and electromagnetic wave propagation near caustics. \\
Morphology:
\begin{verbatim}
EBNF:
"f.special.airy", "(", ( int | real ), ")"
\end{verbatim}
\begin{itemize}
\item Operands: \texttt{int} or \texttt{real}.
\item Result: list \texttt{[real, real, real, real]} --- $[\text{Ai}(x), \text{Ai}'(x), \text{Bi}(x), \text{Bi}'(x)]$.
\item For complex arguments, see \texttt{f.special.airyC}.
\item A \texttt{Type Error} is raised if the argument is not numeric.
\end{itemize}
Examples:
\begin{verbatim}
f.special.airy(0)        #c# Returns [0.3551, -0.2588, 0.6149, 0.4483]
f.special.airy(1)        #c# Returns [0.1353, -0.1591, 1.2019, 0.9324]
f.special.airy(-1)       #c# Returns [0.5356, 0.1040, 0.1039, -0.4643]
\end{verbatim}
\begin{itemize}
\item $\text{Ai}(0) \approx 0.3551$, $\text{Bi}(0) \approx 0.6149$.
\item $\text{Ai}(x) \to 0$ and $\text{Bi}(x) \to \infty$ as $x \to +\infty$.
\item For large negative $x$, both functions oscillate with decaying amplitude.
\end{itemize}

\subsubsection{Airy Function (Complex)}
Computes the Airy functions for complex arguments. \\
Morphology:
\begin{verbatim}
EBNF:
"f.special.airyC", "(", complex, ")"
\end{verbatim}
\begin{itemize}
\item Operands: \texttt{int}, \texttt{real}, or \texttt{complex}.
\item Result: list \texttt{[complex, complex, complex, complex]} --- $[\text{Ai}(z), \text{Ai}'(z), \text{Bi}(z), \text{Bi}'(z)]$.
\item For real arguments, prefer \texttt{f.special.airy} which returns \texttt{real} values.
\item A \texttt{Type Error} is raised for non-numeric arguments.
\end{itemize}
Examples:
\begin{verbatim}
f.special.airyC(0)          #c# Returns [0.3551+0i, -0.2588+0i, 0.6149+0i, 0.4483+0i]
f.special.airyC(1+1i)       #c# Returns complex Airy values
\end{verbatim}

\subsubsection{Scaled Airy Function}
Computes exponentially scaled Airy functions $e^{\frac{2}{3} x^{3/2}} \text{Ai}(x)$ and $e^{-\frac{2}{3} |x|^{3/2}} \text{Bi}(x)$ to avoid overflow for large arguments. \\
Morphology:
\begin{verbatim}
EBNF:
"f.special.airyScaled", "(", ( int | real ), ")"
\end{verbatim}
\begin{itemize}
\item Operands: \texttt{int} or \texttt{real}.
\item Result: list \texttt{[real, real, real, real]} --- $[\text{AiS}(x), \text{AiS}'(x), \text{BiS}(x), \text{BiS}'(x)]$.
\item For $x \geq 0$: $\text{AiS}(x) = e^{\frac{2}{3} x^{3/2}} \text{Ai}(x)$.
\item For $x < 0$: $\text{AiS}(x) = \text{Ai}(x)$ (no scaling needed).
\item \text{BiS}(x) = $e^{-|\frac{2}{3} x^{3/2}|} \text{Bi}(x)$ for all $x$.
\item Useful when $|x|$ is large and unscaled \text{Ai} or \text{Bi} would overflow/underflow.
\item A \texttt{Type Error} is raised if the argument is not numeric.
\end{itemize}
Examples:
\begin{verbatim}
f.special.airyScaled(0)     #c# Returns [0.3551, -0.2588, 0.6149, 0.4483]
f.special.airyScaled(5)     #c# Returns scaled values avoiding overflow
f.special.airyScaled(-5)    #c# Returns same as f.special.airy(-5) (no scaling for x < 0)
\end{verbatim}

\subsubsection{Airy Function Zeros}
Computes zeros and values of the Airy functions $\text{Ai}$ and $\text{Bi}$ and their derivatives. \\
Morphology:
\begin{verbatim}
EBNF:
"f.special.airyZerosAi", "(", int, ")"
\end{verbatim}
\begin{itemize}
\item Operands: count $n$ (\texttt{int} $\geq 1$).
\item Result: list \texttt{[list, list]} --- $[a_n, a'_n]$ where $a_n$ are zeros of $\text{Ai}(x)$ and $a'_n$ are zeros of $\text{Ai}'(x)$.
\item A \texttt{Value Error} is raised if $n < 1$.
\item A \texttt{Type Error} is raised if the argument is not \texttt{int}.
\end{itemize}
Examples:
\begin{verbatim}
f.special.airyZerosAi(1)    #c# Returns [[-2.3381], [-1.0186]]
f.special.airyZerosAi(3)    #c# Returns [[-2.3381, -4.0879, -5.5206], [-1.0186, -3.2482, -4.8201]]
\end{verbatim}
\begin{itemize}
\item Zeros of $\text{Ai}(x)$ are all negative (oscillatory regime).
\item Zeros of $\text{Ai}'(x)$ interlace with zeros of $\text{Ai}(x)$.
\end{itemize}

\subsubsection{Bi Airy Function Zeros}
Computes zeros and values of the Airy functions $\text{Bi}$ and $\text{Bi}'$. \\
Morphology:
\begin{verbatim}
EBNF:
"f.special.airyZerosBi", "(", int, ")"
\end{verbatim}
\begin{itemize}
\item Operands: count $n$ (\texttt{int} $\geq 1$).
\item Result: list \texttt{[list, list]} --- $[b_n, b'_n]$ where $b_n$ are zeros of $\text{Bi}(x)$ and $b'_n$ are zeros of $\text{Bi}'(x)$.
\item A \texttt{Value Error} is raised if $n < 1$.
\item A \texttt{Type Error} is raised if the argument is not \texttt{int}.
\end{itemize}
Examples:
\begin{verbatim}
f.special.airyZerosBi(1)    #c# Returns [[-1.1737], [-2.2945]]
f.special.airyZerosBi(3)    #c# Returns [[-1.1737, -3.2711, -4.8307], [-2.2945, -4.1034, -5.5172]]
\end{verbatim}
\begin{itemize}
\item Zeros of $\text{Bi}(x)$ and $\text{Bi}'(x)$ are all negative.
\item Used in quantum mechanical bound state problems with linear potentials.
\end{itemize}

\subsubsection{Airy Integrals}
Computes integrals of the Airy functions: $\int_0^x \text{Ai}(t)\,dt$ and $\int_0^x \text{Bi}(t)\,dt$. \\
Morphology:
\begin{verbatim}
EBNF:
"f.special.airyIntegral", "(", ( int | real ), ")"
\end{verbatim}
\begin{itemize}
\item Operands: \texttt{int} or \texttt{real}.
\item Result: list \texttt{[real, real]} --- $[\int_0^x \text{Ai}(t)\,dt,\; \int_0^x \text{Bi}(t)\,dt]$.
\item $\int_0^0 = [0.0, 0.0]$.
\item A \texttt{Type Error} is raised if the argument is not numeric.
\end{itemize}
Examples:
\begin{verbatim}
f.special.airyIntegral(0)     #c# Returns [0.0, 0.0]
f.special.airyIntegral(1)     #c# Returns approximately [0.2018, 0.9640]
f.special.airyIntegral(-1)    #c# Returns approximately [0.2310, -0.1199]
\end{verbatim}
\begin{itemize}
\item Useful in diffraction theory and quantum scattering calculations.
\end{itemize}

\subsubsection{Scorer's Function Hi}
Computes the Scorer's function $\text{Hi}(z)$, a particular solution to the inhomogeneous Airy equation $y'' - xy = \frac{1}{\pi}$. Scorer's functions arise in diffraction theory, wave scattering, and electromagnetic propagation near caustics. \\
Morphology:
\begin{verbatim}
EBNF:
"f.special.scorerHi", "(", ( int | real ), ")"
\end{verbatim}
\begin{itemize}
\item Operands: \texttt{x} (\texttt{int} or \texttt{real}).
\item Result: \texttt{real} --- value of $\text{Hi}(x)$.
\item For complex arguments, see \texttt{f.special.scorerHiC}.
\end{itemize}
\begin{verbatim}
f.special.scorerHi(0)       #c# Returns approximately 0.5
f.special.scorerHi(1)       #c# Returns approximately 0.3016
f.special.scorerHi(-1)      #c# Returns approximately 0.4643
\end{verbatim}

\subsubsection{Scorer's Function Gi}
Computes the Scorer's function $\text{Gi}(z)$, the complementary particular solution to the inhomogeneous Airy equation. Related to $\text{Hi}$ by $\text{Gi}(z) = \text{Bi}(z) \int_0^z \text{Ai}(t)\,dt - \text{Ai}(z) \int_0^z \text{Bi}(t)\,dt$. \\
Morphology:
\begin{verbatim}
EBNF:
"f.special.scorerGi", "(", ( int | real ), ")"
\end{verbatim}
\begin{itemize}
\item Operands: \texttt{x} (\texttt{int} or \texttt{real}).
\item Result: \texttt{real} --- value of $\text{Gi}(x)$.
\item For complex arguments, see \texttt{f.special.scorerGiC}.
\end{itemize}
\begin{verbatim}
f.special.scorerGi(0)       #c# Returns approximately 0.5
f.special.scorerGi(1)       #c# Returns approximately 0.1606
f.special.scorerGi(-1)      #c# Returns approximately 0.4643
\end{verbatim}

\subsubsection{Scorer's Function Hi (Complex)}
Computes the Scorer's function $\text{Hi}(z)$ for complex arguments. \\
Morphology:
\begin{verbatim}
EBNF:
"f.special.scorerHiC", "(", complex, ")"
\end{verbatim}
\begin{itemize}
\item Operands: \texttt{z} (\texttt{complex}).
\item Result: \texttt{complex} --- value of $\text{Hi}(z)$.
\item For real arguments, prefer \texttt{f.special.scorerHi} which returns \texttt{real} values.
\end{itemize}
\begin{verbatim}
f.special.scorerHiC(0)          #c# Returns approximately 0.5+0i
f.special.scorerHiC(1+1i)       #c# Returns complex Scorer Hi value
\end{verbatim}

\subsubsection{Scorer's Function Gi (Complex)}
Computes the Scorer's function $\text{Gi}(z)$ for complex arguments. \\
Morphology:
\begin{verbatim}
EBNF:
"f.special.scorerGiC", "(", complex, ")"
\end{verbatim}
\begin{itemize}
\item Operands: \texttt{z} (\texttt{complex}).
\item Result: \texttt{complex} --- value of $\text{Gi}(z)$.
\item For real arguments, prefer \texttt{f.special.scorerGi} which returns \texttt{real} values.
\end{itemize}
\begin{verbatim}
f.special.scorerGiC(0)          #c# Returns approximately 0.5+0i
f.special.scorerGiC(1+1i)       #c# Returns complex Scorer Gi value
\end{verbatim}

\subsubsection{Align Up Function}
Rounds $x$ up to the next multiple of \texttt{alignment}. \\
Morphology:
\begin{verbatim}
EBNF:
"f.int.alignUp", "(", int, ",", int, ")"
\end{verbatim}
\begin{itemize}
\item Operands: $x$ (\texttt{int}), alignment (\texttt{int}).
\item Result: \texttt{int}.
\item A \texttt{Value Error} is raised if alignment $\leq 0$.
\item A \texttt{Type Error} is raised if arguments are not \texttt{int}.
\end{itemize}
Examples:
\begin{verbatim}
f.alignUp(7, 8)       #c# Returns 8
f.alignUp(9, 8)       #c# Returns 16
f.alignUp(8, 8)       #c# Returns 8 (already aligned)
\end{verbatim}

\subsubsection{Align Down Function}
Rounds $x$ down to the previous multiple of \texttt{alignment}. \\
Morphology:
\begin{verbatim}
EBNF:
"f.int.alignDown", "(", int, ",", int, ")"
\end{verbatim}
\begin{itemize}
\item Operands: $x$ (\texttt{int}), alignment (\texttt{int}).
\item Result: \texttt{int}.
\item A \texttt{Value Error} is raised if alignment $\leq 0$.
\item A \texttt{Type Error} is raised if arguments are not \texttt{int}.
\end{itemize}
Examples:
\begin{verbatim}
f.alignDown(7, 8)     #c# Returns 0
f.alignDown(15, 8)    #c# Returns 8
\end{verbatim}

\subsubsection{Is Power of Two Function}
Returns \texttt{True} if $n$ is a power of two. \\
Morphology:
\begin{verbatim}
EBNF:
"f.int.isPowerOfTwo", "(", int, ")"
\end{verbatim}
\begin{itemize}
\item Operands: \texttt{int}.
\item Result: \texttt{bool}.
\item A \texttt{Type Error} is raised if the argument is not \texttt{int}.
\end{itemize}
Examples:
\begin{verbatim}
f.isPowerOfTwo(8)      #c# Returns True
f.isPowerOfTwo(7)      #c# Returns False
\end{verbatim}

\subsubsection{Next Power of Two Function}
Returns the smallest power of two $\geq n$. \\
Morphology:
\begin{verbatim}
EBNF:
"f.int.nextPowerOfTwo", "(", int, ")"
\end{verbatim}
\begin{itemize}
\item Operands: \texttt{int} (must be $> 0$).
\item Result: \texttt{int}.
\item A \texttt{Value Error} is raised for $n \leq 0$.
\item A \texttt{Type Error} is raised if the argument is not \texttt{int}.
\end{itemize}
Examples:
\begin{verbatim}
f.nextPowerOfTwo(7)    #c# Returns 8
f.nextPowerOfTwo(8)    #c# Returns 8 (already a power of two)
\end{verbatim}

\subsubsection{Popcount Function}
Returns the number of set bits (1-bits) in the two's complement representation of $n$. \\
Morphology:
\begin{verbatim}
EBNF:
"f.int.popcount", "(", int, ")"
\end{verbatim}
\begin{itemize}
\item Operands: \texttt{int}.
\item Result: \texttt{int}.
\item For negative numbers, operates on infinite-width two's complement (returns $\infty$).
\item A \texttt{Type Error} is raised if the argument is not \texttt{int}.
\end{itemize}
Examples:
\begin{verbatim}
f.int.popcount(7)       #c# Returns 3  (binary 111)
f.int.popcount(8)       #c# Returns 1  (binary 1000)
\end{verbatim}

\subsubsection{Parity Function}
Returns 0 if the number of 1-bits is even, 1 if odd. \\
Morphology:
\begin{verbatim}
EBNF:
"f.int.parity", "(", int, ")"
\end{verbatim}
\begin{itemize}
\item Operands: \texttt{int}.
\item Result: \texttt{int} (0 or 1).
\item For negative numbers, operates on infinite-width two's complement.
\item A \texttt{Type Error} is raised if the argument is not \texttt{int}.
\end{itemize}
Examples:
\begin{verbatim}
f.int.parity(7)         #c# Returns 1  (odd: three 1-bits)
f.int.parity(15)        #c# Returns 0  (even: four 1-bits)
\end{verbatim}

\subsubsection{Bit Length Function}
Returns the number of bits required to represent $n$ in binary (excluding the sign bit). \\
Morphology:
\begin{verbatim}
EBNF:
"f.int.bitLength", "(", int, ")"
\end{verbatim}
\begin{itemize}
\item Operands: \texttt{int}.
\item Result: \texttt{int}.
\item \texttt{f.int.bitLength(0)} returns 0.
\item For negative numbers, returns the bit length of the absolute value.
\item A \texttt{Type Error} is raised if the argument is not \texttt{int}.
\end{itemize}
Examples:
\begin{verbatim}
f.int.bitLength(0)      #c# Returns 0
f.int.bitLength(8)      #c# Returns 4
f.int.bitLength(-8)     #c# Returns 4
\end{verbatim}

\subsubsection{Integer Log2 Function}
Returns the integer floor of the base-2 logarithm. Equivalent to $\lfloor \log_2(n) \rfloor$. \\
Morphology:
\begin{verbatim}
EBNF:
"f.int.ilog2", "(", int, ")"
\end{verbatim}
\begin{itemize}
\item Operands: \texttt{int} (must be $> 0$).
\item Result: \texttt{int}.
\item A \texttt{Value Error} is raised if the argument is $\leq 0$.
\item A \texttt{Type Error} is raised if the argument is not \texttt{int}.
\end{itemize}
Examples:
\begin{verbatim}
f.int.ilog2(1)       #c# Returns 0   (2^0 = 1)
f.int.ilog2(8)       #c# Returns 3   (2^3 = 8)
f.int.ilog2(9)       #c# Returns 3   (floor of 3.17)
f.int.ilog2(1023)    #c# Returns 9   (2^9 = 512)
f.int.ilog2(1024)    #c# Returns 10  (2^10 = 1024)
\end{verbatim}

\subsubsection{Ceiling Log2 Function}
Returns the smallest integer $k$ such that $2^k \geq n$. Equivalent to $\lceil \log_2(n) \rceil$. \\
Morphology:
\begin{verbatim}
EBNF:
"f.int.ceilLog2", "(", int, ")"
\end{verbatim}
\begin{itemize}
\item Operands: \texttt{int} (must be $> 0$).
\item Result: \texttt{int}.
\item A \texttt{Value Error} is raised if the argument is $\leq 0$.
\item A \texttt{Type Error} is raised if the argument is not \texttt{int}.
\end{itemize}
Examples:
\begin{verbatim}
f.int.ceilLog2(1)       #c# Returns 0   (2^0 = 1)
f.int.ceilLog2(8)       #c# Returns 3   (2^3 = 8)
f.int.ceilLog2(9)       #c# Returns 4   (2^4 = 16 >= 9)
f.int.ceilLog2(1023)    #c# Returns 10  (2^10 = 1024)
f.int.ceilLog2(1024)    #c# Returns 10  (2^10 = 1024)
\end{verbatim}

\subsubsection{ToBase Function}
Converts an integer to a string representation in the given base (2--36). \\
Morphology:
\begin{verbatim}
EBNF:
"f.int.toBase", "(", int, ",", int, ")"
\end{verbatim}
\begin{itemize}
\item Operands: value (\texttt{int}), base (\texttt{int}, 2--36).
\item Result: \texttt{string} (lowercase digits).
\item Negative values produce a \texttt{"-"} prefix.
\item A \texttt{Value Error} is raised for base outside 2--36.
\item A \texttt{Type Error} is raised if arguments are not \texttt{int}.
\end{itemize}
Examples:
\begin{verbatim}
f.toBase(255, 16)     #c# Returns "ff"
f.toBase(255, 2)      #c# Returns "11111111"
f.toBase(-10, 16)     #c# Returns "-a"
\end{verbatim}

\subsubsection{FromBase Function}
Parses a string in the given base to an integer. \\
Morphology:
\begin{verbatim}
EBNF:
"f.int.fromBase", "(", string, ",", int, ")"
\end{verbatim}
\begin{itemize}
\item Operands: string (\texttt{string}), base (\texttt{int}, 2--36).
\item Result: \texttt{int}.
\item The \texttt{"-"} prefix is handled for negative values.
\item A \texttt{Value Error} is raised for invalid characters or base outside 2--36.
\item A \texttt{Type Error} is raised if the first argument is not a \texttt{string} or the second is not an \texttt{int}.
\end{itemize}
Examples:
\begin{verbatim}
f.fromBase("ff", 16)  #c# Returns 255
f.fromBase("1010", 2) #c# Returns 10
f.fromBase("-a", 16)  #c# Returns -10
\end{verbatim}

\subsubsection{Bin Function}
Converts an integer to a binary string (base 2). \\
Morphology:
\begin{verbatim}
EBNF:
"f.int.bin", "(", int, ")"
\end{verbatim}
\begin{itemize}
\item Operands: \texttt{int}.
\item Result: \texttt{string}.
\item Equivalent to \texttt{f.toBase(x, 2)}.
\item A \texttt{Type Error} is raised if the argument is not \texttt{int}.
\end{itemize}
Examples:
\begin{verbatim}
f.int.bin(10)          #c# Returns "1010"
f.int.bin(255)         #c# Returns "11111111"
\end{verbatim}

\subsubsection{Oct Function}
Converts an integer to an octal string (base 8). \\
Morphology:
\begin{verbatim}
EBNF:
"f.int.oct", "(", int, ")"
\end{verbatim}
\begin{itemize}
\item Operands: \texttt{int}.
\item Result: \texttt{string}.
\item Equivalent to \texttt{f.toBase(x, 8)}.
\item A \texttt{Type Error} is raised if the argument is not \texttt{int}.
\end{itemize}
Examples:
\begin{verbatim}
f.int.oct(10)          #c# Returns "12"
\end{verbatim}

\subsubsection{Hex Function}
Converts an integer to a hexadecimal string (base 16). \\
Morphology:
\begin{verbatim}
EBNF:
"f.int.hex", "(", int, ")"
\end{verbatim}
\begin{itemize}
\item Operands: \texttt{int}.
\item Result: \texttt{string} (lowercase).
\item Equivalent to \texttt{f.toBase(x, 16)}.
\item A \texttt{Type Error} is raised if the argument is not \texttt{int}.
\end{itemize}
Examples:
\begin{verbatim}
f.int.hex(255)         #c# Returns "ff"
\end{verbatim}

\subsubsection{FromBin Function}
Parses a binary string to an integer. \\
Morphology:
\begin{verbatim}
EBNF:
"f.int.fromBin", "(", string, ")"
\end{verbatim}
\begin{itemize}
\item Operands: \texttt{string}.
\item Result: \texttt{int}.
\item Equivalent to \texttt{f.fromBase(s, 2)}.
\item A \texttt{Type Error} is raised if the argument is not a \texttt{string}.
\end{itemize}
Examples:
\begin{verbatim}
f.fromBin("1010")  #c# Returns 10
\end{verbatim}

\subsubsection{FromOct Function}
Parses an octal string to an integer. \\
Morphology:
\begin{verbatim}
EBNF:
"f.int.fromOct", "(", string, ")"
\end{verbatim}
\begin{itemize}
\item Operands: \texttt{string}.
\item Result: \texttt{int}.
\item Equivalent to \texttt{f.fromBase(s, 8)}.
\item A \texttt{Type Error} is raised if the argument is not a \texttt{string}.
\end{itemize}
Examples:
\begin{verbatim}
f.fromOct("12")    #c# Returns 10
\end{verbatim}

\subsubsection{FromHex Function}
Parses a hexadecimal string to an integer. \\
Morphology:
\begin{verbatim}
EBNF:
"f.int.fromHex", "(", string, ")"
\end{verbatim}
\begin{itemize}
\item Operands: \texttt{string}.
\item Result: \texttt{int}.
\item Equivalent to \texttt{f.fromBase(s, 16)}.
\item A \texttt{Type Error} is raised if the argument is not a \texttt{string}.
\end{itemize}
Examples:
\begin{verbatim}
f.fromHex("ff")    #c# Returns 255
\end{verbatim}

\subsubsection{Choice and Majority Functions}
\paragraph{Choice (Ch)}
\begin{verbatim}
EBNF:
"f.int.ch", "(", ( bool | int ), ",", ( bool | int ), ",", ( bool | int ), ")"
\end{verbatim}
\begin{itemize}
\item Operands: three values of the same type
\item For \texttt{bool}: $\text{Ch}(x, y, z) = (x \land y) \oplus (\lnot x \land z)$ -- if $x$ then $y$ else $z$
\item For \texttt{int}: bitwise version -- $\text{Ch}(x, y, z) = (x \mathbin{\&} y) \mathbin{\hat{}\hat{}} (\mathbin{b!\&}(x) \mathbin{\&} z)$
\item A \texttt{Type Error} is raised if the three arguments are not the same type, or if the type is not \texttt{bool} or \texttt{int}
\end{itemize}
Examples:
\begin{verbatim}
f.int.ch(True, True, False)      #c# Returns True
f.int.ch(False, True, False)     #c# Returns False
f.int.ch(True, False, True)      #c# Returns False
f.int.ch(0xFF, 0x0F, 0x00)       #c# Returns 0x0F (bitwise)
\end{verbatim}

\paragraph{Majority (Maj)}
\begin{verbatim}
EBNF:
"f.int.maj", "(", ( bool | int ), ",", ( bool | int ), ",", ( bool | int ), ")"
\end{verbatim}
\begin{itemize}
\item Operands: three values of the same type
\item For \texttt{bool}: $\text{Maj}(x, y, z) = (x \land y) \oplus (x \land z) \oplus (y \land z)$ -- majority vote
\item For \texttt{int}: bitwise version -- $\text{Maj}(x, y, z) = (x \mathbin{\&} y) \mathbin{\hat{}\hat{}} (x \mathbin{\&} z) \mathbin{\hat{}\hat{}} (y \mathbin{\&} z)$
\item A \texttt{Type Error} is raised if the three arguments are not the same type, or if the type is not \texttt{bool} or \texttt{int}
\end{itemize}
Examples:
\begin{verbatim}
f.int.maj(True, True, False)     #c# Returns True
f.int.maj(True, False, False)    #c# Returns False
f.int.maj(False, True, True)     #c# Returns True
f.int.maj(0xFF, 0x0F, 0x33)      #c# Returns 0x33 (bitwise)
\end{verbatim}

\subsubsection{Is Prime}
Tests whether an integer is prime. \\
Morphology:
\begin{verbatim}
EBNF:
"f.int.isPrime", "(", int, ")"
\end{verbatim}
\begin{itemize}
\item Operand: \texttt{int}.
\item Result: \texttt{bool}. Returns \texttt{False} for $n < 2$.
\end{itemize}
Examples:
\begin{verbatim}
f.int.isPrime(2)        #c# Returns True
f.int.isPrime(17)       #c# Returns True
f.int.isPrime(1)        #c# Returns False
f.int.isPrime(-5)       #c# Returns False
\end{verbatim}

\subsubsection{Next Prime}
Returns the smallest prime strictly greater than $n$. \\
Morphology:
\begin{verbatim}
EBNF:
"f.int.nextPrime", "(", int, ")"
\end{verbatim}
\begin{itemize}
\item Operand: \texttt{int}.
\item Result: \texttt{int}. Returns $2$ if $n < 2$.
\item A \texttt{Value Error} is raised if $n < 0$.
\end{itemize}
Examples:
\begin{verbatim}
f.int.nextPrime(1)        #c# Returns 2
f.int.nextPrime(10)       #c# Returns 11
f.int.nextPrime(17)       #c# Returns 19
\end{verbatim}

\subsubsection{Previous Prime}
Returns the largest prime strictly less than $n$. \\
Morphology:
\begin{verbatim}
EBNF:
"f.int.previousPrime", "(", int, ")"
\end{verbatim}
\begin{itemize}
\item Operand: \texttt{int}.
\item Result: \texttt{int}.
\item A \texttt{Value Error} is raised if $n \leq 2$ (no prime less than 2).
\end{itemize}
Examples:
\begin{verbatim}
f.int.previousPrime(10)     #c# Returns 7
f.int.previousPrime(17)     #c# Returns 13
\end{verbatim}

\subsubsection{Nearest Prime}
Returns the prime nearest to $n$. For ties (e.g.\ $n = 6$, equidistant from 5 and 7), returns the smaller prime. \\
Morphology:
\begin{verbatim}
EBNF:
"f.int.nearestPrime", "(", int, ")"
\end{verbatim}
\begin{itemize}
\item Operand: \texttt{int}.
\item Result: \texttt{int}. Returns $2$ if $n < 2$.
\item A \texttt{Value Error} is raised if $n < 0$.
\end{itemize}
Examples:
\begin{verbatim}
f.int.nearestPrime(10)     #c# Returns 11
f.int.nearestPrime(6)      #c# Returns 5  (tie, returns smaller)
f.int.nearestPrime(12)     #c# Returns 11
\end{verbatim}

\subsubsection{Nth Prime}
Returns the $n$-th prime (1-indexed: $p_1 = 2$, $p_2 = 3$, $p_3 = 5$). \\
Morphology:
\begin{verbatim}
EBNF:
"f.int.nthPrime", "(", int, ")"
\end{verbatim}
\begin{itemize}
\item Operand: \texttt{int} ($\geq 1$).
\item Result: \texttt{int}.
\item A \texttt{Value Error} is raised if $n < 1$.
\end{itemize}
Examples:
\begin{verbatim}
f.int.nthPrime(1)        #c# Returns 2
f.int.nthPrime(10)       #c# Returns 29
f.int.nthPrime(100)      #c# Returns 541
\end{verbatim}

\subsubsection{Prime Pi}
Computes $\pi(n)$, the number of primes less than or equal to $n$. \\
Morphology:
\begin{verbatim}
EBNF:
"f.int.primePi", "(", int, ")"
\end{verbatim}
\begin{itemize}
\item Operand: \texttt{int} ($\geq 0$).
\item Result: \texttt{int}.
\item A \texttt{Value Error} is raised if $n < 0$.
\end{itemize}
Examples:
\begin{verbatim}
f.int.primePi(10)       #c# Returns 4  (2, 3, 5, 7)
f.int.primePi(100)      #c# Returns 25
f.int.primePi(1)        #c# Returns 0
\end{verbatim}

\subsubsection{Prime Factors}
Returns the prime factorization of $n$ as a sorted list (with repetition). \\
Morphology:
\begin{verbatim}
EBNF:
"f.int.primeFactors", "(", int, ")"
\end{verbatim}
\begin{itemize}
\item Operand: \texttt{int} ($\geq 1$).
\item Result: \texttt{list[int]}. Returns $[]$ for $n = 1$.
\item A \texttt{Value Error} is raised if $n < 0$.
\end{itemize}
Examples:
\begin{verbatim}
f.int.primeFactors(12)      #c# Returns [2, 2, 3]
f.int.primeFactors(100)     #c# Returns [2, 2, 5, 5]
f.int.primeFactors(1)       #c# Returns []
f.int.primeFactors(17)      #c# Returns [17]
\end{verbatim}

\subsubsection{Distinct Prime Factors}
Returns the distinct prime factors of $n$ as a sorted list. \\
Morphology:
\begin{verbatim}
EBNF:
"f.int.distinctPrimeFactors", "(", int, ")"
\end{verbatim}
\begin{itemize}
\item Operand: \texttt{int} ($\geq 1$).
\item Result: \texttt{list[int]}. Returns $[]$ for $n = 1$.
\item A \texttt{Value Error} is raised if $n < 0$.
\end{itemize}
Examples:
\begin{verbatim}
f.int.distinctPrimeFactors(12)      #c# Returns [2, 3]
f.int.distinctPrimeFactors(100)     #c# Returns [2, 5]
f.int.distinctPrimeFactors(1)       #c# Returns []
\end{verbatim}

\subsubsection{Omega Function}
Computes $\omega(n)$, the number of distinct prime factors of $n$. \\
Morphology:
\begin{verbatim}
EBNF:
"f.int.omega", "(", int, ")"
\end{verbatim}
\begin{itemize}
\item Operand: \texttt{int} ($\geq 1$). $\omega(1) = 0$.
\item Result: \texttt{int}.
\item A \texttt{Value Error} is raised if $n < 0$.
\end{itemize}
Examples:
\begin{verbatim}
f.int.omega(1)          #c# Returns 0
f.int.omega(12)         #c# Returns 2  (factors: 2, 3)
f.int.omega(30)         #c# Returns 3  (factors: 2, 3, 5)
\end{verbatim}

\subsubsection{Big Omega Function}
Computes $\Omega(n)$, the total number of prime factors of $n$ (with multiplicity). \\
Morphology:
\begin{verbatim}
EBNF:
"f.int.bigOmega", "(", int, ")"
\end{verbatim}
\begin{itemize}
\item Operand: \texttt{int} ($\geq 1$). $\Omega(1) = 0$.
\item Result: \texttt{int}.
\item A \texttt{Value Error} is raised if $n < 0$.
\end{itemize}
Examples:
\begin{verbatim}
f.int.bigOmega(1)       #c# Returns 0
f.int.bigOmega(12)      #c# Returns 3  (2 * 2 * 3)
f.int.bigOmega(30)      #c# Returns 3  (2 * 3 * 5)
\end{verbatim}

\subsubsection{Euler's Totient Function}
Computes $\varphi(n)$, the number of integers in $[1, n]$ that are coprime to $n$. \\
Morphology:
\begin{verbatim}
EBNF:
"f.int.totient", "(", int, ")"
\end{verbatim}
\begin{itemize}
\item Operand: \texttt{int} ($\geq 1$). $\varphi(1) = 1$.
\item Result: \texttt{int}.
\item A \texttt{Value Error} is raised if $n < 0$.
\end{itemize}
Examples:
\begin{verbatim}
f.int.totient(1)        #c# Returns 1
f.int.totient(9)        #c# Returns 6  (1, 2, 4, 5, 7, 8)
f.int.totient(12)       #c# Returns 4  (1, 5, 7, 11)
\end{verbatim}

\subsubsection{M\"obius Function}
Computes the M\"obius function $\mu(n)$: $0$ if $n$ has a squared prime factor, $(-1)^k$ if $n$ is a product of $k$ distinct primes, $1$ if $n = 1$. \\
Morphology:
\begin{verbatim}
EBNF:
"f.int.mobius", "(", int, ")"
\end{verbatim}
\begin{itemize}
\item Operand: \texttt{int} ($\geq 1$). $\mu(1) = 1$.
\item Result: \texttt{int} ($-1$, $0$, or $1$).
\item A \texttt{Value Error} is raised if $n \leq 0$.
\end{itemize}
Examples:
\begin{verbatim}
f.int.mobius(1)         #c# Returns 1
f.int.mobius(6)         #c# Returns 1  (2 * 3, two distinct primes)
f.int.mobius(12)        #c# Returns 0  (has squared factor 2^2)
f.int.mobius(30)        #c# Returns -1  (2 * 3 * 5, three distinct primes)
\end{verbatim}

\subsubsection{Liouville Function}
Computes $\lambda(n) = (-1)^{\Omega(n)}$, where $\Omega(n)$ is the total number of prime factors. \\
Morphology:
\begin{verbatim}
EBNF:
"f.int.liouville", "(", int, ")"
\end{verbatim}
\begin{itemize}
\item Operand: \texttt{int} ($\geq 1$). $\lambda(1) = 1$.
\item Result: \texttt{int} ($-1$ or $1$).
\item A \texttt{Value Error} is raised if $n < 0$.
\end{itemize}
Examples:
\begin{verbatim}
f.int.liouville(1)      #c# Returns 1
f.int.liouville(6)      #c# Returns 1  (2*3, two factors → (-1)^2)
f.int.liouville(8)      #c# Returns -1  (2^2*2, three factors → (-1)^3)
\end{verbatim}

\subsubsection{Jacobi Symbol}
Computes the Jacobi symbol $(a/n)$, a generalization of the Legendre symbol to composite moduli. \\
Morphology:
\begin{verbatim}
EBNF:
"f.int.jacobiSymbol", "(", int, ",", int, ")"
\end{verbatim}
\begin{itemize}
\item Operands: $a$ (\texttt{int}), $n$ (\texttt{int}, odd $> 0$).
\item Result: \texttt{int} ($-1$, $0$, or $1$).
\item A \texttt{Value Error} is raised if $n \leq 0$ or $n$ is even.
\end{itemize}
Examples:
\begin{verbatim}
f.int.jacobiSymbol(2, 3)     #c# Returns -1
f.int.jacobiSymbol(3, 5)     #c# Returns -1
f.int.jacobiSymbol(2, 7)     #c# Returns 1
\end{verbatim}

\subsubsection{Legendre Symbol}
Computes the Legendre symbol $(a/p)$: $1$ if $a$ is a quadratic residue mod $p$, $-1$ if not, $0$ if $p \mid a$. \\
Morphology:
\begin{verbatim}
EBNF:
"f.int.legendreSymbol", "(", int, ",", int, ")"
\end{verbatim}
\begin{itemize}
\item Operands: $a$ (\texttt{int}), $p$ (\texttt{int}, must be an odd prime).
\item Result: \texttt{int} ($-1$, $0$, or $1$).
\item A \texttt{Value Error} is raised if $p$ is not an odd prime.
\end{itemize}
Examples:
\begin{verbatim}
f.int.legendreSymbol(2, 7)     #c# Returns 1  (2 is QR mod 7: 3^2 ≡ 2)
f.int.legendreSymbol(3, 7)     #c# Returns -1
f.int.legendreSymbol(4, 7)     #c# Returns 1  (4 = 2^2)
\end{verbatim}

\subsubsection{Kronecker Symbol}
Computes the Kronecker symbol $(a/n)$, a generalization of the Jacobi symbol to all integers $n$. \\
Morphology:
\begin{verbatim}
EBNF:
"f.int.kroneckerSymbol", "(", int, ",", int, ")"
\end{verbatim}
\begin{itemize}
\item Operands: two \texttt{int} values.
\item Result: \texttt{int} ($-1$, $0$, or $1$).
\item A \texttt{Value Error} is raised if both arguments are $0$.
\end{itemize}
Examples:
\begin{verbatim}
f.int.kroneckerSymbol(2, 3)     #c# Returns -1
f.int.kroneckerSymbol(5, 1)     #c# Returns 1
f.int.kroneckerSymbol(3, -1)    #c# Returns -1
\end{verbatim}

\subsubsection{Carmichael Lambda Function}
Computes $\lambda(n)$, the smallest positive integer such that $a^{\lambda(n)} \equiv 1 \pmod{n}$ for all $a$ coprime to $n$. \\
Morphology:
\begin{verbatim}
EBNF:
"f.int.carmichaelLambda", "(", int, ")"
\end{verbatim}
\begin{itemize}
\item Operand: \texttt{int} ($\geq 1$).
\item Result: \texttt{int}.
\item A \texttt{Value Error} is raised if $n < 1$.
\end{itemize}
Examples:
\begin{verbatim}
f.int.carmichaelLambda(1)      #c# Returns 1
f.int.carmichaelLambda(7)      #c# Returns 6  (= phi(7))
f.int.carmichaelLambda(8)      #c# Returns 2
f.int.carmichaelLambda(12)     #c# Returns 2
\end{verbatim}

\subsubsection{Continued Fraction}
Computes the continued fraction representation of a real number as a list of integers $[a_0; a_1, a_2, \ldots]$.
\begin{verbatim}
EBNF:
"f.int.continuedFraction", "(", real, [",", int], ")"
\end{verbatim}
\begin{itemize}
\item Operands: \texttt{x} (value to expand), optional \texttt{maxTerms} (\texttt{int}, default 20).
\item Result: \texttt{list[int]} --- continued fraction coefficients $[a_0; a_1, a_2, \ldots]$.
\item A \texttt{Value Error} is raised if \texttt{maxTerms} $\leq 0$.
\end{itemize}
\begin{verbatim}
f.int.continuedFraction(c.math.pi)       #c# Returns [3, 7, 15, 1, 292, ...]
f.int.continuedFraction(3.14159, 5)      #c# Returns [3, 7, 15, 1, 292]
f.int.continuedFraction(1.41421)         #c# Returns [1, 2, 2, 2, 2, ...] (sqrt(2))
\end{verbatim}

\subsubsection{Continued Fraction Value}
Evaluates a continued fraction from a list of coefficients and returns the resulting rational or real value.
\begin{verbatim}
EBNF:
"f.int.continuedFractionValue", "(", list[int], ")"
\end{verbatim}
\begin{itemize}
\item Operands: \texttt{list[int]} --- continued fraction coefficients $[a_0; a_1, a_2, \ldots]$.
\item Result: \texttt{real} --- the value of the continued fraction.
\item A \texttt{Value Error} is raised if the list is empty.
\end{itemize}
\begin{verbatim}
f.int.continuedFractionValue([3, 7, 15, 1])        #c# Returns approximately 3.14159
f.int.continuedFractionValue([1, 2, 2, 2, 2])      #c# Returns approximately 1.41421
f.int.continuedFractionValue([0, 1, 1, 1, 1, 1])   #c# Returns approximately 0.61803 (golden ratio - 1)
\end{verbatim}

\subsubsection{Is Perfect Power}
Checks whether an integer is a perfect power, i.e.\ $n = m^k$ for some integers $m \geq 1$, $k \geq 2$.
\begin{verbatim}
EBNF:
"f.int.isPerfect", "(", int, ")"
\end{verbatim}
\begin{itemize}
\item Operands: \texttt{int}.
\item Result: \texttt{bool} --- \texttt{True} if $n$ is a perfect power (perfect square, cube, etc.), \texttt{False} otherwise.
\item Returns \texttt{True} for $n = 0$ and $n = 1$.
\item A \texttt{Type Error} is raised if the argument is not \texttt{int}.
\end{itemize}
\begin{verbatim}
f.int.isPerfect(4)        #c# True  (2^2)
f.int.isPerfect(8)        #c# True  (2^3)
f.int.isPerfect(9)        #c# True  (3^2)
f.int.isPerfect(10)       #c# False
f.int.isPerfect(27)       #c# True  (3^3)
f.int.isPerfect(100)      #c# True  (10^2)
\end{verbatim}

\subsubsection{Is Abundant}
Checks whether an integer is an abundant number (sum of proper divisors exceeds $n$).
\begin{verbatim}
EBNF:
"f.int.isAbundant", "(", int, ")"
\end{verbatim}
\begin{itemize}
\item Operands: \texttt{int} (must be $\geq 1$).
\item Result: \texttt{bool}.
\item A \texttt{Value Error} is raised if the argument is $< 1$.
\end{itemize}
\begin{verbatim}
f.int.isAbundant(12)      #c# True  (1 + 2 + 3 + 4 + 6 = 16 > 12)
f.int.isAbundant(6)       #c# False (1 + 2 + 3 = 6)
f.int.isAbundant(18)      #c# True  (1 + 2 + 3 + 6 + 9 = 21 > 18)
\end{verbatim}

\subsubsection{Is Deficient}
Checks whether an integer is a deficient number (sum of proper divisors is less than $n$).
\begin{verbatim}
EBNF:
"f.int.isDeficient", "(", int, ")"
\end{verbatim}
\begin{itemize}
\item Operands: \texttt{int} (must be $\geq 1$).
\item Result: \texttt{bool}.
\item A \texttt{Value Error} is raised if the argument is $< 1$.
\end{itemize}
\begin{verbatim}
f.int.isDeficient(8)      #c# True  (1 + 2 + 4 = 7 < 8)
f.int.isDeficient(6)      #c# False (1 + 2 + 3 = 6, perfect)
f.int.isDeficient(10)     #c# True  (1 + 2 + 5 = 8 < 10)
\end{verbatim}

\subsubsection{Is Square-Free}
Checks whether an integer is square-free (not divisible by any perfect square other than 1).
\begin{verbatim}
EBNF:
"f.int.isSquareFree", "(", int, ")"
\end{verbatim}
\begin{itemize}
\item Operands: \texttt{int}.
\item Result: \texttt{bool} --- \texttt{True} if no prime squared divides $n$.
\item Returns \texttt{True} for $n = 0$ and $n = 1$.
\end{itemize}
\begin{verbatim}
f.int.isSquareFree(6)     #c# True  (2 * 3)
f.int.isSquareFree(12)    #c# False (divisible by 4)
f.int.isSquareFree(30)    #c# True  (2 * 3 * 5)
f.int.isSquareFree(1)     #c# True
\end{verbatim}

\subsubsection{Is Power}
Checks whether an integer is a perfect $k$-th power for a given exponent.
\begin{verbatim}
EBNF:
"f.int.isPower", "(", int, ",", int, ")"
\end{verbatim}
\begin{itemize}
\item Operands: \texttt{n} (value), \texttt{k} (exponent, must be $\geq 2$).
\item Result: \texttt{bool} --- \texttt{True} if $n = m^k$ for some integer $m$.
\item A \texttt{Value Error} is raised if \texttt{k $<$ 2}.
\end{itemize}
\begin{verbatim}
f.int.isPower(8, 3)       #c# True  (2^3 = 8)
f.int.isPower(9, 2)       #c# True  (3^2 = 9)
f.int.isPower(10, 2)      #c# False
f.int.isPower(16, 4)      #c# True  (2^4 = 16)
f.int.isPower(1024, 10)   #c# True  (2^10 = 1024)
\end{verbatim}

\subsubsection{Integer Square Root}
Returns the floor of the square root of a non-negative integer. Equivalent to $\lfloor\sqrt{n}\rfloor$ computed exactly using integer arithmetic.
\begin{verbatim}
EBNF:
"f.int.integerSqrt", "(", int, ")"
\end{verbatim}
\begin{itemize}
\item Operands: \texttt{int} (must be $\geq 0$).
\item Result: \texttt{int}.
\item A \texttt{Value Error} is raised if the argument is $< 0$.
\end{itemize}
\begin{verbatim}
f.int.integerSqrt(0)       #c# Returns 0
f.int.integerSqrt(1)       #c# Returns 1
f.int.integerSqrt(8)       #c# Returns 2  (floor of 2.828)
f.int.integerSqrt(9)       #c# Returns 3
f.int.integerSqrt(100)     #c# Returns 10
f.int.integerSqrt(99)      #c# Returns 9
\end{verbatim}

\subsubsection{Valuation}
Returns the $p$-adic valuation of $n$: the exponent of the highest power of $p$ that divides $n$.
\begin{verbatim}
EBNF:
"f.int.valuation", "(", int, ",", int, ")"
\end{verbatim}
\begin{itemize}
\item Operands: \texttt{n} (value), \texttt{p} (prime base, must be $> 1$).
\item Result: \texttt{int} --- largest $k$ such that $p^k \mid n$.
\item Returns 0 if $p$ does not divide $n$.
\item A \texttt{Value Error} is raised if \texttt{p $\leq$ 1}.
\end{itemize}
\begin{verbatim}
f.int.valuation(8, 2)        #c# Returns 3  (2^3 | 8)
f.int.valuation(12, 2)       #c# Returns 2  (2^2 | 12, but 2^3 does not)
f.int.valuation(12, 3)       #c# Returns 1  (3^1 | 12)
f.int.valuation(100, 10)     #c# Returns 2  (10^2 | 100)
f.int.valuation(7, 2)        #c# Returns 0  (2 does not divide 7)
\end{verbatim}

\subsubsection{Chinese Remainder}
Solves a system of simultaneous congruences using the Chinese Remainder Theorem. Finds $x$ such that $x \equiv a_i \pmod{m_i}$ for all $i$.
\begin{verbatim}
EBNF:
"f.int.chineseRemainder", "(", list[int], ",", list[int], ")"
\end{verbatim}
\begin{itemize}
\item Operands: \texttt{remainders} ($[a_0, a_1, \ldots]$), \texttt{moduli} ($[m_0, m_1, \ldots]$).
\item Result: \texttt{int} --- smallest non-negative solution $x$.
\item Moduli must be pairwise coprime.
\item A \texttt{Value Error} is raised if lists are empty, have different lengths, or moduli are not pairwise coprime.
\end{itemize}
\begin{verbatim}
f.int.chineseRemainder([2, 3, 2], [3, 5, 7])    #c# Returns 23  (23 % 3 = 2, 23 % 5 = 3, 23 % 7 = 2)
f.int.chineseRemainder([1, 2, 3], [3, 4, 5])    #c# Returns 58  (58 % 3 = 1, 58 % 4 = 2, 58 % 5 = 3)
\end{verbatim}

\subsubsection{Primitive Root}
Finds the smallest primitive root modulo $n$, or returns 0 if none exists. A primitive root $g$ modulo $n$ generates the multiplicative group: $g^1, g^2, \ldots, g^{\phi(n)}$ produce all integers coprime to $n$.
\begin{verbatim}
EBNF:
"f.int.primitiveRoot", "(", int, ")"
\end{verbatim}
\begin{itemize}
\item Operands: \texttt{n} (modulus, must be $> 1$).
\item Result: \texttt{int} --- smallest primitive root modulo $n$, or 0 if none exists.
\item A \texttt{Value Error} is raised if \texttt{n $\leq$ 1}.
\end{itemize}
\begin{verbatim}
f.int.primitiveRoot(2)       #c# Returns 1
f.int.primitiveRoot(3)       #c# Returns 2
f.int.primitiveRoot(5)       #c# Returns 2
f.int.primitiveRoot(7)       #c# Returns 3
f.int.primitiveRoot(8)       #c# Returns 0  (no primitive root mod 8)
f.int.primitiveRoot(11)      #c# Returns 2
\end{verbatim}

% ──────────────────────────────────────────────────────────────
%  Additional Special Functions
% ──────────────────────────────────────────────────────────────

\subsubsection{Polygamma Function}
Computes the polygamma function $\psi^{(n)}(x) = \frac{d^{n+1}}{dx^{n+1}} \ln \Gamma(x)$. For $n=0$ this is the digamma function. \\
Morphology:
\begin{verbatim}
EBNF:
"f.special.polygamma", "(", int, ",", ( int | real ), ")"
\end{verbatim}
\begin{itemize}
\item Operands: $n$ (\texttt{int}, order $\geq 0$), $x$ (\texttt{int} or \texttt{real}).
\item Result: \texttt{real}.
\item A \texttt{Value Error} is raised if $x$ is a non-positive integer (pole of $\Gamma$).
\item A \texttt{Type Error} is raised if arguments are not \texttt{int}/\texttt{real}.
\end{itemize}
Examples:
\begin{verbatim}
f.special.polygamma(0, 1)        #c# Returns -0.5772... (Euler-Mascheroni constant, negated)
f.special.polygamma(0, 5)        #c# Returns 2.0256...
f.special.polygamma(1, 1)        #c# Returns 1.6449... (pi^2 / 6)
f.special.polygamma(2, 1)        #c# Returns -2.4041...
\end{verbatim}

\subsubsection{IsNaN}
Returns \texttt{True} if the value is Not-a-Number (\texttt{NaN}). \\
Morphology:
\begin{verbatim}
EBNF:
"f.math.isNaN", "(", ( int | real ), ")"
\end{verbatim}
\begin{itemize}
\item Operand: $x$ (\texttt{int} or \texttt{real}).
\item Result: \texttt{bool}.
\end{itemize}
Examples:
\begin{verbatim}
f.math.isNaN(0/0)               #c# Returns True  (0/0 = NaN in real division)
f.math.isNaN(1)                 #c# Returns False
\end{verbatim}

\subsubsection{IsInf}
Returns \texttt{True} if the value is positive or negative infinity. \\
Morphology:
\begin{verbatim}
EBNF:
"f.math.isInf", "(", ( int | real ), ")"
\end{verbatim}
\begin{itemize}
\item Operand: $x$ (\texttt{int} or \texttt{real}).
\item Result: \texttt{bool}.
\end{itemize}
Examples:
\begin{verbatim}
f.math.isInf(1/0)               #c# Returns True  (1/0 = Inf)
f.math.isInf(1)                 #c# Returns False
\end{verbatim}

\subsubsection{Remap}
Linearly maps a value from one range to another. \\
Morphology:
\begin{verbatim}
EBNF:
"f.math.remap", "(", ( int | real ), ",", ( int | real ), ",", ( int | real ), ",",
                     ( int | real ), ",", ( int | real ), ")"
\end{verbatim}
\begin{itemize}
\item Operands: $value$, $from\_low$, $from\_high$, $to\_low$, $to\_high$ (all \texttt{int} or \texttt{real}).
\item Result: \texttt{real}.
\end{itemize}
Examples:
\begin{verbatim}
f.math.remap(5, 0, 10, 0, 100)  #c# Returns 50.0
f.math.remap(0, 0, 10, 0, 100)  #c# Returns 0.0
\end{verbatim}

\subsubsection{Covariance}
Computes the population covariance of two lists of equal length. \\
Morphology:
\begin{verbatim}
EBNF:
"f.stat.covariance", "(", list, ",", list, ")"
\end{verbatim}
\begin{itemize}
\item Operands: Two \texttt{list[int]} or \texttt{list[real]} of equal length.
\item Result: \texttt{real} (population covariance).
\item A \texttt{Value Error} is raised if lists are empty or unequal length.
\end{itemize}
Examples:
\begin{verbatim}
f.stat.covariance([1,2,3], [4,5,6])   #c# Returns 0.6666...
\end{verbatim}

\subsubsection{Correlation}
Computes the Pearson correlation coefficient of two lists of equal length. \\
Morphology:
\begin{verbatim}
EBNF:
"f.stat.correlation", "(", list, ",", list, ")"
\end{verbatim}
\begin{itemize}
\item Operands: Two \texttt{list[int]} or \texttt{list[real]} of equal length.
\item Result: \texttt{real} ($-1$ to $1$).
\item A \texttt{Value Error} is raised if lists are empty, unequal length, or have zero variance.
\end{itemize}
Examples:
\begin{verbatim}
f.stat.correlation([1,2,3], [2,4,6])   #c# Returns 1.0  (perfect positive)
f.stat.correlation([1,2,3], [6,4,2])   #c# Returns -1.0 (perfect negative)
\end{verbatim}

\subsubsection{Linear Regression}
Performs simple linear regression on two lists: $y = mx + b$. Returns \texttt{[slope, intercept]}. \\
Morphology:
\begin{verbatim}
EBNF:
"f.stat.linearRegression", "(", list, ",", list, ")"
\end{verbatim}
\begin{itemize}
\item Operands: $x$ (\texttt{list[int]} or \texttt{list[real]}), $y$ (\texttt{list[int]} or \texttt{list[real]}).
\item Result: \texttt{list[real]} with \texttt{[slope, intercept]}.
\item A \texttt{Value Error} is raised if lists are empty or have fewer than 2 points.
\end{itemize}
Examples:
\begin{verbatim}
f.stat.linearRegression([1,2,3], [2,4,6])   #c# Returns [2.0, 0.0]
f.stat.linearRegression([0,1,2], [1,3,5])   #c# Returns [2.0, 1.0]
\end{verbatim}

% ──────────────────────────────────────────────────────────────
%  Additional Special Functions (batch 2)
% ──────────────────────────────────────────────────────────────

\subsubsection{Dawson Function}
Computes the Dawson function $F(x) = e^{-x^2} \int_0^x e^{t^2}\,dt$. \\
Morphology:
\begin{verbatim}
EBNF:
"f.special.dawson", "(", ( int | real ), ")"
\end{verbatim}
\begin{itemize}
\item Operand: \texttt{int} or \texttt{real}.
\item Result: \texttt{real}.
\item Odd function: $F(-x) = -F(x)$.
\end{itemize}
Examples:
\begin{verbatim}
f.special.dawson(0)        #c# Returns 0.0
f.special.dawson(1)        #c# Returns approximately 0.5381
f.special.dawson(-1)       #c# Returns approximately -0.5381
\end{verbatim}

\subsubsection{Fresnel Sine Integral}
Computes the Fresnel sine integral $S(x) = \int_0^x \sin\!\left(\frac{\pi t^2}{2}\right) dt$. \\
Morphology:
\begin{verbatim}
EBNF:
"f.special.fresnelS", "(", ( int | real ), ")"
\end{verbatim}
\begin{itemize}
\item Operand: \texttt{int} or \texttt{real}.
\item Result: \texttt{real}.
\item Odd function: $S(-x) = -S(x)$.
\end{itemize}
Examples:
\begin{verbatim}
f.special.fresnelS(0)     #c# Returns 0.0
f.special.fresnelS(1)     #c# Returns approximately 0.4383
\end{verbatim}

\subsubsection{Fresnel Cosine Integral}
Computes the Fresnel cosine integral $C(x) = \int_0^x \cos\!\left(\frac{\pi t^2}{2}\right) dt$. \\
Morphology:
\begin{verbatim}
EBNF:
"f.special.fresnelC", "(", ( int | real ), ")"
\end{verbatim}
\begin{itemize}
\item Operand: \texttt{int} or \texttt{real}.
\item Result: \texttt{real}.
\item Odd function: $C(-x) = -C(x)$.
\end{itemize}
Examples:
\begin{verbatim}
f.special.fresnelC(0)     #c# Returns 0.0
f.special.fresnelC(1)     #c# Returns approximately 0.7799
\end{verbatim}

\subsubsection{Exponential Integral}
Computes the exponential integral $\text{Ei}(x) = -\int_{-x}^{\infty} \frac{e^{-t}}{t}\,dt$. \\
Morphology:
\begin{verbatim}
EBNF:
"f.special.expIntegral", "(", ( int | real ), ")"
\end{verbatim}
\begin{itemize}
\item Operand: \texttt{int} or \texttt{real}.
\item Result: \texttt{real}.
\item A \texttt{Value Error} is raised at $x = 0$ (singularity).
\end{itemize}
Examples:
\begin{verbatim}
f.special.expIntegral(1)   #c# Returns approximately 1.8951
f.special.expIntegral(-1)  #c# Returns approximately -0.2194
\end{verbatim}

\subsubsection{Logarithmic Integral}
Computes the logarithmic integral $\text{li}(x) = \int_0^x \frac{dt}{\ln t}$. \\
Morphology:
\begin{verbatim}
EBNF:
"f.special.logIntegral", "(", ( int | real ), ")"
\end{verbatim}
\begin{itemize}
\item Operand: \texttt{int} or \texttt{real}.
\item Result: \texttt{real}.
\item A \texttt{Value Error} is raised for $x \leq 0$ (domain restriction).
\end{itemize}
Examples:
\begin{verbatim}
f.special.logIntegral(2)   #c# Returns approximately 1.0452
f.special.logIntegral(10)  #c# Returns approximately 5.1200
\end{verbatim}

\subsubsection{Cosine Integral}
Computes the cosine integral $\text{Ci}(x) = -\int_x^{\infty} \frac{\cos t}{t}\,dt$. \\
Morphology:
\begin{verbatim}
EBNF:
"f.special.ci", "(", ( int | real ), ")"
\end{verbatim}
\begin{itemize}
\item Operand: \texttt{int} or \texttt{real}.
\item Result: \texttt{real}.
\item A \texttt{Value Error} is raised at $x = 0$ (singularity).
\end{itemize}
Examples:
\begin{verbatim}
f.special.ci(1)            #c# Returns approximately 0.3374
f.special.ci(5)            #c# Returns approximately -0.1900
\end{verbatim}

\subsubsection{Hyperbolic Sine Integral}
Computes the hyperbolic sine integral $\text{Shi}(x) = \int_0^x \frac{\sinh t}{t}\,dt$. \\
Morphology:
\begin{verbatim}
EBNF:
"f.special.shi", "(", ( int | real ), ")"
\end{verbatim}
\begin{itemize}
\item Operand: \texttt{int} or \texttt{real}.
\item Result: \texttt{real}.
\item Odd function: $\text{Shi}(-x) = -\text{Shi}(x)$.
\end{itemize}
Examples:
\begin{verbatim}
f.special.shi(0)           #c# Returns 0.0
f.special.shi(1)           #c# Returns approximately 1.0573
\end{verbatim}

\subsubsection{Hyperbolic Cosine Integral}
Computes the hyperbolic cosine integral $\text{Chi}(x) = \gamma + \ln x + \int_0^x \frac{\cosh t - 1}{t}\,dt$. \\
Morphology:
\begin{verbatim}
EBNF:
"f.special.chi", "(", ( int | real ), ")"
\end{verbatim}
\begin{itemize}
\item Operand: \texttt{int} or \texttt{real}.
\item Result: \texttt{real}.
\item A \texttt{Value Error} is raised at $x = 0$ (singularity).
\end{itemize}
Examples:
\begin{verbatim}
f.special.chi(1)           #c# Returns approximately 0.8379
f.special.chi(2)           #c# Returns approximately 1.3060
\end{verbatim}

\subsubsection{Hankel Function of the First Kind}
Computes the Hankel function of the first kind $H_n^{(1)}(x) = J_n(x) + i Y_n(x)$. \\
Morphology:
\begin{verbatim}
EBNF:
"f.special.hankelH1", "(", int, ",", ( int | real ), ")"
\end{verbatim}
\begin{itemize}
\item Operands: order $n$ (\texttt{int}), argument $x$ (\texttt{int} or \texttt{real}).
\item Result: \texttt{complex}.
\end{itemize}
Examples:
\begin{verbatim}
f.special.hankelH1(0, 1)   #c# Returns approximately 0.7652+0.0883i
\end{verbatim}

\subsubsection{Hankel Function of the Second Kind}
Computes the Hankel function of the second kind $H_n^{(2)}(x) = J_n(x) - i Y_n(x)$. \\
Morphology:
\begin{verbatim}
EBNF:
"f.special.hankelH2", "(", int, ",", ( int | real ), ")"
\end{verbatim}
\begin{itemize}
\item Operands: order $n$ (\texttt{int}), argument $x$ (\texttt{int} or \texttt{real}).
\item Result: \texttt{complex}.
\end{itemize}
Examples:
\begin{verbatim}
f.special.hankelH2(0, 1)   #c# Returns approximately 0.7652-0.0883i
\end{verbatim}

\subsubsection{Kelvin Function ber}
Computes the Kelvin function $\text{ber}_n(x) = \text{Re}\left[J_n\!\left(x e^{i 3\pi/4}\right)\right]$. \\
Morphology:
\begin{verbatim}
EBNF:
"f.special.kelvinBer", "(", int, ",", ( int | real ), ")"
\end{verbatim}
\begin{itemize}
\item Operands: order $n$ (\texttt{int}), argument $x$ (\texttt{int} or \texttt{real}).
\item Result: \texttt{real}.
\end{itemize}
Examples:
\begin{verbatim}
f.special.kelvinBer(0, 0)  #c# Returns 1.0
f.special.kelvinBer(0, 1)  #c# Returns approximately 0.9844
\end{verbatim}

\subsubsection{Kelvin Function bei}
Computes the Kelvin function $\text{bei}_n(x) = \text{Im}\left[J_n\!\left(x e^{i 3\pi/4}\right)\right]$. \\
Morphology:
\begin{verbatim}
EBNF:
"f.special.kelvinBei", "(", int, ",", ( int | real ), ")"
\end{verbatim}
\begin{itemize}
\item Operands: order $n$ (\texttt{int}), argument $x$ (\texttt{int} or \texttt{real}).
\item Result: \texttt{real}.
\end{itemize}
Examples:
\begin{verbatim}
f.special.kelvinBei(0, 0)  #c# Returns 0.0
f.special.kelvinBei(0, 1)  #c# Returns approximately -0.1196
\end{verbatim}

\subsubsection{Kelvin Function ker}
Computes the Kelvin function $\text{ker}_n(x) = \text{Re}\left[K_n\!\left(x e^{i 3\pi/4}\right)\right]$. \\
Morphology:
\begin{verbatim}
EBNF:
"f.special.kelvinKer", "(", int, ",", ( int | real ), ")"
\end{verbatim}
\begin{itemize}
\item Operands: order $n$ (\texttt{int}), argument $x$ (\texttt{int} or \texttt{real}).
\item Result: \texttt{real}.
\item A \texttt{Value Error} is raised for $x \leq 0$ (singularity).
\end{itemize}
Examples:
\begin{verbatim}
f.special.kelvinKer(0, 1)  #c# Returns approximately 0.1159
\end{verbatim}

\subsubsection{Kelvin Function kei}
Computes the Kelvin function $\text{kei}_n(x) = \text{Im}\left[K_n\!\left(x e^{i 3\pi/4}\right)\right]$. \\
Morphology:
\begin{verbatim}
EBNF:
"f.special.kelvinKei", "(", int, ",", ( int | real ), ")"
\end{verbatim}
\begin{itemize}
\item Operands: order $n$ (\texttt{int}), argument $x$ (\texttt{int} or \texttt{real}).
\item Result: \texttt{real}.
\end{itemize}
Examples:
\begin{verbatim}
f.special.kelvinKei(0, 1)  #c# Returns approximately -0.7895
\end{verbatim}

\subsubsection{Struve Function H}
Computes the Struve function $\mathbf{H}_n(x)$. \\
Morphology:
\begin{verbatim}
EBNF:
"f.special.struveH", "(", int, ",", ( int | real ), ")"
\end{verbatim}
\begin{itemize}
\item Operands: order $n$ (\texttt{int}), argument $x$ (\texttt{int} or \texttt{real}).
\item Result: \texttt{real}.
\end{itemize}
Examples:
\begin{verbatim}
f.special.struveH(0, 0)    #c# Returns 0.0
f.special.struveH(0, 1)    #c# Returns approximately 0.7102
\end{verbatim}

\subsubsection{Struve Function L}
Computes the modified Struve function $\mathbf{L}_n(x)$. \\
Morphology:
\begin{verbatim}
EBNF:
"f.special.struveL", "(", int, ",", ( int | real ), ")"
\end{verbatim}
\begin{itemize}
\item Operands: order $n$ (\texttt{int}), argument $x$ (\texttt{int} or \texttt{real}).
\item Result: \texttt{real}.
\end{itemize}
Examples:
\begin{verbatim}
f.special.struveL(0, 0)    #c# Returns 0.0
f.special.struveL(0, 1)    #c# Returns approximately 0.7102
\end{verbatim}

\subsubsection{Riemann Zeta Function}
Computes the Riemann zeta function $\zeta(s) = \sum_{n=1}^{\infty} n^{-s}$. \\
Morphology:
\begin{verbatim}
EBNF:
"f.special.riemannZeta", "(", ( int | real ), ")"
\end{verbatim}
\begin{itemize}
\item Operand: \texttt{int} or \texttt{real}.
\item Result: \texttt{real}.
\item A \texttt{Value Error} is raised at $s = 1$ (simple pole).
\item $\zeta(0) = -0.5$, $\zeta(2) = \pi^2 / 6$, $\zeta(4) = \pi^4 / 90$.
\end{itemize}
Examples:
\begin{verbatim}
f.special.riemannZeta(2)   #c# Returns approximately 1.6449 (pi^2/6)
f.special.riemannZeta(4)   #c# Returns approximately 1.0823 (pi^4/90)
f.special.riemannZeta(0)   #c# Returns -0.5
\end{verbatim}

\subsubsection{Hurwitz Zeta Function}
Computes the Hurwitz zeta function $\zeta(s, q) = \sum_{n=0}^{\infty} (n+q)^{-s}$. For $q = 1$ this is the Riemann zeta function. \\
Morphology:
\begin{verbatim}
EBNF:
"f.special.hurwitzZeta", "(", ( int | real ), ",", ( int | real ), ")"
\end{verbatim}
\begin{itemize}
\item Operands: $s$ (\texttt{int} or \texttt{real}), $q$ (\texttt{int} or \texttt{real}).
\item Result: \texttt{real}.
\item A \texttt{Value Error} is raised if $s \leq 1$ or $q \leq 0$.
\end{itemize}
Examples:
\begin{verbatim}
f.special.hurwitzZeta(2, 1)   #c# Returns approximately 1.6449 (pi^2/6)
f.special.hurwitzZeta(2, 0.5) #c# Returns approximately 4.9348 (pi^2/2)
\end{verbatim}

\subsubsection{Dirichlet Eta Function}
Computes the Dirichlet eta function $\eta(s) = (1 - 2^{1-s})\,\zeta(s)$. Also known as the alternating zeta function. \\
Morphology:
\begin{verbatim}
EBNF:
"f.special.dirichletEta", "(", ( int | real ), ")"
\end{verbatim}
\begin{itemize}
\item Operand: \texttt{int} or \texttt{real}.
\item Result: \texttt{real}.
\item $\eta(0) = 0.5$, $\eta(1) = \ln 2$.
\end{itemize}
Examples:
\begin{verbatim}
f.special.dirichletEta(0)   #c# Returns 0.5
f.special.dirichletEta(1)   #c# Returns approximately 0.6931 (ln 2)
f.special.dirichletEta(2)   #c# Returns approximately 0.8225 (pi^2/12)
\end{verbatim}

\subsubsection{Lerch Transcendent}
Computes the Lerch transcendent $\Phi(z, s, q) = \sum_{n=0}^{\infty} \frac{z^n}{(n+q)^s}$. \\
Morphology:
\begin{verbatim}
EBNF:
"f.special.lerchPhi", "(", ( int | real | complex ), ",", ( int | real ), ",", ( int | real ), ")"
\end{verbatim}
\begin{itemize}
\item Operands: $z$ (\texttt{int}, \texttt{real}, or \texttt{complex}), $s$ (\texttt{int} or \texttt{real}), $q$ (\texttt{int} or \texttt{real}).
\item Result: \texttt{complex}.
\item A \texttt{Value Error} is raised if $q \leq 0$ or if $|z| \geq 1$ and $s \leq 1$.
\end{itemize}
Examples:
\begin{verbatim}
f.special.lerchPhi(0.5, 2, 1)   #c# Returns approximately 1.1447
\end{verbatim}

\subsubsection{Confluent Hypergeometric Function 0F1}
Computes ${}_0F_1(; b; z) = \sum_{n=0}^{\infty} \frac{z^n}{(b)_n\, n!}$. \\
Morphology:
\begin{verbatim}
EBNF:
"f.special.hyp0F1", "(", ( int | real ), ",", ( int | real | complex ), ")"
\end{verbatim}
\begin{itemize}
\item Operands: $b$ (\texttt{int} or \texttt{real}), $z$ (\texttt{int}, \texttt{real}, or \texttt{complex}).
\item Result: \texttt{complex}.
\item A \texttt{Value Error} is raised if $b$ is a non-positive integer.
\end{itemize}
Examples:
\begin{verbatim}
f.special.hyp0F1(1.5, -0.25)   #c# Returns approximately 0.9048
\end{verbatim}

\subsubsection{Kummer's Confluent Hypergeometric Function 1F1}
Computes ${}_1F_1(a; b; z) = \sum_{n=0}^{\infty} \frac{(a)_n\, z^n}{(b)_n\, n!}$. Also known as Kummer's function $M(a, b, z)$. \\
Morphology:
\begin{verbatim}
EBNF:
"f.special.hyp1F1", "(", ( int | real ), ",", ( int | real ), ",", ( int | real | complex ), ")"
\end{verbatim}
\begin{itemize}
\item Operands: $a$ (\texttt{int} or \texttt{real}), $b$ (\texttt{int} or \texttt{real}), $z$ (\texttt{int}, \texttt{real}, or \texttt{complex}).
\item Result: \texttt{complex}.
\item A \texttt{Value Error} is raised if $b$ is a non-positive integer.
\end{itemize}
Examples:
\begin{verbatim}
f.special.hyp1F1(1, 1, 1)   #c# Returns approximately 2.7183 (e)
\end{verbatim}

\subsubsection{Gauss Hypergeometric Function 2F1}
Computes ${}_2F_1(a, b; c; z) = \sum_{n=0}^{\infty} \frac{(a)_n (b)_n}{(c)_n} \frac{z^n}{n!}$. \\
Morphology:
\begin{verbatim}
EBNF:
"f.special.hyp2F1", "(", ( int | real ), ",", ( int | real ), ",", ( int | real ), ",", ( int | real | complex ), ")"
\end{verbatim}
\begin{itemize}
\item Operands: $a$, $b$, $c$ (\texttt{int} or \texttt{real}), $z$ (\texttt{int}, \texttt{real}, or \texttt{complex}).
\item Result: \texttt{complex}.
\item A \texttt{Value Error} is raised if $c$ is a non-positive integer.
\item Converges for $|z| < 1$; analytic continuation for $|z| \geq 1$.
\end{itemize}
Examples:
\begin{verbatim}
f.special.hyp2F1(1, 1, 2, 0.5)   #c# Returns approximately 1.3863
\end{verbatim}

\subsubsection{Generalized Hypergeometric Function pFq}
Computes ${}_pF_q(a_1, \ldots, a_p; b_1, \ldots, b_q; z)$. \\
Morphology:
\begin{verbatim}
EBNF:
"f.special.hypPFQ", "(", list, ",", list, ",", ( int | real | complex ), ")"
\end{verbatim}
\begin{itemize}
\item Operands: numerator parameters (\texttt{list}), denominator parameters (\texttt{list}), argument (\texttt{int}, \texttt{real}, or \texttt{complex}).
\item Result: \texttt{complex}.
\item A \texttt{Type Error} is raised if the lists do not contain numeric values.
\item A \texttt{Value Error} is raised if any denominator parameter is a non-positive integer.
\end{itemize}
Examples:
\begin{verbatim}
f.special.hypPFQ([1], [2], 1)          #c# Returns approximately 1.7183 (e-1)
f.special.hypPFQ([1, 1], [2], 0.5)     #c# Returns approximately 1.2730
\end{verbatim}

\subsubsection{Meijer G-Function}
Computes the Meijer G-function $G_{p,q}^{m,n}\!\left(z \;\middle|\; \begin{smallmatrix} a_1, \ldots, a_p \\ b_1, \ldots, b_q \end{smallmatrix}\right)$. A very general function that includes most common special functions as special cases. \\
Morphology:
\begin{verbatim}
EBNF:
"f.special.meijerG", "(", list, ",", list, ",", ( int | real | complex ), ")"
\end{verbatim}
\begin{itemize}
\item Operands: upper parameters (\texttt{list}), lower parameters (\texttt{list}), argument (\texttt{int}, \texttt{real}, or \texttt{complex}).
\item Result: \texttt{complex}.
\item A \texttt{Type Error} is raised if the lists do not contain numeric values.
\end{itemize}
Examples:
\begin{verbatim}
f.special.meijerG([0], [0], 1)   #c# Returns 1.0  (delta function case)
\end{verbatim}

\subsubsection{Fox H-Function}
Computes the Fox H-function, a generalization of the Meijer G-function allowing non-integer parameters. \\
Morphology:
\begin{verbatim}
EBNF:
"f.special.foxH", "(", list, ",", list, ",", ( int | real | complex ), ")"
\end{verbatim}
\begin{itemize}
\item Operands: numerator parameters (\texttt{list}), denominator parameters (\texttt{list}), argument (\texttt{int}, \texttt{real}, or \texttt{complex}).
\item Result: \texttt{complex}.
\item A \texttt{Type Error} is raised if the lists do not contain numeric values.
\end{itemize}
Examples:
\begin{verbatim}
f.special.foxH([[1,1]], [[1]], 1)   #c# Returns approximately 1.0
\end{verbatim}

\subsubsection{Jacobi Theta Function}
Computes the Jacobi theta function $\vartheta_n(\tau, z)$ for $n \in \{1, 2, 3, 4\}$. \\
Morphology:
\begin{verbatim}
EBNF:
"f.special.jacobiTheta", "(", int, ",", ( int | real ), ",", ( int | real ), ")"
\end{verbatim}
\begin{itemize}
\item Operands: index $n$ (\texttt{int}, must be 1--4), $\tau$ (\texttt{int} or \texttt{real}), $z$ (\texttt{int} or \texttt{real}).
\item Result: \texttt{complex}.
\item A \texttt{Value Error} is raised if $n \notin \{1, 2, 3, 4\}$.
\end{itemize}
Examples:
\begin{verbatim}
f.special.jacobiTheta(1, 0, 0)   #c# Returns 0.0
f.special.jacobiTheta(3, 0, 0)   #c# Returns approximately 1.0
\end{verbatim}

\subsubsection{Complete Elliptic Integral of the First Kind}
Computes $K(m) = \int_0^{\pi/2} \frac{d\theta}{\sqrt{1 - m \sin^2\theta}}$. \\
Morphology:
\begin{verbatim}
EBNF:
"f.special.ellipticK", "(", ( int | real ), ")"
\end{verbatim}
\begin{itemize}
\item Operand: parameter $m$ (\texttt{int} or \texttt{real}).
\item Result: \texttt{real}.
\item A \texttt{Value Error} is raised for $m \geq 1$ (singularity).
\end{itemize}
Examples:
\begin{verbatim}
f.special.ellipticK(0)     #c# Returns approximately 1.5708 (pi/2)
f.special.ellipticK(0.5)   #c# Returns approximately 1.8541
\end{verbatim}

\subsubsection{Complete Elliptic Integral of the Second Kind}
Computes $E(m) = \int_0^{\pi/2} \sqrt{1 - m \sin^2\theta}\,d\theta$. \\
Morphology:
\begin{verbatim}
EBNF:
"f.special.ellipticE", "(", ( int | real ), ")"
\end{verbatim}
\begin{itemize}
\item Operand: parameter $m$ (\texttt{int} or \texttt{real}).
\item Result: \texttt{real}.
\end{itemize}
Examples:
\begin{verbatim}
f.special.ellipticE(0)     #c# Returns approximately 1.5708 (pi/2)
f.special.ellipticE(0.5)   #c# Returns approximately 1.3506
\end{verbatim}

\subsubsection{Complete Elliptic Integral of the Third Kind}
Computes $\Pi(n, m) = \int_0^{\pi/2} \frac{d\theta}{(1 - n \sin^2\theta)\sqrt{1 - m \sin^2\theta}}$. \\
Morphology:
\begin{verbatim}
EBNF:
"f.special.ellipticPi", "(", ( int | real ), ",", ( int | real ), ")"
\end{verbatim}
\begin{itemize}
\item Operands: characteristic $n$ (\texttt{int} or \texttt{real}), parameter $m$ (\texttt{int} or \texttt{real}).
\item Result: \texttt{real}.
\item A \texttt{Value Error} is raised for $m \geq 1$ or $n \geq 1$ (singularities).
\end{itemize}
Examples:
\begin{verbatim}
f.special.ellipticPi(0, 0.5)   #c# Returns approximately 1.8541 (= K(0.5))
f.special.ellipticPi(0.5, 0.5) #c# Returns approximately 2.1565
\end{verbatim}

\subsubsection{Jacobi Elliptic Functions}
Computes the Jacobi elliptic functions $\text{sn}(u, m)$, $\text{cn}(u, m)$, $\text{dn}(u, m)$. \\
Morphology:
\begin{verbatim}
EBNF:
"f.special.jacobiSN", "(", ( int | real ), ",", ( int | real ), ")"
"f.special.jacobiCN", "(", ( int | real ), ",", ( int | real ), ")"
"f.special.jacobiDN", "(", ( int | real ), ",", ( int | real ), ")"
\end{verbatim}
\begin{itemize}
\item Operands: argument $u$ (\texttt{int} or \texttt{real}), parameter $m$ (\texttt{int} or \texttt{real}).
\item Result: \texttt{real}.
\item Identity: $\text{sn}^2 + \text{cn}^2 = 1$, $m\,\text{sn}^2 + \text{dn}^2 = 1$.
\end{itemize}
Examples:
\begin{verbatim}
f.special.jacobiSN(0, 0.5)   #c# Returns 0.0
f.special.jacobiCN(0, 0.5)   #c# Returns 1.0
f.special.jacobiDN(0, 0.5)   #c# Returns 1.0
f.special.jacobiSN(1, 0.5)   #c# Returns approximately 0.7997
\end{verbatim}

\subsubsection{Weierstrass Elliptic Function}
Computes the Weierstrass elliptic function $\wp(z; g_2, g_3)$, the inverse of the elliptic integral. \\
Morphology:
\begin{verbatim}
EBNF:
"f.special.weierstrassP", "(", ( int | real | complex ), ",", ( int | real | complex ), ",", ( int | real | complex ), ")"
\end{verbatim}
\begin{itemize}
\item Operands: $z$ (\texttt{int}, \texttt{real}, or \texttt{complex}), invariants $g_2$ and $g_3$ (\texttt{int}, \texttt{real}, or \texttt{complex}).
\item Result: \texttt{complex}.
\end{itemize}
Examples:
\begin{verbatim}
f.special.weierstrassP(0, 1, 0)   #c# Returns infinity (pole at z=0)
f.special.weierstrassP(1, 1, 0)   #c# Returns approximately 1.1674
\end{verbatim}

\subsubsection{Lower Incomplete Gamma Function}
Computes $\gamma(s, x) = \int_0^x t^{s-1} e^{-t}\,dt$. \\
Morphology:
\begin{verbatim}
EBNF:
"f.special.incompleteGamma", "(", ( int | real ), ",", ( int | real ), ")"
\end{verbatim}
\begin{itemize}
\item Operands: $s$ (\texttt{int} or \texttt{real}), $x$ (\texttt{int} or \texttt{real}).
\item Result: \texttt{real}.
\item $\gamma(s, 0) = 0$, $\gamma(s, \infty) = \Gamma(s)$.
\item A \texttt{Value Error} is raised if $s \leq 0$ and $x = 0$.
\end{itemize}
Examples:
\begin{verbatim}
f.special.incompleteGamma(1, 0)   #c# Returns 0.0
f.special.incompleteGamma(1, 1)   #c# Returns approximately 0.6321 (1 - 1/e)
\end{verbatim}

\subsubsection{Upper Incomplete Beta Function}
Computes $B(a, b; x) = \int_x^1 t^{a-1} (1-t)^{b-1}\,dt$, the complement of the regularized incomplete beta function. \\
Morphology:
\begin{verbatim}
EBNF:
"f.special.incompleteBeta", "(", ( int | real ), ",", ( int | real ), ",", ( int | real ), ")"
\end{verbatim}
\begin{itemize}
\item Operands: $a$ (\texttt{int} or \texttt{real}), $b$ (\texttt{int} or \texttt{real}), $x$ (\texttt{int} or \texttt{real}).
\item Result: \texttt{real}.
\item $B(a, b; 0) = B(a, b)$, $B(a, b; 1) = 0$.
\item A \texttt{Value Error} is raised if $a$ or $b$ is a non-positive integer.
\end{itemize}
Examples:
\begin{verbatim}
f.special.incompleteBeta(1, 1, 0)   #c# Returns 1.0  (= B(1,1) = 1)
f.special.incompleteBeta(1, 1, 1)   #c# Returns 0.0
f.special.incompleteBeta(2, 3, 0.5) #c# Returns approximately 0.3125
\end{verbatim}

\section{User-defined Functions}
User-defined functions allow creation of reusable code blocks with local scoping.

\subsection{Function Definition}
Functions are defined using the \texttt{function} keyword with the \texttt{f.u.} prefix:
\begin{verbatim}
EBNF:
function_definition = "function", "f.u.", id, 
  "(", [ { id, "," }, id], ")", "::", "{", 
  f_statement, {f_statement}, "}"

f_statement = statement | "return", [data]

Example:
function f.u.add( $a, $b )::{
    $retval [int]= $a + $b
}
\end{verbatim}

Function parameters are local variables scoped to the function body. Parameters may optionally include type annotations using the \texttt{[type]} syntax, consistent with lambda parameters. To return a value, set the \texttt{\$retval} variable.

\subsection{Function Calls}
User-defined functions are called using the \texttt{f.u.} prefix:
\begin{verbatim}
EBNF:
function_call_ud = "f.u.", id, "(", 
                   [{ data, "," }, data], ")"

Example:
$result [int]= f.u.add( 3, 5 )   #c# $result = 8
\end{verbatim}

\subsection{Return Values}
Functions return values by setting the \texttt{\$retval} variable. A bare \texttt{return} without a value, or a function that never sets \texttt{\$retval}, returns \texttt{None}.
\begin{verbatim}
function f.u.square( $x )::{
    $retval [int]= $x ^ 2
}

function f.u.greet( $name )::{
    f.util.print( f.str.str( "Hello, ", $name, "!" ) )
    #c# No return value
}

$sq [int]= f.u.square( 4 )     #c# $sq = 16
f.u.greet( "World" )          #c# Prints: Hello, World!
\end{verbatim}

\subsection{Local Scoping}
\label{sec:function-scope}
Functions have local scoping. Variables declared inside a function are only accessible within that function. Parameters are also local variables.

\paragraph{Scope Rules}
\begin{itemize}
\item Parameters (\$a, \$b, etc.) are local to the function body
\item Variables declared inside the function with \texttt{\$var [type]= val} are local
\item Local variables shadow global variables of the same name
\item The \texttt{\$retval} variable is always local to the function
\item Functions cannot access local variables from other function invocations
\end{itemize}

\paragraph{Return Values}
Functions can return values using two equivalent mechanisms:
\begin{itemize}
\item Assign to \texttt{\$retval} and let execution fall through the closing \texttt{\}}
\item Use the \texttt{return} keyword (see Chapter 3, \texttt{f\_statement} rule), which is syntactic sugar that assigns its argument to \texttt{\$retval} and immediately exits the function
\end{itemize}
Both \texttt{return \$value} and \texttt{\$retval [type]= \$value} followed by natural function exit are valid and equivalent. If both are used, \texttt{return} takes precedence.

\paragraph{Examples}
\begin{verbatim}
$x [int]= 10      #c# Global variable

function f.u.doublen()::{
    $x [int]= $x * 2  #c# Local $x shadows global
    f.util.print($x)            #c# Prints: 20
}

f.u.doublen()
f.util.print($x)          #c# Prints: 10 (global unchanged)

function f.u.addLocal( $a, $b )::{
    $sum [int]= $a + $b   #c# $sum is local
    $retval [int]= $sum
}

f.u.addLocal( 3, 4 )   #c# Returns 7
#c# $sum is not accessible here
\end{verbatim}

\paragraph{Recursion}
Local scoping enables recursion since each call has its own copy of parameters:
\begin{verbatim}
function f.u.factorial( $n )::{
    if $n <= 1 ::{
        $retval [int]= 1
    } else ::{
        $retval [int]= $n * 
            f.u.factorial( $n - 1 )
    }
}

$result [int]= f.u.factorial( 5 )  #c# = 120
\end{verbatim}




\part{Formal Semantics}
